% Options for packages loaded elsewhere
\PassOptionsToPackage{unicode}{hyperref}
\PassOptionsToPackage{hyphens}{url}
\documentclass[
  11pt,
]{article}
\usepackage{xcolor}
\usepackage[margin=1in]{geometry}
\usepackage{amsmath,amssymb}
\setcounter{secnumdepth}{-\maxdimen} % remove section numbering
\usepackage{iftex}
\ifPDFTeX
  \usepackage[T1]{fontenc}
  \usepackage[utf8]{inputenc}
  \usepackage{textcomp} % provide euro and other symbols
\else % if luatex or xetex
  \usepackage{unicode-math} % this also loads fontspec
  \defaultfontfeatures{Scale=MatchLowercase}
  \defaultfontfeatures[\rmfamily]{Ligatures=TeX,Scale=1}
\fi
\usepackage{lmodern}
\ifPDFTeX\else
  % xetex/luatex font selection
\fi
% Use upquote if available, for straight quotes in verbatim environments
\IfFileExists{upquote.sty}{\usepackage{upquote}}{}
\IfFileExists{microtype.sty}{% use microtype if available
  \usepackage[]{microtype}
  \UseMicrotypeSet[protrusion]{basicmath} % disable protrusion for tt fonts
}{}
\makeatletter
\@ifundefined{KOMAClassName}{% if non-KOMA class
  \IfFileExists{parskip.sty}{%
    \usepackage{parskip}
  }{% else
    \setlength{\parindent}{0pt}
    \setlength{\parskip}{6pt plus 2pt minus 1pt}}
}{% if KOMA class
  \KOMAoptions{parskip=half}}
\makeatother
\setlength{\emergencystretch}{3em} % prevent overfull lines
\providecommand{\tightlist}{%
  \setlength{\itemsep}{0pt}\setlength{\parskip}{0pt}}
\usepackage{mathrsfs}
\usepackage{bookmark}
\IfFileExists{xurl.sty}{\usepackage{xurl}}{} % add URL line breaks if available
\urlstyle{same}
\hypersetup{
  pdftitle={The retained cusp{,} spatial refinement{,} and the spectral measure of the actual magnetic vacuum translation},
  hidelinks,
  pdfcreator={LaTeX via pandoc}}

\title{The retained cusp, spatial refinement, and the spectral measure
of the actual magnetic vacuum translation}
\author{}
\date{8 September 2026}

\begin{document}
\maketitle

\subsection{1. Objects and scope of the
calculation}\label{objects-and-scope-of-the-calculation}

We calculate simultaneous geometric and coupling sequences in the
original full spatial SU(2) theory. The first retains an arbitrary fixed
positive coupling. For that sequence we prove an explicit vanishing
bound for the actual, unmodified, centered magnetic state. The second
retains a strictly positive, explicitly varying coupling at every
regulator. For the second sequence we prove convergence of the entire
excitation spectral probability after recording its energy dilation,
escape of that probability from every bounded physical energy interval,
convergence of configuration measures under exact refinement maps, and a
failure of strong continuity for the resulting unscaled temporal
correlations. The estimates use the actual vacuum at every regulator. In
particular, a constant function is used as a comparison vector in
inequalities, never as a replacement vacuum.

These results do not construct a four-dimensional interacting continuum
counterexample. The vanishing bound alone does not determine the
fixed-coupling spectral probability after division by its actual nonzero
squared norm. Sections 9--11 continue the calculation: a volume-uniform
gap proves full magnetic spectral escape on an explicit fixed-coupling
interval; an all-positive-coupling argument proves loss of high-energy
mass and failure of direct strongly continuous dynamics for fixed
physical Wilson states; explicit decreasing-coupling trajectories retain
their full coefficients. Sections 12--13 derive actual weak-coupling
energy and magnetic-state bounds and construct a further exact
pure-gauge vacuum configuration limit. Sections 14--16 prove explicit
nonzero covariance and native-mass lower bounds, then construct a
corrected cusp on which the entire native physical spectral probability
approaches the actual electric covariance probability in total variation
on two decreasing-coupling paths. The second sequence provides a proved
obstruction to the specific
small-interaction-relative-to-the-total-electric-operator route, with an
exact map to its limiting configuration measure. Sections 17--18 then
prove the actual weighted-electric graph limit, compute its full
physical pair spectrum, and construct a selected weak-coupling diagonal
on which the finite gaps close while the corrected native probability
escapes finite energies. All constants in the Hamiltonian, including the
scalar Wilson constant, remain displayed.

Section 26 proves a further original-state result without changing the
cusp depth: nonlinear separation of the full plaquette potential,
arbitrary potential moments on the actual physical finite-energy space,
and a polynomial raw-norm lower bound give native finite-energy
probability smaller than every coupling power at fixed nonzero angle and
box. A common actual coupling sequence retains \(T_j=j^2\) and all
earlier convergence tests. Its native probability escapes, and becomes
maximally separated in total variation from its electric covariance
probability; their raw vectors become orthogonal after their explicitly
recorded norm ratios. The actual relative angular Taylor remainder tends
to one on that sequence. These are results for the stated sequence, with
the prescribed paths left open.

For an integer \(L\ge2\), the spatial vertices are
\(\{-L,\ldots,L\}^3\). We retain every positive edge whose two endpoints
lie in the box and every elementary face. Let \[
 N_L=3(2L)(2L+1)^2,
 \qquad M_L=3(2L)^2(2L+1)=12L^2(2L+1).
\] The configuration space is \(Q_L=SU(2)^{N_L}\), with product Haar
probability \(\lambda_L\). Reverse traversal means the inverse of the
same link. A face is traversed as \[
 U_p=U_i(n)U_j(n+e_i)U_i(n+e_j)^{-1}U_j(n)^{-1},
 \qquad W_p=\operatorname{tr}U_p,\quad i<j.
\] With \(T_a=-i\sigma_a/2\), let \(X_{e,a}\) be differentiation under
left multiplication by \(\exp(tT_a)\), and set \[
 E_e=-\sum_{a=1}^3X_{e,a}^2,\quad H_0=\sum_eE_e,
 \quad \mathcal W=\sum_pW_p,
\] \[
 H_{L,a,g}=\kappa H_0+b\sum_p(2-W_p)
 =\kappa(H_0-\xi\mathcal W)+2bM_L I,
\quad \kappa=\frac{2g^2}{a},\quad b=\frac1{2g^2a},\quad
 \xi=\frac1{4g^4}.
\tag{1}
\] The physical Lie metric is \(c(X,Y)=-\operatorname{tr}(XY)/2\), so
\(\Delta^c=4\Delta^T\). Equation (1) is therefore also the Hamiltonian
with electric part \(-g^2(2a)^{-1}\sum_e\Delta_e^c\). The positive real
numbers \(a,g\) have not been absorbed into any new metric.

The vertex group acts by \(U_e\mapsto g_{s(e)}U_eg_{t(e)}^{-1}\). Its
Haar average \(\Pi_L\) is an orthogonal projection because Haar measure
is invariant under inversion and multiplication. The operator \(H\)
commutes with it: left and right multiplication preserve each link
Laplacian, and each face word is conjugated at its base vertex. Its
range is denoted \(\mathcal H_L^{\mathrm{inv}}\).

The operator \(H_0\) is the product compact-group Laplacian, with
operator domain \(H^2(Q_L)\) and form domain \(H^1(Q_L)\). The potential
is smooth, bounded and real, with \(0\le b\sum_p(2-W_p)\le4bM_L\). Hence
\(H\) has the same domains and compact resolvent. A ground minimizer can
be made nonnegative by the gradient inequality for its absolute value;
the smooth approximants \((|f|^2+\varepsilon^2)^{1/2}\) justify that
inequality at zeros. Elliptic regularity and the strong maximum
principle give a smooth strictly positive ground vector \(\psi\), with
\(\int\psi^2d\lambda_L=1\). For smooth \(F\), the product rule and Haar
integration by parts give \[
 \langle\psi F,(H-E)\psi F\rangle
 =\kappa\sum_{e,a}\int\psi^2|X_{e,a}F|^2d\lambda_L.
\tag{2}
\] If another ground vector existed, its quotient by \(\psi\) would have
every derivative zero by (2), and thus would be constant on connected
\(Q_L\). This proves simplicity. Gauge transformations preserve
positivity and the chosen norm, so \(\Pi_L\psi=\psi\). All quotients by
\(\psi\) below are valid on the finite compact manifold, where \(\psi\)
has a positive minimum.

\subsection{2. Original cusp and the complete physical
dictionary}\label{original-cusp-and-the-complete-physical-dictionary}

The canonical cusp supplies holomorphic functions
\(\mu(t_c),h(t_c),b_{\mathrm{per}}(t_c)\), with the exact period
identity \[
 Z(s)=s\begin{pmatrix}0&1\\-1&0\end{pmatrix}
 +\begin{pmatrix}6\mu&h\\b_{\mathrm{per}}-h&\mu\end{pmatrix},
 \qquad t_c=e^{2\pi i s}.
\] The notation \(b_{\mathrm{per}}\) records the canonical period
function called \(b\) in the cusp source; it is distinct from the
Hamiltonian coefficient in (1). On \(s=iT\), write \[
 \mu_T=\alpha_T+im_T,\quad h_T=\eta_T+iH_T,
 \quad b_{\mathrm{per},T}=\zeta_T+iB_T^{\mathrm{per}},
\] \[
 L_T=T+H_T,\quad q_T=L_T-B_T^{\mathrm{per}},\quad
 c_T=\zeta_T-\eta_T,\quad D_T=L_Tq_T+6m_T^2>0.
\] In original real coordinate order
\((x_1,x_2,x_3,x_4)=(\operatorname{Re}z_1,\operatorname{Im}z_1,\operatorname{Re}z_2,\operatorname{Im}z_2)\),
the exact matrix is \[
 P_T=\begin{pmatrix}
 6\alpha_T&\eta_T&1&0\\6m_T&L_T&0&0\\
 c_T&\alpha_T&0&1\\-q_T&m_T&0&0
 \end{pmatrix},\qquad \det P_T=-D_T.
\tag{3}
\] Every bounded period correction remains in this matrix. Holomorphy
gives convergence of its bounded entries with errors \(O(e^{-2\pi T})\).
In particular, with \[
 A_T=\begin{pmatrix}6\alpha_T&\eta_T\\c_T&\alpha_T\end{pmatrix},
 \quad B_T=\begin{pmatrix}6m_T&L_T\\-q_T&m_T\end{pmatrix}
 =TJ+C_T,\quad J=\begin{pmatrix}0&1\\-1&0\end{pmatrix},
\] there is a finite \(C=\sup_{T\ge T_0}\|C_T\|\), after restricting to
the specified cusp ray. Thus \[
 D_T\ge(T-C)^2\quad(T>C),\qquad D_T=T^2+O(T).
\tag{4}
\] For the first inequality, \(|B_Tv|\ge(T-C)|v|\) bounds both singular
values from below; their product is the positive determinant \(D_T\).
The second follows by multiplying the retained expressions for
\(L_T,q_T,m_T\).

For a cover integer \(M_{\mathrm{cov}}\), set
\(S_M=\operatorname{diag}(1,1,M_{\mathrm{cov}},M_{\mathrm{cov}})\). The
torus covering is induced by \[
 P_TS_M\mathbb Z^4\subset P_T\mathbb Z^4.
\] It has fibre degree \(M_{\mathrm{cov}}^2\), determinant
\(-M_{\mathrm{cov}}^2D_T\), metric \(G=(P_TS_M)^{\mathsf T}P_TS_M\), and
positive volume density \(M_{\mathrm{cov}}^2D_T\). In the recorded
permutation \((x_1,x_3;x_2,x_4)\), its period matrix and inverse are \[
 Q=\begin{pmatrix}A_T&M_{\mathrm{cov}}I_2\\B_T&0\end{pmatrix},
 \quad Q^{-1}=\begin{pmatrix}0&B_T^{-1}\\
 M_{\mathrm{cov}}^{-1}I_2&-M_{\mathrm{cov}}^{-1}A_TB_T^{-1}\end{pmatrix},
 \quad B_T^{-1}=D_T^{-1}\begin{pmatrix}m_T&-L_T\\q_T&6m_T\end{pmatrix}.
\tag{5}
\] Multiplication verifies both inverse identities. The permutation in
(5) has determinant \(-1\); it records coordinates and does not alter
the determinant of (3). Formula (5) retains the complete inverse metric
as \(Q^{-1}Q^{-\mathsf T}\).

The original marked monodromy is \[
 M_0=I+R_0,\qquad
 R_0=\begin{pmatrix}0&0&0&0\\0&0&0&0\\0&1&0&0\\-1&0&0&0\end{pmatrix}.
\] Indeed it sends \(\widehat\gamma\) to
\(\widehat\gamma-\widehat\delta\), \(\widehat u\) to
\(\widehat u+\widehat w\), and fixes \(\widehat w,\widehat\delta\).
Direct multiplication gives \[
 R_0^2=0,\qquad S_M^{-1}M_0^{M_{\mathrm{cov}}}S_M=M_0.
\tag{6}
\] Thus the degree-\(M_{\mathrm{cov}}^2\) fibre cover descends after the
base change \(t_c=v^{M_{\mathrm{cov}}}\). The proof is that continuation
around \(v=0\) acts on the original lattice by
\(M_0^{M_{\mathrm{cov}}}\), which (6) carries to an integral
automorphism in the cover basis. No cover of the entire original base or
completed compact threefold is required here.

A nonzero covered period, in (5)'s order, is
\((A_Tn+M_{\mathrm{cov}}z,B_Tn)\), for \(n,z\in\mathbb Z^2\). If
\(n\ne0\), its length is at least \(T-C\); if \(n=0\), its length is at
least \(M_{\mathrm{cov}}\). Therefore \[
 \operatorname{sys}\ge\min(T-C,M_{\mathrm{cov}}).
\tag{7}
\] When this bound exceeds \(4R\), the radius-\(R\) Euclidean ball is
mapped isometrically to the torus: for two points in that ball and a
nonzero period \(\lambda\), \(|x-y+\lambda|\ge|\lambda|-2R>2R\ge|x-y|\).

The physical time and spatial coordinates used for the Hamiltonian are
\[
 (y_0,y_1,y_2,y_3)=(x_3,x_2,x_4,x_1).
\tag{8}
\] This permutation has determinant \(+1\) and preserves the exact
Euclidean metric. The physical magnetic plane is \((y_1,y_2)\). From
(5), \[
 u_1=D_T^{-1}(m_Ty_1-L_Ty_2),\quad
 u_2=D_T^{-1}(q_Ty_1+6m_Ty_2),\quad
 du_1\wedge du_2=D_T^{-1}dy_1\wedge dy_2.
\] The original bundle connection \(d+2\pi u_1H_cdu_2\), with
\(H_c=i\sigma_3=-2T_3\), therefore has the local expression \[
 A_{\mathrm{orig}}=\frac{2\pi H_c}{D_T^2}
 (m_Ty_1-L_Ty_2)(q_Tdy_1+6m_Tdy_2).
\] Put \[
 f_T(y)=\frac{2\pi}{D_T^2}
 \left(\frac{m_Tq_T}{2}y_1^2-L_Tq_Ty_1y_2-3m_TL_Ty_2^2\right).
\] Differentiation, with all signs retained, proves \[
 A_{\mathrm{orig}}=\frac{2\pi}{D_T}H_cy_1dy_2+H_cdf_T.
\] The two gauge maps \(\exp(-H_cf_T)\) and
\(\exp(-2\pi H_cy_1y_2/D_T)\), under \(A^h=h^{-1}Ah+h^{-1}dh\), give
exactly \[
 A_1=-\frac{2\pi}{D_T}H_cy_2dy_1,\qquad
 F_{12}=\frac{2\pi}{D_T}H_c.
\tag{9}
\] The original transitions are \(\exp(-2\pi k_1u_2H_c)\); on a
nonwrapping embedded ball they are represented by the displayed local
frames. The flat reference is chosen as the identity transport in the
final frame. Subsequent frame changes act on the reference and magnetic
transport simultaneously. For any such pair \(M_e,R_e\), their relative
transport \(h_e=M_eR_e^{-1}\) changes by conjugation at the source,
proving covariance of the translated state under changes of vertex
frames.

The inverse transport of (9) on the edge from \(an\) to \(a(n+e_i)\) is
\[
 h_1(n)=\exp(-\theta n_2H_c),\quad h_2(n)=h_3(n)=I,
 \qquad \theta=\frac{2\pi a^2}{D_T}.
\tag{10}
\] Integration is exact because the coefficients commute along the edge.
The face word is \(\exp(\theta H_c)\) in the 12 plane and identity in
the other two planes.

For general connections, the map to the period coordinates is the usual
exact pullback:
\(\widetilde A_i=\sum_\alpha(P_TS_M)_{\alpha i}A_\alpha(P_TS_My)\), and
\(\widetilde F\) is its exterior-square pullback. The chain rule
preserves \(dA\), and bilinearity preserves every term of \(A\wedge A\).
Its action is \[
 \frac{M_{\mathrm{cov}}^2D_T}{2g^2}
 \int_{[0,1)^4}\sum_{i<j,\,k<l}
 \bigl[(G^{-1})_{ik}(G^{-1})_{jl}-(G^{-1})_{il}(G^{-1})_{jk}\bigr]
 c(\widetilde F_{ij},\widetilde F_{kl})\,d^4y.
\tag{11}
\] This follows by the Jacobian formula for volume and the Gram
determinant for two covectors. It includes all mixed terms. Pullback
carries the parallel-transport differential equation to the identical
equation on each path; uniqueness gives the holonomy map, including for
noncommuting fields. On the embedded balls in (7), this pullback is a
local isometry and preserves the full action of every compactly
supported connection, extended by zero through a collar. Thus no
nonlinear interaction is lost in passing to a physical cubical patch.

\subsubsection{2.1 A simultaneous geometric exhaustion in unchanged
physical
units}\label{a-simultaneous-geometric-exhaustion-in-unchanged-physical-units}

Let \(j=2^r\), for integers \(r\) sufficiently large, and set \[
 T_j=j^2,\quad t_{c,j}=e^{-2\pi j^2},\quad
 M_{\mathrm{cov},j}=j,\quad v_j=e^{-2\pi j},
\] \[
 a_j=\frac1{100j},\quad L_j=j^2,\quad
 \ell_j=2a_jL_j=\frac j{50},\quad
 D_j=L_{j^2}q_{j^2}+6m_{j^2}^2,\quad
 \theta_j=\frac{2\pi}{10^4j^2D_j}.
\tag{12}
\] Here subscripts on \(L_{j^2},q_{j^2},m_{j^2}\) denote the cusp
functions, whereas \(L_j=j^2\) on the preceding line is the integer
half-width of the spatial graph. In particular \(D_j\) is never replaced
by \(j^4\). The exact identity \(v_j^j=t_{c,j}\) and (6) prove monodromy
compatibility. The fibre degree is \(j^2\) and the base-change degree is
\(j\).

The four-dimensional box \(|y_\mu|\le j/100\), with time coordinate (8),
has radius \(R_j=j/50\). For \(j^2-C\ge j\), (7) gives
\(\operatorname{sys}\ge j>8R_j\). Hence the spatial box at \(y_0=0\),
the growing time interval, and a collar all embed in the actual covered
cusp. The mesh tends to zero while the physical box tends to all of
\(\mathbb R^3\), and the time interval tends to \(\mathbb R\).

The fixed physical reference length is the length \(\ell_*\) of the
original real period vector \(e_1\) in (3). Its value in the original
Euclidean metric is exactly 1. We use its reciprocal \(E_*=\ell_*^{-1}\)
as a reference inverse length for energies. Neither \(\ell_*\), the
metric, nor time (8) varies with \(j\). Thus \(a_j/\ell_*=1/(100j)\),
the reference momentum \(p_*=2\pi/\ell_*\) stays fixed, and the finite
energy interval \([0,\Omega E_*]\) means the same physical interval at
every regulator. These are reference units, not an assertion that the
quantum theory has a particle of mass \(E_*\).

The exact reciprocal-vector approximation can also retain this momentum.
For a physical covector \(p\), choose each integer component of \(k_j\)
by rounding \((P_{j^2}S_j)^{\mathsf T}p/(2\pi)\). Then \[
 p_j=2\pi(P_{j^2}S_j)^{-\mathsf T}k_j,\qquad
 |p_j-p|\le\pi\sum_{i=1}^4|(P_{j^2}S_j)^{-\mathsf T}e_i|.
\] Formula (5) bounds the right side by a constant times
\(j^{-2}+j^{-1}\), with the original \(A_T,B_T\) retained. This proves
that fixed physical spatial or temporal momenta are geometrically
resolvable. It is not a dispersion relation for the interacting
Hamiltonian.

The magnetic field itself supplies another exact physical length and
inverse length, \[
 \ell_{B,j}=\sqrt{\frac{D_j}{2\pi}},\qquad
 E_{B,j}=\ell_{B,j}^{-1}=\sqrt{\frac{2\pi}{D_j}}.
\] They are defined by the retained curvature norm in (9), so \[
 \frac{a_j}{\ell_{B,j}}=\sqrt{\theta_j},\qquad
 \frac{\ell_j}{\ell_{B,j}}=\frac{j\sqrt{2\pi}}{50\sqrt{D_j}}
 \sim\frac{\sqrt{2\pi}}{50j},\qquad
 E_{B,j}\sim\frac{\sqrt{2\pi}}{j^2}.
\] Thus the box exhausts physical space while remaining small relative
to the growing magnetic length. The flux through its full magnetic
square is exactly \(2\pi\ell_j^2/D_j\), which tends to zero. These
equations identify the geometric soft scale without equating it to a
quantum excitation energy. Every finite cover with its pullback metric
preserves the local value \(2\pi/D_j\) of the curvature, since its
differential is a local isometry. Consequently no choice of cover degree
can hold this original nonzero magnetic scale fixed while
\(T_j\to\infty\). A varying unit of length would have the additional
metric and energy effects calculated in (61); it cannot be silently used
to keep this field scale constant.

\subsection{3. Exact spectral measure and a bound at every positive
coupling}\label{exact-spectral-measure-and-a-bound-at-every-positive-coupling}

Let \((T_hF)(U)=F((h_e^{-1}U_e)_e)\). Haar invariance makes this a
unitary operator; its electric commutator is zero by bi-invariance. Set
\[
 c_\theta=\langle\psi,T_h\psi\rangle,
 \quad \chi_\theta=\Pi_LT_h\psi-c_\theta\psi,
 \quad d_\theta=\|\chi_\theta\|^2.
\tag{13}
\] The positive overlap is real, and \(\chi_\theta\) is invariant and
orthogonal to \(\psi\). The full magnetic translation has at most one
nonidentity translated edge at each source vertex. Conjugating it by
gauge transformations therefore gives an independent conjugacy average
on each such link. On invariant inputs, \[
 \Pi_LT_h\Pi_L=K_\theta\Pi_L,\qquad
 K_\theta=\prod_{e=(n,1)}C_{e,\theta n_2}.
\tag{14}
\] Here \(C_{e,s}\) averages the left translations by the conjugacy
class of \(\exp(sH_c)\). Schur's lemma applied to spin \(q\) gives the
real multiplier \[
 c_q(s)=\frac1{2q+1}\sum_{m=-q}^q e^{2ims}
 =\frac{\sin((2q+1)s)}{(2q+1)\sin s},\qquad q\in\tfrac12\mathbb Z_{\ge0}.
\tag{15}
\] The finite sum defines the value when \(\sin s=0\). Inversion
preserves the conjugacy class, so each \(C\) and their commuting product
are self-adjoint contractions. They preserve constants and Haar
integrals, and commute with every \(E_e\) and with \(\Pi_L\). Equations
(13)-(14) thus give the exact state
\(\chi_\theta=(K_\theta-c_\theta)\psi\), with no angular expansion.

Write \(h_{L,\xi}=H_0-\xi\mathcal W\), let \(e_{L,\xi}\) be its ground
energy, and define \(A_{L,\xi}=h_{L,\xi}-e_{L,\xi}\). Then \[
 E=\kappa e_{L,\xi}+2bM_L,\qquad H-E=\kappa A_{L,\xi}.
\tag{16}
\] The same \(\psi\), \(\chi_\theta\), and \(d_\theta\) appear on both
sides. For \(d_\theta>0\), the two spectral probabilities are \[
 \rho_{L,\xi,\theta}(B)=d_\theta^{-1}
 \langle\chi_\theta,\mathbf1_B(A_{L,\xi})\chi_\theta\rangle,
\] \[
 \nu_{L,a,g,\theta}(B)=d_\theta^{-1}
 \langle\chi_\theta,\mathbf1_B(H-E)\chi_\theta\rangle
 =\rho_{L,\xi,\theta}(\kappa^{-1}B).
\tag{17}
\] This is the exact pushforward by \(\lambda\mapsto\kappa\lambda\). It
follows first for spectral intervals from the spectral resolution of a
scalar multiple, then for all Borel sets by countable additivity. The
finite measure before division is \(d_\theta\nu\); the state and its
amplitude have not been changed. In particular \[
 \langle\chi_\theta,e^{-t(H-E)}\chi_\theta\rangle
 =d_\theta\int e^{-t\kappa\lambda}d\rho(\lambda).
\tag{18}
\]

\subsubsection{3.1 An all-angle operator
inequality}\label{an-all-angle-operator-inequality}

Put \(\Gamma=\sum_{e=(n,1)}n_2^2E_e\) and \(J_\theta=K_\theta-I\). On a
joint spin block, (15) and \(1-\cos x\le x^2/2\) give \[
 |1-c_q(s)|=1-c_q(s)
 \le\frac{2s^2}{2q+1}\sum_{m=-q}^qm^2
 =\frac23s^2q(q+1).
\] The last finite-sum formula follows by pairing the weights and
summing consecutive squares, also for half-integral \(q\). For real
numbers \(|z_i|\le1\), telescoping the product gives
\(|1-\prod_i z_i|\le\sum_i|1-z_i|\). Therefore, on every joint spin
block, \[
 |1-k(\theta)|\le\frac23\theta^2\sum_{e=(n,1)}n_2^2q_e(q_e+1).
\] Parseval's identity and completeness of the joint matrix coefficients
prove, for every \(F\in D(H_0)\), \[
 \|J_\theta F\|\le\frac23\theta^2\|\Gamma F\|
 \le\frac23\theta^2L^2\|H_0F\|.
\tag{19}
\] The second inequality follows on the same joint blocks from
\(0\le\Gamma\le L^2H_0\); these are simultaneously diagonal positive
operators. This is a graph-domain estimate, not a bounded-operator
Taylor expansion.

The constant trial vector has zero electric energy and \(\int W_p=0\),
so \(0\le E/\kappa\le2\xi M_L\). The eigen-equation, with
\(0\le\sum_p(2-W_p)\le4M_L\), consequently gives \[
 \|H_0\psi\|
 =\|[E/\kappa-\xi\sum_p(2-W_p)]\psi\|
 \le4\xi M_L.
\tag{20}
\] Centering is an orthogonal projection off \(\psi\), so (19)-(20)
imply the all-coupling estimate \[
 \boxed{\quad\|\chi_\theta\|\le\frac83\theta^2L^2\xi M_L.\quad}
\tag{21}
\] This estimate concerns the full interacting vacuum and is valid
without any restriction on \(\xi M_L\).

For the fixed-coupling version of (12), set \(g_j=g_0>0\), hence
\(\xi_0=1/(4g_0^4)\), \(\kappa_j=200g_0^2j\), and \(b_j=50j/g_0^2\).
Once \(j^2\ge2C\), (4) yields \(D_j\ge j^4/4\). Since
\(M_{L_j}\le36j^6\), direct substitution in (21) gives \[
 \boxed{\quad\|\chi_j\|\le
 \frac{6144\pi^2\xi_0}{10^8}\,j^{-2},\qquad
 d_j\le\left(\frac{6144\pi^2\xi_0}{10^8}\right)^2j^{-4}.\quad}
\tag{22}
\] In deriving this bound, the exact angle satisfies
\(\theta_j^2L_j^2=4\pi^2/(10^8D_j^2)\); the determinant estimate is
applied only afterwards. Thus every unscaled excitation spectral measure
\(d_j\nu_j\) converges to the zero measure in total mass. Division by
\(d_j\) in (17) prevents this fact from deciding \(\nu_j\).

For completeness, the uncentered projected vectors also converge to
their respective vacua. Indeed \(|c_j-1|\le\|J_{\theta_j}\psi_j\|\), and
(19)-(20) give exactly the same right side as (22). If
\(F_j=K_{\theta_j}\psi_j\), then for each multiplication observable
\(f\) with \(\|f\|_\infty\le1\), \[
 \left|\int f(F_j^2-\psi_j^2)d\lambda_j\right|
 \le\|F_j-\psi_j\|(\|F_j\|+\|\psi_j\|)
 \le2\|J_{\theta_j}\psi_j\|.
\tag{23}
\] Here \(F_j^2d\lambda_j\) is the actual finite measure of mass
\(\|F_j\|^2\), which tends to 1; no mass convention has been changed.

\subsection{4. A volume-explicit estimate for the whole excitation
measure}\label{a-volume-explicit-estimate-for-the-whole-excitation-measure}

We now prove a relative estimate that remains useful as both the angle
and the volume vary. Its constants do not divide by an unspecified small
angular remainder.

\subsubsection{4.1 The exact first physical electric
eigenvalue}\label{the-exact-first-physical-electric-eigenvalue}

On one link, spin \(q\) has \(E_e\) eigenvalue \(q(q+1)\); for the
fundamental representation this is \(3/4\), obtained directly from
\(-\sum_aT_a^2=3I/4\). At a vertex, a gauge invariant tensor involving
exactly one nontrivial incident representation is impossible: that
irreducible representation has no fixed vector. Hence the support graph
of nonzero edge spins in any nonconstant gauge invariant joint-spin
block has no vertex of degree one. Any nonempty finite graph with this
property contains a cycle, as follows by following edges without
immediately reversing and stopping at the first repeated vertex. The
open cubical graph is simple and bipartite and has no triangular cycles,
so a cycle has at least four edges. Consequently its electric sum is at
least \(4(3/4)=3\). Equality is attained by the fundamental trace around
any elementary square, whose four links have spin \(1/2\).

Peter--Weyl decomposition and vertex averaging preserve each finite
joint-spin block and give a complete invariant basis. The preceding
argument therefore proves the form inequality \[
 H_0\big|_{\mathcal H_L^{\mathrm{inv}}\ominus\mathbb C1}\ge3I.
\tag{24}
\] The constant is an exact physical-sector gap for the comparison
operator \(H_0\), not an assumed gap for \(H\).

For distinct elementary faces, integrating a link occurring in just one
trace gives zero, since multiplication of that link by \(-I\) changes
the sign. Two distinct elementary squares have distinct four-edge
supports. A single face product is Haar, and the trace has second moment
1: in the unit-quaternion model the scalar coordinate has second moment
\(1/4\). Hence \[
 \int W_p=0,\quad\langle W_p,W_q\rangle=\delta_{pq},
 \quad H_0W_p=3W_p,\quad\|\mathcal W\|_2=\sqrt{M_L},
 \quad\|\mathcal W\|_\infty=2M_L.
\tag{25}
\] The last equality is attained at the identity configuration and holds
for the essential supremum by continuity and positive Haar measure of
neighborhoods.

The upper endpoint \(4bM_L\) of the full potential is attained as well.
Define central links \(z_i(n)=(-1)^{\sum_{k<i}n_k}I\). On a face
\(i<j\), shifting in direction \(i\) changes the sign of \(z_j\), while
shifting in direction \(j\) does not change \(z_i\). The product of the
four central signs is therefore \(-I\). Every face trace is \(-2\), so
the potential equals \(4bM_L\). Continuity again identifies this maximum
with the multiplication-operator norm. This proves the norm assertion
used in (39) without dropping any face.

Since \(h_{L,\xi}=H_0-\xi\mathcal W\), its ground eigenvalue satisfies
\(-2\xi M_L\le e\le0\). The min-max principle and (24)-(25) give the
first eigenvalue above its ground at least \(3-2\xi M_L\), before
subtraction of \(e\). Since \(e\le0\), \[
 \boxed{\quad A_{L,\xi}\big|_{\psi^\perp\cap\mathcal H_L^{\mathrm{inv}}}
 \ge(3-2\xi M_L)I.\quad}
\tag{26}
\] This lower bound is useful when its right side is positive. It is
valid for the entire invariant excitation space, independently of the
chosen magnetic state.

\subsubsection{4.2 A controlled comparison with the actual
vacuum}\label{a-controlled-comparison-with-the-actual-vacuum}

Write \(\psi=s1+\zeta\), with \(s=\int\psi>0\), \(\int\zeta=0\), and
\(s^2+\|\zeta\|^2=1\). Let \(Q\) be the projection off the constant
function. Put \(\varepsilon=\xi M_L\), and suppose for the bounds in
this subsection that \(0<\varepsilon\le1/2\). The exact eigen-equation
projected by \(Q\) is \[
 (H_0-e)\zeta=\xi s\mathcal W+\xi Q\mathcal W\zeta.
\] By (24), \(\|(H_0-e)^{-1}Q\|\le1/(3-e)\). Applying (25), rearranging,
and using \(s\le1,e\le0\) proves \[
 \|\zeta\|\le\frac{\xi\sqrt{M_L}}{3-2\varepsilon}
 \le\frac{\xi\sqrt{M_L}}2.
\tag{27}
\] Set \(A=\mathcal W/3\) and retain the exact remainder
\(r=\psi-1-\xi A\). The same equation gives \[
 H_0r=\xi(s-1)\mathcal W+e\zeta+\xi Q\mathcal W\zeta.
\] Since \(1-s=\|\zeta\|^2/(1+s)\le\|\zeta\|^2\),
\(|e|\le2\varepsilon\), and \(\xi\le1\), equation (27) yields \[
 \boxed{\quad\|H_0r\|\le
 \frac{\xi^3M_L^{3/2}}4+2\xi^2M_L^{3/2}
 \le\frac94\xi^2M_L^{3/2}.\quad}
\tag{28}
\] No convergence of a perturbation series, no unknown Sobolev constant,
and no approximation of the true probability density are used in
(27)-(28).

Define the explicit finite vector \[
 V_\theta=J_\theta A=\frac13\sum_p(\alpha_p(\theta)-1)W_p.
\tag{29}
\] The coefficients follow from (14)-(15) on the fundamental
representation of the two direction-1 edges of each face: \[
 \alpha_{n;12}=\cos(n_2\theta)\cos((n_2+1)\theta),\quad
 \alpha_{n;13}=\cos^2(n_2\theta),\quad\alpha_{n;23}=1.
\tag{30}
\] The two translated edges of such a face have distinct source
vertices, which justifies the product of the two averages in this
particular background. Thus \[
 \|V_\theta\|^2=\mathcal D_L(\theta)/9,\quad
 H_0V_\theta=3V_\theta,\quad\int V_\theta=0,
\] \[
 \mathcal D_L(\theta)=2L(2L+1)\sum_{m=-L}^{L-1}
 [1-\cos(m\theta)\cos((m+1)\theta)]^2
 +4L^2\sum_{m=-L}^{L}\sin^4(m\theta).
\tag{31}
\]

The actual centered state is nonzero for every \(b>0\) and
\(\theta\notin2\pi\mathbb Z\), without a small-coupling restriction. To
prove this, suppose \(\chi_\theta=0\). Haar integration in
\(K_\theta\psi=c_\theta\psi\) gives \(c_\theta=1\), since \(\int\psi>0\)
and \(K_\theta\) preserves that integral. Unitarity then gives
\(\|T_h\psi-\psi\|^2=2-2c_\theta=0\). Apply \(T_h\) to the exact
eigen-equation, use its electric commutation, and subtract the original
equation. The result is \(b(T_h\mathcal W-\mathcal W)\psi=0\). Strict
positivity and smoothness imply \(T_h\mathcal W=\mathcal W\) pointwise.
Its gauge average, (30), and the orthogonality (25) imply \(\alpha_p=1\)
for every face. A 12 face based at \(n_2=0\) exists and has
\(\alpha_p=\cos\theta\), so \(\theta\in2\pi\mathbb Z\). Conversely, at
these angles every link (10) is identity and \(\chi_\theta=0\). This
proves the exact zero set and, in particular, legitimizes (17) on both
sequences (12) for all sufficiently large indices.

We need a lower estimate with a completely displayed domain. Let
\(r_m=m^2+m+1/2\). The elementary inequality \[
 0\le t^2/2-(1-\cos t)\le t^4/24
\] follows from \(|\sin(t/2)|\le|t/2|\) and the Taylor integral
remainder. Applying it to \(x+y\) and \(x-y\) gives, for
\(u=1-\cos x\cos y\), \(q=(x^2+y^2)/2\), \[
 0\le u\le q,\quad 0\le q-u\le q^2/3,
 \quad0\le q^2-u^2\le2q^3/3.
\] The same last estimate holds for \(u=\sin^2x,q=x^2\). Substitution
into (31) gives \[
 \mathcal D_L(\theta)\ge\mathcal A_L\theta^4-16L^9\theta^6,
\] \[
 \mathcal A_L=
 \frac{L^2(2L+1)}{15}(24L^4+24L^3+8L^2-4L+3)
 \ge\frac{16}{5}L^7.
\tag{32}
\] For the exact polynomial, expand \(r_m^2\), sum, and use \[
 \sum_{m=-L}^{L-1}r_m^2=\frac{L(4L^4+1)}{10},\qquad
 \sum_{m=-L}^Lm^4=\frac{L(L+1)(2L+1)(3L^2+3L-1)}{15}.
\] Each identity is verified at \(L=1\); subtracting its value at \(L\)
from that at \(L+1\) gives exactly the two new endpoint summands,
proving it for every integer. The error coefficient is bounded by
\((2/3)[2L(2L+1)(2L)L^6+4L^2(2L+1)L^6]\le16L^9\), because
\(r_m\le L^2\). Expansion of \(\mathcal A_L\) gives
\((16/5)L^7+(24/5)L^6+(8/3)L^5+(2/15)L^3+(1/5)L^2\), proving its lower
bound.

Consequently, whenever \(0<\theta^2L^2\le1/10\), \[
 \|V_\theta\|\ge\frac13\theta^2L^{7/2}>0.
\tag{33}
\] Combining (19),(28),(33), and \(M_L\le36L^3\) gives \[
 B_\theta:=\|J_\theta\psi-\xi V_\theta\|
 \le\frac32\theta^2L^2\xi^2M_L^{3/2},
 \qquad\frac{B_\theta}{\xi\|V_\theta\|}\le27\varepsilon.
\tag{34}
\] Since \(J_\theta\) preserves Haar integrals and annihilates 1, \[
 c_\theta-1=\langle\zeta,J_\theta\psi\rangle.
\] Thus (13) and (34) give \[
 \|\chi_\theta-\xi V_\theta\|
 \le B_\theta+\|\zeta\|(\xi\|V_\theta\|+B_\theta).
\] For \(\varepsilon\le1/2\), (27) implies
\(\|\zeta\|\le\varepsilon/2\le1/4\). We have proved the explicit
relative estimate \[
 \boxed{\quad
 \|\chi_\theta-\xi V_\theta\|\le
 40\varepsilon\,\xi\|V_\theta\|,
 \qquad 0<\varepsilon\le\tfrac12,
 \quad0<\theta^2L^2\le\tfrac1{10}.
 \quad}
\tag{35}
\] The constant follows from \((1+1/4)27+1/2=137/4<40\). In particular,
when \(40\varepsilon<1\), the actual state is nonzero and \[
 (1-40\varepsilon)^2\frac{\xi^2\mathcal D_L(\theta)}9
 \le d_\theta\le
 (1+40\varepsilon)^2\frac{\xi^2\mathcal D_L(\theta)}9.
\tag{36}
\] This is a simultaneous bound, not an iterated fixed-volume
coefficient.

\subsubsection{4.3 Concentration of the full excitation
probability}\label{concentration-of-the-full-excitation-probability}

For \(\delta>0\), let \(P_\delta\) be the spectral projection of
\(A_{L,\xi}\) on \(\{\lambda:|\lambda-3|>\delta\}\). From (25),(29), \[
 \|(A_{L,\xi}-3)V_\theta\|
 =\|(-\xi\mathcal W-e)V_\theta\|
 \le4\varepsilon\|V_\theta\|.
\] The spectral theorem gives
\(\|P_\delta V_\theta\|\le(4\varepsilon/\delta)\|V_\theta\|\). Equation
(35) bounds its actual-state counterpart by \[
 \|P_\delta\chi_\theta\|
 \le(40+4/\delta)\varepsilon\xi\|V_\theta\|.
\] When \(\varepsilon\le1/100\), (36) makes
\(\|\chi_\theta\|\ge(1-40\varepsilon)\xi\|V_\theta\|\ge\tfrac12\xi\|V_\theta\|\).
Hence \[
 \boxed{\quad
 \rho_{L,\xi,\theta}(\{\lambda:|\lambda-3|>\delta\})
 \le(80+8/\delta)^2\varepsilon^2.
 \quad}
\tag{37}
\] This estimate treats all spectral channels of the interacting
operator. It imposes no cutoff on edge spins and discards none of the
off-diagonal Wilson matrix elements.

\subsection{5. Explicit joint sequence and its physical spectral
consequences}\label{explicit-joint-sequence-and-its-physical-spectral-consequences}

Keep all geometric data (12), and define the second coupling sequence by
\[
 M_j:=M_{L_j}=12j^4(2j^2+1),\qquad
 \xi_j=\frac1{j^2M_j},\qquad
 g_j^4=\frac{j^2M_j}{4},\qquad
 g_j^2=\frac j2\sqrt{M_j}.
\tag{38}
\] Every \(g_j\) is strictly positive and finite. Its actual Hamiltonian
coefficients are \[
 \kappa_j=100j^2\sqrt{M_j},\qquad
 b_j=\frac{100}{\sqrt{M_j}},\qquad
 2b_jM_j=200\sqrt{M_j},\qquad
 \xi_jM_j=j^{-2}.
\tag{39}
\] In particular the complete magnetic potential still has norm
\(4b_jM_j=400\sqrt{M_j}\), which diverges. It was not removed from the
physical operator; the norm of its nonscalar part relative to
\(\kappa_j\) is \(2/j^2\).

The angular condition of (35) holds for all sufficiently large \(j\),
because \[
 \theta_j^2L_j^2=\frac{4\pi^2}{10^8D_j^2}\longrightarrow0.
\] For \(j\ge10\) satisfying that condition, (36)-(37) apply to the
actual \(\psi_j\) and \(\chi_j\). They prove \[
 \rho_j\Longrightarrow\delta_3,\qquad
 \rho_j(|\lambda-3|>\delta)\le(80+8/\delta)^2j^{-4}.
\tag{40}
\] For a bounded continuous test function, continuity near 3 and the
tail bound in (40) prove weak convergence by splitting its integral into
\(|\lambda-3|\le\delta\) and its complement. The same estimate proves
tightness of these probabilities in the recorded variable
\(\lambda=\omega/\kappa_j\).

The physical probabilities have a different computed behavior. Equations
(17),(26),(39) give \[
 \operatorname{supp}\nu_j\subset
 [\kappa_j(3-2/j^2),\infty),\qquad
 \kappa_j\sim100\sqrt{24}\,j^5.
\tag{41}
\] For every fixed finite physical \(\Omega\), therefore, \[
 \boxed{\quad \nu_j([0,\Omega])=0\quad\hbox{for all sufficiently large }j.\quad}
\tag{42}
\] No subsequence of \(\nu_j\) is tight on \([0,\infty)\): every compact
subset is bounded and eventually has zero mass. The measures converge
vaguely to zero against continuous compactly supported functions, while
each retains total mass 1. This is spectral escape, not a zero-energy
limit.

The original state amplitude has a simultaneous limit as well. Equations
(31)-(32), with \(L_j=j^2\), \(D_j=j^4+O(j^2)\), and (12), give \[
 j^{10}\mathcal D_{L_j}(\theta_j)
 \longrightarrow\frac{16}{5}\left(\frac{2\pi}{10^4}\right)^4.
\] Indeed the degree-seven term of \(\mathcal A_{j^2}\) times
\(\theta_j^4\) has this limit, and the error in (32), after
multiplication by \(j^{10}\), is bounded by
\(16j^{28}\theta_j^6=O(j^{-8})\). Since \(j^8\xi_j\to1/24\), the
relative bounds (36) prove the genuine joint limit \[
 \boxed{\quad
 j^{26}d_j\longrightarrow
 \frac1{1620}\left(\frac{2\pi}{10^4}\right)^4.
 \quad}
\tag{43}
\] The raw measure remains \(d_j\nu_j\). Equations (41)-(43) display
both its vanishing amplitude and its energy escape without exchanging
the two limits.

The physical time scale in (18) is also explicit. For every fixed
\(t>0\), \[
 0\le d_j^{-1}\langle\chi_j,e^{-t(H_j-E_j)}\chi_j\rangle
 \le e^{-t\kappa_j(3-2/j^2)}\longrightarrow0.
\tag{44}
\] At the varying physical time \(t_j=s/\kappa_j\), with \(s\ge0\)
fixed, (18),(40) instead give \[
 d_j^{-1}\langle\chi_j,e^{-(s/\kappa_j)(H_j-E_j)}\chi_j\rangle
 \longrightarrow e^{-3s}.
\tag{45}
\] The test function \(e^{-s\lambda}\) is bounded and continuous, so
(45) follows directly from the full spectral convergence. Time \(s\) is
the explicitly dilated time; \(t_j\) tends to zero in the physical
coordinate (8). Holding \(\kappa_j\) at a fixed positive \(\kappa_*\)
while keeping the coupling (38) would require \[
 a_j=\frac{2g_j^2}{\kappa_*}=\frac{j\sqrt{M_j}}{\kappa_*},
\] which diverges. Conversely, holding \(\kappa_*\) with
\(a_j\downarrow0\) forces \(g_j^2=\kappa_*a_j/2\) and \[
 \xi_jM_j=\frac{M_j}{\kappa_*^2a_j^2}\longrightarrow\infty.
\tag{46}
\] Thus the controlled sequence (38) cannot be corrected into a fixed
electric energy scale by retaining its coupling and also refining the
mesh.

A broader precise obstruction follows directly from (26). On any
sequence with \(a_j\downarrow0\) and \(0<\xi_jM_j\le\eta<3/2\), the
original dictionary gives \[
 \kappa_j=\frac1{a_j\sqrt{\xi_j}}
 \ge\frac{\sqrt{M_j}}{a_j\sqrt\eta},
\quad \operatorname{gap}(H_j)\ge
 \frac{(3-2\eta)\sqrt{M_j}}{a_j\sqrt\eta}\longrightarrow\infty.
\tag{47}
\] Every vacuum-orthogonal invariant spectral probability then escapes
bounded physical energies, not only the magnetic one. This statement has
the displayed coupling-volume domain. A fixed positive coupling on the
growing box has \(\xi_jM_j\to\infty\) and is outside that domain.
Equation (47) makes no assertion of spectral escape there. The
independent volume-uniform proof in Section 9 establishes escape on the
fixed interval \(g^4\ge12288\), without this total-volume restriction.

\subsection{6. Exact refinement maps, measures, and their
limits}\label{exact-refinement-maps-measures-and-their-limits}

\subsubsection{6.1 Configuration, gauge, field and electric-operator
maps}\label{configuration-gauge-field-and-electric-operator-maps}

Because \(j\) is dyadic, every edge of the graph in (12) at \(j\) is the
concatenation of two fine edges at \(2j\); the old physical box is
contained in the new one. Define \[
 p_{j,2j}:Q_{L_{2j}}\longrightarrow Q_{L_j},\qquad
 (p_{j,2j}U)_{(n,i)}=U_i(2n)U_i(2n+e_i).
\tag{48}
\] It ignores all other fine links. Concatenations for distinct coarse
edges use disjoint fine edges. The map is continuous and surjective:
assign the first link in each pair the desired coarse value and the
second identity, then assign unused links identity. This gives a
continuous section in the fixed frames. Fine vertex gauge
transformations cancel at the intermediate vertex, leaving exactly the
coarse endpoint action. Thus (48) induces a map of gauge quotients as
well. Associativity proves \(p_{j,4j}=p_{j,2j}p_{2j,4j}\).

For each pair, the variable change \[
 (U^{(1)},U^{(2)})\longleftrightarrow
 (W=U^{(1)}U^{(2)},Z=U^{(1)}),\qquad
 U^{(2)}=Z^{-1}W
\tag{49}
\] is a diffeomorphism preserving product Haar measure. To prove the
measure assertion, integrate first over \(U^{(2)}\) at fixed \(Z\) and
use left invariance. Applying (49) to each disjoint pair proves \[
 (p_{j,2j})_*\lambda_{2j}=\lambda_j.
\tag{50}
\] Thus pullback is an exact Haar-Hilbert-space isometry and an algebra
homomorphism, preserving complex conjugation and the sup norm. Inverse
parallel transport along two consecutive segments multiplies in the
order in (48), so every smooth connection's link assignment obeys this
same map, with its full non-Abelian path ordering.

For smooth \(f\) on the coarse configuration, each of the two fine edge
Casimirs differentiating \(f(U^{(1)}U^{(2)})\) gives the coarse Casimir.
For the first edge this is direct left differentiation. For the second,
differentiation is conjugated by \(U^{(1)}\); the orthogonality of its
adjoint action preserves the sum of three squares. All unused-edge
derivatives annihilate the pullback. Consequently \[
 H_{0,2j}p_{j,2j}^*f=2p_{j,2j}^*H_{0,j}f.
\tag{51}
\] For the full Hamiltonians the exact defect, with
\(V_j=b_j\sum_p(2-W_p)\), is \[
 H_{2j}p^*f-p^*H_jf
 =(2\kappa_{2j}-\kappa_j)p^*H_{0,j}f
 +(V_{2j}-p^*V_j)p^*f.
\tag{52}
\] The fine potential and the coarse face word in the second term are
both retained; (52) is not set to zero. Smooth cylinders are within both
operator domains, so it is an equality of Hilbert-space vectors on the
displayed domain.

The changing cusp determinant also appears in the background map. On a
coarse direction-1 edge, the product of the two fine background links
equals \[
 \exp\left(-\frac{2\pi a_j^2}{D_{2j}}n_2H_c\right).
\] The original coarse background instead has \(D_j\) in the
denominator. Their exact relative factor is \[
 \exp\left[-2\pi a_j^2n_2
 \left(\frac1{D_{2j}}-\frac1{D_j}\right)H_c\right].
\tag{53}
\] Even for a fixed determinant, arbitrary quantum links cannot be
translated and then blocked by merely multiplying the background
translations. For two fine links the actual translated product is \[
 (h_1^{-1}U_1)(h_2^{-1}U_2)
 =h_1^{-1}(U_1h_2^{-1}U_1^{-1})(U_1U_2).
\tag{54}
\] The effective left insertion is therefore
\(h_1^{-1}\operatorname{Ad}_{U_1}(h_2^{-1})\), with its exact dependence
on the integrated link. Formula (54) is the morphism between the two
presentations; no independence of that conjugation is assumed.

\subsubsection{6.2 The exact vacuum and translated-state measure
maps}\label{the-exact-vacuum-and-translated-state-measure-maps}

Let \(d\mu_{2j}=\psi_{2j}^2d\lambda_{2j}\). In the coordinates (49),
collect all split-link \(Z\)'s and all unused links into
\(Z_{\mathrm{all}}\). Write \(\Phi(W,Z_{\mathrm{all}})\) for the inverse
change of variables. The exact marginal density is \[
 r_j^{(2j)}(W)=\int
 \psi_{2j}(\Phi(W,Z_{\mathrm{all}}))^2dZ_{\mathrm{all}},
 \qquad (p_{j,2j})_*\mu_{2j}=r_j^{(2j)}d\lambda_j.
\tag{55}
\] This is positive and smooth on the compact coarse configuration, by
differentiation under the finite compact integral. It is gauge invariant
by (48) and invariance of \(\mu_{2j}\). Equation (55) does not identify
\(r_j^{(2j)}\) with \(\psi_j^2\).

For the exact centered state write \(f_{2j}=\chi_{2j}/\psi_{2j}\),
retaining its amplitude. Its pushed-forward finite measure is \[
 (p_{j,2j})_*(|\chi_{2j}|^2d\lambda_{2j})
 =r_j^{(2j)}(W)
 \left[\frac{\int |f_{2j}(\Phi)|^2\psi_{2j}(\Phi)^2dZ_{\mathrm{all}}}
 {r_j^{(2j)}(W)}\right]d\lambda_j(W).
\tag{56}
\] The conditional mean \[
 \bar f_j(W)=\frac{\int f_{2j}(\Phi)\psi_{2j}(\Phi)^2dZ_{\mathrm{all}}}
 {r_j^{(2j)}(W)}
\] defines the orthogonal projection of \(f_{2j}\) onto coarse
observables in \(L^2(\mu_{2j})\). Indeed subtracting its pullback has
zero inner product with every coarse function, directly by Fubini. Its
exact norm decomposition is \[
 \|\chi_{2j}\|^2
 =\int |\bar f_j|^2r_j^{(2j)}d\lambda_j
 +\int|f_{2j}-p^*\bar f_j|^2d\mu_{2j}.
\tag{57}
\] The second term is retained. Equations (54)-(57) give explicit maps
of quantum link variables, measures and gauge-invariant states through
refinement, including the interaction-dependent conditional law. They do
not require a proposed continuum vacuum.

\subsubsection{6.3 Actual configuration compactness at fixed
coupling}\label{actual-configuration-compactness-at-fixed-coupling}

Let \(\mathcal Q_\infty\) be the inverse limit of the countable compact
spaces \(Q_{L_j}\) under (48). It is a closed subset of their compact
metrizable product, since each equality \(p_{j,2j}U_{2j}=U_j\) is
closed; hence it is compact and metrizable. Surjectivity of the
coordinate maps follows by successively using the section after (48).
Cylindrical continuous functions form a unital self-adjoint algebra
separating its points; the Stone--Weierstrass theorem makes them
uniformly dense in \(C(\mathcal Q_\infty)\).

For either actual sequence of vacua, and any fixed coarse index \(k\),
the pushforwards \((p_{k,j})_*\mu_j\), \(j\ge k\), are probability
measures on compact \(Q_{L_k}\). Their weak sequential compactness
follows by taking a countable uniformly dense set of continuous
functions, choosing a diagonal subsequence of their bounded integrals,
and extending the limiting positive functional by uniform continuity;
the Riesz representation theorem supplies the limiting probability.
Diagonalize once more over the countably many \(k\). The resulting
probabilities \(\mu_k^\infty\) are consistent, since \[
 (p_{k,l})_* (p_{l,j})_*\mu_j=(p_{k,j})_*\mu_j
\] at every finite regulator, and continuous pullback permits passage to
the weak limit. They define a positive functional on the cylindrical
algebra; the same density and Riesz argument give a unique probability
\(\mu^\infty\) on \(\mathcal Q_\infty\) with these marginals. Gauge
invariance passes to this limit. This proves actual subsequential
configuration compactness for the fixed-coupling vacua. It is a compact
topology of holonomies on the displayed dyadic path family, not a
claimed topology of distributional curvature or convergence of
Hamiltonians.

Equation (23) shows that the uncentered projected magnetic finite
measures have exactly the same subsequential limits as their
corresponding vacua on these cylinders. Their centered finite measures
have total mass tending to zero by (22). These are proved convergence
statements at fixed positive coupling, while these configuration
estimates alone do not control the probability obtained from the
centered state after division by \(d_j\). Section 9 proves its escape on
a fixed strong-coupling interval.

\subsubsection{6.4 The explicit configuration limit of
(38)}\label{the-explicit-configuration-limit-of-38}

For the controlled sequence, (27) gives \[
 \|\psi_j-1\|^2=2(1-s_j)\le2\|\zeta_j\|^2,
 \quad\|\psi_j-1\|\le\frac{\sqrt2}{2}\xi_j\sqrt{M_j}.
\] For the total variation convention
\(\|\mu-\lambda\|_{\mathrm{var}}=\sup_{\|f\|_\infty\le1}|\mu(f)-\lambda(f)|\),
Cauchy--Schwarz therefore proves \[
 \|\mu_j-\lambda_j\|_{\mathrm{var}}
 \le\int|\psi_j^2-1|d\lambda_j
 \le2\|\psi_j-1\|
 \le\sqrt2\xi_j\sqrt{M_j}=\frac{\sqrt2}{j^2\sqrt{M_j}}.
\tag{58}
\] Pushforward cannot increase this norm, because composition preserves
the bound on a test function. Equations (50),(58) give, for every fixed
\(k\), \[
 \|(p_{k,j})_*\mu_j-\lambda_k\|_{\mathrm{var}}
 \le\frac{\sqrt2}{j^2\sqrt{M_j}}\longrightarrow0.
\tag{59}
\] Thus the whole sequence, not only a subsequence, converges on all
cylinders to the projective Haar probability. Existence and uniqueness
follow from the construction in the preceding subsection and the exact
consistency (50). This is the restriction of the projective Haar
generalized-connection construction to the present countable cubical
path family. The measure is a computed limit of the actual interacting
vacua in (38); no equality between a finite \(\psi_j\) and the constant
vector is asserted.

\subsection{7. The computed temporal obstruction and an exact change of
scales}\label{the-computed-temporal-obstruction-and-an-exact-change-of-scales}

Take a fixed elementary square of one coarse graph, and let \(W\) be its
trace as a cylinder observable on every refinement. Haar integration on
that graph gives \(\lambda(W)=0\) and \(\lambda(W^2)=1\), by the
four-edge argument in (25). Equation (59) proves \[
 \mu_j(W)\longrightarrow0,\qquad
 v_j:=\|(W-\mu_j(W))\psi_j\|^2\longrightarrow1.
\] Let \[
 C_j(t)=\langle(W-\mu_j(W))\psi_j,
 e^{-t(H_j-E_j)}(W-\mu_j(W))\psi_j\rangle.
\] The vector is invariant and orthogonal to the actual vacuum. Applying
(26),(39) to it proves \[
 0\le C_j(t)\le v_j e^{-t\kappa_j(3-2/j^2)}.
\] Hence the pointwise limiting function satisfies \[
 C(0)=1,\qquad C(t)=0\quad(t>0).
\tag{60}
\] There is no Hilbert-space vector \(v\) of squared norm 1 and
nonnegative self-adjoint operator \(A\) whose strongly continuous
semigroup has \(C(t)=\langle v,e^{-tA}v\rangle\). In fact spectral
calculus and dominated convergence give
\(\|e^{-tA}v-v\|^2=\int|e^{-t\omega}-1|^2d\langle v,\mathbf1_{d\omega}(A)v\rangle\to0\),
forcing \(C(t)\to1\) as \(t\downarrow0\). This contradicts (60). Thus
the direct configuration limit in (59), with these fixed physical Wilson
observables and unscaled time, does not reconstruct such an evolution.

This is a specific dynamical obstruction in addition to spectral escape.
It does not forbid other coupling trajectories or differently
constructed observables. The exact comparison between them is (17),
(46), (52), and (54)-(57), rather than a declaration that different
presentations have no relation.

For an actual uniform spatial dilation \(x'=sx\) of the same finite
graph, retaining its link variables and \(g\), the pulled connection is
\(A'(x')=s^{-1}A(x'/s)\); its curvature is \(s^{-2}F(x'/s)\), including
the commutator, since both differentiated and bracket terms acquire
\(s^{-2}\). The dilated metric volume is \(s^4dx\), so the
four-dimensional action (11) is unchanged. The lattice dictionary gives
\[
 a'=sa,\quad\kappa'=\kappa/s,\quad b'=b/s,
 \quad H'=H/s,\quad E'=E/s,\quad
 \nu'(B)=\nu(sB).
\tag{61}
\] These are identities on the same finite configuration Hilbert space.
If the time coordinate is also dilated by \(s\), then
\(t'(H'-E')=t(H-E)\). Thus changes of coordinate scale have an exact
spectral dictionary; none of (40)-(45) is a spectral disproof obtained
by concealing such a dilation. In (12), the original physical metric and
\(\ell_*\) stay fixed throughout.

\subsection{8. Relation to temporal transfer and continuum
literature}\label{relation-to-temporal-transfer-and-continuum-literature}

The retained finite-spatial-regulator temporal Wilson coefficients are
\[
 b_t(\epsilon)=\frac{a}{2g^2\epsilon},\quad
 b_s(\epsilon)=\frac{\epsilon}{2g^2a},\quad
 Z(b_t)=\frac{I_1(2b_t)}{b_t}.
\] The exact temporal-link integral is \[
 \mathcal T_\epsilon^{\mathrm{raw}}
 =[e^{-2b_t}Z(b_t)]^{N_L}
 M_\epsilon K_\epsilon M_\epsilon\Pi_L,
 \quad M_\epsilon=e^{-\epsilon V/2},
\] where the single-edge multiplier of \(K_\epsilon\) is
\(I_{2q+1}(2b_t)/I_1(2b_t)\).

The companion temporal-transfer proof, retained in the sources
directory, proves its strong power limit to \(e^{-tH}\Pi_L\), with the
scalar energy shift \(N_L[2b_t-\log Z(b_t)]/\epsilon\) retained. Its
correlation-ratio identity applies directly to (18) and to \(C_j(t)\):
at each fixed \(j\), the raw scalar cancels exactly between numerator
and vacuum denominator. Thus all spectral statements here are statements
about the actual temporal-limit Hamiltonian of that full finite Wilson
system. No uniform-in-\(j\) error bound for replacing \(H_j\) by a
simultaneous finite-time-step transfer is used or asserted.

The literature comparison is tied to the proved maps above:

\begin{itemize}
\item
  K. Osterwalder and E. Seiler, \emph{Gauge field theories on a
  lattice}, Annals of Physics \textbf{110} (1978), 440-471,
  \href{https://doi.org/10.1016/0003-4916(78)90039-8}{doi:10.1016/0003-4916(78)90039-8},
  establish transfer positivity and strongly coupled infinite-volume
  lattice results. Their lattice transfer framework supports comparison
  with the finite temporal operator just displayed; it does not identify
  a single-time Wilson density with the Hamiltonian density \(\psi_j^2\)
  in (55).
\item
  A. Ashtekar and J. Lewandowski, \emph{Projective techniques and
  functional integration for gauge theories}, Journal of Mathematical
  Physics \textbf{36} (1995), 2170-2191,
  \href{https://doi.org/10.1063/1.531037}{doi:10.1063/1.531037},
  \href{https://arxiv.org/abs/gr-qc/9411046}{arXiv:gr-qc/9411046},
  Theorem 2 and Sections 2.3, 3.3, give the compact projective-measure
  framework. The precise hypotheses here are continuous surjective maps
  (48), their composition identity, and consistent limiting
  probabilities proved after (57). The finite vacuum measures are not
  presumed consistent. Our countable cubical construction is proved
  directly, and (59)-(60) calculate both its limiting measure and its
  temporal limitation. Their source TeX is retained under
  \texttt{sources/Ashtekar\_Lewandowski\_1994}.
\item
  S. Chatterjee, \emph{Yang-Mills for probabilists}, in
  \emph{Probability and Analysis in Interacting Physical Systems},
  Springer Proceedings in Mathematics and Statistics \textbf{283}
  (2019), 1-16,
  \href{https://doi.org/10.1007/978-3-030-15338-0_1}{doi:10.1007/978-3-030-15338-0\_1},
  \href{https://arxiv.org/abs/1803.01950}{arXiv:1803.01950}, Sections 3
  and 6-7, distinguishes finite Wilson measures and the scaling of loop
  observables needed for a continuum limit. His lattice action
  \(S_\Lambda=\sum_p\operatorname{Re}\operatorname{Tr}(I-U_p)\) is
  exactly \(\sum_p(2-W_p)\) for SU(2). Therefore his coefficient
  \(\beta\) equals \(b_t\) on our temporal faces and \(b_s\) on spatial
  faces of the anisotropic finite integral; it is not the Hamiltonian
  \(b=b_s/\epsilon\). This coefficient dictionary is exact and prevents
  substituting a fixed isotropic Wilson parameter for (38)-(39). Source
  TeX is retained under \texttt{sources/Chatterjee\_2018}.
\end{itemize}

The local primary corpus \texttt{Tong\_Gauge\_Theory\_20260904/gt.txt}
was also consulted for the lattice Hamiltonian and continuum discussion;
it supplies context rather than a vacuum estimate used in any proof
here. Every spectral estimate in Sections 3-7 is proved above with the
original finite operator and original coefficients.

A current primary construction can be compared by an exact density map.
S. Chatterjee, \emph{A scaling limit of SU(2) lattice Yang-Mills-Higgs
theory}, \href{https://arxiv.org/abs/2401.10507v2}{arXiv:2401.10507v2},
Theorem 3.2, treats the unitary-gauge stereographic field
\(\sqrt2\sigma_3(\tau(U_e))/g_H\). With \(\alpha\to\infty\),
\(g_H\to0\), \(\alpha g_H=c\epsilon\), and \(g_H=O(\epsilon^{50d})\),
its specified spatially dilated distributional fields converge to three
independent Proca fields of mass \(c/\sqrt2\). Its equation (5.2) gives
\(\|I-U\|_{\mathrm{HS}}^2=4-2\operatorname{Re}\operatorname{tr}U\).
Therefore, on any one fixed finite graph, its gauge-fixed density (3.3)
is related to the pure isotropic Wilson probability with coefficient
\(\beta=1/g_H^2\) by the exact Radon-Nikodym derivative \[
 \frac{d\mu_H}{d\mu_W}(U)
 =\frac{\exp[-(\alpha^2/2)\sum_e(2-\operatorname{tr}U_e)]}
 {\mu_W(\exp[-(\alpha^2/2)\sum_e(2-\operatorname{tr}U_e)])}.
\] This follows simply by dividing the two positive densities and
integrating to determine the denominator. The numerator is nonconstant
for \(\alpha>0\), since varying one link with all others fixed changes
its trace. It retains the added Higgs-induced interaction explicitly.
The correspondence does not remove that term or provide the missing
pure-vacuum estimate in (17). The retained source TeX records the
theorem's boundary conditions, observable map and full hypotheses.

\subsection{9. A volume-uniform gap and the fixed-coupling magnetic
sequence}\label{a-volume-uniform-gap-and-the-fixed-coupling-magnetic-sequence}

The comparison (47) required a small total interaction \(\xi M_L\). We
now prove a different estimate with no volume factor in its coupling
range: \[
 0<\xi\le\frac1{49152},\qquad
 (H-E)|_{\psi^\perp}\ge
 \frac{3\kappa}{4}\left(1-\frac{512\xi}{3}\right)I
 \ge\frac{287\kappa}{384}I .
 \tag{62}
\] The perpendicular complement here is in the full link Hilbert space;
hence the inequality also holds for the gauge-invariant centered
magnetic state. The proof below retains the infinite-dimensional link
spaces. It develops the creation-coordinate argument of D. A. Yarotsky,
\emph{Quasi-particles in weak perturbations of non-interacting quantum
lattice systems},
\href{https://arxiv.org/abs/math-ph/0411042v1}{arXiv:math-ph/0411042v1},
Section 2, with all numerical constants calculated for the present
four-link interactions. The complete source TeX is retained. The
companion volume-uniform calculation first supplied these explicit
constants; we give the proof here to make the continuum consequence
independently readable.

\subsubsection{9.1 Exact creation coordinates and the interacting
eigenvector}\label{exact-creation-coordinates-and-the-interacting-eigenvector}

In this section write \(\mathbb H_0=\kappa H_0\), where \(H_0\) remains
the operator in (1). Put \[
 V_p=-bW_p,\quad J=2b,\quad s=4,\quad d_{\rm inc}=4,\quad
 \gamma_*=\frac{3\kappa}{4}.
\] Thus \(H=\mathbb H_0+2bM_L+\sum_pV_p\), \(\|V_p\|\le J\), every
support has \(s\) distinct links and every link meets at most
\(d_{\rm inc}\) supports. None of these definitions changes a physical
coefficient.

Let \(\Omega=1\). For \(I\subset\mathsf E_L\), let \(\mathcal H'_I\) be
the tensor product of the mean-zero single-link spaces on \(I\), with
the empty product \(\mathbb C\). The orthogonal sector projection
\(\mathsf P_I\) applies the mean-zero projection on \(I\) and the
constant projection on \(I^c\). Then \[
 \mathcal H_L=\bigoplus_{I\subset\mathsf E_L}\mathcal H'_I,\qquad
 \mathbb H_I=\kappa\sum_{e\in I}E_e,\qquad
 \mathbb H_I\ge\gamma_*|I|,\qquad
 \|\mathbb H_I^{-1}\|\le(\gamma_*|I|)^{-1}\quad(I\ne\varnothing).
 \tag{63}
\] The isomorphism inserts constants in \(I^c\), and its inverse is the
family of the displayed projections. Parseval proves both inverse
identities. The lower bound follows because each nonconstant factor has
a nonzero spin and single-link energy at least \(3\kappa/4\).

For \(z_I\in\mathcal H'_I\), define \[
 \widehat z_I=|z_I\rangle\langle1_I|\otimes I_{I^c}.
\] It is bounded of norm \(\|z_I\|\). Integrating a common link shows
that \[
 \widehat z_I\widehat y_K=0\quad(I\cap K\ne\varnothing),
 \qquad
 \widehat z_I\widehat y_K=\widehat{z_I\otimes y_K}
       \quad(I\cap K=\varnothing).
 \tag{64}
\] For entangled vectors this follows by approximating them with finite
sums of elementary tensors; the defining operators converge in norm. In
particular all creation operators commute, and any product of more than
\(N_L\) nonempty creation operators vanishes.

Use the finite graph whose vertices are physical links and whose
adjacency means sharing a face. For nonempty \(I\), let \(t(I)\) be the
minimum number of edges in a connected subgraph containing \(I\),
allowing additional vertices, and put \(w(I)=2^{t(I)+1}\). This graph is
connected. A face support forms a four-vertex clique and has
\(w(p)=16\). Inclusion is monotone. Joining minimizing connected
subgraphs to a tree on a face proves \[
 w\left(p\cup\bigcup_{j=1}^n I_j\right)
 \le16\prod_{j=1}^n w(I_j)\quad\hbox{when every }I_j\cap p\ne\varnothing.
 \tag{65}
\] Indeed their connected union has at most \(3+\sum_jt(I_j)\) edges;
the extra factors \(2^n\) on the right only increase the upper bound.

Consider the complete space of collections \(z_I\in D(\mathbb H_I)\),
\(I\ne\varnothing\), with \[
 \|z\|_*=\max_x\sum_{I\ni x}
       w(I)\frac{\|\mathbb H_Iz_I\|}{\gamma_*}.
 \tag{66}
\] Completeness follows from (63) and the finite number of sectors. If
\(a_I=w(I)\|z_I\|\) and \(R=\|z\|_*\), then \[
 \sup_x\sum_{I\ni x}|I|a_I\le R,\qquad
 \sum_{I:I\cap p\ne\varnothing}a_I\le sR .
\] Set \(\mathscr C(z)=\sum_{I\ne\varnothing}\widehat z_I\). At each
finite volume it is bounded and nilpotent, so \(e^{\mathscr C}\) and
\(e^{-\mathscr C}\) are mutually inverse finite polynomials. On
\(D(\mathbb H_0)\), \[
 [\mathbb H_0,\widehat z_I]=\widehat{\mathbb H_Iz_I},
 \qquad
 e^{-\mathscr C}\mathbb H_0e^{\mathscr C}
       =\mathbb H_0+\sum_I\widehat{\mathbb H_Iz_I}.
 \tag{67}
\] The first identity holds on finite spectral tensors: exterior
energies commute and the local energy annihilates the constant bra. Its
bounded right side proves extension to the graph domain, and
preservation of that domain by both exponentials. The second commutator
vanishes by (64), proving the second identity exactly.

The nonempty components of the eigen-equation are the fixed-point
equation \[
 \mathcal F(z)_I=-\mathbb H_I^{-1}\mathsf P_I
       e^{-\mathscr C(z)}\left(\sum_pV_p\right)e^{\mathscr C(z)}\Omega .
 \tag{68}
\] The inverse in (63) places every output in its stated domain. We
calculate the contraction instead of postulating a convergent expansion.

For one \(p\), expand its conjugation in nested commutators with
\(\mathscr C\). A creation index disjoint from \(p\) contributes zero
since it commutes with both \(V_p\) and all other creation operators. A
commutator of degree \(n\) has \(2^n\) ordered words. On either side of
\(V_p\), supports in a nonzero word are disjoint by (64); at most four
of them can meet \(p\). Thus the exact expansion ends at degree eight.

Outside \(p\), the potential acts as identity. If two creation supports
share an outside link, the constant bra and mean-zero projection
annihilate that word. Otherwise its excited outside set is fixed to be
\((\bigcup_jI_j)\setminus p\). Only the subsets of the four inside links
vary. There are at most \(16\) orthogonal output sectors. Consequently
Cauchy--Schwarz gives, for each word, \[
 \sum_K\|\mathsf P_K(\text{word})\Omega\|
       \le4J\prod_{j=1}^n\|z_{I_j}\| .
 \tag{69}
\] All arguments use projections, so remain valid for entangled vectors
by the norm approximation establishing (64).

For an anchor \(x\) in a nonempty output \(K\), either \(x\in p\) or
\(x\in I_j\) for some \(j\). In the first case there are at most
\(d_{\rm inc}\) faces and the index sums contribute \((sR)^n\). In the
second choose \(j\) and \(I_j\ni x\); at most \(d_{\rm inc}|I_j|\) faces
meet it, and the factor \(|I_j|\) is absorbed by (66). The two upper
bounds are \(d_{\rm inc}(sR)^n\) and \(nd_{\rm inc}R(sR)^{n-1}\), with
the second zero at \(n=0\). Combining (65), (68) and (69), cancelling
\(\mathbb H_I\) against its inverse, proves \[
 \|\mathcal F(z)\|_*
 \le\frac{2048}{3}\xi(1+2R)e^{8R}.
\] Replacing one input at a time in each multilinear word gives the
Lipschitz bound on the radius-\(R\) ball: \[
 \|\mathcal F(z)-\mathcal F(y)\|_*
 \le\frac{2048}{3}\xi(10+16R)e^{8R}\|z-y\|_* .
 \tag{70}
\] For the second inequality, the degree-\(n\) positive majorant
acquires the factor \(nR^{n-1}\) in place of \(R^n\). Differentiation of
the preceding positive majorant gives precisely \(10+16R\). This counts
anchors in the differing position as well as the other positions.

Take \(R_*=1/16\). The geometric-series comparison \(e^{1/2}<2\) proves,
throughout the range of (62), \[
 \|\mathcal F(z)\|_*\le1536\xi\le1/32,\qquad
 \operatorname{Lip}(\mathcal F)\le(45056/3)\xi\le11/36.
 \tag{71}
\] Iteration from zero therefore converges: successive differences are
bounded by a geometric series, and completeness and (70) give the
fixed-point equation and uniqueness in the ball. Conjugation commutes
with every factor of (68), so the solution and its eigenvalue are real.
Equations (67)--(68) yield \[
 \phi=e^{\mathscr C(z)}\Omega,\qquad H\phi=\varepsilon\phi,\qquad
 \varepsilon=2bM_L+
 \mathsf P_\varnothing e^{-\mathscr C(z)}
       \left(\sum_pV_p\right)e^{\mathscr C(z)}\Omega .
 \tag{72}
\] Every nonempty creation product has zero Haar mean, so
\(\langle\Omega,\phi\rangle=1\). The scalar \(2bM_L\) has remained in
the eigenvalue.

\subsubsection{9.2 Resolvent exclusion in the full
domain}\label{resolvent-exclusion-in-the-full-domain}

Fix the constructed \(z\). Define \[
 \widetilde H=e^{-\mathscr C}(H-\varepsilon)e^{\mathscr C}
       =\mathbb H_0+\mathscr K,\qquad \widetilde H\Omega=0 .
\] This is a bounded similarity preserving \(D(\mathbb H_0)\); it is not
asserted unitary. Equip the same finite-volume space with the equivalent
sector norm \[
 \|f\|_{\oplus,1}=\sum_I\|\mathsf P_If\|,\qquad
 \|f\|\le\|f\|_{\oplus,1}\le2^{N_L/2}\|f\| .
\] For \(u_I\in D(\mathbb H_I)\), commute the similarity with
\(\widehat u_I\), use \(\widetilde H\Omega=0\), and subtract the
\(\mathbb H_Iu_I\) term to get \[
 \mathscr K\widehat u_I\Omega
 =e^{-\mathscr C}\left[\sum_pV_p,\widehat u_I\right]
       e^{\mathscr C}\Omega .
\] Only the at most \(d_{\rm inc}|I|\) faces meeting \(I\) contribute.
In the conjugation expansion the same outside-face sector argument
proving (69) applies, now including the additional creation support
\(I\). The commutator adds two words. Since
\(\sum_{K:K\cap p\ne\varnothing}\|z_K\|\le sR_*\), it follows that \[
 \|\mathscr K\widehat u_I\Omega\|_{\oplus,1}
 \le8Jd_{\rm inc}|I|e^{8R_*}\|u_I\| .
\] Sum over \(I\), use (63), and use \(e^{1/2}<2\). This proves \[
 \|\mathscr Kf\|_{\oplus,1}
 \le c_\xi\|\mathbb H_0f\|_{\oplus,1},\qquad
 c_\xi=\frac{512\xi}{3}\le\frac1{288},
 \quad f\in D(\mathbb H_0).
 \tag{73}
\] The empty component gives zero. The estimates first hold on finite
spectral tensors and extend by graph closure.

For real \(t<0\), (73) gives
\(\|\mathscr K(\mathbb H_0-t)^{-1}\|_{\oplus,1}\le c_\xi<1\). For
\(0<t<(1-c_\xi)\gamma_*\), it gives \[
 \|\mathscr K(\mathbb H_0-t)^{-1}\|_{\oplus,1}
 \le\frac{c_\xi\gamma_*}{\gamma_*-t}<1 .
 \tag{74}
\] Indeed \(\lambda/(\lambda-t)\) on \(\lambda\ge\gamma_*\) has this
supremum; the empty sector is annihilated by \(\mathscr K\). The Neumann
inverse followed by \((\mathbb H_0-t)^{-1}\) is a bounded inverse of
\(\widetilde H-t\), with image in its graph domain. Norm equivalence
makes it a bounded Hilbert-space inverse at each finite volume. No
uniform estimate on the equivalence constant is used in excluding the
real intervals. The bounded similarity transfers these intervals to the
resolvent of \(H-\varepsilon\). Since \(H\) is self-adjoint and
\(\varepsilon\) is a real eigenvalue, the exclusion of all \(t<0\)
proves \(\varepsilon=E\). Ground-state simplicity from Section 1
identifies the unchanged physical vacuum: \[
 \psi=\langle\Omega,\psi\rangle e^{\mathscr C}\Omega,\qquad
 \langle\Omega,\psi\rangle>0,\qquad
 \langle\Omega,\psi\rangle^2\|e^{\mathscr C}\Omega\|^2=1 .
 \tag{75}
\] The positive real interval excluded in (74), and the spectral
theorem, now prove (62).

\subsubsection{9.3 Consequence for the original fixed-coupling cusp
state}\label{consequence-for-the-original-fixed-coupling-cusp-state}

The range in original coefficients is \(g^4\ge12288\). Retain exactly
the geometry (12), \(D_j\), the state \(\chi_j\), and a fixed \(g=g_0\)
in that range. For all sufficiently large \(j\), its angle is nonzero
and not a multiple of \(2\pi\), so Section 3 proves \(d_j>0\). Equation
(62) gives \[
 \operatorname{supp}\nu_j\subset[\Delta_j,\infty),\qquad
 \Delta_j=
 \left(150g_0^2-\frac{6400}{g_0^2}\right)j
 \ge\frac{7175}{48}g_0^2j .
 \tag{76}
\] Thus \(\nu_j([0,\Omega])=0\) eventually for every fixed finite
physical \(\Omega\), and
\(\langle\chi_j,e^{-t(H_j-E_j)}\chi_j\rangle/d_j
 \le e^{-t\Delta_j}\) for all \(t\ge0\). This is a conclusion about the
actual centered magnetic spectral probability, despite the vanishing
\(d_j\). It extends beyond (47), since now \(\xi_0M_j\) tends to
infinity.

More generally every \(a_j\to0\) trajectory that stays in
\(g_j^4\ge12288\) has \(\Delta_j\ge287g_j^2/(192a_j)\to\infty\). All
vacuum-orthogonal vectors have the corresponding spectral exclusion.
Finite positive energy filters, refinements, and choices of centering
cannot create spectral weight below this support for those same physical
operators. This result concerns the displayed strong-coupling interval.
Section 10 gives a different all-positive-coupling conclusion for
specified loop states.

\subsection{10. Fixed physical loop states at every positive
coupling}\label{fixed-physical-loop-states-at-every-positive-coupling}

\subsubsection{10.1 A local density bound from the nonlinear vacuum
equation}\label{a-local-density-bound-from-the-nonlinear-vacuum-equation}

Write \(u=\log\psi\) and \(\rho=\psi^2\) for the same vacuum. In the
link metric with orthonormal \(T_a\), \(SU(2)\) is the round
three-sphere of radius 2: in quaternion coordinates \(T_a\) has
Euclidean length \(1/2\), and left multiplication preserves the
quaternion metric. Its Ricci tensor is one half its metric and its
diameter is \(2\pi\). These factors preserve \(\Delta^c=4\Delta^T\).

Let \(\Delta=\sum_e\Delta_e=-H_0\) and \(V=b\sum_p(2-W_p)\). Dividing
the full equation by its positive vacuum gives \[
 \Delta u+|\nabla u|^2=(V-E)/\kappa .
\] For \(w_e=|\nabla_eu|^2\), differentiation on the product manifold
gives the partial Bochner identity \[
 (\Delta+2\nabla u\cdot\nabla)w_e
 =2\sum_f|\nabla_f\nabla_eu|_{\rm HS}^2+w_e
   +\frac2\kappa\langle\nabla_eu,\nabla_eV\rangle .
 \tag{77}
\] For completeness, differentiate the three squared gradient components
in orthonormal geodesic frames. The product Laplacian produces the first
Hessian sum and \(2\langle\nabla_eu,\nabla_e\Delta u\rangle\). Commuting
the two factor-\(e\) derivatives adds
\(2\operatorname{Ric}_e(\nabla_eu,\nabla_eu)=w_e\); mixed factors have
zero curvature. Substituting the differentiated vacuum equation produces
a term \(-4\nabla^2u(\nabla u,\nabla_eu)\), which cancels exactly with
the drift of \(w_e\). This proves (77), retaining every mixed Hessian.

If \(r_e\le4\) is the actual number of faces incident to \(e\), then
\(|\nabla_eV|\le br_e\). To verify its unit constant, write a face word
as \(AUB\) or \(AU^{-1}B\). Bi-invariance and inversion reduce its trace
derivative to that of \(2u_0\) on the radius-2 sphere, whose squared
gradient is \(1-u_0^2\le1\). At a global maximum of \(w_e\) on the
compact full product, its gradient is zero and its Laplacian
nonpositive. Equation (77) gives \(0\ge w_e-2r_e\xi\sqrt{w_e}\), and
hence \[
 \|\nabla_e\log\psi\|_\infty\le2r_e\xi,\qquad
 \|\nabla_e\log\rho\|_\infty\le4r_e\xi .
 \tag{78}
\] This is valid at all \(g>0\) and every volume, with no expansion in
\(\xi\).

For an exterior configuration \(y\), let \[
 p_e(U\mid y)=\frac{\rho(U,y)}{\int\rho(U',y)\,dU'},\qquad
 \alpha(g)=e^{-32\pi\xi}=e^{-8\pi/g^4}.
\] The geodesic diameter and (78) bound the oscillation of \(\log p_e\)
by \(8\pi r_e\xi\). Since its integral is 1, its minimum and maximum
straddle 1; the ratio estimate therefore proves \[
 e^{-8\pi r_e\xi}\le p_e(U\mid y)
       \le e^{8\pi r_e\xi},\qquad
 \alpha(g)\le p_e(U\mid y)\le\alpha(g)^{-1}.
 \tag{79}
\] These are the actual conditional densities; correlations among other
links are unchanged.

\subsubsection{10.2 Raw spectral mass of a simple
loop}\label{raw-spectral-mass-of-a-simple-loop}

Let \(C\) be a closed lattice path with \(\ell\) distinct physical
edges, each traversed once, and let \(W_C=\operatorname{tr}U_C\) with
its full ordered product and inverse orientations. Define \[
 m_C=\mu(W_C),\quad
 v_C=\psi(W_C-m_C),\quad d_C=\|v_C\|^2,\quad
 \eta_C(B)=\langle v_C,\mathbf1_B(H-E)v_C\rangle,\quad
 \nu_C=\eta_C/d_C .
\] The loop is gauge invariant. On any one of its edges, at fixed
exterior links, the change of variable from that link to \(U_C\)
preserves Haar, including for inverse traversal. Haar on \(SU(2)\) has
\(\int\operatorname{tr}U=0\) and \(\int(\operatorname{tr}U)^2=1\). One
can calculate these directly from \(U=u_0I+i\mathbf u\cdot\sigma\):
symmetry gives \(\int u_0=0\), \(\int u_0^2=1/4\). Consequently, for
every real \(c\), \[
 \int(W_C-c)^2p_e\,dU_e\ge\alpha(1+c^2).
\] The conditional variance is at least \(\alpha\). The identity of
total variance, obtained by subtracting and adding the conditional mean
and integrating the zero cross term, now gives \(d_C\ge\alpha\). Also
\(d_C\le4\) because \(|W_C|\le2\). In the same conditional integral,
\(\int(1-W_C^2/4)\,dU_e=3/4\). Thus the ground-state form (2) and the
link trace derivative calculation give exactly \[
 d_C\in[\alpha,4],\qquad
 A_{1,C}:=\int\omega\,d\eta_C(\omega)
 =\kappa\ell\left(1-\frac{\mu(W_C^2)}4\right)
 \ge\frac{3\alpha}{4}\kappa\ell .
 \tag{80}
\] There are no cross-link derivatives in this form, and each of the
\(\ell\) distinct links contributes the same squared gradient
\(1-W_C^2/4\).

Multiplication by \(\psi\) is the exact unitary from \(L^2(\mu)\) to
\(L^2(\lambda_L)\). Under it, \[
 \mathscr L=\psi^{-1}(H-E)\psi
       =-\kappa(\Delta+2\nabla u\cdot\nabla).
\] The operator acts on the same \(H^2\) domain, since the finite-volume
positive smooth \(\psi\) and its inverse preserve Sobolev spaces. Since
\(-\Delta W_C=(3\ell/4)W_C\), (78) implies the pointwise bound \[
 |\mathscr L W_C|\le\kappa\ell B(g),\qquad
 B(g)=\frac32+16\xi=\frac32+\frac4{g^4}.
\] Every reverse traversal gives the same Casimir by inversion
invariance. The derivative contribution is bounded by
\(2\sum_{e\in C}(2r_e\xi)|\nabla_eW_C|\le16\xi\ell\). The full second
spectral moment is therefore bounded by \[
 A_{2,C}:=\int\omega^2\,d\eta_C(\omega)
       =\|\mathscr L W_C\|_{L^2(\mu)}^2
       \le B(g)^2\kappa^2\ell^2 .
 \tag{81}
\]

Put \[
 R_C=\frac{3\alpha(g)}{32}\kappa\ell,\qquad
 c_{\rm raw}(g)=\frac{9\alpha(g)^2}{64B(g)^2},\qquad
 c_{\rm prob}(g)=\frac{9\alpha(g)^2}{256B(g)^2}.
\] Then the full spectral measure, with no discarded channels, satisfies
\[
 \eta_C([R_C,\infty))\ge c_{\rm raw}(g),\qquad
 \nu_C([R_C,\infty))\ge c_{\rm prob}(g)>0 .
 \tag{82}
\] Indeed split the first moment below and above \(R_C\).
Cauchy--Schwarz for the second piece gives \[
 A_{1,C}\le R_Cd_C+
    \sqrt{A_{2,C}\eta_C([R_C,\infty))}.
\] Equations (80)--(81) and \(d_C\le4\) show
\(A_{1,C}-R_Cd_C\ge3\alpha\kappa\ell/8\). Squaring proves the raw bound
in (82), and division by \(d_C\le4\) proves its probability version.
Keeping the raw norm in this calculation avoids an unnecessary third
power of the conditional-density constant.

\subsubsection{10.3 Exact spatial refinement and failure of energy
tightness}\label{exact-spatial-refinement-and-failure-of-energy-tightness}

Fix a square of one coarse graph \(Q_{L_k}\), of physical side
\(s=a_k>0\). For each subsequent dyadic \(j\), the map \(p_{k,j}\) in
Section 6 sends its four coarse edges to chains of \(m_j=a_k/a_j=j/k\)
fine edges. Thus \(W_j=p_{k,j}^*W\) is the same physical Wilson
observable and \[
 \ell_j=4m_j=\frac{4s}{a_j},\qquad
 R_j=\frac{3\alpha(g_j)g_j^2s}{4a_j^2}.
 \tag{83}
\] This follows from the exact ordered edge products; no conditional
vacuum is substituted. Its state is precisely
\(v_j=\psi_j(W_j-\mu_j(W_j))\), and its raw measure is that in (82).

At fixed \(g_0>0\), \(R_j\) tends to infinity and (82) retains a
positive raw mass \(c_{\rm raw}(g_0)\) and probability mass
\(c_{\rm prob}(g_0)\) above \(R_j\). Hence neither family of these
measures is tight in physical energy. This proves loss of a positive
spectral fraction; it does not prove that the entire loop probability
escapes, or that its low-energy part is zero outside the range of
Section 9.

The same statement holds uniformly on any subsequence with
\(g_j\ge g_{\min}>0\). In fact \(\alpha(g_j)\ge\alpha(g_{\min})\),
\(B(g_j)\le B(g_{\min})\), and \(g_j^2\ge g_{\min}^2\). Equations
(82)--(83) then give uniform positive mass beyond a threshold diverging
as \(a_j^{-2}\). It follows that tightness for this family of unchanged
fixed physical Wilson states along a full \(a_j\to0\) sequence requires
\[
 g_j\longrightarrow0.
 \tag{84}
\] This is a proved necessary coupling consequence for the displayed
observables. It is not a sufficiency claim or a theorem about every
possible renormalized observable.

There is also a direct dynamical consequence for the configuration
limits already constructed. Along a subsequence in Section 6.3,
continuity of \(W\) and \(W^2\) on \(Q_{L_k}\) gives
\(d_j\to d_\infty\ge\alpha(g_0)\). For every fixed physical \(t>0\),
spectral calculus and (82) give \[
 \langle v_j,e^{-t(H_j-E_j)}v_j\rangle
 \le d_j-c_{\rm raw}(g_0)(1-e^{-tR_j}).
 \tag{85}
\] There is an actual subsequence on which the correlations converge for
every \(t>0\), and its lost spectral mass has an exact representation.
Use the unchanged reference length \(\ell_*=1\) from Section 2.1 to
define the dimensionless variable \(x=\exp(-\ell_*\omega)\). Let
\(\zeta_j\) be the pushforward of the raw measure \(\eta_j\) onto
\((0,1]\), regarded as a finite measure on compact \([0,1]\). Its total
mass is \(d_j\le4\). The countable-test-function and Riesz argument used
in Section 6.3 gives a weakly convergent further subsequence
\(\zeta_j\to\zeta_\infty\), of total mass \(d_\infty\).

For each fixed \(0<\delta<1\), eventually \(e^{-\ell_*R_j}<\delta\), so
(82) gives \(\zeta_j([0,\delta])\ge c_{\rm raw}(g_0)\). The closed-set
weak-convergence inequality yields
\(\zeta_\infty([0,\delta])\ge c_{\rm raw}(g_0)\). One can derive that
inequality by decreasing continuous functions bounded by 1 to the
indicator of the closed interval and then using dominated convergence.
Letting \(\delta\downarrow0\) proves \[
 m_\infty:=\zeta_\infty(\{0\})\ge c_{\rm raw}(g_0)>0.
 \tag{85a}
\] The point \(x=0\) represents infinite physical energy. For every
\(t>0\), the function \(x^{t/\ell_*}\), defined to be zero at \(x=0\),
is continuous on the whole compact interval. The same subsequence
therefore has the exact correlation limit \[
 C(t)=\int_{[0,1]}x^{t/\ell_*}\,d\zeta_\infty(x)\quad(t>0),
 \qquad C(0)=d_\infty,\qquad
 C(0+)=d_\infty-m_\infty .
 \tag{85b}
\] The last identity follows by dominated convergence on \((0,1]\). Thus
its positive-time limit loses at least \(c_{\rm raw}(g_0)\) of the
original squared state norm at time zero.

The exact Hilbert-space map recording this loss is \[
 L^2(\zeta_\infty)
 =L^2(\zeta_\infty|_{(0,1]})\oplus L^2(m_\infty\delta_0),
 \qquad P_{\rm fin}f=\mathbf1_{(0,1]}f .
 \tag{85c}
\] On the first summand define \(A_{\rm fin}\) by multiplication by
\(-\ell_*^{-1}\log x\), with domain consisting exactly of the functions
for which that product is square integrable. Bounded-energy truncations
show the domain is dense. Multiplication by
\((-\ell_*^{-1}\log x-z)^{-1}\) is a bounded inverse for every nonreal
\(z\), with image in the stated domain; symmetry and these resolvents
prove self-adjointness. Its nonnegative semigroup is multiplication by
\(x^{t/\ell_*}\), and dominated convergence proves strong continuity.
Its constant vector has squared norm \(d_\infty-m_\infty\) and
reproduces \(C(t)\) for \(t>0\).

On the whole space (85c), the exact limiting operators are
\(T_t=e^{-tA_{\rm fin}}\oplus0\) for \(t>0\) and \(T_0=I\). They obey
the semigroup law, but \(T_t\) tends strongly to \(P_{\rm fin}\), not to
\(I\), as \(t\downarrow0\). The original constant vector has squared
norm \(d_\infty\), and its lost component has squared norm exactly
\(m_\infty\). Consequently no strongly continuous nonnegative
self-adjoint semigroup with that original squared norm can reproduce
(85b). This gives the precise projection and residual scalar spectral
dynamics associated with the obstruction. It does not construct the
interacting spatial observable algebra from one loop's scalar spectral
measure. The configuration limit has nowhere been identified with Haar.

\subsubsection{10.4 Shrinking loops and the retained field
map}\label{shrinking-loops-and-the-retained-field-map}

For any elementary face, integrating one of its links and using (79)
gives \[
 \mu(W_p\le0)\ge\frac{\alpha(g)}2,\qquad
 \mu\bigl(\|I-U_p\|_{\rm HS}^2\bigr)\ge4\alpha(g).
 \tag{86}
\] The Haar probability of \(\operatorname{tr}U\le0\) is \(1/2\), by the
measure-preserving map \(U\mapsto-U\). Also
\(\|I-U\|_{\rm HS}^2=4-2\operatorname{tr}U\), whose Haar mean is 4.
These prove both inequalities. The expressions are invariant under the
conjugation of the based holonomy, so the statements descend to the
original gauge quotient.

For comparison, the original cusp background on a magnetic elementary
face has \[
 \|I-U_p^{\rm bg}\|_{\rm HS}^2
       =4-4\cos\theta_j\le2\theta_j^2\longrightarrow0 .
 \tag{87}
\] Its value follows directly from the retained links \(h_i\) and is
unchanged by conjugation or the sign of the face orientation. This exact
comparison explains why the smooth background limit does not itself
determine the fluctuation law of the actual vacuum. At a fixed positive
coupling the elementary quantum holonomy fails to approach identity in
probability: on \(W_p\le0\) its squared distance is at least 4 and that
event has uniformly positive probability.

The continuum consequence has precisely this scope. A topology or
reconstruction that requires these unmodified shrinking holonomies to
converge to identity cannot realize that fixed-coupling sequence. A
distributional field reconstruction need not impose that pointwise
property. Equations (79), (83), (85) and (87) give the actual measure,
refinement, spectral and field maps; no general nonrelationship or
universal continuum impossibility is inferred.

\subsection{11. Explicit coupling continuations outside the proved
exclusions}\label{explicit-coupling-continuations-outside-the-proved-exclusions}

The deductions (76) and (84) give concrete information for the next
interacting calculation. They do not fix a continuum coupling trajectory
by an unproved renormalization rule. For example retain the entire
geometry (12), take \(j\ge2\) dyadic, and choose the strictly positive
coupling \[
 g_j^2=\frac1{\log j},\quad
 \xi_j=\frac{(\log j)^2}{4},\quad
 \kappa_j=\frac{200j}{\log j},\quad
 b_j=50j\log j,\quad
 2b_jM_j=100jM_j\log j .
 \tag{88}
\] All plaquettes and the unchanged physical reference length remain
present. The calculation (22), which is uniform in \(\xi\) before
specialization, now gives \[
 \|\chi_j\|\le
 \frac{1536\pi^2}{10^8}\frac{(\log j)^2}{j^2},\qquad
 d_j\le\left(\frac{1536\pi^2}{10^8}\right)^2
             \frac{(\log j)^4}{j^4}.
 \tag{89}
\] The nonzero-state proof still applies. Thus the exact
configuration-compactness and uncentered-state comparison argument in
Section 6.3 applies to this trajectory too: it needs compact
configurations, gauge invariance and the vanishing norm just proved, not
a constant coupling. The centered finite measures have zero total-mass
limit. Their spectral probabilities remain the exact pushforward (17)
and are not determined by (89).

For this trajectory the constants in the all-coupling loop result are
explicitly \[
 \alpha_j=e^{-8\pi(\log j)^2},\quad B_j=\frac32+4(\log j)^2,\quad
 R_j=\frac{7500s\,j^2}{\log j}e^{-8\pi(\log j)^2}\longrightarrow0 .
 \tag{90}
\] The lower masses in (82) tend to zero as well. Thus that proved
estimate does not exclude tightness on (88). It does not establish it
either. The cluster-gap range is eventually violated and \(\xi_jM_j\)
diverges, so neither earlier whole-spectrum exclusion applies.

A different exact coefficient choice keeps the original electric
coefficient equal to a fixed physical inverse length \(\kappa_*>0\): \[
 g_j^2=\frac{\kappa_*}{200j},\quad
 \xi_j=\frac{10000j^2}{\kappa_*^2},\quad
 \kappa_j=\kappa_*,\quad
 b_j=\frac{10000j^2}{\kappa_*},\quad
 2b_jM_j=\frac{20000j^2M_j}{\kappa_*}.
 \tag{91}
\] This changes the coupling, with the full magnetic term displayed; it
is not the coordinate dilation (61). Substitution in the same all-angle
state bound gives only \[
 \|\chi_j\|\le\frac{6144\pi^2}{10000\kappa_*^2}.
 \tag{92}
\] This constant upper bound alone gives no vanishing conclusion or
limiting Hamiltonian. The independent full-vacuum moment calculation in
Section 12 proves the stronger vanishing estimate (104) for this same
trajectory. The next target is a spectral estimate for the actual
interacting states along a coupling trajectory outside (76), such as
(88) or (91), with its full limiting dynamics constructed. None of the
unknown convergence statements has been installed as an assumption.

\subsection{12. A weak-coupling variational bound for the actual
vacuum}\label{a-weak-coupling-variational-bound-for-the-actual-vacuum}

\subsubsection{12.1 Exact compact-group comparison with every physical
coefficient
retained}\label{exact-compact-group-comparison-with-every-physical-coefficient-retained}

Write \(N=N_L\), \(M=M_L\), and \(r=N/M=1+1/(2L)\). These counts retain
every boundary link and face. In quaternion coordinates
\(U=xI+i\sum_{\alpha=1}^3y_\alpha\sigma_\alpha\), the Haar marginal of
\(x\in[-1,1]\) is \[
 d\lambda_x=\frac2\pi(1-x^2)^{1/2}\,dx,\qquad
 \operatorname{tr}U=2x.
 \tag{93}
\] To calculate this marginal, parameterize the unit three-sphere by
\((\cos\varphi,\sin\varphi\,\mathbf n)\), with \(\mathbf n\in S^2\) and
\(0\le\varphi\le\pi\). Its area element is
\(\sin^2\varphi\,d\varphi\,dS^2\). The substitution \(x=\cos\varphi\),
followed by the integral \(\int_{-1}^1\sqrt{1-x^2}\,dx=\pi/2\), gives
(93). The radius-2 link metric multiplies total area by the constant
eight, which cancels from this same Haar probability convention; it
remains in the derivatives below.

For \(t>0\), define the exact integral and unit trial factor \[
 Z_t=\int_{SU(2)}e^{2tx}\,d\lambda,\qquad
 f_t(U)=Z_t^{-1/2}e^{tx},\qquad
 q_t=Z_t^{-1}\int x e^{2tx}\,d\lambda .
 \tag{94}
\] This is a comparison vector, not the physical vacuum. The scalar
\(Z_t\) is retained. Pairing \(x\) with \(-x\) shows \(0<q_t<1\), since
the positive-\(x\) contribution to the numerator is \(2x\sinh(2tx)>0\).

There is the following elementary integral bound, with no large-\(t\)
asymptotics: \[
 0<1-q_t\le\frac{3}{4t}.
 \tag{95}
\] Put \(z=1-x\). Its tilted Haar density is proportional to
\(e^{-2tz}z^{1/2}\sqrt{2-z}\,\mathbf1_{[0,2]}(z)\). On \([0,\infty)\),
let \(\gamma_t\) have density proportional to \(e^{-2tz}z^{1/2}\);
integration by parts of \(z^{3/2}e^{-2tz}\) gives its mean \(3/(4t)\).
The factor \(w(z)=\sqrt{2-z}\mathbf1_{[0,2]}(z)\) is nonnegative and
decreasing. For two independent variables \(z,z'\) of law \(\gamma_t\),
\[
 2\left(\int zw\,d\gamma_t-\int z\,d\gamma_t\int w\,d\gamma_t\right)
 =\iint(z-z')(w(z)-w(z'))\,d\gamma_t(z)d\gamma_t(z')\le0 .
\] Division by the positive mean of \(w\) gives (95). The integrations
are absolutely convergent; the boundary term in the displayed
integration by parts vanishes at zero and infinity.

The metric in Section 10 gives \[
 |\nabla x|^2=(1-x^2)/4,\qquad \Delta x=-3x/4 .
\] The first follows from the radius-2 sphere, and the second also
follows by tracing \(\sum_\alpha T_\alpha^2U=-3U/4\). Integrating
\(\operatorname{div}(e^{2tx}\nabla x)\) over the compact group gives \[
 \frac{t}{2}\langle1-x^2\rangle_t=\frac34q_t,\qquad
 \langle f_t,E_e f_t\rangle=\frac{t^2}{4}\langle1-x^2\rangle_t
 =\frac{3tq_t}{8}\le\frac{3t}{8}.
 \tag{96}
\] Here \(\langle\cdot\rangle_t\) uses precisely \(f_t^2d\lambda\).

Let \(F_t(U)=\prod_{e\in\mathsf E_L}f_t(U_e)\). Its norm is 1. Under its
product density, conjugation invariance and symmetry in \(\mathbf y\)
give \(\langle U_e\rangle_t=q_t I\) and
\(\langle U_e^{-1}\rangle_t=q_t I\). Entrywise integration of the
ordered four-link product therefore gives, with its original
orientations, \[
 \langle F_t,W_pF_t\rangle=2q_t^4,\qquad
 2-2q_t^4=2(1-q_t)(1+q_t+q_t^2+q_t^3)\le\frac6t .
 \tag{97}
\] Independence is used only in this explicit comparison vector.

The exact trial energy and its bound are \[
 \mathfrak q_H[F_t]
 =\frac{3\kappa Ntq_t}{8}+2bM(1-q_t^4)
 \le\frac{3\kappa N}{8}t+\frac{6bM}{t}.
 \tag{98}
\] Every scalar and every face in \(V=b\sum_p(2-W_p)\) remains here. The
physical vacuum from Section 1 is the full ground state and is gauge
invariant, so the variational principle directly gives
\(E\le\mathfrak q_H[F_t]\). One can also make the comparison gauge
invariant without changing the potential integral. Define \[
 \bar\rho_t(U)=\int_{\mathcal G_L}F_t(h\cdot U)^2\,dh,\qquad
 \bar F_t(U)=\sqrt{\bar\rho_t(U)}.
\] This is a smooth positive gauge-invariant unit vector.
Differentiating under the compact integral and applying Cauchy--Schwarz
yields \[
 |\nabla_e\bar F_t|^2
 =\frac{|\int F_t(h\cdot U)\nabla_e(F_t(h\cdot U))\,dh|^2}
       {\int F_t(h\cdot U)^2\,dh}
 \le\int|\nabla_e(F_t(h\cdot U))|^2\,dh.
\] Gauge transformations act by link isometries, and the potential is
gauge invariant. Haar integration proves
\(\mathfrak q_H[\bar F_t]\le\mathfrak q_H[F_t]\). This supplies the
exact gauge map for the comparison; neither vector is substituted for
\(\psi\).

The choice \(t=4\sqrt{bM/(\kappa N)}=2/(g^2\sqrt r)>0\) makes the two
terms in the upper bound (98) equal. Since \(\kappa b=a^{-2}\), and
since the constant comparison vector gives \(E\le2bM\), we obtain \[
 \begin{gathered}
 0\le E\le\mathcal B_{L,a,g}
 :=\min\left\{2bM,\frac{3\sqrt{NM}}{a}\right\},
 \\
 \frac1M\sum_p\mu(2-W_p)
 \le\min\{2,6\sqrt r\,g^2\}
 \le\min\{2,3\sqrt5\,g^2\}.
 \end{gathered}
 \tag{99}
\] The expectation bound follows from \(b\sum_p\mu(2-W_p)=\mu(V)\le E\).
The last inequality uses \(1<r\le5/4\), valid for every \(L\ge2\). Thus
the full actual vacuum has an explicitly controlled mean plaquette
defect at decreasing coupling. This is an average over the actual faces;
no translation invariance of the open box is assumed.

\subsubsection{12.2 An improved full electric second
moment}\label{an-improved-full-electric-second-moment}

The operator equation still reads \(\kappa H_0\psi=(E-V)\psi\).
Positivity and the full potential bound give \[
 0\le\mu(V)\le E\le2bM,\qquad
 \mu(V^2)\le4bM\mu(V).
\] Consequently the exact second moment satisfies \[
 \begin{split}
 \kappa^2\|H_0\psi\|^2
 &=E^2-2E\mu(V)+\mu(V^2)\\
 &\le E^2+(4bM-2E)\mu(V)\\
 &\le4bME-E^2\\
 &\le4bM\mathcal B_{L,a,g}-\mathcal B_{L,a,g}^2 .
 \end{split}
 \tag{100}
\] The coefficient \(4bM-2E\) is nonnegative, which justifies the third
line. The function \(x\mapsto4bMx-x^2\) increases on \([0,2bM]\), which
justifies the last line using (99). All vectors lie in the same \(H^2\)
domain; these equalities and inequalities concern the unchanged
eigenvector.

For an additional explicit upper bound, retain (100) and then bound it
above by \(4bM(3\sqrt{NM}/a)\). The original coefficients give \[
 \|H_0\psi\|
 \le\frac{\sqrt{4bM\mathcal B_{L,a,g}-\mathcal B_{L,a,g}^2}}{\kappa}
 \le\frac{M}{g^3}\sqrt{\frac32}\,r^{1/4}.
 \tag{101}
\] Combining this with the exact central-convolution graph estimate (19)
and orthogonal centering yields \[
 \|\chi_\theta\|
 \le\frac{2\theta^2L^2}{3\kappa}
       \sqrt{4bM\mathcal B_{L,a,g}-\mathcal B_{L,a,g}^2}
 \le\frac23\theta^2L^2\frac{M}{g^3}
       \sqrt{\frac32}\,r^{1/4}.
 \tag{102}
\] No expansion in \(g\), \(\theta\), or total volume has been inserted.
The earlier all-angle bound remains valid independently.

On the original geometric sequence (12), retain \(D_j\) in \(\theta_j\)
and then use \(D_j\ge j^4/4\), \(M_j\le36j^6\), and \(r_j\le5/4\). This
proves \[
 \|\chi_j\|\le C_0\,g_j^{-3}j^{-2},\qquad
 C_0=\frac{768\pi^2\sqrt{3\sqrt5}}{10^8}.
 \tag{103}
\] For \(g_j^2=1/\log j\), it gives
\(\|\chi_j\|\le C_0(\log j)^{3/2}j^{-2}\), strengthening (89). For the
fixed physical electric coefficient trajectory (91), it gives the new
vanishing estimate \[
 \|\chi_j\|\le
 \frac{48\pi^2\sqrt{6\sqrt5}}{3125\,\kappa_*^{3/2}}j^{-1/2},
 \qquad
 d_j\le\frac{13824\sqrt5\,\pi^4}
 {9765625\,\kappa_*^3}j^{-1}.
 \tag{104}
\] The original state remains nonzero at each sufficiently large finite
\(j\). The bound is on its raw squared norm; it does not assert a limit
for the probability after division by \(d_j\).

More generally, (103) proves this vanishing for every positive coupling
sequence satisfying \(g_j^3j^2\to\infty\). The exact
compact-configuration argument in Section 6.3 therefore applies both to
(88) and to (91). The actual vacuum cylinder measures have consistent
subsequential limits, the uncentered projected magnetic measures have
the same limits, and the centered finite measures have total mass
tending to zero. The proof of the uncentered statement is exactly (23),
applied to the improved bound on \(\|(K_{\theta_j}-I)\psi_j\|\) from
(102), before centering. No identification of the limiting Hamiltonian
is required for or inferred from these measure statements.

\subsection{13. Exact plaquette-to-loop maps and an explicit flat
configuration
limit}\label{exact-plaquette-to-loop-maps-and-an-explicit-flat-configuration-limit}

\subsubsection{13.1 Ordered non-Abelian rectangular
filling}\label{ordered-non-abelian-rectangular-filling}

Let \(C\) be the oriented boundary of an \(m\)-by-\(m\) planar square of
fine faces, based at a specified corner. Its original ordered holonomy
has an exact factorization \[
 U_C=\prod_{\nu=1}^{m^2}
       A_\nu U_{p_\nu}^{\varepsilon_\nu}A_\nu^{-1},
 \qquad \varepsilon_\nu\in\{1,-1\},
 \tag{105}
\] for a definite ordering of the faces and connecting path holonomies
\(A_\nu\). Each face occurs once; all conjugating products depend on the
actual link variables.

Here is a constructive proof and specification of the factors. First
take the boundary base at the lower left corner and the positive
orientation in the chosen coordinate plane. Starting with the lower left
face, attach the remaining faces from left to right in the first row,
then from left to right in each successive row. Every added face meets
the existing boundary in a connected path of one or two edges, and the
accumulated union is a disk containing this base corner. If its current
boundary word is \(AqB\), with \(q\) the shared boundary segment, and
the new boundary replaces that segment by \(r\), then \[
 (ArB)(AqB)^{-1}=A(rq^{-1})A^{-1}.
\] The word \(rq^{-1}\) is the boundary of precisely the added face.
Moving its starting vertex to the source convention of Section 1 is a
cyclic conjugation, absorbed into \(A_\nu\); reversing its orientation
gives \(\varepsilon_\nu=-1\). This defines the next factor with no
Abelian rearrangement. Starting from the first face, induction gives
(105), with the new factor multiplied on the left at each step,
equivalently with the displayed product in reverse attachment order. For
a different specified corner, write the full boundary as \(PQ\), where
\(P\) ends at that corner. The boundary based there is
\(QP=P^{-1}(PQ)P\), so conjugate every completed factor by \(P^{-1}\).
Reversing the full boundary inverts the product, reverses its factor
order and changes each sign \(\varepsilon_\nu\). This algorithm fixes
every word for the original base and orientations.

Unitary invariance of the Hilbert--Schmidt norm, the identity
\(I-XY=(I-X)+X(I-Y)\), and then Cauchy--Schwarz imply the deterministic
estimates \[
 \|I-U_C\|_{\rm HS}
 \le\sum_{\nu=1}^{m^2}\|I-U_{p_\nu}\|_{\rm HS},\qquad
 \|I-U_C\|_{\rm HS}^2
 \le2m^2\sum_{\nu=1}^{m^2}(2-W_{p_\nu}).
 \tag{106}
\] The norm of \(I-U_p^{-1}\) equals that of \(I-U_p\). These facts
justify the estimates for every orientation, with all quantum
conjugations in (105) retained.

Applying (99) to the actual vacuum, and using positivity to bound this
subset by the sum over all faces, gives \[
 \mu(\|I-U_C\|_{\rm HS}^2)
 \le12m^2g^2\sqrt{NM}.
 \tag{107}
\] This global upper bound uses no independence and no translation
symmetry of the vacuum. The deterministic subset estimate (106) remains
available when sharper local expectations are known.

\subsubsection{13.2 A full sequence of exact vacuum cylinder
limits}\label{a-full-sequence-of-exact-vacuum-cylinder-limits}

Retain (12), and take the strictly positive coupling sequence \[
 g_j^2=j^{-10},\qquad
 \xi_j=j^{20}/4,\qquad
 \kappa_j=200j^{-9},\qquad
 b_j=50j^{11},\qquad
 2b_jM_j=100j^{11}M_j .
 \tag{108}
\] This changes the coupling and retains the original metric and all
interactions. The vanishing of \(\kappa_j\) does not remove its electric
operator from any finite Hamiltonian, and no inference about an
excitation energy is made from that coefficient alone.

Fix one coarse index \(k\), and push the actual vacuum \(\mu_j\) to
\(Q_{L_k}\) by the exact map \(p_{k,j}\). Every coarse elementary square
has \(m_j=j/k\) links per side on the fine graph. Since
\(\sqrt{N_jM_j}=\sqrt{r_j}M_j\le18\sqrt5\,j^6\), (107) gives \[
 \int\|I-U_p\|_{\rm HS}^2\,d(p_{k,j})_*\mu_j
 \le\frac{216\sqrt5}{k^2}\,g_j^2j^8
 =\frac{216\sqrt5}{k^2}\,j^{-2}.
 \tag{109}
\] All \(M_{L_k}\) coarse faces are fixed in number when \(k\) is fixed.
Their summed defect therefore tends to zero in mean and probability.
More generally the same conclusion holds for any coupling sequence with
\(g_j^2j^8\to0\).

Let \[
 \mathcal F_k=\{U\in Q_{L_k}: U_p=I
                     \text{ for every coarse elementary face }p\}.
\] It is a compact subset. We identify it and its gauge-invariant
probability exactly. Let \(o=(-L_k,-L_k,-L_k)\). For each coarse vertex
\(n\), choose the canonical path from \(o\) to \(n\) that first
increases coordinate 1, then coordinate 2, then coordinate 3, and write
its holonomy as \(q_n\), with \(q_o=I\). Flatness of each elementary
square gives \[
 U_i(n)U_h(n+e_i)=U_h(n)U_i(n+e_h)
\] whenever the square is present, including the inverse form for
backward steps. Moving the last step of a canonical path past the
coordinate-3 and coordinate-2 segments one elementary square at a time
proves \(q_nU_i(n)=q_{n+e_i}\) for each edge. All these paths remain in
the original open box. Consequently \[
 U_e=q_{s(e)}^{-1}q_{t(e)},\qquad
 \mathfrak f_k:
 \prod_{n\ne o}SU(2)\longrightarrow\mathcal F_k,\quad
 (q_n)\longmapsto(q_{s(e)}^{-1}q_{t(e)})_e
 \tag{110}
\] is a continuous bijection whose inverse is the stated canonical path
product. Products telescope, proving both inverse identities. In
particular \(\mathcal F_k\) is exactly the gauge orbit of the identity
configuration, and its gauge quotient is one point.

Push independent Haar variables \(q_n\), \(n\ne o\), through (110); call
the resulting probability \(\sigma_k\). It is invariant under the full
original vertex gauge action. Indeed
\(U_e\mapsto h_{s(e)}U_eh_{t(e)}^{-1}\) corresponds to
\(q_n\mapsto h_o q_nh_n^{-1}\); it preserves the independent Haar law
and keeps \(q_o=I\). Moreover it is the unique invariant probability on
\(\mathcal F_k\). For any continuous \(F\), its gauge average
\(\int F(h\cdot U)\,dh\) has the same value for every \(U\) in this
transitive orbit, by Haar invariance. Integration against any invariant
probability gives that value, which is also its integral under
\(\sigma_k\). Continuous functions determine measures.

Compactness gives weakly convergent subsequences of
\((p_{k,j})_*\mu_j\). Equation (109) and continuity of the nonnegative
summed face defect force every such limit to have support in
\(\mathcal F_k\). Gauge invariance passes to the limit because the
finite action and its pullback on continuous functions commute with the
finite measures. Uniqueness just proved forces every subsequential limit
to equal \(\sigma_k\). If the full sequence did not converge, a
continuous test function and a subsequence of separated integrals would
have a weakly convergent further subsequence, a contradiction. Thus the
full sequence satisfies \[
 (p_{k,j})_*\mu_j\Longrightarrow\sigma_k
 \quad\text{for every fixed }k
 \quad\text{whenever }g_j^2j^8\to0 .
 \tag{111}
\]

The measures \(\sigma_k\) are consistent under refinement. A fine
product of links in (110) telescopes exactly to \(q_s^{-1}q_t\). For
strict refinement \(k'>k\), the retained vertices in fine indices are
\((k'/k)n\), with coordinates bounded by \(kk'<(k')^2\). Thus none is
the fixed fine root \((-(k')^2,-(k')^2,-(k')^2)\), and their \(q\)
variables are all independent Haar. Rebase at the coarse root \(o'\) by
replacing \(q_n\) with \(q_{o'}^{-1}q_n\). Conditional on \(q_{o'}\),
the other retained variables are independent Haar, so the rebased
variables have exactly the law defining \(\sigma_k\). For \(k'=k\) the
map is the identity. This proves \((p_{k,k'})_*\sigma_{k'}=\sigma_k\).
The compact inverse-limit construction of Section 6.3 consequently gives
their unique consistent probability \(\sigma_\infty\), supported on the
compatible pure-gauge configurations. This is an explicit quantum-vacuum
configuration limit of the interacting finite Hamiltonians, with its
parameter trajectory, gauge map and entire limiting cylinder measure
proved.

For every gauge-invariant continuous cylinder \(F\), its value on the
flat orbit is exactly \(F(I)\). Equation (111) therefore yields \[
 \mu_j(F)\to F(I),\qquad
 \|\psi_j(F-\mu_j(F))\|^2
 =\mu_j(|F|^2)-|\mu_j(F)|^2\longrightarrow0 .
 \tag{112}
\] In particular every fixed physical Wilson square has limiting mean 2
and variance zero. This differs from the projective-Haar limit (59),
whose corresponding mean is 0 and variance 1; both limits and their
exact refinement maps have now been constructed. On the gauge quotient
the limit (111) has only constant cylinder observables. It does not
provide the nonzero interacting state required by the research target.
Quantitative control of rescaled curvature observables and the native
magnetic spectral probability remains a separate calculation; neither
follows from weak convergence of these bounded holonomies.

The shrinking curvature background (87) and the limiting quantum
holonomies now both approach the flat configuration on this particular
trajectory, with their exact quantitative bounds (87) and (109). This
relation uses the full quantum measure rather than replacing it by the
deterministic background. On (108) the all-angle magnetic bound (103)
does not tend to zero, so (111)--(112) make no claim that the centered
native translated state's spectral probability or its uncentered measure
has the same limit.

\subsection{14. Quantitative covariance and magnetic mass at every
positive
coupling}\label{quantitative-covariance-and-magnetic-mass-at-every-positive-coupling}

\subsubsection{14.1 Objects, domains and the local score
estimate}\label{objects-domains-and-the-local-score-estimate}

For the complete open box with vertices \(\{-L,\ldots,L\}^3\),
\(L\ge2\), retain \[
 N=3(2L)(2L+1)^2,\qquad M=12L^2(2L+1),
 \qquad H_0=\sum_eE_e,\quad E_e=-\sum_aX_{e,a}^2,
\] \[
 \mathcal W=\sum_p W_p,\qquad
 H=\kappa(H_0-\xi\mathcal W)+2bMI,
 \quad \kappa=\frac{2g^2}{a},\quad b=\frac1{2g^2a},
 \quad \xi=\frac b\kappa=\frac1{4g^4}>0.
\] The positive unit vacuum is the same \(\psi\) at every occurrence.
Write \[
 (H_0-\xi\mathcal W)\psi=e\psi,\quad
 E=\kappa e+2bM,\quad d\mu=\psi^2d\lambda,
 \quad u=\log\psi.
\] Every function and operator in the commutators below is smooth on the
finite product of compact groups when applied to \(\psi\). Thus all
displayed commutators and integrations by parts are defined without an
unbounded-operator domain inference.

Let \(r_e\) be the exact number of incident faces. It satisfies
\(2\le r_e\le4\), including boundary edges. In the metric whose
orthonormal basis is \(T_a=-i\sigma_a/2\), the group is the radius-2
sphere, with diameter \(2\pi\). The full-vacuum partial Bochner identity
in Section 10 proves \[
 \lVert\nabla_eu\rVert_\infty\le2r_e\xi,
 \qquad
 \lVert\nabla_e\log\psi^2\rVert_\infty\le4r_e\xi.
 \tag{113}
\] For the actual conditional density of link \(f\), at arbitrary fixed
exterior links, this gives \[
 p_f(U_f\mid U_{f^c})\ge e^{-8\pi r_f\xi}
 \ge\alpha,\qquad \alpha=e^{-32\pi\xi}.
 \tag{114}
\] Indeed, the logarithmic oscillation is at most \(8\pi r_f\xi\), and
the conditional density integrates to one. The same proof gives, for any
finite edge set \(S\), \[
 p_S(U_S\mid U_{S^c})\ge
 \exp\!\left[-8\pi\xi\sum_{f\in S}r_f\right].
 \tag{115}
\] For (115), join configurations by changing their links one at a time,
apply (113) to each segment, then use the normalization of that
particular conditional density. No conditional product measure is
asserted.

The original weighted electric operator is \[
 \Gamma=\sum_{e=(n,1)}n_2^2E_e.
 \tag{116}
\] In particular, links at \(n_2=0\) have weight zero; they are still
present in \(H_0\), \(\mathcal W\), \(\psi\) and every expectation
below.

\subsubsection{14.2 A local lower bound for the full Wilson
gradient}\label{a-local-lower-bound-for-the-full-wilson-gradient}

Fix an edge \(e\), and choose an incident elementary face \(p\). Let
\(f\) be the other edge of \(p\) parallel to \(e\). This edge exists for
every incident face, even at the boundary. Among the faces incident to
\(e\), only \(p\) contains \(f\): the direction and sign of the
displacement from \(e\) to \(f\) determine that face uniquely.

Fix all links other than \(f\), including \(U_e\), and write \[
 \nabla_e\mathcal W=B+\nabla_eW_p,
 \qquad B=\sum_{q\ni e,\ q\ne p}\nabla_eW_q.
\] The vector \(B\) is independent of \(U_f\). The oriented plaquette
word contains \(U_f\) or its inverse exactly once. The change from that
variable to the full plaquette product preserves Haar, including inverse
traversal. Therefore \[
 \int \nabla_eW_p\,dU_f=0,\qquad
 \int |\nabla_eW_p|^2dU_f=\frac34.
 \tag{117}
\] For the first equality, differentiate the zero Haar mean of the
fundamental trace. For the second, bi-invariance and inversion give the
exact pointwise identity \[
 |\nabla_eW_p|^2=1-\frac{W_p^2}{4},
\] and the Haar second moment of the fundamental trace is one.
Equivalently, for \(U=u_0I+i\mathbf u\cdot\sigma\), the trace is
\(2u_0\), and uniform measure on the unit three-sphere has
\(\int u_0^2=1/4\).

Applying (114) to the nonnegative squared norm and retaining its cross
term until Haar integration gives \[
 \int |\nabla_e\mathcal W|^2 p_f\,dU_f
 \ge e^{-8\pi r_f\xi}
       \int|B+\nabla_eW_p|^2dU_f
 =e^{-8\pi r_f\xi}\left(|B|^2+\frac34\right).
\] The exterior is then integrated against its actual marginal. Hence \[
 \mu(|\nabla_e\mathcal W|^2)
 \ge\frac34 e^{-8\pi r_f\xi}\ge\frac34\alpha.
 \tag{118}
\] There is no cancellation assumption about the other plaquettes. They
remain in \(B\), and its squared norm has nonnegative sign after the
exact integral.

For a sharper finite-box value define \[
 a_e=\max_{p\ni e}\exp[-8\pi r_{f(e,p)}\xi].
\] The same result is valid with \(\alpha\) replaced by \(a_e\),
separately for each edge. In particular, \[
 S_\Gamma:=\mu\!\left(\sum_{e=(n,1)}n_2^2
                     |\nabla_e\mathcal W|^2\right)
 \ge\frac34\sum_{e=(n,1)}n_2^2a_e
 \ge\frac34\alpha\,\mathscr S_L,
 \tag{119}
\] where the exact weighted edge count is \[
 \begin{split}
 \mathscr S_L
 &= (2L)(2L+1)\sum_{m=-L}^Lm^2\\
 &=\frac23L^2(L+1)(2L+1)^2,\qquad
 \frac{\mathscr S_L}{M}=\frac{(L+1)(2L+1)}{18}.
 \end{split}
 \tag{120}
\]

\subsubsection{14.3 A commutator lower bound for the actual electric
covariance}\label{a-commutator-lower-bound-for-the-actual-electric-covariance}

Use a new symbol for the skew operator, to keep it distinct from the
nonnegative excitation operator \(A_{L,\xi}\) in Section 3: \[
 \mathscr B=[H_0,\mathcal W]
 =3\mathcal W-2\sum_e\nabla_e\mathcal W\cdot\nabla_e.
 \tag{121}
\] The coefficient three is exact because every elementary trace has
four spin-\(1/2\) links, each of Casimir \(3/4\), so \(H_0W_p=3W_p\).
Haar integration by parts gives \(\mathscr B^*=-\mathscr B\) on smooth
functions. All relevant vectors are real, so
\(\langle\psi,\mathscr B\psi\rangle=0\).

The separate link Casimirs commute, hence \([\Gamma,H_0]=0\). The Jacobi
identity and the actual eigen-equation imply \[
 \begin{split}
 \langle\psi,[\Gamma,\mathscr B]\psi\rangle
 &=\langle\psi,[H_0,[\Gamma,\mathcal W]]\psi\rangle\\
 &=\xi\langle\psi,[\mathcal W,[\Gamma,\mathcal W]]\psi\rangle
 =2\xi S_\Gamma.
 \end{split}
 \tag{122}
\] To check the middle equality directly, substitute
\(H_0\psi=(e+\xi\mathcal W)\psi\) in the two outer matrix elements; the
scalar \(e\) cancels. To check the final equality, apply the product
rule to any smooth \(F\): \[
 [\Gamma,\mathcal W]F=(\Gamma\mathcal W)F
       -2\sum_{e=(n,1)}n_2^2\nabla_e\mathcal W\cdot\nabla_eF,
\] \[
 [\mathcal W,[\Gamma,\mathcal W]]F
       =2\sum_{e=(n,1)}n_2^2|\nabla_e\mathcal W|^2F.
\] This calculation retains all nonlinear Wilson contributions and all
cross terms inside each full gradient \(\nabla_e\mathcal W\).

Write \[
 \gamma=\langle\psi,\Gamma\psi\rangle,
 \qquad v_\Gamma=(\Gamma-\gamma)\psi.
\] The adjoint identities turn the left side of (122) into
\(2\operatorname{Re}\langle v_\Gamma,\mathscr B\psi\rangle\). Thus \[
 \boxed{\quad
 \lVert v_\Gamma\rVert\ge
 \frac{\xi S_\Gamma}{\lVert\mathscr B\psi\rVert}.
 \quad}
 \tag{123}
\] Its denominator cannot vanish: otherwise (122) would contradict
\(S_\Gamma>0\), already proved in (119). The following explicit upper
bounds remove that denominator from the numerical estimate.

Let \[
 R_2=\sum_er_e^2=384L^3+48L^2-24L\le16M.
 \tag{124}
\] For this count, set \(r(m)=1\) at \(m=\pm L\) and \(r(m)=2\)
otherwise. For a direction-one edge, \(r_e=r(n_2)+r(n_3)\). Use
\(\sum r(m)=4L\), \(\sum r(m)^2=8L-2\), sum \((r(n_2)+r(n_3))^2\),
multiply by \(2L\), and then by three orientations. This gives exactly
(124).

The derivative norm of a single trace is at most one, so
\(|\nabla_e\mathcal W|\le r_e\). Equation (113), followed by the
triangle inequality in (121), gives \[
 \lVert\mathscr B\psi\rVert
 \le6M+4\xi R_2\le M(6+64\xi).
 \tag{125}
\] Consequently, at every \(\xi>0\) and every \(L\ge2\), \[
 \boxed{\quad
 \lVert(\Gamma-\langle\Gamma\rangle_\psi)\psi\rVert
 \ge
 \frac{\xi e^{-32\pi\xi}(L+1)(2L+1)}
      {24(6+64\xi)} >0.
 \quad}
 \tag{126}
\] The exponent in this estimate is local in \(\xi\); it does not
contain \(M\). No substitute vacuum or assumption about the spatial
correlation length was used.

One can keep a stronger denominator if useful. Put
\(\mathcal E_0=\langle\psi,H_0\psi\rangle\). Cauchy--Schwarz pointwise
in the full product tangent space, then in \(L^2\), gives \[
 \lVert\mathscr B\psi\rVert
 \le6M+2\sqrt{R_2\mathcal E_0}.
 \tag{127}
\] Both (125) and (127) concern the same skew operator acting on the
same vacuum. With physical units restored,
\([H,\mathcal W]=\kappa\mathscr B\), so \[
 \lVert[H,\mathcal W]\psi\rVert
 \le\kappa(6M+4\xi R_2),
 \qquad
 \lVert[H,\mathcal W]\psi\rVert
 \le\kappa(6M+2\sqrt{R_2\mathcal E_0}).
 \tag{128}
\] The constant \(2bM\) commutes with \(\mathcal W\); it is retained in
\(H\) and in the expression for \(E\).

\subsubsection{14.4 An exact all-angle lower bound for the magnetic
state}\label{an-exact-all-angle-lower-bound-for-the-magnetic-state}

Retain \[
 K_\theta=\prod_{e=(n,1)}C_{e,\theta n_2},\quad
 c_\theta=\langle\psi,K_\theta\psi\rangle,
 \quad\chi_\theta=(K_\theta-c_\theta)\psi,
 \quad d_\theta=\lVert\chi_\theta\rVert^2.
\] By Section 3, \(K_\theta\) is a self-adjoint Markov contraction
commuting with \(H_0\). Denote its Markov kernel by \(P_\theta(U,dV)\);
independently on each direction-one link, \(V_e=k_e^{-1}U_e\), where
\(k_e\) is uniformly conjugate to \(\exp(\theta n_2H_c)\). Other links
are unchanged. Symmetry means \(\lambda(dU)P_\theta(U,dV)\) is invariant
under swapping \(U,V\).

Define the exact positive quantity \[
 I_\theta=\int\lambda(dU)P_\theta(U,dV)\,
       \psi(U)\psi(V)
       (\mathcal W(U)-\mathcal W(V))^2.
 \tag{129}
\] The same Jacobi calculation as above, now with \(K_\theta\), proves
\[
 \begin{split}
 \langle\psi,[K_\theta,\mathscr B]\psi\rangle
 &=\xi\langle\psi,[\mathcal W,[K_\theta,\mathcal W]]\psi\rangle\\
 &=-\xi I_\theta.
 \end{split}
 \tag{130}
\] Indeed the kernel of the nested multiplication commutator is
\(-[\mathcal W(U)-\mathcal W(V)]^2P_\theta(U,dV)\). Since
\(\langle\psi,\mathscr B\psi\rangle=0\), Cauchy--Schwarz yields the
exact all-angle comparison \[
 \boxed{\qquad
 \lVert\chi_\theta\rVert\ge
 \frac{\xi I_\theta}{2\lVert\mathscr B\psi\rVert}.
 \qquad}
 \tag{131}
\]

Here is a fully explicit lower bound for (129) which avoids the global
Haar overlap of \(\psi\). For \(s\in\mathbb R\), define \[
 \delta(s)=\min_{m\in\mathbb Z}|s-2\pi m|\in[0,\pi].
\] In the preserved radius-2 metric, left translation by a conjugate of
\(\exp(sH_c)\) moves a point by geodesic distance \(2\delta(s)\). By
(113), changing the translated links one at a time gives \[
 \psi(V)\ge e^{-\mathscr R_\theta}\psi(U),\qquad
 \mathscr R_\theta=4\xi\sum_{e=(n,1)}r_e\delta(\theta n_2).
 \tag{132}
\] This formula keeps the original link angles, including all
resonances. A convenient weaker bound is \[
 \mathscr R_\theta
 \le4\xi|\theta|\sum_{e=(n,1)}r_e|n_2|
 =16\xi|\theta|L^3(4L+3).
 \tag{133}
\] For the count in (133), \(\sum|m|=L(L+1)\), \(\sum r(m)|m|=2L^2\),
and \(\sum r(m)=4L\), so the weighted incidence sum is
\(2L[(2L+1)2L^2+4L^2(L+1)]
 =4L^3(4L+3)\).

Choose any elementary face \(p\), and let \(S_p\) be its four physical
links and \(R_p=\sum_{f\in S_p}r_f\le16\). At fixed exterior and fixed
conjugating variables in the kernel, Haar integration in \(S_p\) makes
the function \(W_p(U)-W_p(V)\) orthogonal to every other face
difference. To prove this, two different elementary faces share at most
one physical edge. There is therefore an edge of \(p\) absent from the
other face. The fundamental trace has zero integral in that edge, and
its translated trace has the same zero integral. Each term from the
other face is independent of that edge. This proves each cross-term
vanishes; the square of the sum of all other face differences remains
nonnegative.

It follows, after the kernel average, that \[
 \int dU_{S_p}\,dP_\theta\,
    (\mathcal W(U)-\mathcal W(V))^2
 \ge 2(1-\alpha_p(\theta)),
 \tag{134}
\] where the exact multiplier of \(K_\theta\) on \(W_p\) is \[
 \alpha_p(\theta)=
 \begin{cases}
 \cos(\theta n_2)\cos(\theta(n_2+1)),&p\text{ has type }12,\\
 \cos^2(\theta n_2),&p\text{ has type }13,\\
 1,&p\text{ has type }23.
 \end{cases}
\] For (134), both face-trace squares have Haar integral one. Their
cross term, after the independent conjugacy average, is
\(\langle W_p,K_\theta W_p\rangle=\alpha_p\); this also proves that the
bound is independent of the fixed exterior.

Apply (115) to the nonnegative integrand in (134), and then integrate
the true exterior marginal. Combining with (132) gives \[
 I_\theta\ge
 2e^{-\mathscr R_\theta-8\pi\xi R_p}
       (1-\alpha_p(\theta)).
 \tag{135}
\] Thus every selected face furnishes the all-positive-coupling,
all-angle bound \[
 \boxed{\quad
 d_\theta\ge
 \left[
 \frac{\xi e^{-\mathscr R_\theta-8\pi\xi R_p}
              (1-\alpha_p(\theta))}
      {6M+4\xi R_2}
 \right]^2.
 \quad}
 \tag{136}
\] The denominator can instead be the sharper expression in (127), with
any proved upper bound for \(\mathcal E_0\). A type-12 face based at
\(n_2=0\) has \(1-\alpha_p=1-\cos\theta\), strictly positive precisely
when \(\theta\notin2\pi\mathbb Z\). One can also choose the definite
corner face of type 12 based at \((-L,-L,-L)\). Its four incidence
numbers are \(2,3,3,2\), so \(R_p=10\), while \[
 \alpha_p=\cos(L\theta)\cos((L-1)\theta)
 \le\frac{1+\cos\theta}{2}.
\] Thus a single fully explicit choice in (136) gives \[
 d_\theta\ge
 \left[
 \frac{\xi e^{-\mathscr R_\theta-80\pi\xi}
       [1-\cos(L\theta)\cos((L-1)\theta)]}
      {6M+4\xi R_2}
 \right]^2
 \ge
 \left[
 \frac{\xi e^{-\mathscr R_\theta-80\pi\xi}(1-\cos\theta)}
      {2(6M+4\xi R_2)}
 \right]^2.
 \tag{137}
\] Therefore (136) is a quantitative version of the exact nonzero-state
theorem, valid also on the original cusp sequence; it retains its actual
transport penalty rather than declaring it uniform. Choosing a test face
never removes any other face from the theory.

\subsection{15. Relative comparison of the full physical spectral
probabilities}\label{relative-comparison-of-the-full-physical-spectral-probabilities}

\subsubsection{15.1 Exact joint spin calculation at every
angle}\label{exact-joint-spin-calculation-at-every-angle}

Let \(L\ge2\), let \(Q_L=SU(2)^{N_L}\) with product Haar probability,
and retain \(T_a=-i\sigma_a/2\), \(E_e=-\sum_aX_{e,a}^2\) and

\[
 H_0=\sum_eE_e,\qquad
 \Gamma=\sum_{e=(n,1)}n_2^2E_e.
\]

The sum defining \(H_0\) includes every edge; the displayed weighted sum
defining \(\Gamma\) includes every edge in direction 1, with its
original integer coordinate \(n_2\), including those with zero weight.
Write

\[
 K_\theta=\prod_{e=(n,1)}C_{e,\theta n_2},\qquad
 c_q(s)=\frac1{2q+1}\sum_{m=-q}^{q}e^{2ims},
 \quad q\in\tfrac12\mathbb Z_{\ge0}.
\]

The finite sum is the definition at all angles, including
\(s\in\pi\mathbb Z\). On a joint spin block, put

\[
 \lambda_e=q_e(q_e+1),\qquad w_e=n_2^2,\qquad
 \gamma=\sum_{e=(n,1)}w_e\lambda_e.
\]

For this calculation only, give each translated edge an independent
random variable \(m_e\) uniformly distributed on its exact spin weights

\[
 \{-q_e,-q_e+1,\ldots,q_e\}.
\]

These finite auxiliary random variables express the exact central
multiplier; they do not change the quantum state or the Haar convention.
The signed integer coordinate remains in

\[
 Y=2\sum_{e=(n,1)}n_2m_e.
\]

Independence, symmetry and the finite-product identity give

\[
 k(\theta)=\prod_ec_{q_e}(\theta n_2)
 =\mathbb E e^{i\theta Y}=\mathbb E\cos(\theta Y).
\tag{138}
\]

The exact moments of each weight variable are

\[
 \mathbb E m_e=\mathbb E m_e^3=0,\qquad
 \mathbb E m_e^2=\frac{\lambda_e}{3},\qquad
 \mathbb E m_e^4=\frac{\lambda_e(3\lambda_e-1)}{15}.
\tag{139}
\]

For an explicit verification at integral and half-integral spin, set

\[
 d=2q\in\mathbb Z_{\ge0},\qquad m=k-d/2,
 \quad k=0,\ldots,d.
\]

Expanding \((k-d/2)^2\) and \((k-d/2)^4\), then using

\[
 \sum_{k=0}^dk=\frac{d(d+1)}2,\quad
 \sum_{k=0}^dk^2=\frac{d(d+1)(2d+1)}6,\quad
 \sum_{k=0}^dk^3=\frac{d^2(d+1)^2}4,
\]

\[
 \sum_{k=0}^dk^4
 =\frac{d(d+1)(2d+1)(3d^2+3d-1)}{30},
\]

gives (139) after division by \(d+1\), with

\[
 \lambda=q(q+1)=\frac{d(d+2)}4.
\]

The consecutive-power sum identities hold at \(d=0\); subtracting their
right sides at \(d\) and \(d-1\) gives \(d,d^2,d^3,d^4\), respectively,
which proves them by induction. Symmetry gives the two odd moments.

Thus the entire joint second and fourth moments are

\[
 \mathbb E Y^2=\frac43\gamma,
\]

\[
 \begin{aligned}
 \mathbb E Y^4
 &=16\left[
 \sum_ew_e^2\frac{\lambda_e(3\lambda_e-1)}{15}
 +6\sum_{e<f}w_ew_f\frac{\lambda_e\lambda_f}{9}
 \right]\\
 &=\frac{16}{3}\gamma^2
 -\frac{32}{15}\sum_ew_e^2\lambda_e^2
 -\frac{16}{15}\sum_ew_e^2\lambda_e
 \le\frac{16}{3}\gamma^2.
 \end{aligned}
\tag{140}
\]

No cross term has been omitted: terms having an odd multiplicity vanish
by (139), and each unordered pair of distinct edges has precisely six
placements in the fourth power. The signed coordinates produce exactly

\[
 n_2^4=w_e^2,\qquad n_2^2(n'_2)^2=w_ew_f
\]

in the surviving terms.

For all real \(x\),

\[
 0\le\cos x-1+\frac{x^2}{2}\le\frac{x^4}{24}.
\tag{141}
\]

Indeed \(1-\cos x\le x^2/2\), because for \(x\ge0\) one has

\[
 1-\cos x=\int_0^x\sin t\,dt\le\int_0^xt\,dt,
\]

and both sides are even. The function on the middle of (141) has value
and first derivative zero at zero and second derivative \(1-\cos x\),
which lies between zero and \(x^2/2\). Integrating these bounds twice
for

\[
 x\ge0
\]

gives (141), and evenness gives the other half-line.

Combining (138)--(141) proves on every joint spin block, for every real

\[
 \theta\in\mathbb R,
\]

the exact scalar bounds

\[
 0\le k(\theta)-1+\frac23\theta^2\gamma
 \le\frac{\theta^4}{24}\mathbb E Y^4
 \le\frac29\theta^4\gamma^2.
\tag{142}
\]

The full product matrix coefficients give an orthogonal Hilbert sum of
these joint spin blocks, with their complete multiplicities. For

\[
 F\in D(\Gamma^2),
\]

square (142), multiply by the squared norm of the corresponding block
coefficient, and sum over all blocks. Parseval's identity yields

\[
 \boxed{
 \left\|\left(K_\theta-I+\frac23\theta^2\Gamma\right)F\right\|
 \le\frac29\theta^4\|\Gamma^2F\|.}
\tag{143}
\]

This is an inequality on the stated graph domain. Every block is
included; there is no finite spin truncation and no assertion of a
bounded-operator Taylor expansion. Since the exact eigenvalues satisfy

\[
 0\le\gamma\le L^2\sum_e\lambda_e,
\]

one also has, for \(F\in D(H_0^2)\),

\[
 \|\Gamma^2F\|\le L^4\|H_0^2F\|.
\tag{144}
\]

\subsubsection{15.2 Fourth electric moment of the actual
vacuum}\label{fourth-electric-moment-of-the-actual-vacuum}

Retain the full physical Hamiltonian and scalar Wilson term:

\[
 H=\kappa H_0+b\sum_p(2-W_p),\qquad
 \kappa=\frac{2g^2}{a},\quad b=\frac1{2g^2a},\quad
 \xi=\frac b\kappa=\frac1{4g^4}.
\]

Let \(\psi>0\) be its actual smooth ground vector with

\[
 \|\psi\|=1,\qquad H\psi=E\psi.
\]

Define only the exact displayed quantities

\[
 e_0=\frac E\kappa,\qquad
 v=\xi\sum_p(2-W_p)=2\xi M_L-\xi\mathcal W.
\]

In particular \(e_0\) is not the shifted energy \(e_{L,\xi}\): the
relation is \(e_0=e_{L,\xi}+2\xi M_L\). The vacuum equation is

\[
 H_0\psi=(e_0-v)\psi.
\tag{145}
\]

Every face has four distinct link occurrences on this open cubic box,
and each occurrence carries fundamental spin \(1/2\). The exact link
Casimir is \(3/4\), hence

\[
 H_0W_p=3W_p,\qquad H_0v=-3\xi\mathcal W.
\tag{146}
\]

The product formula for \(H_0=-\Delta\) is

\[
 H_0(f\psi)=fH_0\psi+(H_0f)\psi
             -2\sum_{e,a}(X_{e,a}f)(X_{e,a}\psi).
\]

Apply it to \(f=e_0-v\), use (145)--(146), and retain all derivatives to
get

\[
 \boxed{H_0^2\psi=(e_0-v)^2\psi
          +3\xi\mathcal W\psi
          +2\sum_{e,a}(X_{e,a}v)(X_{e,a}\psi).}
\tag{147}
\]

Smoothness on the finite compact product makes every term and every
operator domain in (147) legitimate.

The constant trial vector gives \(E\le2bM_L\), because its electric
energy is zero and the Haar mean of each fundamental face trace is zero.
The latter follows by integrating one of the four distinct links. Also
the original potential is nonnegative and bounded by \(4bM_L\).
Therefore

\[
 0\le e_0\le2\xi M_L,\qquad
 0\le v\le4\xi M_L,\qquad
 |e_0-v|\le4\xi M_L.
\tag{148}
\]

For completeness, if a face is written with a particular link as

\[
 AUB\quad\hbox{or}\quad AU^{-1}B,
\]

bi-invariance and inversion show that the squared norm of its derivative
in that link equals the squared gradient norm of the trace on \(SU(2)\).
Write \(U=u_0I+i\mathbf u\cdot\boldsymbol\sigma\). Direct
differentiation under \(\exp(tT_a)U\) gives
\(X_a\operatorname{tr}U=u_a\), and hence

\[
 \sum_a|X_a\operatorname{tr}U|^2
 =\sum_au_a^2=1-u_0^2\le1.
\]

Thus if \(r_e\le4\) is the actual number of incident faces,

\[
 |\nabla_ev|\le\xi r_e\le4\xi,
 \qquad \||\nabla v|\|_\infty\le4\xi\sqrt{N_L}.
\tag{149}
\]

Haar integration by parts and (145) give

\[
 \int|\nabla\psi|^2d\lambda_L
 =\langle\psi,H_0\psi\rangle
 =e_0-\int v\psi^2d\lambda_L
 \le2\xi M_L.
\tag{150}
\]

Pointwise Cauchy--Schwarz in the complete \((e,a)\) index, followed by
(149)--(150), therefore proves

\[
 \left\|2\sum_{e,a}(X_{e,a}v)(X_{e,a}\psi)\right\|
 \le8\sqrt2\,\xi^{3/2}\sqrt{N_LM_L}.
\tag{151}
\]

No mixed electric contribution has been removed. The first two terms of
(147) have norms at most \(16(\xi M_L)^2\) and \(6\xi M_L\).
Consequently

\[
 \|H_0^2\psi\|
 \le16(\xi M_L)^2+6\xi M_L
      +8\sqrt2\,\xi^{3/2}\sqrt{N_LM_L}.
\tag{152}
\]

Now retain the exact finite-box counts

\[
 N_L=3(2L)(2L+1)^2,\qquad M_L=12L^2(2L+1),
\]

so that \(N_L/M_L=1+1/(2L)\le5/4\). For \(\xi\ge1\), \(L\ge2\), one has
\(\xi M_L\ge1\) and \(\sqrt\xi M_L\ge1\). Thus (152) gives

\[
 \boxed{
 \|H_0^2\psi\|
 \le(22+4\sqrt{10})(\xi M_L)^2
 <35(\xi M_L)^2.}
\tag{153}
\]

The last strict numerical inequality follows from \(4\sqrt{10}<13\),
equivalently \(160<169\).

There is also a slightly stronger optional bound with the same
hypotheses. Because every face has four incident edges,

\[
 \sum_er_e=4M_L,\qquad \sum_er_e^2\le4\sum_er_e=16M_L.
\]

Using the first inequality in (149) with this sum replaces the last term
of (152) by \(8\sqrt2\,\xi^{3/2}M_L\), and proves

\[
 \|H_0^2\psi\|\le(22+8\sqrt2)(\xi M_L)^2.
\tag{154}
\]

Both bounds concern the same actual vacuum, with its full magnetic
interaction and original physical coefficients.

\paragraph{15.2.1 Improved exponent using the proved weak-coupling
energy
estimate}\label{improved-exponent-using-the-proved-weak-coupling-energy-estimate}

We next prove the stronger bound

\[
 B_2=15\xi^{7/4}M_L^2\qquad(\xi\ge1,\ L\ge2).
\]

This follows from the exact vacuum equation (147) and the already proved
variational energy and electric second-moment estimates (99)--(101)
above. Here is the complete passage between those quantities. Write
\(r=N_L/M_L\), retaining its exact value \(1+1/(2L)\). The proved energy
inequality gives

\[
 e_0=\frac E\kappa
 \le\frac{3\sqrt{N_LM_L}/a}{2g^2/a}
 =3\sqrt{N_LM_L}\sqrt\xi.
\tag{155}
\]

Moreover, without discarding the exact moment identity,

\[
 \begin{aligned}
 \|H_0\psi\|^2
 &=e_0^2-2e_0\mu(v)+\mu(v^2)\\
 &\le e_0^2+(4\xi M_L-2e_0)\mu(v)\\
 &\le4\xi M_Le_0-e_0^2
 \le4\xi M_Le_0,
 \end{aligned}
\]

where \(\mu(v)\le e_0\), \(\mu(v^2)\le4\xi M_L\mu(v)\), and

\[
 4\xi M_L-2e_0\ge0
\]

follow from (148) and the nonnegative electric form. Substitute (155) to
obtain

\[
 \|H_0\psi\|\le2\sqrt3\,r^{1/4}\xi^{3/4}M_L.
\tag{156}
\]

The first term of (147) can now be bounded through (145) itself:

\[
 \|(e_0-v)^2\psi\|
 \le\|e_0-v\|_\infty\|(e_0-v)\psi\|
 \le8\sqrt3\,r^{1/4}\xi^{7/4}M_L^2.
\]

The second term remains \(6\xi M_L\). For the third, retain the full
incidence count from (154), so

\[
 \||\nabla v|\|_\infty\le4\xi\sqrt{M_L},\qquad
 \||\nabla\psi|\|_2\le\sqrt{e_0}
 \le\sqrt3\,r^{1/4}\xi^{1/4}\sqrt{M_L}.
\]

Therefore (147) gives, with the explicit coefficient

\[
 A_r=8\sqrt3\,r^{1/4},
\]

the full upper bound

\[
 \boxed{\|H_0^2\psi\|
 \le A_r\xi^{7/4}M_L^2+6\xi M_L+A_r\xi^{5/4}M_L.}
\tag{157}
\]

For \(L\ge2\), \(M_L\ge M_2=240\) and \(r\le5/4\). The exact rational
comparisons

\[
 \sqrt3\le\frac74,\qquad
 r^{1/4}\le\frac{17}{16}
\]

follow respectively from \(3\le49/16\) and

\[
 \frac54=\frac{81920}{65536}
 <\frac{83521}{65536}=\left(\frac{17}{16}\right)^4.
\]

Thus \(A_r\le119/8\). Divide the right side of (157) by the positive
quantity \(\xi^{7/4}M_L^2\); for \(\xi\ge1\) its ratio is at most

\[
 A_r+\frac{6}{\xi^{3/4}M_L}
       +\frac{A_r}{\xi^{1/2}M_L}
 \le\frac{119}{8}+\frac{6+119/8}{240}
 =\frac{28727}{1920}<15.
\]

This proves the stronger explicit result

\[
 \boxed{\|H_0^2\psi\|\le15\xi^{7/4}M_L^2
 \qquad(\xi\ge1,\ L\ge2).}
\tag{158}
\]

The strict gap to the last integer constant is \(73/1920\). All earlier
all-angle and relative-state estimates therefore remain valid with this
improved \(B_2\). No replacement ground vector or weak-coupling
expansion has been used.

\subsubsection{15.3 Exact centering and relative
error}\label{exact-centering-and-relative-error}

Let

\[
 P_\psi^\perp=I-|\psi\rangle\langle\psi|,\qquad
 v_\Gamma=P_\psi^\perp\Gamma\psi,
 \qquad\sigma_\Gamma=\|v_\Gamma\|,
\]

\[
 c_\theta=\langle\psi,K_\theta\psi\rangle,\qquad
 \chi_\theta=(K_\theta-c_\theta)\psi,
 \qquad d_\theta=\|\chi_\theta\|^2.
\]

In fact \(\sigma_\Gamma>0\) for every \(\xi>0\), \(L\ge2\). To prove
this, suppose \(v_\Gamma=0\). Then \(\Gamma\psi=A\psi\), where

\[
 A=\langle\psi,\Gamma\psi\rangle\ge0.
\]

Haar integration of \(\Gamma\psi\) is zero, whereas

\[
 \int\psi\,d\lambda_L>0.
\]

It follows that \(A=0\). The quadratic form of \(\Gamma\) is a sum of

\[
 n_2^2\sum_a\|X_{e,a}\psi\|^2,
\]

so \(\psi\) is independent of each direction-1 link with \(n_2\ne0\).
Choose such a link \(e\) incident to at least one face, which exists in
the stated box. The function \(H_0\psi\) is also independent of \(U_e\):
all other link derivatives preserve this independence and the \(e\)-link
derivatives vanish. Equation (145), positivity of \(\psi\), and its
independence of \(U_e\) imply that \(v\) is independent of \(U_e\). But
when all other links are the identity, every incident face word is

\[
 U_e\quad\hbox{or}\quad U_e^{-1},
\]

and every nonincident face word is the identity. Hence exactly

\[
 v(U_e,I_{\ne e})=\xi r_e(2-\operatorname{tr}U_e),
\]

which is nonconstant since \(\xi r_e>0\). This contradiction proves

\[
 \sigma_\Gamma>0.
\tag{159}
\]

Write \(R_\theta=K_\theta-I+(2/3)\theta^2\Gamma\). Exact centering
yields

\[
 \boxed{\chi_\theta=-\frac23\theta^2v_\Gamma
                   +P_\psi^\perp R_\theta\psi.}
\tag{160}
\]

For any proved upper bound \(B_2\ge\|H_0^2\psi\|\), (143)--(144) imply

\[
 \left\|\chi_\theta+\frac23\theta^2v_\Gamma\right\|
 \le\frac29\theta^4L^4B_2.
\tag{161}
\]

Here \(B_2\) can be taken to be the explicit right side of (152) at
every positive \(\xi\), or the improved \(15\xi^{7/4}M_L^2\) from (158),
or the earlier \(35(\xi M_L)^2\), on the completely verified range

\[
 \xi\ge1,\qquad L\ge2.
\]

For \(\theta\ne0\), define the displayed scalar bound

\[
 \eta=\frac{\theta^2L^4B_2}{3\sigma_\Gamma},\qquad
 t_\theta=\frac23\theta^2\sigma_\Gamma.
\tag{162}
\]

Thus the actual relative vector error is at most \(\eta\):

\[
 \left\|\chi_\theta+\frac23\theta^2v_\Gamma\right\|
 \le\eta t_\theta.
\tag{163}
\]

The raw state and its raw squared norm remain unchanged. In particular,

\[
 \max(0,1-\eta)^2\frac49\theta^4\sigma_\Gamma^2
 \le d_\theta
 \le(1+\eta)^2\frac49\theta^4\sigma_\Gamma^2.
\tag{164}
\]

When the explicit bound \(\eta<1\), (164) alone proves that the actual
translated state is nonzero. The exact all-angle zero set from Section 4
above is \(\theta\in2\pi\mathbb Z\); on its complement the spectral
probability below exists regardless of the size of \(\eta\).

\subsubsection{15.4 Full physical spectral probability and raw finite
measures}\label{full-physical-spectral-probability-and-raw-finite-measures}

Set \(A=H-E\), with the exact physical coefficient and scalar shift in
the preceding sections. Its spectral projections are denoted

\[
 P(B)=\mathbf1_B(H-E),\qquad B\subset\mathbb R\text{ Borel}.
\]

On \(\theta\notin2\pi\mathbb Z\), define the two actual spectral
probabilities by their displayed raw denominators:

\[
 \nu_\theta(B)=\frac{\langle\chi_\theta,P(B)\chi_\theta\rangle}{d_\theta},
 \qquad
 \zeta_\Gamma(B)=\frac{\langle v_\Gamma,P(B)v_\Gamma\rangle}
                          {\sigma_\Gamma^2}.
\tag{165}
\]

Both vectors are orthogonal to the same actual vacuum. No spectral
projection onto a finite spin subspace or finite energy window has been
used in their construction.

Here is a direct finite-rank argument proving the following quantitative
bound. Let

\[
 x=\chi_\theta,\qquad y=-\frac23\theta^2v_\Gamma,
 \qquad r=\|x\|,\qquad t=\|y\|=t_\theta.
\]

For this proof only, use the unit directions \(u=x/r\) and \(w=y/t\) in
the same Hilbert space; they are not replacement physical states. Since

\[
 (I-|u\rangle\langle u|)x=0,
\]

\begin{enumerate}
\def\labelenumi{(\arabic{enumi})}
\setcounter{enumi}{162}
\tightlist
\item
  gives
\end{enumerate}

\[
 \sqrt{1-|\langle u,w\rangle|^2}
 =\|(I-|u\rangle\langle u|)w\|
 \le\frac{\|y-x\|}{t}\le\eta.
\tag{166}
\]

Let \(s=\sqrt{1-|\langle u,w\rangle|^2}\), so \(0\le s\le\min(1,\eta)\).
The self-adjoint finite-rank operator

\[
 D=|u\rangle\langle u|-|w\rangle\langle w|
\]

vanishes on the orthogonal complement of their span and has trace zero.
Direct multiplication gives

\[
 \operatorname{Tr}D^2=2(1-|\langle u,w\rangle|^2)=2s^2.
\]

If \(s=0\), then \(D=0\). Otherwise its two eigenvalues are \(s,-s\),
because their sum is zero and the sum of their squares is \(2s^2\). For
associated orthonormal eigenvectors \(e_+,e_-\), any projection \(P\)
therefore satisfies

\[
 |\operatorname{Tr}(DP)|
 =s\,|\langle e_+,Pe_+\rangle-\langle e_-,Pe_-\rangle|
 \le s,
\]

because the two expectations lie in \([0,1]\). Apply this to every full
spectral projection \(P(B)\) to obtain

\[
 \boxed{\sup_{B\text{ Borel}}|\nu_\theta(B)-\zeta_\Gamma(B)|
         \le\min(1,\eta).}
\tag{167}
\]

To remove any convention ambiguity, define the total variation mass norm
of a finite real signed measure \(m\) by

\[
 \|m\|_{\mathrm{TV},1}
 =\sup_{f\text{ real Borel},\,|f|\le1}\left|\int f\,dm\right|.
\]

The same two-eigenvalue calculation with the self-adjoint contraction

\[
 f(H-E)
\]

gives

\[
 \boxed{\|\nu_\theta-\zeta_\Gamma\|_{\mathrm{TV},1}
 \le2\min(1,\eta).}
\tag{168}
\]

The probability-distance convention is the supremum over Borel sets in
(167), which is one half this mass norm.

The corresponding raw finite measures retain their exact masses. Put

\[
 m_\theta(B)=\langle\chi_\theta,P(B)\chi_\theta\rangle
            =d_\theta\nu_\theta(B),
\]

\[
 m_{\Gamma,\theta}(B)
 =\frac49\theta^4\langle v_\Gamma,P(B)v_\Gamma\rangle
 =\frac49\theta^4\sigma_\Gamma^2\zeta_\Gamma(B).
\tag{169}
\]

For every bounded operator \(F\) of norm at most one,

\[
 \begin{aligned}
 |\langle x,Fx\rangle-\langle y,Fy\rangle|
 &\le\|x-y\|(\|x\|+\|y\|)\\
 &\le\eta(2+\eta)t_\theta^2.
 \end{aligned}
\]

Consequently, applying this to the same real bounded spectral functions,

\[
 \boxed{\|m_\theta-m_{\Gamma,\theta}\|_{\mathrm{TV},1}
 \le\eta(2+\eta)\frac49\theta^4\sigma_\Gamma^2.}
\tag{170}
\]

Taking \(F=I\) gives the corresponding raw mass error; (164) gives the
sharper two-sided mass statement. The same physical spectral operators
appear throughout. Thus no vanishing raw mass has been silently removed.

For example, \(e^{-t(H-E)}\) for \(t\ge0\) is a contraction because the
actual excitation operator is nonnegative. Equation (168) applies to its
bounded spectral function at every physical time, while (170) applies to
the unscaled matrix elements. Neither estimate by itself bounds an
unbounded energy moment or proves low-energy weight: those require
further information about the same exact electric covariance state

\[
 v_\Gamma=(\Gamma-\langle\psi,\Gamma\psi\rangle)\psi.
\]

\subsection{16. An explicit cusp giving vanishing relative spectral
error}\label{an-explicit-cusp-giving-vanishing-relative-spectral-error}

\subsubsection{16.1 A bound with all finite parameters
displayed}\label{a-bound-with-all-finite-parameters-displayed}

Write \(M=M_L\) only within this subsection. At every \(L\ge2\) and
\(\xi\ge1\), the preceding lower covariance bound and fourth electric
moment give the two explicit quantities \[
 s_{L,\xi}:=
 \frac{\xi e^{-32\pi\xi}(L+1)(2L+1)}{24(6+64\xi)}
 \le\sigma_\Gamma,
 \qquad B_{2,L,\xi}:=15\xi^{7/4}M^2\ge\|H_0^2\psi\|.
 \tag{171}
\] In particular their quotient contains no unspecified lower variance.
For the actual angle \(\theta=2\pi a^2/D_T\), the relative error in
Section 15 is bounded by \[
 \begin{split}
 \eta_{L,a,g,T}
 &:=\frac{\theta^2L^4B_{2,L,\xi}}{3\sigma_\Gamma}
 \le\bar\eta_{L,a,g,T},\\
 \bar\eta_{L,a,g,T}
 &:=\frac{120\theta^2L^4\xi^{3/4}M^2(6+64\xi)}
 {e^{-32\pi\xi}(L+1)(2L+1)} .
 \end{split}
 \tag{172}
\] This is an inequality for the complete actual eigenvector of (1). The
factor \(e^{-32\pi\xi}\) was obtained by integration in one physical
link, with the full exterior marginal retained. It has no volume in its
exponent.

\subsubsection{16.2 The retained periods on a corrected simultaneous
sequence}\label{the-retained-periods-on-a-corrected-simultaneous-sequence}

Keep the same dyadic integers \(j\), spatial boxes, mesh, cover degree,
coordinate permutation, physical reference length and physical time as
in (12). For either of the positive coupling trajectories specified
below, set \[
 \begin{gathered}
 L_j=j^2,\qquad a_j=\frac1{100j},\qquad
 M_j=12j^4(2j^2+1),\qquad M_{\mathrm{cov},j}=j,\\
 \widetilde T_j=2C+j^4e^{8\pi\xi_j},\qquad
 \widetilde v_j=e^{-2\pi\widetilde T_j/j},\qquad
 \widetilde t_{c,j}=e^{-2\pi\widetilde T_j},\\
 \widetilde D_j=L_{\widetilde T_j}q_{\widetilde T_j}
                    +6m_{\widetilde T_j}^{\,2},\qquad
 \widetilde\theta_j=\frac{2\pi}{10^4j^2\widetilde D_j}.
 \end{gathered}
 \tag{173}
\] Here \(C\ge0\) is exactly the fixed period-correction bound in (4).
Restrict to the tail where \(\widetilde T_j\ge T_0\), as required for
that original cusp chart. This tail exists since
\(\widetilde T_j\ge j^4\). Every original real period in
\(P_{\widetilde T_j}\), including
\(\alpha_{\widetilde T_j},\eta_{\widetilde T_j},c_{\widetilde T_j}\), is
retained. The cover still has determinant \(-j^2\widetilde D_j\),
positive volume density \(j^2\widetilde D_j\), and its exact pullback
metric. Equation (6) applies with the same \(S_j\); the identity
\(\widetilde v_j^{\,j}=\widetilde t_{c,j}\) proves the required base
change.

The lower period bound now gives \[
 \begin{gathered}
 \widetilde T_j-C=C+j^4e^{8\pi\xi_j}
                       \ge j^4e^{8\pi\xi_j},\\
 \widetilde D_j\ge j^8e^{16\pi\xi_j},\qquad
 \widetilde\theta_j^{\,2}
 \le\frac{4\pi^2e^{-32\pi\xi_j}}{10^8j^{20}},\qquad
 \operatorname{sys}\ge\min(j,\widetilde T_j-C)=j.
 \end{gathered}
 \tag{174}
\] For the last equality \(j\ge2\) suffices. Thus the same
four-dimensional box \(|y_\mu|\le j/100\), radius \(j/50\), and a collar
embed isometrically. The original physical side length remains \(j/50\),
while the reference length \(\ell_*=1\), reference energy
\(E_*=1/\ell_*\), and time \(y_0\) are unchanged. No geometric bound
here is a quantum dispersion relation.

There is an exact map between the two local presentations. In the
original real coordinate order, let \[
 \mathcal P_j=P_{j^2}S_j,\qquad
 \widetilde{\mathcal P}_j=P_{\widetilde T_j}S_j,
 \qquad
 F_j^{\mathrm{chart}}: u\longmapsto
             \widetilde{\mathcal P}_j^{-1}\mathcal P_j u.
 \tag{175}
\] The domain and codomain are the lifts of the same embedded physical
box to the two period-coordinate charts, using its nonwrapping
trivializations. Both inverses are obtained from (5) by restoring its
recorded row permutation. Explicitly, let \(\mathscr P(x_1,x_2,x_3,x_4)
=(x_1,x_3,x_2,x_4)\); then \(Q_T=\mathscr P P_TS_j\) and
\((P_TS_j)^{-1}=Q_T^{-1}\mathscr P\), evaluated at each retained period.
For this comparison keep the old tail \(j^2\ge T_0\), \(j^2-C\ge j\)
from (12), so both local charts contain the stated physical box and
collar. On these domains, \[
 \widetilde{\mathcal P}_j F_j^{\mathrm{chart}}(u)=\mathcal P_j u,
 \qquad
 (F_j^{\mathrm{chart}})^*
  (\widetilde{\mathcal P}_j^{\mathsf T}\widetilde{\mathcal P}_j)
  =\mathcal P_j^{\mathsf T}\mathcal P_j,
 \qquad
 \det DF_j^{\mathrm{chart}}=D_j/\widetilde D_j>0 .
 \tag{176}
\] Matrix multiplication proves the first two identities, and the two
original negative determinants prove the third. The pullback acts on
one-forms by \(DF_j^{\mathsf T}\) and on their curvature by its exterior
square. The chain rule gives \(d(F_j^*A)=F_j^*(dA)\); matrix
multiplication of the coefficient functions gives
\((F_j^*A)\wedge(F_j^*A)=F_j^*(A\wedge A)\). Hence the entire local
nonlinear connection and action map is retained. The map is asserted on
the displayed local charts; no descent of this real linear matrix to an
isometry between the entire differently sized tori is needed.

In the final magnetic frames of (9), the two backgrounds are related by
the exact local formula \[
 \begin{split}
 \widetilde A_1-A_1
 &=-2\pi\left(\frac1{\widetilde D_j}-\frac1{D_j}\right)
                       H_cy_2\,dy_1,\\
 \widetilde h_1(n)h_1(n)^{-1}
 &=\exp\left[-2\pi a_j^2 n_2
       \left(\frac1{\widetilde D_j}-\frac1{D_j}\right)H_c\right],
 \qquad \widetilde h_2=\widetilde h_3=h_2=h_3=I.
 \end{split}
 \tag{177}
\] Both connections are obtained from their original bundle connection
by the explicitly proved gauge maps in Section 2, evaluated at the
respective periods. Their flat references are the identity in those same
final frames. The exponentials in (177) commute because their generator
is the retained \(H_c=i\sigma_3=-2T_3\). This proves the displayed
relation without any replacement of the original connection or its sign.

The physical link endpoints and their inverse traversals are identical
in the two boxes. Their identification therefore induces the identity
map on \(Q_{L_j}\), its Haar measure, vertex gauge action, electric
operators and all ordered Wilson words. The complete finite Hamiltonian
(1) is the same operator for the two angles, since \(a_j,L_j,g_j\) are
unchanged. Its unique actual vacuum and energy \(\psi_j,E_j\)
consequently coincide. Only the chosen background transport, and thus
the actual translated vector, changes. This establishes the exact
relation between the two presentations used in the spectral comparison.

For completeness their raw vectors also obey a direct comparison. In the
joint weight representation of Section 15, \[
 |\cos(\theta Y)-\cos(\phi Y)|
 \le\tfrac12|\theta^2-\phi^2|Y^2
 \quad(\theta,\phi\ge0).
\] Indeed the function \(s\mapsto\cos(Y\sqrt s)\) on \(s\ge0\) has
derivative of absolute value at most \(Y^2/2\), with the same bound at
zero by its limit. Integrate that derivative between \(\phi^2\) and
\(\theta^2\). Averaging, using \(\mathbb E Y^2=4\gamma/3\), and then
applying Parseval and the same orthogonal centering proves \[
 \|\widetilde\chi_j-\chi_j\|
 \le\frac23|\widetilde\theta_j^{\,2}-\theta_j^2|
                                      \|\Gamma_j\psi_j\|.
 \tag{178}
\] This raw estimate does not by itself compare the old and new
probabilities after their different squared norms are used. The relative
spectral comparison below concerns the corrected angle and its own
actual norm.

The corrected background retains its exact magnetic length \[
 \widetilde\ell_{B,j}=\sqrt{\widetilde D_j/(2\pi)},\qquad
 \widetilde E_{B,j}=\sqrt{2\pi/\widetilde D_j},\qquad
 \widetilde\ell_{B,j}\sim\widetilde T_j/\sqrt{2\pi}.
 \tag{179}
\] The last relation follows by applying (4) along (173). In particular
this length diverges while its inverse tends to zero; it has not been
identified with an excitation mass. For any fixed physical covector the
rounding map of Section 2.1 uses \(\widetilde{\mathcal P}_j\) and gives
error bounded by a constant times \(j^{-1}+\widetilde T_j^{-1}\). To see
this directly in (5), its four inverse blocks have norms \(0\), at most
\((\widetilde T_j-C)^{-1}\), \(j^{-1}\), and at most
\(j^{-1}\|A_{\widetilde T_j}\|(\widetilde T_j-C)^{-1}\), respectively.
The real correction matrix \(A_T\) is bounded on the cusp ray. The same
componentwise rounding proof therefore retains the fixed physical
momentum and energy reference.

\subsubsection{16.3 Explicit error rates on both decreasing-coupling
paths}\label{explicit-error-rates-on-both-decreasing-coupling-paths}

Substitute the exact face count and (174) in (172). Every factor is
positive, so the inequality gives \[
 \begin{split}
 \bar\eta_j
 &\le\frac{69120\pi^2}{10^8}\,
       \xi_j^{3/4}(6+64\xi_j)j^{-4}
       \frac{2j^2+1}{j^2+1}\\
 &\le A_{\mathrm{rel}}\,
       \xi_j^{3/4}(6+64\xi_j)j^{-4},\qquad
 A_{\mathrm{rel}}:=\frac{138240\pi^2}{10^8}.
 \end{split}
 \tag{180}
\] Here the intermediate multiplication is \(120\cdot4\cdot144=69120\):
\(L_j^4=j^8\), \(M_j^2=144j^8(2j^2+1)^2\), and the denominator is
\((j^2+1)(2j^2+1)\). Thus no asymptotic replacement of the face count is
required. The second line follows from \(2j^2+1<2(j^2+1)\).

First retain the logarithmic path of (88), with its entire physical
dictionary \[
 g_j^2=\frac1{\log j},\quad
 \xi_j=\frac{(\log j)^2}{4},\quad
 \kappa_j=\frac{200j}{\log j},\quad b_j=50j\log j,
 \quad 2b_jM_j=100j(\log j)M_j .
 \tag{181}
\] For the dyadic tail where \(\log j\ge2\), the hypotheses
\(\xi_j\ge1\) hold. Equations (173) and (180) give, explicitly, \[
 \widetilde T_j=2C+j^4e^{2\pi(\log j)^2},\qquad
 \bar\eta_j\le\frac{A_{\mathrm{rel}}}{2\sqrt2}
 \left[6(\log j)^{3/2}+16(\log j)^{7/2}\right]j^{-4}
 \longrightarrow0.
 \tag{182}
\] For the convergence, put \(x=\log j\). For any integer \(n>7/2\), the
exponential series gives \(e^{4x}\ge(4x)^n/n!\) for \(x>0\), and hence
\(x^{7/2}e^{-4x}\le n!4^{-n}x^{7/2-n}\to0\); the smaller power follows
by the same argument.

Second fix the original physical electric coefficient \(\kappa_*>0\) and
retain the path (91), namely \[
 g_j^2=\frac{\kappa_*}{200j},\quad
 \xi_j=\frac{10000j^2}{\kappa_*^2},\quad
 \kappa_j=\kappa_*,\quad b_j=\frac{10000j^2}{\kappa_*},
 \quad 2b_jM_j=\frac{20000j^2M_j}{\kappa_*}.
 \tag{183}
\] For the dyadic tail \(j\ge\kappa_*/100\), \(\xi_j\ge1\). Then \[
 \begin{gathered}
 \widetilde T_j=2C+j^4e^{80000\pi j^2/\kappa_*^2},\\
 \bar\eta_j\le A_{\mathrm{rel}}
 \left[\frac{6000}{\kappa_*^{3/2}}j^{-5/2}
       +\frac{640000000}{\kappa_*^{7/2}}j^{-1/2}\right]
 \longrightarrow0 .
 \end{gathered}
 \tag{184}
\] Both constants follow by inserting
\(\xi_j^{3/4}=1000j^{3/2}/\kappa_*^{3/2}\) into (180). Every finite
magnetic interaction and the full scalar Wilson term remain in (183).
The convergence is an ordinary limit of the two displayed negative
powers.

\subsubsection{16.4 Full spectral consequences with the raw amplitudes
retained}\label{full-spectral-consequences-with-the-raw-amplitudes-retained}

For either path, let \[
 \begin{gathered}
 \widetilde\chi_j=(K_{\widetilde\theta_j}
     -\langle\psi_j,K_{\widetilde\theta_j}\psi_j\rangle)\psi_j,
 \qquad \widetilde d_j=\|\widetilde\chi_j\|^2,\\
 v_{\Gamma,j}=(\Gamma_j-\langle\psi_j,\Gamma_j\psi_j\rangle)\psi_j,
 \qquad \sigma_{\Gamma,j}=\|v_{\Gamma,j}\|,\qquad
 P_j(B)=\mathbf1_B(H_j-E_j),\\
 \widetilde\nu_j(B)=
   \frac{\langle\widetilde\chi_j,P_j(B)\widetilde\chi_j\rangle}
        {\widetilde d_j},\qquad
 \zeta_{\Gamma,j}(B)=
   \frac{\langle v_{\Gamma,j},P_j(B)v_{\Gamma,j}\rangle}
        {\sigma_{\Gamma,j}^2}.
 \end{gathered}
 \tag{185}
\] The projections use the complete original physical operator. Equation
(174) gives \(0<\widetilde\theta_j<2\pi\) on the specified tail, and
Section 14 gives \(\sigma_{\Gamma,j}>0\) and \(\widetilde d_j>0\). Thus
both probabilities exist at every index considered. Equations
(180)--(184) give a further tail on which \(\bar\eta_j<1\), and the
proved raw mass and probability estimates are \[
 \begin{gathered}
 (1-\bar\eta_j)^2\frac49\widetilde\theta_j^{\,4}
       \sigma_{\Gamma,j}^{\,2}
 \le\widetilde d_j\le
 (1+\bar\eta_j)^2\frac49\widetilde\theta_j^{\,4}
       \sigma_{\Gamma,j}^{\,2},\\
 \widetilde d_j\ge
 (1-\bar\eta_j)^2\frac49\widetilde\theta_j^{\,4}s_{L_j,\xi_j}^{\,2},\\
 \sup_{B\text{ Borel}}|\widetilde\nu_j(B)-\zeta_{\Gamma,j}(B)|
       \le\min(1,\bar\eta_j)\longrightarrow0,\\
 \|\widetilde\nu_j-\zeta_{\Gamma,j}\|_{\mathrm{TV},1}
       \le2\min(1,\bar\eta_j)\longrightarrow0.
 \end{gathered}
 \tag{186}
\] For the first line the function \((1-x)^2\) decreases on \([0,1]\),
and \((1+x)^2\) increases there; substitute \(\eta_j\le\bar\eta_j<1\) in
Section 15. The second line then uses (171). The last two lines use the
same comparison on the full Hilbert space, with no energy or spin
truncation.

The corresponding raw measures satisfy, in the explicitly defined mass
norm of Section 15, \[
 \left\|\widetilde d_j\widetilde\nu_j
 -\frac49\widetilde\theta_j^{\,4}\sigma_{\Gamma,j}^{\,2}
                 \zeta_{\Gamma,j}\right\|_{\mathrm{TV},1}
 \le\bar\eta_j(2+\bar\eta_j)
       \frac49\widetilde\theta_j^{\,4}\sigma_{\Gamma,j}^{\,2}.
 \tag{187}
\] The polynomial \(x(2+x)\) increases for \(x\ge0\), so (187) follows
directly from the raw comparison in Section 15. This displays exactly
what division by the shrinking squared norms does to the two measures.

These estimates also give actual vanishing of the corrected raw norm.
From (101), \(\sigma_{\Gamma,j}\le\|\Gamma_j\psi_j\|
\le L_j^2\|H_0\psi_j\|\), and \(r_j\le5/4\), \(M_j\le36j^6\), equations
(174) and (186) imply on \(\bar\eta_j<1\) that \[
 \|\widetilde\chi_j\|
 \le\frac{384\sqrt3(5/4)^{1/4}\pi^2}{10^8}
          \xi_j^{3/4}e^{-32\pi\xi_j}j^{-10}
 \longrightarrow0 .
 \tag{188}
\] Indeed \(\|\widetilde\chi_j\|\le(4/3)\widetilde\theta_j^2
\sigma_{\Gamma,j}\), and the preceding estimates multiply to the stated
constant. For \(x\ge1\), \(x^{3/4}\le x\) and \(e^{32\pi x}\ge32\pi x\),
so \(x^{3/4}e^{-32\pi x}\le1/(32\pi)\); this proves the limit in (188)
without discarding its original amplitude.

For every fixed physical energy \(\Omega>0\), (186) proves \[
 \big|\widetilde\nu_j([0,\Omega E_*])
              -\zeta_{\Gamma,j}([0,\Omega E_*])\big|\longrightarrow0.
 \tag{189}
\] It follows that their lower and upper limiting weights on this
interval are equal, because two bounded real sequences whose difference
tends to zero have identical liminf and limsup. In particular a positive
lower limiting weight for either displayed sequence is equivalent to one
for the other; this is a proved relation, not an assumed positive
weight.

The two families are tight on the unchanged physical half-line
\([0,\infty)\) simultaneously. To prove this assertion, suppose one
family is tight and fix \(\varepsilon>0\). Its tail beyond a
sufficiently large \(R\) has mass below \(\varepsilon/2\) at every
index. Choose \(J\) so that \(\bar\eta_j<\varepsilon/2\) for \(j\ge J\).
Equation (186) bounds the other family's tail by \(\varepsilon\) for
these indices. For the finitely many earlier probability measures,
continuity from above gives tails tending to zero as \(R\to\infty\);
increasing \(R\) handles all of them. The reverse implication uses the
same symmetric estimate. Moreover, along any subsequence, weak
convergence of either probability to a probability on \([0,\infty)\)
implies weak convergence of the other to the same limit: apply the
mass-norm estimate to each bounded continuous test function divided by
its supremum norm.

For every \(t\ge0\), the function \(e^{-t\omega}\) is a bounded spectral
function of the same \(H_j-E_j\). Equations (186)--(187) therefore also
compare its probability integrals and its raw matrix elements, with
errors uniform in this physical time parameter. This argument does not
apply to unbounded energy powers, and it does not determine the spectral
probability of \(v_{\Gamma,j}\). Establishing its finite-energy mass and
dynamical continuum is the remaining calculation. The comparison (186)
already proves the exact link to the original native magnetic states on
the corrected cusp without assuming either of those conclusions.

\subsection{17. Strong graph convergence of the actual weighted electric
state}\label{strong-graph-convergence-of-the-actual-weighted-electric-state}

Throughout this section fix the original integer \(L\ge2\) and physical
spacing \(a>0\), before sending the strictly positive \(g\) to zero. The
complete companion proofs are retained in
sources/weak\_coupling\_state\_inputs: the finite-box physical
eigenvalue proof, the actual vacuum and spectral-projection proof, and
the electric-state graph proof. The coordinate definitions and
derivative passage needed here are given explicitly below. The fixed-box
constants are not asserted uniform in volume.

\subsubsection{17.1 Original operator, exact quotient, and graph
domain}\label{original-operator-exact-quotient-and-graph-domain}

Let \(\mathsf E,\mathsf P\) be every positive edge and every elementary
oriented face in the box \(\{-L,\ldots,L\}^3\). Write

\[
 N=3(2L)(2L+1)^2,\qquad M=12L^2(2L+1).
\]

The original product Haar probability space is \(Q_L=SU(2)^N\). With
\(T_\alpha=-i\sigma_\alpha/2\), define

\[
 X_{e,\alpha}F(U)
 =\left.\frac{d}{dt}F(\ldots,e^{tT_\alpha}U_e,\ldots)\right|_{t=0},
 \quad E_e=-\sum_{\alpha=1}^3X_{e,\alpha}^2,\quad H_0=\sum_eE_e.
\]

For each face use its original oriented word and write

\[
 W_p=\operatorname{tr}U_p,\qquad
 W=\sum_{p\in\mathsf P}(2-W_p).
\]

The full physical operator and its exact constant multiple are

\[
 H_g=\frac{2g^2}{a}H_0+\frac{W}{2g^2a},\qquad
 P_g=\frac a2H_g=g^2H_0+\frac{W}{4g^2}.
\tag{190}
\]

Thus the scalar term \(2bM\), \(b=1/(2g^2a)\), remains in \(H_g\). Let
\(\psi_g>0\) be its actual unit vacuum, with physical energy \(E_g\).
Put

\[
 P_g\psi_g=\epsilon_g\psi_g,\qquad
 \epsilon_g=\frac a2E_g,\qquad
 0\le\epsilon_g\le\epsilon_*:=\frac32\sqrt{NM}.
\tag{191}
\]

The last inequality is the already proved actual-vacuum variational
bound; it holds with the full Wilson potential. The vacuum is smooth on
the finite compact product, invariant under every vertex gauge
transformation, and its norm is exactly one.

Retain the weighted operator

\[
 w_e=
 \begin{cases}
 n_2^2,&e=(n,1),\\
 0,&e\text{ is in direction }2\text{ or }3,
 \end{cases}
 \qquad
 \Gamma=\sum_ew_eE_e.
\tag{192}
\]

All original edges, including those of zero weight, remain in \(H_0,W\)
and \(\psi_g\). The commuting nonnegative operators \(E_e\) have the
complete product matrix-coefficient decomposition with eigenvalues
\(q_e(q_e+1)\). On every joint block,

\[
 0\le\sum_ew_eq_e(q_e+1)\le L^2\sum_eq_e(q_e+1).
\]

Squaring and applying Parseval to the entire Hilbert sum proves

\[
 \|\Gamma F\|\le L^2\|H_0F\|,\qquad F\in D(H_0).
\tag{193}
\]

There is no spin truncation in (193). Each \(E_e\) commutes with the
vertex gauge action, so the same statement holds on the exact invariant
subspaces.

The maximal tree is rooted at \(o=(-L,-L,-L)\); the parent of a nonroot
vertex decreases the first coordinate, in order \(1,2,3\), that exceeds
\(-L\). Denote its ordered tree holonomies by \(t_v\), and its chords by
\(\mathcal C\), of cardinality

\[
 r=N-(2L+1)^3+1=2(2L)^3+3(2L)^2.
\]

The exact chord coordinates and inverse are

\[
 Z_c=t_{s(c)}U_ct_{t(c)}^{-1},\qquad
 U_c=t_{s(c)}^{-1}Z_ct_{t(c)}.
\tag{194}
\]

Tree links are retained as independent variables in the inverse. For
fixed tree links, each chord change is a left and right Haar
translation; Fubini proves exact product Haar preservation. Based gauge
averaging eliminates the tree variables. The remaining physical
condition is simultaneous conjugation of all chords by the root gauge
element.

Let \(\iota F(U)=F(Z(U))\). The map \(\iota\) is an isometry from chord
Haar \(L^2\) onto the based-invariant subspace, by the exact
product-Haar and compact-gauge averaging proof in the retained
finite-box companion. We identify a based-invariant function with its
chord representative only through this isometry. In particular no value
of an \(L^2\) class on a Haar-null tree slice is used.

For the path-incidence numbers \(\varepsilon(v,e)\in\{0,1\}\), the exact
chord fields are

\[
 \mathcal X_{e,\alpha}=
 \begin{cases}
 L_{e,\alpha},&e\in\mathcal C,\\
 \displaystyle\sum_{c\in\mathcal C}
 [\varepsilon(s(c),e)L_{c,\alpha}
       -\varepsilon(t(c),e)R_{c,\alpha}],&e\text{ in the tree}.
 \end{cases}
\tag{195}
\]

Here \(L_{c,\alpha}\) and \(R_{c,\alpha}\) differentiate left and right
multiplication of \(Z_c\) by \(e^{tT_\alpha}\). They are divergence-free
for product Haar. Differentiation of (194), followed by invariance of
the sum over the three generator components under the adjoint rotations,
intertwines the original edge forms with

\[
 H_0^{\mathrm q}=-\sum_{e,\alpha}\mathcal X_{e,\alpha}^2,\qquad
 \Gamma^{\mathrm q}=-\sum_{e,\alpha}w_e\mathcal X_{e,\alpha}^2.
\tag{196}
\]

The same argument works with the weights because each weight multiplies
the complete three-component sum at that edge. Integration against the
exact product Haar measure makes (196) the represented operators of
these forms. Chord fields alone span the tangent space, hence the full
kinetic operator is elliptic, with domain \(H^2\) on the compact chord
manifold. All compactly supported comparison vectors below belong to
this domain. Formula (193) transfers through \(\iota\) without changing
its constant.

\subsubsection{17.2 Exact Haar/logarithm/dilation map and local
operators}\label{exact-haarlogarithmdilation-map-and-local-operators}

Outside the Haar-null set where some chord equals \(-I\), write uniquely

\[
 Z_c=\exp(y_c^\alpha T_\alpha),\qquad |y_c|<2\pi.
\]

The exact Haar density and rescaled domain are

\[
 \mathcal J(y)=(16\pi^2)^{-r}
       \prod_c\left[\frac{\sin(|y_c|/2)}{|y_c|/2}\right]^2,
 \quad \Omega_g=\{x:|x_c|<2\pi/g\ \forall c\}.
\]

The isometry into Euclidean \(L^2\) is

\[
 (\mathcal B_gF)(x)=
 \begin{cases}
 g^{3r/2}\mathcal J(gx)^{1/2}F(\exp(gx)),&x\in\Omega_g,\\
 0,&x\notin\Omega_g.
 \end{cases}
\tag{197}
\]

Its image is exactly the closed subspace supported in \(\Omega_g\); its
adjoint is zero on the orthogonal complement of that subspace. Its
angular rotation action is the exact residual physical action. The power
of \(g\), the Haar factor, and the zero extension are all retained.

Write \(\mathcal X_{e,\alpha}=a_{e,\alpha}^i(y)\partial_{y_i}\) on the
chart, with \(i=(c,\beta)\). These are the full smooth coefficients of
(195). On every fixed Euclidean ball contained in \(\Omega_g\), define
the exact first-order operators

\[
 \mathcal D_{g,e,\alpha}
 =a_{e,\alpha}^i(gx)
   \left[\partial_{x_i}
        -\frac g2(\partial_{y_i}\log\mathcal J)(gx)\right].
\tag{198}
\]

Direct differentiation of the inverse of (197) proves

\[
 \mathcal B_g(g\mathcal X_{e,\alpha})\mathcal B_g^*
 =\mathcal D_{g,e,\alpha}.
\]

The identity is local on compactly supported chart functions and on the
restriction of the smooth vacuum. It makes no assertion of regularity of
the zero extension across \(\partial\Omega_g\).

The exact local second-order operators are

\[
 K_g=-\sum_{e,\alpha}\mathcal D_{g,e,\alpha}^2,\qquad
 G_g=-\sum_{e,\alpha}w_e\mathcal D_{g,e,\alpha}^2,\qquad
 V_g(x)=\frac{W(gx)}{4g^2}.
\tag{199}
\]

In particular \(K_g\) is the local conjugate of \(g^2H_0^{\mathrm q}\),
and \(G_g\) is the local conjugate of \(g^2\Gamma^{\mathrm q}\).

Retain the additive tree map \(T:\mathbb R^{\mathsf E}\to
\mathbb R^{\mathcal C}\),

\[
 (Tz)_c=p_{s(c)}(z)+z_c-p_{t(c)}(z),
 \quad G=TT^*,\quad C=d_1j,\quad
 G_w=T\operatorname{diag}(w_e)T^*.
\tag{200}
\]

Here \(p_v\) is the original additive tree-path integral and \(j\)
inserts chord entries with zero tree entries. The exact tangent map of
(194) is \(T\) in each color component. Consequently the coefficient
limits on every fixed ball are

\[
 K_g\longrightarrow K_0
 =-\sum_{\alpha,c,d}G_{cd}
       \partial_{x_c^\alpha}\partial_{x_d^\alpha},
 \qquad
 G_g\longrightarrow\Gamma_{\mathrm{osc}}
 =-\sum_{\alpha,c,d}(G_w)_{cd}
       \partial_{x_c^\alpha}\partial_{x_d^\alpha}.
\tag{201}
\]

Convergence here means all coefficients required in these second-order
differential expressions converge uniformly on the fixed ball. It
follows by differentiating the smooth coefficients at \(gx\), with the
factor \(g\) on their spatial derivatives, and retaining the density
term in (198). In particular that density term tends to zero; it was not
omitted before taking the limit.

The original full word expansion, with its proved Taylor remainder, is

\[
 W(y)=\frac14\sum_\alpha\|Cy^\alpha\|^2+O_L(|y|^3).
\]

Its smooth differentiated version gives

\[
 V_g\longrightarrow V_0=\frac1{16}\sum_\alpha\|Cx^\alpha\|^2
\tag{202}
\]

uniformly with every fixed number of derivatives on each fixed ball.
This is a statement about the exact potential there, not replacement of
the potential outside the ball.

The complete retained companion proofs establish

\[
 f_g:=\mathcal B_g\psi_g\longrightarrow\Phi_0
       \quad\hbox{strongly in }L^2(\mathbb R^{3r}),\qquad
 \epsilon_g\longrightarrow
       \epsilon_0=\frac34\sum_{\nu=1}^r\sigma_\nu.
\tag{203}
\]

Here \(\Phi_0\) is the actual limiting positive direction selected by
the positive finite-coupling vacua:

\[
 \Phi_0(x)=(\det S)^{3/4}(4\pi)^{-3r/4}
      \exp\left[-\frac18\sum_\alpha(x^\alpha)^TSx^\alpha\right],
\]

\[
 S=G^{-1/2}(G^{1/2}C^*CG^{1/2})^{1/2}G^{-1/2}.
\tag{204}
\]

The positive matrices and every scalar are as in the inputs. The source
proves \((K_0+V_0)\Phi_0=\epsilon_0\Phi_0\) and unit norm, using the
actual vacuum equation and compactness. We use (203), not a trial-state
identification, as the starting vector convergence.

\subsubsection{17.3. Exact original operator and potential-gradient
inequality}\label{exact-original-operator-and-potential-gradient-inequality}

Use the finite-box definitions and exact coordinate maps of the
companion manuscripts \emph{The original finite-box SU(2) Hamiltonian at
small positive coupling} and \emph{Actual weak-coupling vacuum,
observable and spectral maps}. In particular \[
 \begin{gathered}
 H=\kappa H_0+bW,\quad H_0=\sum_eE_e,\qquad
 \kappa=2g^2/a,\quad b=1/(2g^2a),\\
 W=\sum_p w_p,\quad w_p=2-\operatorname{tr}U_p,\qquad
 E_e=-\sum_{\alpha=1}^3X_{e,\alpha}^2.
 \end{gathered}                                            \tag{205}
\] The anti-Hermitian generators are \(T_\alpha=-i\sigma_\alpha/2\). All
positive contained links and all contained faces are included. The
strictly positive unit vacuum is \(\psi\), with actual energy
\(\mathcal E\), so \(H\psi=\mathcal E\psi\). The scalar \(2b|P|\) is
part of bW and is not removed in this equation.

For any face p and any of its four links e, \[
 \sum_\alpha|X_{e,\alpha}w_p|^2=w_p-w_p^2/4\le w_p.       \tag{206}
\] To check every factor, hold the other links fixed. Its oriented face
word is a Haar translate of U\_e or U\_e\^{}\{-1\}; conjugations rotate
the three generators orthogonally. Write the resulting SU(2) element as
\(u_0I+i\boldsymbol u\cdot\boldsymbol\sigma\), with
\(u_0^2+|\boldsymbol u|^2=1\). Its trace is 2u\_0, and
\(\operatorname{tr}(T_\alpha U)=u_\alpha\). The three squared
derivatives sum to \(|\boldsymbol u|^2\). Since \(w_p=2-2u_0\), this is
exactly (206). The inverse-link case has the same sum by inversion and
bi-invariance, not a changed orientation.

Let \(r_e\) be the actual number of faces incident to e; on this open
box \(r_e\le4\), including every boundary edge. Cauchy--Schwarz in that
finite incident set, followed by summation in e, gives \[
 \sum_{e,\alpha}|X_{e,\alpha}W|^2
  \le\sum_e r_e\sum_{p\ni e}\sum_\alpha|X_{e,\alpha}w_p|^2
  \le16W.                                                   \tag{207}
\] Thus the constant 16 is independent of the box; W itself still
contains every face. In particular this is not a local replacement
potential.

\subsubsection{17.4. All potential moments from the actual vacuum
equation}\label{all-potential-moments-from-the-actual-vacuum-equation}

Write \(M_k=\int W^k\psi^2\,dU\), with \(M_0=1\). For every integer
k\textgreater=1, \[
 \boxed{bM_{k+1}\le\mathcal E M_k+4\kappa k^2M_{k-1},
 \qquad M_1\le\mathcal E/b.}                              \tag{208}
\] Here is a proof that retains the derivative terms. The exact
ground-state identity for a real smooth f is \[
 q_H[f\psi]-\mathcal E\|f\psi\|^2
   =\kappa\int\psi^2\sum_{e,\alpha}|X_{e,\alpha}f|^2dU.
                                                               \tag{209}
\] It follows by differentiating the product twice in (205), integrating
with the original Haar measure, and substituting the full
eigen-equation. The kinetic term of \(q_H[f\psi]\) is nonnegative. For
k\textgreater=2 take \(f=W^{k/2}\), or its smooth positive
regularization if needed. Equation (207) bounds the right side by
\(4\kappa k^2M_{k-1}\); the left side is at least
\(bM_{k+1}-\mathcal E M_k\), proving (208). For k=1 use
\(f=(W+\epsilon)^{1/2}\). Its squared gradient is at most
\(4W/(W+\epsilon)\le4\). In (209) the lower bound is
\(b(M_2+\epsilon M_1)-\mathcal E(M_1+\epsilon)\). Letting epsilon
decrease to zero proves the same formula for k=1. For higher odd k the
identical regularization and dominated convergence on the compact group
justify all powers. The bound for M1 follows directly from \(bW\le H\)
in quadratic forms on \(\psi\).

The proved variational bound (99) supplies the explicit constant
\(K=3\sqrt{NM}/a\), with \(\mathcal E\le K\) at every positive \(g\).
The same energy is denoted \(E_g\) above; here \(\mathcal E=E_g\). This
constant is retained with its full box dependence. Substituting the
original coefficients into (208) gives the exact scaled recursion \[
 \frac{M_{k+1}}{g^{2k+2}}
 \le 2aK\frac{M_k}{g^{2k}}+
                 16k^2\frac{M_{k-1}}{g^{2k-2}}.             \tag{210}
\] Define explicitly \(B_0=1\), \(B_1=2aK\), and
\(B_{k+1}=2aK B_k+16k^2 B_{k-1}\). Induction proves
\(M_k\le B_k g^{2k}\) for every fixed k. The recursion keeps the full
vacuum energy estimate and every original g,a factor. For example,
\(B_2=(2aK)^2+16\) and \(B_3=(2aK)^3+80(2aK)\). Only fixed-box
boundedness of K is used.

\subsubsection{17.5. Potential-weighted strong convergence and the total
electric
graph}\label{potential-weighted-strong-convergence-and-the-total-electric-graph}

Retain the exact full-measure map \(\mathcal B_g\), matrices T,G,C, and
Gaussian \(\Phi_0\) from the preceding exact coordinate construction.
Put \[
 V_2(x)=\tfrac14\sum_\alpha\|Cx^\alpha\|^2,
 \qquad D=-\sum_{\alpha,c,d}G_{cd}
                         \partial_{x_c^\alpha}\partial_{x_d^\alpha}.
                                                               \tag{211}
\] The retained actual-state proof establishes
\(\mathcal B_g\psi\to\Phi_0\) strongly. We strengthen it to \[
 \mathcal B_g\bigl((W/g^2)\psi\bigr)\to V_2\Phi_0,
 \qquad
 \mathcal B_g(g^2H_0\psi)\to D\Phi_0
                \quad\hbox{strongly in }L^2.                \tag{212}
\] The first convergence is local by the exact Taylor limit
\(W(gx)/g^2\to V_2\) and the strong vacuum convergence. To prove that no
weighted tail is lost, choose the same fixed radius \(\rho\) as in the
companion: \(W(y)\ge c|y|^2\) inside it and \(W\ge c_\rho\) outside it.
On the subset \(|x|>R\) with \(|gx|<\rho\), the variable \(W/g^2\) is at
least \(cR^2\). Outside the \(\rho\) neighborhood it is at least
\(c_\rho/g^2\). Therefore (210) at \(k=3\) gives \[
 \int_{|x|>R}|\mathcal B_g((W/g^2)\psi)|^2dx
  \le \frac{B_3}{cR^2}+\frac{g^2B_3}{c_\rho}.             \tag{213}
\] The limiting polynomial times Gaussian has an integrable squared
tail. Let R tend to infinity after the local convergence to obtain the
first part of (212).

For the second, the actual eigen-equation gives the exact identity \[
 g^2H_0\psi=\frac{a\mathcal E}{2}\psi-\frac{W}{4g^2}\psi.
                                                               \tag{214}
\] The first limit in (212) and \(\mathcal E\to\mu_0\) show that its
image tends to \((a\mu_0/2-V_2/4)\Phi_0\). The exact oscillator equation
\([(2/a)D+V_2/(2a)]\Phi_0=\mu_0\Phi_0\) identifies this as \(D\Phi_0\).
This proves (212) from full potential moments, not merely from
convergence in L2 of the vacuum.

\subsubsection{17.6. Every retained electric weight and its graph
limit}\label{every-retained-electric-weight-and-its-graph-limit}

For any fixed real edge weights w define \[
 \Gamma_w=\sum_e w_eE_e,\qquad
 G_w=T\operatorname{diag}(w)T^*,\qquad
 D_w=-\sum_{\alpha,c,d}(G_w)_{cd}
                        \partial_{x_c^\alpha}\partial_{x_d^\alpha}.
                                                               \tag{215}
\] No positivity of w is needed. The electric operators on distinct
links strongly commute by their product-group spectral decomposition.
Their eigenvalues are nonnegative, so on \(\operatorname{Dom}H_0\) \[
 \|\Gamma_w F\|\le\|w\|_\infty\|H_0F\|.                \tag{216}
\] Their individual Casimir spectra and all signs of w remain in this
inequality. Restriction to the physical space is valid because every
E\_e commutes with the full vertex gauge projection.

We claim the actual graph limit \[
 \boxed{\mathcal B_g(g^2\Gamma_w\psi)\longrightarrow D_w\Phi_0.}
                                                               \tag{217}
\] For a smooth invariant compactly supported u(x), its reciprocal map
\(\mathcal B_g^*u\) is a smooth physical function supported in the
logarithm chart for all sufficiently small g. Differentiating the exact
tree fields and density gives \[
 \mathcal B_g(g^2H_0\mathcal B_g^*u)\to Du,
 \qquad
 \mathcal B_g(g^2\Gamma_w\mathcal B_g^*u)\to D_wu          \tag{218}
\] strongly, since every coefficient and derivative converges uniformly
on the fixed compact support. In detail the principal coefficients of
gX\_e are its derivative columns of T at y=0; derivatives of these
coefficients and of the Haar density each carry an extra factor g.
Summing their exact squared fields gives G and Gw. The lower-order terms
tend to zero on this compact support, not by deletion from the original
operator.

Choose smooth invariant radial cutoffs \(u_R=\chi_R\Phi_0\). Their
Gaussian decay gives \(u_R\to\Phi_0\) in the graph norms of both
constant-coefficient operators D and Dw. The first limit in (218),
combined with the total electric graph limit (212), shows \[
 \lim_{g\to0}\|g^2H_0(\psi-\mathcal B_g^*u_R)\|
        =\|D(\Phi_0-u_R)\|.                                \tag{219}
\] Apply (216) to the difference, then use the second limit in (218).
The limsup error in (217) is at most \(\|w\|_\infty\|D(\Phi_0-u_R)\|+
\|D_w(u_R-\Phi_0)\|\), which tends to zero as R grows. This proves
(217). Every limit here fixes L,a,w before g tends to zero.

Inner products with the strongly converging unit vacua now prove \[
 g^2\langle\Gamma_w\rangle_\psi\to
          \langle\Phi_0,D_w\Phi_0\rangle,
\quad
 \mathcal B_g\bigl(g^2(\Gamma_w-\langle\Gamma_w\rangle)\psi\bigr)
     \to V_w:=(D_w-\langle D_w\rangle)\Phi_0.              \tag{220}
\] This is the raw actual covariance-state map whose derivative control
was not supplied by the strong vacuum convergence (203) alone.

\subsubsection{17.7 Full physical spectral projections and the
finite-box
scope}\label{full-physical-spectral-projections-and-the-finite-box-scope}

The actual-state companion uses the same \(\mathcal B_g\), with its
adjoint zero outside \(\Omega_g\), and proves convergence in operator
norm of every finite physical spectral projection whose endpoints avoid
the oscillator spectrum. Its proof can be recalled directly from the
maps above. A family of actual eigenvectors of bounded absolute energy
has rescaled potential tails bounded by a constant times \(R^{-2}+g^2\):
the full potential is bounded below by \(c|y|^2\) in a fixed small
identity chart, and by a positive constant outside it. On each fixed
rescaled ball the full elliptic form and the retained density bound the
local \(H^1\) norm. Rellich compactness there and the tail bound give
strong global \(L^2\) subsequences. For finitely many vectors their
inner products persist.

Test an actual eigen-equation against the inverse image of a smooth
compactly supported physical function. All local coefficients and their
required derivatives converge under the exact map above; weak local
\(H^1\) convergence and strong local \(L^2\) convergence therefore
identify every subsequential limit with the corresponding oscillator
eigenvector. The eigenvalue convergence, with multiplicity, is proved by
the full localization and min--max argument in the retained finite-box
source. It gives equal finite ranks below every chosen cutoff for
sufficiently small \(g\). Thus a converging orthonormal eigenbasis of
the actual finite spectral range gives a complete orthonormal basis of
the limiting spectral range. The finite sums of their rank-one
projections converge in operator norm. Every subsequence has the same
projection limit, proving the full-family conclusion. Subtracting the
actual convergent ground energy yields the same assertion for excitation
projections.

In particular the fixed-box physical gap tends to \[
 \Delta_{\mathrm{osc}}(L,a)=
     \frac{4\sqrt2}{a}\sin\frac{\pi}{4L+2}.
\] The full frequency calculation, including the open boundary
conditions, is also spelled out in (236) below. Its smallest positive
frequency is \(\sqrt8\sin[\pi/(4L+2)]\). Residual simultaneous color
rotation has no invariant one-quantum vector. Any physical excitation
therefore has at least two quanta, and an invariant contraction of two
creators in a lowest mode achieves exactly twice that frequency divided
by \(a\). This proves the displayed comparison gap from the original
residual gauge constraint. Its convergence as an actual finite-box gap
uses the just-stated full physical eigenvalue theorem. Neither that
convergence nor the strong graph limit above specifies an error uniform
in a simultaneous sequence \(L\to\infty\).

\subsection{18. Complete covariance spectral weights and a weak-coupling
spatial
diagonal}\label{complete-covariance-spectral-weights-and-a-weak-coupling-spatial-diagonal}

\subsubsection{18.1 The original coordinate weights in the transverse
kinetic
matrix}\label{the-original-coordinate-weights-in-the-transverse-kinetic-matrix}

Use the exact maximal-tree maps and the actual-vacuum convergence of
Section 17. We first calculate their limiting covariance state at each
fixed \(L,a\), retaining every original edge weight. Write \[
 \begin{gathered}
 m=2L,\qquad r=2m^3+3m^2=16L^3+12L^2,\\
 w_{(n,1)}=n_2^2,\quad w_{(n,2)}=w_{(n,3)}=0,\qquad
 \mathsf W=\operatorname{diag}(w_e).
 \end{gathered}
 \tag{221}
\] Here \(\mathsf W\) is a matrix on real edge cochains, distinct from
the full nonlinear Wilson potential \(\mathcal V=\sum_p(2-W_p)\). Let
\(T\) be the exact additive tree-to-chord map, \(G=TT^{\mathsf T}\),
\(C=d_1j_{\mathcal C}\), and \[
 O^{\mathsf T}G^{1/2}C^{\mathsf T}CG^{1/2}O=\Lambda^2,
 \qquad \Lambda=\operatorname{diag}(\sigma_1,\ldots,\sigma_r).
\] The chord insertion \(j_{\mathcal C}\) has zero entries on the
original tree and retains every chord; it is unrelated to the regulator
index used later. Each \(\sigma_\nu>0\) is the exact positive open-box
curl frequency. Define \[
 \mathscr R=T^{\mathsf T}G^{-1/2}O:
       \mathbb R^r\longrightarrow \ker d_0^{\mathsf T},
 \qquad D=\mathscr R^{\mathsf T}\mathsf W\mathscr R.
 \tag{222}
\] The map \(\mathscr R\) is an isometry onto the indicated original
transverse subspace, by \(TT^{\mathsf T}=G\) and
\(\ker T=\operatorname{im}d_0\). Thus \(0\le D\le L^2I_r\). No
independence of the entries of \(D\) is asserted. The original
coordinate origin remains in \(n_2^2\).

Under the exact logarithm and Haar map followed by the constant-Jacobian
change \(x=G^{1/2}Oz\) in each color, the limit of \(g^2\Gamma\) on the
actual vacuum is \[
 \Gamma_{\mathrm{osc}}=
   -\sum_{\alpha=1}^3\sum_{\nu,\mu=1}^r
       D_{\nu\mu}\partial_{z_\nu^\alpha}\partial_{z_\mu^\alpha}.
 \tag{223}
\] Indeed the principal coefficient of the original weighted electric
operator at the identity is \(T\mathsf W T^{\mathsf T}\) in chord
coordinates. Conjugating that matrix by \(O^{\mathsf T}G^{-1/2}\) gives
exactly (222). The strong graph convergence proved in Section 17 is what
justifies applying this principal coefficient to the actual vacuum
limit; a coefficient calculation alone would not justify it.

The comparison Hamiltonian and its unit ground function in these
coordinates are \[
 H_{\mathrm{osc}}=\sum_{\nu=1}^r
   \left[-\frac2a\Delta_{z_\nu}
                 +\frac{\sigma_\nu^2}{8a}|z_\nu|^2\right],
 \qquad
 \Phi_0(z)=\prod_{\nu=1}^r
       \left(\frac{\sigma_\nu}{4\pi}\right)^{3/4}
       e^{-\sigma_\nu|z_\nu|^2/8},
 \qquad
 \mu_0=\frac3{2a}\sum_\nu\sigma_\nu .
 \tag{224}
\] The unitary change from an \(x\)-function to a \(z\)-function is
\(f(x)\mapsto(\det G)^{3/4}f(G^{1/2}Oz)\). This determinant and the
nonconstant Haar density in Section 17 are both retained. Residual gauge
invariance is simultaneous adjoint rotation of the color components of
all \(z_\nu\); it does not remove separate modes.

\subsubsection{18.2 The entire centered pair state and its raw spectral
measure}\label{the-entire-centered-pair-state-and-its-raw-spectral-measure}

For one real coordinate define \[
 a_{\nu\alpha}=
   \sqrt{\frac2{\sigma_\nu}}\,\partial_{z_\nu^\alpha}
       +\sqrt{\frac{\sigma_\nu}{8}}\,z_\nu^\alpha,\qquad
 a_{\nu\alpha}^{\dagger}=
   -\sqrt{\frac2{\sigma_\nu}}\,\partial_{z_\nu^\alpha}
       +\sqrt{\frac{\sigma_\nu}{8}}\,z_\nu^\alpha .
 \tag{225}
\] Integration by parts on polynomials times the Gaussian proves the
adjoint relation. The identity \([\partial_z,z]=1\) proves the
commutator \([a_{\nu\alpha},a_{\mu\beta}^{\dagger}]
=\delta_{\nu\mu}\delta_{\alpha\beta}\); all other such commutators
vanish. Also \(a_{\nu\alpha}\Phi_0=0\), and \[
 H_{\mathrm{osc}}-\mu_0
   =\sum_{\nu,\alpha}\frac{\sigma_\nu}{a}
                      a_{\nu\alpha}^{\dagger}a_{\nu\alpha}.
 \tag{226}
\] These identities follow by direct multiplication of (225); they fix
the physical coefficient and all excitation energies used below.

Since \(\partial_{z_\nu^\alpha}
=\sqrt{\sigma_\nu/8}(a_{\nu\alpha}-a_{\nu\alpha}^{\dagger})\), applying
(223) to \(\Phi_0\) gives \[
 \gamma_{\mathrm{osc}}=
 \langle\Phi_0,\Gamma_{\mathrm{osc}}\Phi_0\rangle
       =\frac38\operatorname{tr}(\Lambda D),\qquad
 v_{\mathrm{osc}}
   :=(\Gamma_{\mathrm{osc}}-\gamma_{\mathrm{osc}})\Phi_0
   =-\frac18\sum_{\alpha,\nu,\mu}
       \sqrt{\sigma_\nu\sigma_\mu}\,D_{\nu\mu}
       a_{\nu\alpha}^{\dagger}a_{\mu\alpha}^{\dagger}\Phi_0.
 \tag{227}
\] For verification, expand
\((a_\nu-a_\nu^\dagger)(a_\mu-a_\mu^\dagger)\Phi_0
=-\delta_{\nu\mu}\Phi_0+
a_\nu^\dagger a_\mu^\dagger\Phi_0\) at each color. Multiplication by the
minus sign in (223) gives the positive scalar in (227) and its negative
pair term. Each color contraction is invariant under simultaneous
adjoint rotation, so every displayed state is in the full physical
subspace.

For \(\nu=\mu\), the three vectors \((a_{\nu\alpha}^{\dagger})^2\Phi_0\)
are orthogonal, each of squared norm two. This follows by commuting the
two annihilators across the two creators, using (225), and then
annihilating the vacuum. Their sum has squared norm six. For
\(\nu<\mu\), the three vectors
\(a_{\nu\alpha}^{\dagger}a_{\mu\alpha}^{\dagger}\Phi_0\) are orthogonal
unit vectors; their sum has squared norm three. Different unordered
spatial pairs are orthogonal because at least one coordinate occupation
differs. In (227) the two ordered occurrences of a distinct pair give
the coefficient \(-\sqrt{\sigma_\nu\sigma_\mu}D_{\nu\mu}/4\).
Consequently its complete raw excitation measure is \[
 \begin{split}
 m^{\Gamma}_{L,a}
 &:=\langle v_{\mathrm{osc}},
       \mathbf1_{(\cdot)}(H_{\mathrm{osc}}-\mu_0)v_{\mathrm{osc}}\rangle\\
 &=\frac3{32}\sum_{\nu=1}^r
        \sigma_\nu^2D_{\nu\nu}^2\,\delta_{\,2\sigma_\nu/a}\\
 &\quad+\frac3{16}\sum_{\nu<\mu}
        \sigma_\nu\sigma_\mu D_{\nu\mu}^2\,
                      \delta_{\,(\sigma_\nu+\sigma_\mu)/a}.
 \end{split}
 \tag{228}
\] When different pairs have the same energy their displayed weights
add; there are no cross terms because the corresponding vectors are
already orthogonal. The exact raw mass is \[
 C_L^{\Gamma}:=\|v_{\mathrm{osc}}\|^2
       =\frac3{32}\operatorname{tr}(\Lambda D\Lambda D).
 \tag{229}
\] Thus the normalized probability, using its displayed nonzero raw
mass, is \(\rho_{L,a}^{\Gamma}=m_{L,a}^{\Gamma}/C_L^\Gamma\). Positivity
of this mass will be proved quantitatively below.

These formulas are limits of the actual electric covariance measures.
Let \[
 v_{\Gamma,g}=(\Gamma-\langle\psi_g,\Gamma\psi_g\rangle)\psi_g,
 \qquad
 m_{L,a,g}^{\Gamma}(B)
 =g^4\langle v_{\Gamma,g},
           \mathbf1_B(H_g-E_g)v_{\Gamma,g}\rangle.
 \tag{230}
\] Section 17 proves strong convergence of the images of
\(g^2v_{\Gamma,g}\) to (227), including convergence of its squared norm.
The same full finite-projection argument proves \[
 m_{L,a,g}^{\Gamma}\Longrightarrow m_{L,a}^{\Gamma},\qquad
 g^4\|v_{\Gamma,g}\|^2\longrightarrow C_L^\Gamma,\qquad
 \zeta_{\Gamma,L,a,g}\Longrightarrow\rho_{L,a}^{\Gamma}
 \quad(g\downarrow0;\ L,a\ \hbox{fixed}).
 \tag{231}
\] Here weak convergence means convergence against every bounded
continuous function on the unchanged physical half-line. To give the
complete passage, choose an energy cutoff avoiding the discrete
oscillator spectrum. Norm convergence of its full physical spectral
projection and strong convergence of the vectors give convergence of the
mass inside that cutoff. Total mass converges by the strong vector
limit. Taking increasing cutoffs makes the remaining oscillator mass
vanish by completeness, and hence also bounds the actual tail uniformly
for small \(g\). Inside a cutoff there are finitely many eigenvalue
clusters; their eigenvalues and finite spectral projections converge, so
uniform continuity of a bounded continuous test function on that compact
interval gives convergence of its integrals. The tail is controlled by
its supremum norm. Finally divide by the positive masses tending to
\(C_L^\Gamma>0\). This proves the probability assertion with all
original raw factors retained. In particular no norm rescaling of the
original physical vector has been silently made.

\subsubsection{18.3 A uniform lower raw mass and the fraction at finite
physical
energy}\label{a-uniform-lower-raw-mass-and-the-fraction-at-finite-physical-energy}

Let \(r_e\) again be the number of incident elementary faces. The
original oriented curl matrix has diagonal
\((d_1^{\mathsf T}d_1)_{ee}=r_e\), because each incidence coefficient is
\(1\) or \(-1\). It vanishes on \(\operatorname{im}d_0\), and its
restriction to the transverse space is exactly
\(\mathscr R\Lambda^2\mathscr R^{\mathsf T}\). It follows that \[
 \operatorname{tr}(\Lambda^2D)=\sum_e w_e r_e
   =:\mathcal A_L
   =\frac43 L^2(2L+1)(4L^2+2L+1).
 \tag{232}
\] The trace identity follows by cyclically moving the rectangular
isometry and using that \(d_1^{\mathsf T}d_1\) is zero on its orthogonal
complement. For the last count, put \(r(t)=1\) at \(t=\pm L\) and
\(r(t)=2\) otherwise. Then \(r_{(n,1)}=r(n_2)+r(n_3)\), \[
 \sum_{t=-L}^L t^2=\frac{L(L+1)(2L+1)}3,\quad
 \sum_{t=-L}^L r(t)=4L,\quad
 \sum_{t=-L}^L r(t)t^2=\frac23L(2L^2+1).
\] The last equality is twice the first sum minus \(2L^2\). Sum over
\(n_3\), then multiply by the \(2L\) possible direction-one sources in
the first coordinate. The result is \(2L[(2L+1)(2/3)L(2L^2+1)
+4L\,L(L+1)(2L+1)/3]\), which is precisely (232). Every boundary
contribution is present.

The exact frequencies satisfy \(0<\sigma_\nu<\sqrt{12}\): they are
square roots of sums of three numbers less than four. Therefore
\(\Lambda\ge\Lambda^2/\sqrt{12}\). Multiplication by \(D\ge0\) and
taking the trace preserves this inequality, since the trace of a product
of two positive matrices is the squared Hilbert--Schmidt norm of the
product of their square roots. The positive matrix
\(\Lambda^{1/2}D\Lambda^{1/2}\) has at most \(r\) eigenvalues.
Cauchy--Schwarz on that finite list proves \[
 \operatorname{tr}(\Lambda D\Lambda D)
 \ge\frac{\operatorname{tr}(\Lambda D)^2}{r}
 \ge\frac{\mathcal A_L^2}{12r}
 \ge\frac{128}{297}L^7.
 \tag{233}
\] For the last inequality use \(\mathcal A_L\ge(32/3)L^5\) in (232) and
\(r=16L^3+12L^2\le22L^3\) for \(L\ge2\). For an upper bound, the
matrix-entry formula for the trace gives
\(\operatorname{tr}(\Lambda D\Lambda D)
=\sum_{\nu,\mu}\sigma_\nu\sigma_\mu D_{\nu\mu}^2
\le12\operatorname{tr}(D^2)\le12rL^4\). In particular the exact mass in
(229) satisfies \[
 \frac4{99}L^7\le C_L^\Gamma\le\frac{99}{4}L^7 .
 \tag{234}
\] These inequalities prove strict positivity for every retained box.

Fix a physical energy \(\Omega\ge0\). A pair in (228) of energy at most
\(\Omega\) has both frequencies at most \(a\Omega\). Let
\(n_{L,a}(\Omega)=\#\{\nu:\sigma_\nu\le a\Omega\}\), counted with their
original transverse multiplicities, and let \(P_\Omega\) be that
coordinate projection in \(\mathbb R^r\). The entry formula and
\(\|D\|\le L^2\) give \[
 \begin{split}
 \rho_{L,a}^{\Gamma}([0,\Omega])
 &\le\frac{(a\Omega)^2
          \operatorname{tr}((P_\Omega D P_\Omega)^2)}
              {\operatorname{tr}(\Lambda D\Lambda D)}\\
 &\le\frac{12r(a\Omega)^2L^4}{\mathcal A_L^2}
                     n_{L,a}(\Omega)\\
 &\le\frac{297}{128}\frac{(a\Omega)^2}{L^3}
                     n_{L,a}(\Omega).
 \end{split}
 \tag{235}
\] The first numerator only enlarges the set of pairs satisfying the
energy condition; it does not discard a term of the original measure.
The factors \(3/32\), with doubled off-diagonal terms, cancel between
that enlarged numerator and the exact mass (229).

Here is a boundary-sensitive count for the remaining dimension. The full
frequency formula on the open box is \[
 \sigma^2=\sum_{i=1}^3
       4\sin^2\frac{\pi j_i}{2(2L+1)},\qquad
 0\le j_i\le2L,\quad k=\#\{i:j_i>0\}\ge2,
 \quad\hbox{multiplicity }k-1.
 \tag{236}
\] For completeness it follows by tensoring the one-dimensional vertex
cosines and edge sines. The vertex basis is \(v_0(t)=(2L+1)^{-1/2}\) and
\(v_j(t)=\sqrt{2/(2L+1)}
\cos[\pi j(t+1/2)/(2L+1)]\), \(0\le t\le2L\). The edge basis is
\(u_j(t)=-\sqrt{2/(2L+1)}
\sin[\pi j(t+1)/(2L+1)]\), \(0\le t<2L\). Finite geometric sums show
orthonormality, and taking the difference of adjacent cosines gives
\(dv_j=2\sin[\pi j/(2(2L+1))]u_j\). At a tensor index there are \(k\)
one-cochain directions and one gradient direction with components
\(s_i=2\sin[\pi j_i/(2(2L+1))]\). The exact oriented curl identity
\(\sum_{i<h}|s_ix_h-s_hx_i|^2
=|s|^2|x|^2-|\langle s,x\rangle|^2\) gives the value \(|s|^2\) on its
\(k-1\) dimensional perpendicular space. These tensor spaces exhaust all
edge cochains. The original index map is \(t=n_i+L\); it does not change
\(n_2^2\) in (221).

Concavity of sine on \([0,\pi/2]\) gives \(\sin t\ge2t/\pi\) there, by
its chord between the endpoints. Each positive coordinate in a frequency
at most \(a\Omega\) therefore satisfies \(2j_i/(2L+1)\le a\Omega\). Put
\[
 K_{L,a}(\Omega)=
 \min\left\{2L,\left\lfloor\frac{(2L+1)a\Omega}{2}\right\rfloor\right\}.
\] There are \(3K^2\) index triples with exactly two positive entries,
of multiplicity one, and \(K^3\) with three positive entries, of
multiplicity two. Hence the complete quantitative result is \[
 \rho_{L,a}^{\Gamma}([0,\Omega])
 \le\min\left\{1,\
     \frac{297}{128}\frac{(a\Omega)^2}{L^3}
       \left[3K_{L,a}(\Omega)^2+2K_{L,a}(\Omega)^3\right]\right\}.
 \tag{237}
\] For \(\Omega=0\), \(K=0\) and the measure has no zero atom, as is
also evident from (228). Along \(L_j=j^2\), \(a_j=1/(100j)\),
\(K_{L_j,a_j}(\Omega)\le3\Omega j/200\) for \(j\ge2\). Thus for each
fixed \(\Omega\) \[
 \rho_{L_j,a_j}^{\Gamma}([0,\Omega])=O_\Omega(j^{-5})
       \longrightarrow0.
 \tag{238}
\] To verify the stated order directly, the two terms in (237) are
bounded by constants times \(j^{-2}j^{-6}j^2=j^{-6}\) and
\(j^{-2}j^{-6}j^3=j^{-5}\), respectively, with all constants already
given by (237). The physical reference remains \(E_*=1\), so this is
escape from every fixed physical energy interval.

\subsubsection{18.4 An actual full-Hamiltonian diagonal and its
corrected native
state}\label{an-actual-full-hamiltonian-diagonal-and-its-corrected-native-state}

We now pass from the fixed-box theorem to a precisely defined sequence
of the original nonlinear Hamiltonians. Its coupling is selected using
the proved finite-box convergence; no uniform dependence on \(L\) is
assumed. Let \(j=2^s\), \(s\ge1\), retain \(L_j=j^2\), \(a_j=1/(100j)\),
and set \[
 d_{\mathrm{BL}}(\nu,\rho)=
 \sup\left\{\left|\int f\,d(\nu-\rho)\right|:
         f:\mathbb R_{\ge0}\to\mathbb R,\
         \|f\|_\infty\le1,\ \operatorname{Lip}(f)\le1\right\}.
 \tag{239}
\] At each fixed \(j\), (231) implies convergence in this distance.
Indeed choose a compact interval carrying all but a prescribed small
tail of both families, which is possible by the tail proof after (231).
The restrictions of the displayed functions to that interval are bounded
and equicontinuous; a finite piecewise-linear approximation on a fine
uniform mesh, with values rounded to a finite mesh in \([-1,1]\), is a
finite net in uniform norm. Convergence of the finitely many continuous
test integrals and the uniform tail bound gives convergence of the
displayed supremum.

Define \(k_j\) to be the least positive integer \(k\) such that every
integer \(q\ge k\) satisfies all four inequalities \[
 \begin{gathered}
 2^{-q}<1/j,\qquad
 d_{\mathrm{BL}}\bigl(
   \zeta_{\Gamma,L_j,a_j,2^{-q}},\rho_{L_j,a_j}^\Gamma\bigr)<1/j,\\
 \left|
   \frac{2^{-4q}\|v_{\Gamma,2^{-q}}\|^2}{C_{L_j}^\Gamma}-1
            \right|<1/j,\\
 \left|
  \operatorname{gap}(H_{L_j,a_j,2^{-q}})
    -\frac{4\sqrt2}{a_j}\sin\frac{\pi}{4L_j+2}
              \right|<1/j.
 \end{gathered}
 \tag{240}
\] The norms and operators in this display all belong to this same
\((L_j,a_j)\). The set defining \(k_j\) is nonempty: (231) proves the
second and third tail assertions, the finite-box physical eigenvalue
theorem in Section 17 proves the fourth, and \(2^{-q}\to0\) proves the
first. Intersect the four proved tails. The well ordering of the
positive integers then defines \(k_j\) without an unproved assumption.
This is an exact selection by finite-operator data; the present theorem
does not give a numerical rate for \(k_j\).

Set \(g_j=2^{-k_j}>0\). Its complete physical coefficients are \[
 g_j^2=2^{-2k_j},\qquad
 \xi_j=2^{4k_j-2},\qquad
 \kappa_j=200j\,2^{-2k_j},\qquad
 b_j=50j\,2^{2k_j},\qquad 2b_jM_j=100j\,2^{2k_j}M_j.
 \tag{241}
\] Every face and electric term of the original Hamiltonian is present.
The inequality \(g_j<1/j\) gives \(\xi_j\ge1\) on this sequence.

To control the actual native state along this selected path, strengthen
only the explicit cusp depth in (173). Put \[
 \begin{gathered}
 F(\xi)=\xi^{3/4}(6+64\xi),\qquad
 T_j^\sharp=2C+j^4e^{8\pi\xi_j}F(\xi_j)^{1/4},\\
 D_j^\sharp=L_{T_j^\sharp}q_{T_j^\sharp}+6m_{T_j^\sharp}^{\,2},
 \qquad
 \theta_j^\sharp=\frac{2\pi}{10^4j^2D_j^\sharp},\\
 v_j^\sharp=e^{-2\pi T_j^\sharp/j},\qquad
 t_{c,j}^\sharp=(v_j^\sharp)^j=e^{-2\pi T_j^\sharp}.
 \end{gathered}
 \tag{242}
\] The cover is still \(j\), with its unchanged pullback metric,
determinant \(-j^2D_j^\sharp\), original real periods, and fixed
physical coordinate permutation. Since \(F(\xi)\ge70\) for \(\xi\ge1\),
the systole and embedding proofs (174)--(177) hold on the same eventual
chart domains with \(T_j^\sharp\) replacing \(\widetilde T_j\). The full
nonlinear pullback, flat reference, and exact background ratio are those
same proved maps evaluated at the displayed original periods. No finite
Hamiltonian or vacuum changes when this external cusp angle is chosen.

In the denominator estimate (174), the fourth power \((T_j^\sharp-C)^4\)
now gains a factor \(F(\xi_j)\). Substitution in the exact count (180)
gives the explicit uniform bound \[
 \bar\eta_j^\sharp
 \le\frac{69120\pi^2}{10^8}\,j^{-4}
          \frac{2j^2+1}{j^2+1}
 \le A_{\mathrm{rel}}j^{-4}\longrightarrow0.
 \tag{243}
\] Thus the full physical native probability \(\nu_j^\sharp\), defined
from its actual centered vector \(\chi_j^\sharp=(K_{\theta_j^\sharp}
-\langle\psi_j,K_{\theta_j^\sharp}\psi_j\rangle)\psi_j\), satisfies \[
 \sup_B|\nu_j^\sharp(B)-\zeta_{\Gamma,L_j,a_j,g_j}(B)|
       \le\bar\eta_j^\sharp.
 \tag{244}
\] All these states are nonzero by Section 14, and the comparison
includes their original raw masses.

For every \(\Omega\ge0\), use the continuous function equal to one on
\([0,\Omega]\), linear from one to zero on \([\Omega,\Omega+1]\), and
zero beyond. It is an admissible test in (239). Equations (240), (244),
and (237) therefore give \[
 \nu_j^\sharp([0,\Omega])
 \le \rho_{L_j,a_j}^{\Gamma}([0,\Omega+1])
                      +j^{-1}+A_{\mathrm{rel}}j^{-4}
 \longrightarrow0.
 \tag{245}
\] This is a theorem about the full interacting Hamiltonians at the
strictly positive couplings (241), with their actual native states
(242). It is not an inference from a comparison Hamiltonian without
state transport. At the same time their actual physical gaps close: \[
 0<\operatorname{gap}(H_j)
 \le400\sqrt2\,j\sin\frac{\pi}{4j^2+2}+j^{-1}
 \le\frac{100\sqrt2\pi+1}{j}\longrightarrow0.
 \tag{246}
\] The first strict inequality follows from compact resolvent and the
proved unique finite vacuum. The last uses \(\sin x\le x\). Thus gap
closure and complete loss of this state's probability from finite energy
intervals occur on the same actual sequence. They are related by the
explicit covariance weights and state maps above, not treated as
independent mathematical presentations.

The raw native amplitudes remain quantitative. Write
\(\varepsilon_j=g_j^4\|v_{\Gamma,g_j}\|^2/C_{L_j}^\Gamma-1\), so
\(|\varepsilon_j|<1/j\) by (240). Then on \(\bar\eta_j^\sharp<1\), \[
 \begin{split}
 &(1-\bar\eta_j^\sharp)^2(1+\varepsilon_j)
       \frac49(\theta_j^\sharp)^4g_j^{-4}C_{L_j}^\Gamma
       \le d_j^\sharp\\
 &\hspace{15mm}\le
 (1+\bar\eta_j^\sharp)^2(1+\varepsilon_j)
       \frac49(\theta_j^\sharp)^4g_j^{-4}C_{L_j}^\Gamma .
 \end{split}
 \tag{247}
\] This follows directly from (186); every quantity has its original
finite value. The raw amplitude estimate (188), applicable here since
\(\xi_j\ge1\) and \(\bar\eta_j^\sharp<1\), with the additional factor
\(F(\xi_j)\) now in the denominator of the squared angle, gives \[
 \|\chi_j^\sharp\|
 \le\frac{384\sqrt3(5/4)^{1/4}\pi^2}{10^8}\,
      \frac{e^{-32\pi\xi_j}}{6+64\xi_j}\,j^{-10}
 \longrightarrow0.
 \tag{248}
\] The finite measures \(d_j^\sharp\nu_j^\sharp\) consequently vanish in
total mass, while (245) records the separate escape of their probability
after division by that mass.

For \(t>0\), split the probability integral at a fixed \(R\):
\(\int e^{-t\omega}d\nu_j^\sharp
\le\nu_j^\sharp([0,R])+e^{-tR}\). First send \(j\) to infinity using
(245), then \(R\) to infinity. The result is zero. At \(t=0\) that
probability integral is exactly one. This gives its exact discontinuous
unscaled correlation limit; the raw correlation also retains the factor
\(d_j^\sharp\) in (247). No strongly continuous nonnegative Hamiltonian
can represent this probability correlation on a unit vector, because its
spectral integral tends to one as \(t\downarrow0\) by bounded
convergence. This obstruction applies to this exact selected
native-state sequence.

The diagonal (240) has not been identified with either explicit running
coupling (181) or (183), and (245) makes no claim about their unresolved
covariance measures. It also does not exclude a different continuum
observable or a different positive-coupling path. The original
interacting four-dimensional construction and nonzero low-energy native
weight outside the proved escape sequences remain the research target.

\subsection{20. A fixed-box bounded low-mode observable and its exact
fate}\label{a-fixed-box-bounded-low-mode-observable-and-its-exact-fate}

The retained fixed-box source proof in the weak-coupling state-inputs
packet contains a second exact state map that is distinct from the
native \(n_2^2\)-weighted translation. It is useful because it produces
genuine low-mode spectral mass at each fixed box without replacing the
actual vacuum. Fix \(L\ge2\) and \(a>0\), use the exact tree/log map of
Section 17, and let \(\mathcal B_g\) be the full Haar-density/dilation
isometry. Strong actual-vacuum convergence and finite
spectral-projection convergence hold as \(g\downarrow0\): \[
 \mathcal B_g\psi_g\longrightarrow\Phi_0,\qquad
 \mathcal B_gP_g(I)\mathcal B_g^*\longrightarrow P_0(I)
 \tag{249}
\] for every bounded interval whose endpoints avoid the oscillator
spectrum. The proof uses the full nonlinear form, compactness from the
retained potential tails, invariant cutoffs, and the exact physical
projection. It does not substitute \(\Phi_0\) for \(\psi_g\) at finite
\(g\).

Choose a transverse mode \(\nu\), a smooth
simultaneous-rotation-invariant cutoff \(\chi\) equal to one near the
identity and supported inside the full logarithm chart, and define the
globally smooth physical multiplier \[
 B_{g,\vartheta}(Z)=
 \begin{cases}
 \exp\!\left(i\vartheta\chi(y)
      |(O^TG^{-1/2}(y/g))_\nu|^2\right),& Z=\exp(y),\\
 1,&\text{outside the chart}.
 \end{cases}
 \tag{250}
\] This is a bounded unitary function for every \(g>0\). Its transported
multiplier converges pointwise to
\(B_{0,\vartheta}(x)=\exp(i\vartheta|z_\nu|^2)\), so the actual centered
vectors \[
 v_g=(B_{g,\vartheta}-\langle\psi_g,B_{g,\vartheta}\psi_g\rangle)\psi_g
 \tag{251}
\] converge strongly under \(\mathcal B_g\) to
\(v_0=(B_{0,\vartheta}-c_0)\Phi_0\). Put \(\sigma=\sigma_\nu\),
\(k=3/2\), and \(\beta=4\vartheta/\sigma\). Direct Gamma integration
gives \[
 c_0=(1-i\beta)^{-3/2},\qquad
 \|v_0\|^2=1-(1+\beta^2)^{-3/2}.
 \tag{252}
\] The centered raw excitation measure is the complete positive series
\[
 \nu_0=
 \sum_{n=1}^{\infty}
 \frac{(3/2)_n}{n!}\,
 \frac{\beta^{2n}}{(1+\beta^2)^{n+3/2}}\,
 \delta_{\,2n\sigma/a}.
 \tag{253}
\] Indeed the radial Gamma density is \(q^{k-1}e^{-q}/\Gamma(k)\), the
invariant radial oscillator eigenvectors are normalized generalized
Laguerre polynomials, and the generating function gives the coefficient
\[
 \left\langle\varphi_n,B_{0,\vartheta}\Phi_0\right\rangle
 =\sqrt{\frac{(k)_n}{n!}}\,
   \frac{(-i\beta)^n}{(1-i\beta)^{n+k}}.
 \tag{254}
\] The \(n=0\) term is removed exactly by (251), and the binomial series
sums (253) to the mass in (252). Thus every term is a raw, nonnegative,
vacuum-orthogonal spectral weight. The corresponding fixed-box
correlation is \[
 \left\langle v_0,e^{-t(H_{\rm osc}-\mu_0)}v_0\right\rangle
 =
 [1+\beta^2(1-e^{-2\sigma t/a})]^{-3/2}
 -(1+\beta^2)^{-3/2}.
 \tag{255}
\]

For the lowest mode \(\sigma=\sqrt8\sin(\pi/(4L+2))\), the choice
\(\vartheta=\sigma/4\) gives \(\beta=1\), total mass \(d=1-2^{-3/2}\),
and first atom \(3/(8\sqrt2)\) at \(4\sqrt2\sin(\pi/(4L+2))/a\). Now
take \(L_j=j^2\), \(a_j=1/(100j)\), for integers \(j\ge2\). For each
fixed \(j\), the finite-box theorem permits a separately selected
strictly positive \(g_j<1/j\) with bounded-Lipschitz distance below
\(1/j\) and centered mass within \(1/j\) of \(d\). Since every
fixed-\(n\) energy \(2n\sigma_j/a_j\) tends to zero and the weights in
(253) are summable and independent of \(j\), \[
 \nu_{g_j,j}\Longrightarrow d\,\delta_0,\qquad
 \langle v_{g_j,j},e^{-t(H_{g_j,j}-\mathcal E_{g_j,j})}v_{g_j,j}\rangle
 \longrightarrow d\quad(t\ge0).
 \tag{256}
\] This is an exact sequence of nonzero physical vacuum-orthogonal
states with positive raw mass, but its limit is a zero-energy atom. Any
nonnegative self-adjoint continuum Hamiltonian \(K\) representing these
limits together with the retained centering would have a unit vacuum
\(\Omega\), a vector \(v\) with \(\langle\Omega,v\rangle=0\) and
\(\|v\|^2=d>0\), and \(\langle v,e^{-tK}v\rangle=d\) for every \(t>0\).
For any such \(t\), the spectral theorem gives
\(\int(1-e^{-tE})\,d\nu_v(E)=0\). The integrand is nonnegative and
strictly positive for \(E>0\), so \(\nu_v((0,\infty))=0\) and \(Kv=0\).
Thus \(v\) is a second independent zero-energy vector, contradicting
uniqueness of the vacuum. Retaining the centering is essential: the
constant autocorrelation alone is also represented by
\(v=\sqrt d\,\Omega\). Hence this particular macroscopic low-mode
identification, with its original centering, cannot be the desired
unique-vacuum interacting continuum. It is a proved degeneration, not a
mass-gap counterexample. This separately selected sequence has not been
identified with the covariance sequence in Section 18.

The construction establishes a favorable finite-box state map and
nonzero low-mode weights before the simultaneous limit. It does not
prove that those weights survive as a nonzero continuum measure on
\((0,\Omega]\), and it does not identify the selected \(g_j\) with the
explicit logarithmic or fixed-electric-coefficient paths. The full
proof, including chart domains, compactness, Laguerre completeness and
projection convergence, is retained in the cited source file.

\subsection{21. Compact fluid concentration and its exact gauge
source}\label{compact-fluid-concentration-and-its-exact-gauge-source}

\subsubsection{21.1. Primary sources and the mathematical
input}\label{primary-sources-and-the-mathematical-input}

The public
\href{https://cdn.openai.com/pdf/32d9f210-8b73-45e0-91bc-82a30aef8a9a/navier-stokes.pdf}{OpenAI
manuscript} states finite-time velocity blowup for every positive
viscosity, with zero initial velocity, a smooth force compactly
supported in space and time, common compact spatial support before time
1, and bounded kinetic energy. This is its Theorem 1.1; the whole-space
and torus conclusions are stated as the Clay C/D alternatives. The
retained 165-page edition and the later 166-page edition have different
hashes. The exact versions and the
\href{https://github.com/openai/NavierStokesAndEuler/tree/8937a8f4cbc7abaab5e9e97d1cc7f5d2319d9538}{pinned
public formal source} are recorded in the source receipt. Equality of
the revisions and a complete independent verification of the claimed
existence proof are not asserted.

\href{https://cims.nyu.edu/~tristanb/statement.pdf}{Tristan Buckmaster's
statement} attributes the smooth-forcing porous-media, Boussinesq and
incompressible Euler results to his personal collaboration with Levent
Alpöge, and credits Diego Córdoba and Luis Martínez-Zoroa for the
preceding forced-blowup programme. Their
\href{https://github.com/tristanbuckmaster/fluid_lean}{public source} is
a separate source object. These attributions do not identify their Euler
theorem with the positive-viscosity Navier--Stokes theorem. No
allegation about access to private material is used in any calculation.

The following results are proved directly for smooth input fields, with
their entire source and coordinate data retained. Thus the identities
and estimates are independently established here; existence of a field
with the newly announced breakdown properties remains the cited source's
claim. This is an exact translation theorem for the displayed class of
inputs, not an assumption of a missing quantum or fluid existence
theorem.

Use the original spatial coordinates \(x\in\mathbb R^3\), time
\(0\le t<1\), positive viscosity \(\nu\), and smooth real fields
\(u,p,f\) satisfying \[
 \partial_tu_i+\sum_{\ell=1}^3u_\ell\partial_\ell u_i
 -\nu\Delta u_i+\partial_i p=f_i,\qquad
 \sum_i\partial_i u_i=0,\qquad u(x,0)=0.
 \tag{257}
\] The supports of \(u(\cdot,t)\) and \(p(\cdot,t)\) lie in one compact
set \(K\subset\mathbb R^3\), and
\(f\in C_c^\infty(\mathbb R^3\times(0,\infty);\mathbb R^3)\). Fix
\(R>0\) with \(K\subset B_R(0)\). The endpoint 1 is the endpoint
specified by the source, not a replacement of an original time
parameter. Write \[
 \omega=\nabla\times u,\qquad
 F_1(t)=\int_0^t\|f(s)\|_{L^2(\mathbb R^3)}\,ds,\qquad
 F_1(1)=\int_0^1\|f(s)\|_2\,ds<\infty .
 \tag{258}
\]

\subsubsection{21.2. Exact energy identity and finite magnetic spacetime
integral}\label{exact-energy-identity-and-finite-magnetic-spacetime-integral}

Multiplication of (257) by \(u_i\), summation, and integration over
\(\mathbb R^3\) give \[
 \frac12\frac{d}{dt}\|u(t)\|_2^2+
 \nu\|\nabla u(t)\|_2^2=\int_{\mathbb R^3}f(x,t)\cdot u(x,t)\,dx .
 \tag{259}
\] All integrations by parts are justified by the common compact support
on every closed interval ending before 1. The transport integral is
\(\int\operatorname{div}(u|u|^2/2)=0\), using \(\operatorname{div}u=0\).
The pressure integral is
\(\int\operatorname{div}(pu)-p\operatorname{div}u=0\). Each viscous
integral is \(\int u_i\partial_\ell^2u_i=-\int|\partial_\ell u_i|^2\).
These equalities prove (259) with its original signs.

For \(\epsilon>0\), set \(Y_\epsilon=(\|u\|_2^2+\epsilon)^{1/2}\).
Equation (259) and Cauchy--Schwarz imply
\(Y_\epsilon'\le\|f\|_2\|u\|_2/Y_\epsilon\le\|f\|_2\). Integrating from
zero and letting \(\epsilon\downarrow0\) proves
\(\|u(t)\|_2\le F_1(t)\). Consequently
\(\int_0^t\langle f,u\rangle\,ds\le
\int_0^tF_1'(s)F_1(s)\,ds=F_1(t)^2/2\). For each fixed time, expansion
of the curl gives \[
 \|\omega\|_2^2
 =\sum_{i,\ell}\int|\partial_i u_\ell|^2
 -\sum_{i,\ell}\int\partial_i u_\ell\,\partial_\ell u_i
 =\|\nabla u\|_2^2 .
 \tag{260}
\] Indeed integration by parts changes the second sum into
\(-\sum_\ell\int u_\ell\partial_\ell(\sum_i\partial_i u_i)=0\).
Integrating (259) and using nonnegative monotone convergence therefore
proves the quantitative endpoint estimate \[
 \frac12\|u(t)\|_2^2+\nu\int_0^t\|\omega(s)\|_2^2\,ds
 \le\frac12F_1(t)^2,\qquad
 \int_0^1\|\omega(s)\|_2^2\,ds\le\frac{F_1(1)^2}{2\nu}.
 \tag{261}
\] No pointwise-in-time bound for vorticity was used or follows from
(261).

\subsubsection{21.3. The entire connection, curvature, current and gauge
law}\label{the-entire-connection-curvature-current-and-gauge-law}

Fix \(c>0\), \(g>0\), and a real calibration \(\lambda\ne0\). Keep
\(T=-i\sigma_3/2\), so \(-2\operatorname{tr}(T^2)=1\). Use \(X^0=ct\),
\(X^i=x^i\) on \(\mathbb R^3\times[0,c)\), with metric
\(\operatorname{diag}(-1,1,1,1)\). Define \[
 A_0=0,\qquad A_i(X)=\lambda u_i(x,X^0/c)T,\qquad
 u_i(x,t)=-\frac{2}{\lambda}\operatorname{tr}(T A_i(ct,x)).
 \tag{262}
\] Thus the map to this specified connection presentation is linear and
injective with the displayed inverse. With velocity in length/time and
\(A_i\) in inverse length, \(\lambda\) has units time/length squared.
Neither \(c,\lambda,g\) nor \(\nu\) is set to one.

All components in (262) commute, so evaluation of the full formula
\(F_{\mu\nu}=\partial_\mu A_\nu-\partial_\nu A_\mu+[A_\mu,A_\nu]\) gives
\[
 F_{0i}=\frac{\lambda}{c}\partial_tu_i\,T,\qquad
 F_{ij}=\lambda(\partial_i u_j-\partial_j u_i)T.
 \tag{263}
\] In particular, define the nonnegative magnetic and electric component
amplitudes by \[
 {\cal M}=\left[-2\sum_{i<j}\operatorname{tr}(F_{ij}^2)\right]^{1/2}
 =|\lambda|\,|\omega|,\qquad
 {\cal E}=\left[-2\sum_i\operatorname{tr}(F_{0i}^2)\right]^{1/2}
 =\frac{|\lambda|}{c}|\partial_tu|.
 \tag{264}
\] These are specified component norms, not the indefinite Lorentzian
action contraction. Under a smooth \(SU(2)\) gauge map \(h(X)\),
\(A^h=h^{-1}Ah+h^{-1}dh\) and \(F^h=h^{-1}Fh\). For verification, acting
on any smooth vector-valued section \(s\) gives
\((d+A^h)s=h^{-1}(d+A)(hs)\); applying the covariant derivative once
more proves the curvature law. Cyclicity of trace then proves that both
amplitudes (264) are gauge invariant, since
\(\operatorname{tr}((h^{-1}Fh)^2)=\operatorname{tr}(F^2)\).

Use exactly \(D^\mu F_{\mu\nu}=g^2j_\nu\). Each commutator with
\(A_\mu\) vanishes on this image. The time component is \[
 j_0=-\frac{\lambda}{g^2c}\partial_t\operatorname{div}u\,T=0,
 \qquad
 j_i=J_iT,\qquad
 J_i=\frac{\lambda}{g^2}
       \left(\Delta u_i-\frac1{c^2}\partial_t^2u_i\right).
 \tag{265}
\] The spatial derivative before using incompressibility is
\(\lambda(\Delta u_i-\partial_i\operatorname{div}u)\); the raised time
derivative contributes \(-\lambda c^{-2}\partial_t^2u_i\). Thus (265)
includes the Lorentzian sign. With \(w_i=\partial_tu_i\),
differentiation of (257) retains the complete fluid dynamics in the
current: \[
 J_i=\frac{\lambda}{g^2}
 \left[\Delta u_i-\frac1{c^2}
 \left(\nu\Delta w_i-\sum_\ell w_\ell\partial_\ell u_i
 -\sum_\ell u_\ell\partial_\ell w_i
 -\partial_i\partial_t p+\partial_t f_i\right)\right].
 \tag{266}
\] In particular \(D^\nu j_\nu=0\): the commutators vanish, \(j_0=0\),
and \(\sum_i\partial_iJ_i=(\lambda/g^2)
(\Delta-c^{-2}\partial_t^2)\operatorname{div}u=0\). This is an exact map
to a conserved gauge source, including viscosity, nonlinear transport,
pressure and forcing.

The physical spacetime volume element in these coordinates is
\(dX^0\,d^3x=c\,dt\,d^3x\). Equations (261) and (264) give \[
 \int_0^c\!\!\int_{\mathbb R^3}{\cal M}^2\,d^3x\,dX^0
 =c\lambda^2\int_0^1\|\omega(t)\|_2^2\,dt
 \le\frac{c\lambda^2}{2\nu}F_1(1)^2 .
 \tag{267}
\] This is a finite magnetic spacetime integral, with the original time
Jacobian. It makes no assertion of finite electric action or of a
quantum vacuum.

\subsubsection{21.4. Common compact support transfers velocity blowup to
curvature}\label{common-compact-support-transfers-velocity-blowup-to-curvature}

For every smooth compactly supported divergence-free vector field, \[
 u(x)=\frac1{4\pi}\int_{\mathbb R^3}
       \frac{\omega(y)\times(x-y)}{|x-y|^3}\,dy .
 \tag{268}
\] Here is the precise reconstruction argument. The function
\(G(x)=1/(4\pi|x|)\) obeys \(-\Delta G=\delta_0\). Away from zero its
Laplacian is zero; integration by parts outside a ball of radius
\(\epsilon\) gives a boundary flux tending to the test function's value
at zero, because \(-\partial_rG=1/(4\pi\epsilon^2)\). The remaining
boundary term containing \(G\partial_r\phi\) tends to zero. This proves
the distribution identity. The vector identity
\(\nabla\times\omega=\nabla\operatorname{div}u-\Delta u=-\Delta u\) then
gives \(G*(\nabla\times\omega)=G*(-\Delta u)=u\), where derivatives may
be moved onto the compactly supported smooth field by integration by
parts. Moving the curl onto \(G\) and using
\(\nabla G(x-y)=-(x-y)/(4\pi|x-y|^3)\) proves (268). The singular kernel
is locally integrable, since its magnitude is bounded by
\(1/(4\pi|x-y|^2)\).

For \(x\in K\), its integration domain lies in \(B_{2R}(x)\), so \[
 |u(x)|\le\frac{\|\omega\|_\infty}{4\pi}
       \int_{B_{2R}(x)}|x-y|^{-2}\,dy
       =2R\|\omega\|_\infty.
 \tag{269}
\] For \(x\notin K\), \(u(x)=0\); hence this is also the global supremum
bound. The same formula proves injectivity of the spatial curvature map
on this divergence-free, compactly supported input space: zero curvature
implies \(\omega=0\), then \(u=0\). This statement is about the
displayed vector-to-curvature map, not an unproved inverse on arbitrary
non-Abelian gauge orbits.

Consequently the compactly supported velocity blowup described by the
source forces \[
 \limsup_{t\uparrow1}\|{\cal M}(t)\|_\infty=\infty,
 \qquad
 \int_0^1\|{\cal E}(t)\|_\infty\,dt=\infty .
 \tag{270}
\] The first assertion follows from
\(\|{\cal M}\|_\infty=|\lambda|\|\omega\|_\infty
\ge|\lambda|\|u\|_\infty/(2R)\). For the second,
\(u(x,t)=\int_0^t\partial_su(x,s)\,ds\) gives
\(\|u(t)\|_\infty\le\int_0^t\|\partial_su(s)\|_\infty\,ds\); a finite
endpoint integral would bound every \(\|u(t)\|_\infty\). Equation (264)
then supplies the exact calibration. In the \(X^0\) time coordinate the
same assertion is \(\int_0^c\|{\cal E}(X^0)\|_\infty\,dX^0
=|\lambda|\int_0^1\|\partial_tu(t)\|_\infty\,dt=\infty\). No claim of
electric \(L^2\) divergence is made. The invariant magnetic divergence
and the finite integral (267) are compatible.

\subsubsection{21.5. The complete source must be singular at the
endpoint}\label{the-complete-source-must-be-singular-at-the-endpoint}

Since the time support of \(f\) is a compact subset of \((0,\infty)\),
it vanishes for \(0\le t\le t_0\) for some \(t_0>0\). The estimate
\(\|u(t)\|_2\le F_1(t)\) proves \(u=0\) on that interval. Thus both
\(u(0)=0\) and \(\partial_tu(0)=0\) hold. Equation (265) is equivalently
the forced wave equation \[
 \partial_t^2u-c^2\Delta u
       =-\frac{g^2c^2}{\lambda}J .
 \tag{271}
\] For every \(t<1\) its exact zero-data solution is \[
 u(t,x)=-\frac{g^2c^2}{4\pi\lambda}
       \int_0^t(t-s)\int_{S^2}
                J(s,x+c(t-s)n)\,dS(n)\,ds .
 \tag{272}
\] For completeness, if \(q\) is smooth and compactly supported, its
spherical mean \(M(r,x)=(4\pi)^{-1}\int_{S^2}q(x+rn)\,dS(n)\) satisfies
\(M_{rr}+2M_r/r=\Delta_xM\). This follows by writing the Laplacian in
spherical coordinates: its angular part integrates to zero on the sphere
by surface integration by parts. Therefore \(S_c(t)q=tM(ct,x)\)
satisfies \((\partial_t^2-c^2\Delta_x)S_c(t)q=0\), \(S_c(0)q=0\), and
\(\partial_tS_c(0)q=q\). Differentiating the integral
\(\int_0^t S_c(t-s)J(s)\,ds\) twice gives its forcing \(J(t)\) and zero
initial data. Multiplication by \(-g^2c^2/\lambda\) proves that the
right side of (272) solves (271). The difference from \(u\) is a
homogeneous wave with zero data. On every finite interval it has compact
spatial support; integration by parts makes its wave energy
\(\frac12\int(|\partial_tz|^2+c^2|\nabla z|^2)\) constant. Zero initial
energy implies \(\partial_tz=\nabla z=0\); the zero initial value then
gives \(z=0\). This proves (272).

Using the Euclidean vector supremum norm in (272) yields \[
 \|u(t)\|_\infty
 \le\frac{g^2c^2}{|\lambda|}
       \int_0^t(t-s)\|J(s)\|_\infty\,ds
 \le\frac{g^2c^2}{|\lambda|}
       \int_0^1(1-s)\|J(s)\|_\infty\,ds .
 \tag{273}
\] In particular finite-time unbounded velocity forces the exact
necessary divergence \[
 \int_0^1(1-s)\|J(s)\|_\infty\,ds=\infty .
 \tag{274}
\] The quantity \(|J|=[-2\sum_i\operatorname{tr}(j_i^2)]^{1/2}\) is
gauge invariant by the same trace calculation as (264). All \(J(s)\)
have support in \(K\). A smooth extension of this current across \(t=1\)
would be bounded on the compact cylinder, making (274) finite, which is
impossible. A zero current is also impossible for a nonzero zero-data
input, directly by (272). Thus smooth compact forcing in the fluid
equation does not produce a smooth gauge current through the singular
endpoint under the exact map (262).

\subsubsection{21.6. The displayed inner profile and its shrinking-loop
holonomy}\label{the-displayed-inner-profile-and-its-shrinking-loop-holonomy}

This calculation uses the explicit inner-circle data in the retained
165-page source edition. Its locators are fixed to that edition:
\(\tau=1-t\), \(A=1/2+h\), \(0<h<1/100\) on page 7; \(q_*>0\),
\(0<X_a<X_{\rm ext}<\infty\), and \(X_{\rm in}\in(0,X_a)\) in Theorem
3.1 on pages 14--15; and \(e_0=E_0(X_{\rm in},0)>0\) in Proposition 9.9,
Step 5, on page 116. The inner value in (3.6) and (10.21), and viscosity
map (10.22), retain \[
 \begin{split}
 u_\theta(\sqrt{2X_{\rm in}\tau},0,1-\tau)
   &=\tau^{-1/2-h}(e_0+\varepsilon(\tau)),\\
 |\varepsilon(\tau)|&\le C_{\rm in}\tau^{2h},
       \quad 0<\tau<\tau_*,\\
 u_\nu(x,t)&=\sqrt{\nu}\,u(x/\sqrt{\nu},t).
 \end{split}
 \tag{275}
\] Here \(C_{\rm in}\ge0\) and \(\tau_*>0\) are the finite remainder
constant and small-time interval supplied by the source's printed
remainder estimate; no numerical value is asserted for either. The
reference \(u\) in (275) is the source's unit-viscosity profile, and
\(u_\nu\) is its displayed map to each prescribed viscosity. All
subsequent expressions retain that prescribed \(\nu\). The annular
corrections vanish for this fixed \(X_{\rm in}<X_a\); the remaining
background is axisymmetric, and its localization cutoff equals one near
the spacetime origin of the blowup. Thus the displayed tangential value
holds around the whole inner circle. A value at only one angular point
would not justify the following calculation.

At time \(t=1-\tau\), parameterize the positively oriented circle by
\(\gamma_{\nu,\tau}(\theta)=
(r_{\nu,\tau}\cos\theta,r_{\nu,\tau}\sin\theta,0)\), where
\(r_{\nu,\tau}=\sqrt{2\nu X_{\rm in}\tau}\) and \(0\le\theta\le2\pi\).
Equations (275), together with
\(\gamma'_{\nu,\tau}=r_{\nu,\tau}e_\theta\), give the exact circulation
and its retained remainder: \[
 \begin{split}
 {\cal C}_{\nu,\tau}
   :=\oint_{\gamma_{\nu,\tau}}u_\nu\cdot dx
   &=2\pi\nu\sqrt{2X_{\rm in}}\,
       \tau^{-h}(e_0+\varepsilon(\tau))
       \longrightarrow+\infty,\\
 \sup_{D_{\nu,\tau}}|\omega_{\nu,3}|
   &\ge\sqrt{\frac{2}{X_{\rm in}}}\,
       \tau^{-1-h}(e_0+\varepsilon(\tau)).
 \end{split}
 \tag{276}
\] The second inequality holds for all sufficiently small \(\tau\). In
fact Stokes' theorem for the disk with normal \(+e_3\) gives
\({\cal C}_{\nu,\tau}=\int_{D_{\nu,\tau}}\omega_{\nu,3}\,dS\). Its area
is \(2\pi\nu X_{\rm in}\tau\); dividing the first line by this original
area proves the second. Positivity of the last factor follows once
\(C_{\rm in}\tau^{2h}\le e_0/2\). When \(C_{\rm in}=0\), it holds
throughout the stated small-time interval. No sign of the orientation or
viscosity factor is dropped.

For the connection (262) formed from \(u_\nu\), ordinary column
transport along this circle solves
\(P'(\theta)=-\lambda T(u_\nu(\gamma(\theta),t)\cdot
\gamma'(\theta))P(\theta)\), \(P(0)=I\). The original lattice endpoint
convention uses inverse transport \(U=P^{-1}\). Differentiating
\(P^{-1}P=I\) gives \(U'=U\lambda T(u_\nu\cdot\gamma')\). Since all
coefficients commute, differentiation of the exponential integral proves
the complete retained pair \[
 \begin{split}
 P_{\nu,\tau}&=\exp(-\lambda T{\cal C}_{\nu,\tau}),\\
 U_{\nu,\tau}=P_{\nu,\tau}^{-1}
    &=\exp(+\lambda T{\cal C}_{\nu,\tau}),\\
 W_{\nu,\tau}=\operatorname{tr}U_{\nu,\tau}
    &=2\cos(\lambda{\cal C}_{\nu,\tau}/2).
 \end{split}
 \tag{277}
\] The eigenvalues of \(T\) are \(-i/2,i/2\), which proves the trace
formula with its factor \(1/2\). For an open path, substitution of
\(h(\gamma(\theta))^{-1}P(\theta)h(\gamma(0))\) into the transport
equation proves \(P^h=h_t^{-1}Ph_s\); taking its inverse gives
\(U^h=h_s^{-1}Uh_t\), exactly the original lattice law. At the closed
endpoint this is conjugation, so \(W_{\nu,\tau}\) is gauge invariant.
The opposite transport signs are both retained; the trace is unchanged
because the cosine is even.

The cluster set of these traces as \(\tau\downarrow0\) is exactly
\([-2,2]\). To prove the nontrivial inclusion, fix \(y\in[-2,2]\) and
set \(\alpha=\arccos(y/2)\). The circulation is continuous on every
regular-time interval and tends to \(+\infty\). The intermediate value
theorem therefore supplies times with
\({\cal C}_{\nu,\tau_n}=(4\pi n+2\alpha)/|\lambda|\) for all
sufficiently large integers \(n\). They tend to zero, since the
circulation is bounded on every closed interval separated from zero.
Formula (277) gives \(W_{\nu,\tau_n}=y\), for either sign of
\(\lambda\). Conversely the cosine formula always lies in \([-2,2]\),
proving equality of the cluster set. This is an exact calculation on the
displayed source core data; it does not independently verify the
source's complete residual construction.

This identifies a genuine invariant curvature concentration and a
quantitative source obstruction. It supplies neither a source-free
solution under this map nor a wave-functional in the physical
vacuum-orthogonal quantum Hilbert space of the preceding sections. The
exact connection-to-link map retained in the source packet may still be
applied on every regular time slice. The source and curvature
calculations here do not change its finite-box vacuum, its raw state
norms, or its spectral inequality.

\subsection{22. A positive Gamma limit and an exact rescaled electric
observable}\label{a-positive-gamma-limit-and-an-exact-rescaled-electric-observable}

\subsubsection{22.1. The radial family with all raw spectral weights
retained}\label{the-radial-family-with-all-raw-spectral-weights-retained}

The retained radial Gamma and amplitude/time proofs extend the
fixed-parameter calculation in Section 20. We first give the full
measure calculation needed for the new operator map below. Keep \[
 L_j=j^2,\quad a_j=\frac1{100j},\quad
 \sigma_j=\sqrt8\sin\frac{\pi}{4j^2+2},\quad
 \lambda_j=\frac{2\sigma_j}{a_j},\quad
 c_j=j\lambda_j,\quad c_\Gamma=100\sqrt2\pi,\quad k=\frac32 .
 \tag{278}
\] Here \(j\ge2\) is an integer, and \(c_\Gamma\) is the physical Gamma
scale, distinct from the spacetime coordinate constant in Section 21.
The full actual finite-box Hamiltonian and vacuum remain those of
(190)--(205). At each fixed box use the same exact tree, Haar, logarithm
and kinetic-coordinate maps. Put \[
 F_j(y)=\chi(y)\frac{\sigma_j}{4}
       |(O^{\mathsf T}G^{-1/2}y)_{\nu_j}|^2,\qquad
 Q_{g,j}=g^{-2}F_j,\qquad
 B_{g,j}(\beta)=e^{i\beta Q_{g,j}} .
 \tag{279}
\] The real invariant cutoff is supported strictly inside the full
logarithm chart and equals one near zero. Extend \(F_j\) by zero off
that chart; it is a globally smooth real physical function. Thus
\(Q_{g,j}\) is a bounded self-adjoint multiplication operator at every
finite \(g>0,j\), even though no bound uniform in the regulator is
asserted. In the transported limit it becomes
\(q=\sigma_j|z_{\nu_j}|^2/4\). Its probability law in the comparison
vacuum is \(d\rho(q)=q^{k-1}e^{-q}dq/\Gamma(k)\). All three colour
coordinates and the original kinetic tensor occur in (279).

Let \(\beta_j=\sqrt j\), subtract the actual mean
\(m_{g,j}=\langle\psi_{g,j},B_{g,j}(\beta_j)\psi_{g,j}\rangle\), and set
\(v_{g,j}=(B_{g,j}(\beta_j)-m_{g,j})\psi_{g,j}\). Its raw norm is
\(1-|m_{g,j}|^2>0\), by the strict Cauchy--Schwarz inequality for a
nonconstant multiplier and the strictly positive vacuum. The finite-box
theorem and (253) give the complete comparison measure \[
 \nu_{0,j}=\sum_{n=1}^{\infty}
 \frac{(k)_n}{n!}\frac{j^n}{(j+1)^{n+k}}\,
 \delta_{\lambda_j n},\qquad
 d_j=\nu_{0,j}([0,\infty))=1-(j+1)^{-k}.
 \tag{280}
\] The missing \(n=0\) coefficient is exactly the removed vacuum
component, not an omitted low-energy mass.

For a direct quantitative limit, let \(Q\) have law \(\rho\), and
conditionally on \(Q=q\) let \(N_j\) be Poisson with mean \(jq\). Gamma
integration gives \[
 {\mathbb P}(N_j=n)
 =\frac{j^n}{n!\Gamma(k)}
   \int_0^\infty q^{n+k-1}e^{-(j+1)q}\,dq
 =\frac{(k)_n}{n!}\frac{j^n}{(j+1)^{n+k}}.
 \tag{281}
\] Thus the law of \(\lambda_jN_j\) is exactly
\(\nu_{0,j}+(j+1)^{-k}\delta_0\). Since
\({\mathbb E}(N_j/j-Q\mid Q)=0\),
\({\mathbb E}((N_j/j-Q)^2\mid Q)=Q/j\), \({\mathbb E}Q=k\) and
\({\mathbb E}Q^2=k(k+1)\), \[
 {\mathbb E}|\lambda_jN_j-c_\Gamma Q|^2
 =\frac{kc_j^2}{j}+k(k+1)(c_j-c_\Gamma)^2.
 \tag{282}
\] The mixed term vanishes by conditional expectation. For
\(x_j=\pi/(4j^2+2)\), integration of \(1-\cos x\le x^2/2\) gives
\(0\le x_j-\sin x_j\le x_j^3/6\), and therefore \[
 0<c_\Gamma-c_j\le
 \frac{c_\Gamma}{2j^2+1}
 +\frac{400\sqrt2\pi^3j^2}{6(4j^2+2)^3}.
 \tag{283}
\] For finite measures, let \(d_{\rm BL}\) use tests of supremum norm at
most one and Lipschitz constant at most one. The coupling,
Cauchy--Schwarz and the removed zero mass yield \[
 \begin{split}
 d_{\rm BL}(\nu_{0,j},\gamma_\Gamma)
 &\le
 \left[\frac{kc_j^2}{j}
       +k(k+1)(c_j-c_\Gamma)^2\right]^{1/2}+(j+1)^{-k}
 \longrightarrow0,\\
 d\gamma_\Gamma(\omega)
 &=\frac{\omega^{1/2}e^{-\omega/c_\Gamma}}
       {\Gamma(3/2)c_\Gamma^{3/2}}\,d\omega .
 \end{split}
 \tag{284}
\] The density follows by the change of variable \(\omega=c_\Gamma q\).
It has total mass one, zero mass at zero, and positive mass on every
\((0,\epsilon)\), \(\epsilon>0\).

At each fixed \(j\), actual-vacuum and finite-projection convergence
give convergence of the actual raw measure to (280), including total
mass. A positive dyadic \(g_j\) can consequently satisfy \[
 g_j<j^{-5},\quad
 d_{\rm BL}(\nu_{g_j,j},\nu_{0,j})<j^{-1},\quad
 |\|v_{g_j,j}\|^2-d_j|<j^{-1},\quad
 |m_{g_j,j}-(1-i\sqrt j)^{-k}|<j^{-1}.
 \tag{285}
\] Below we impose additional proved fixed-box weighted convergence
conditions on this selection. This refines the finite list used to
choose the dyadic; it does not identify it with either prescribed
running-coupling trajectory or with Section 18's native covariance
sequence. All original coefficients remain \(\kappa_j=200jg_j^2\),
\(b_j=50j/g_j^2\), \(\xi_j=1/(4g_j^4)\), including the scalar
\(2b_j|P_j|\). Equations (284)--(285) prove the actual nonlinear raw
limit \(\nu_{g_j,j}\Rightarrow\gamma_\Gamma\) and
\(\|v_{g_j,j}\|^2\to1\). The absence of endpoint atoms gives \[
 \nu_{g_j,j}((0,\epsilon))
 \longrightarrow\gamma_\Gamma((0,\epsilon))>0 .
 \tag{286}
\] This is a positive low-energy limit of this radial observable. It has
not been identified with the original magnetic translation state or a
spatial local interacting continuum.

\subsubsection{22.2. Exact phase conjugation in the full nonlinear
Hamiltonian}\label{exact-phase-conjugation-in-the-full-nonlinear-hamiltonian}

Fix \(j,g>0\), abbreviate \(A=H-\mathcal E\), \(Q=Q_{g,j}\), and retain
every edge and colour derivative \(X=X_{e,\alpha}\). Each \(X\) is
skew-adjoint for the original Haar measure. For a smooth physical
function \(f\), the product rule gives \[
 X(e^{i\beta Q}f)
 =e^{i\beta Q}(Xf+i\beta(XQ)f).
 \tag{287}
\] Applying \(X\) a second time and using the entire Hamiltonian gives
\[
 \begin{split}
 B(\beta)^*AB(\beta)&=A+\beta J+\beta^2R,\\
 J=i[A,Q]&=-i\kappa\sum_{e,\alpha}
       \bigl(2(X_{e,\alpha}Q)X_{e,\alpha}
                    +X_{e,\alpha}^2Q\bigr),\\
 R=\tfrac12[Q,[A,Q]]
     &=\kappa\sum_{e,\alpha}(X_{e,\alpha}Q)^2\ge0 .
 \end{split}
 \tag{288}
\] These identities hold on smooth physical functions and extend to the
original second-order Sobolev operator domain, since each coefficient
and each multiplier is smooth on the compact group. The Wilson
multiplication, its scalar term and \(\mathcal E\) commute with \(B\);
none has been deleted. For the double commutator, \([J,Q]=-2iR\) follows
by applying the displayed first-order expression to \(Qf\), so
\([Q,[A,Q]]=-i[Q,J]=2R\). This checks both signs. Gauge invariance
follows either from the commutator formula or from gauge invariance of
\(A,Q\). In particular \(R\) is an actual nonnegative, bounded,
gauge-invariant multiplication operator at every finite regulator.

For any smooth \(h\), the exact ground-state form identity is
\(q_A[h\psi]=\kappa\int\psi^2\sum|Xh|^2\,dU\). It is obtained by
expanding \(X(h\psi)\), integrating its mixed terms by parts, and
substituting the full equation \(H\psi=\mathcal E\psi\); the potential
cancels against that equation with its complete coefficient. Setting
\(h=e^{i\beta Q}\) in (287) proves \[
 q_A[(B(\beta)-\langle B(\beta)\rangle_\psi)\psi]
 =q_A[B(\beta)\psi]=\beta^2\langle R\rangle_\psi .
 \tag{289}
\] Centering does not change this form because \(A\psi=0\). This is an
exact electric-energy formula for a raw state.

\subsubsection{22.3. The complete fixed-box weighted tail
estimate}\label{the-complete-fixed-box-weighted-tail-estimate}

Continue at fixed \(j\), and abbreviate \(F=F_j\), \(R_g=R_{g,j}\),
\(a=a_j\), and \(\nu=\nu_j\), so \(\sigma_\nu=\sigma_j\). Write
\(W=\sum_p(2-\operatorname{tr}U_p)\) for the complete face potential,
and \(N,M\) for the original edge and face counts. All formulas
involving \(\mathcal B_g\) below are in the original scaled chord
coordinates \(x\); their \(\Phi_0\) is the exact comparison vacuum in
those coordinates. The constant-Jacobian kinetic-coordinate unitary from
(224) then gives the stated radial \(L^2(\rho)\) map.

Let \(\mathcal X_{e,\alpha}\) be the exact quotient fields of Section
17.1, and define the globally smooth nonnegative function \[
 S(y)=\sum_{e,\alpha}(\mathcal X_{e,\alpha}F(y))^2.
\] The quotient identity intertwines the original sums of squared
derivatives, so \[
 R_g=\frac{2}{ag^2}S.
 \tag{290}
\] There is a finite constant \(C_S\), depending on the fixed box, mode
and cutoff, with \[
 S\le C_SW \quad\hbox{on the whole compact chord manifold}.
 \tag{291}
\] Here is the global justification. The original face potential has
expansion \(W(y)=\frac14\sum_\alpha\|Cy^\alpha\|^2+O(|y|^3)\). The
matrix \(C^*C\) is strictly positive on chord coordinates. Hence choose
a chart radius \(\rho>0\) and \(c_*>0\) with \(W(y)\ge c_*|y|^2\) for
\(|y|<\rho\). The smooth function \(F\) has zero value and zero
differential at zero; its differentiated square satisfies
\(S(y)\le C_*|y|^2\) there. For example bounded second derivatives of
\(F\), together with bounded coefficient vectors of the finitely many
\(\mathcal X_{e,\alpha}\), give such a finite \(C_*\) by the fundamental
theorem of calculus along the segment from zero. All faces flat on the
contractible open box imply all based chord holonomies are the identity:
elementary face moves identify the canonical tree paths, and the
original chord words then equal \(I\). Thus \(W\) has precisely one zero
on the chord manifold. Compactness gives \(W\ge c_\rho>0\) off the
radius-\(\rho\) neighborhood. One can take \[
 C_S=\max\left\{\frac{C_*}{c_*},
                 \frac{\max S}{c_\rho}\right\}.
\] This proves (291), including the region where the cutoff changes.

Put \(U_g=W/g^2\) and \(D_S=2C_S/a\). Equations (290)--(291) give \[
 0\le R_g\le D_SU_g.
 \tag{292}
\] The complete actual-vacuum moment proof in Section 17.4 gives, for
every integer \(n\ge0\), \[
 \int U_g^n\psi_g^2\,dU\le B_n,\quad
 B_0=1,\quad B_1=2aK,\quad
 B_{n+1}=2aK B_n+16n^2B_{n-1}\ (n\ge1),\quad
 K=\frac{3\sqrt{NM}}a.
 \tag{293}
\] Here \(N,M\) are the original numbers of edges and faces. These
estimates use the actual energy bound \(\mathcal E_g\le K\). Their
derivation retains the exact identity
\(q_H[f\psi]-\mathcal E_g\|f\psi\|^2
=\kappa\int\psi^2\sum|Xf|^2\), the full-potential bound
\(\sum|XW|^2\le16W\), and consequently
\(bM_{n+1}\le\mathcal E_g M_n+4\kappa n^2M_{n-1}\). Thus (293) is
available for arbitrarily high fixed powers, not merely the first vacuum
moment.

Let \(\mathcal B_g\) be the exact Haar-density and dilation isometry,
zero-extended outside \(\Omega_g\), and \(f_g=\mathcal B_g\psi_g\). For
an integer \(m\ge1\) and radius \(R>0\), split \(|x|>R\) into
\(|gx|<\rho\) and its complement. On the former region
\(U_g\ge c_*R^2\), and on the latter \(U_g\ge c_\rho/g^2\). Equations
(292)--(293) therefore give the explicit weighted squared tail \[
 \boxed{
 \int_{|x|>R}|\mathcal B_g(R_g^m\psi_g)(x)|^2\,dx
 \le D_S^{2m}B_{2m+1}
           \left(\frac1{c_*R^2}+\frac{g^2}{c_\rho}\right).
 }
 \tag{294}
\] In particular \(m=1\) uses \(B_3\), and \(m=2\) uses \(B_5\). No
assumption about a cutoff-region tail has been inserted.

\subsubsection{22.4. Exact coefficient and vacuum-vector
limit}\label{exact-coefficient-and-vacuum-vector-limit}

Write \[
 z^\alpha=O^{\mathsf T}G^{-1/2}x^\alpha,\qquad
 q(x)=\frac{\sigma_\nu}{4}|z_\nu|^2,\qquad
 \lambda=\frac{2\sigma_\nu}{a}.
\] On each fixed ball, \(\chi(gx)=1\) for small enough \(g\). The exact
tangent coefficients of the quotient fields have squared Gram matrix
\(G\otimes I_3\). It follows directly from (290) that the transported
multiplier \(r_g(x)=R_g(\exp(gx))\) satisfies \[
 r_g(x)\longrightarrow
 \frac2a\sum_\alpha(\nabla_{x^\alpha}q)^{\mathsf T}
                       G(\nabla_{x^\alpha}q)
 =\frac2a\frac{\sigma_\nu^2}{4}|z_\nu|^2
 =\lambda q(x).
 \tag{295}
\] This convergence is uniform on each fixed ball. The middle equality
uses the selected row of \(O^{\mathsf T}G^{-1/2}\) and
\(O^{\mathsf T}O=I\); it retains the original quotient kinetic metric.
The Haar-density factor commutes with multiplication, so it introduces
no omitted term into (295).

The actual strong vacuum limit \(f_g\to\Phi_0\), local coefficient
convergence (295), and the tail bound (294) prove \[
 \boxed{\mathcal B_g(R_g^m\psi_g)
            \longrightarrow(\lambda q)^m\Phi_0
       \quad\hbox{strongly in }L^2,\qquad m\ge1.}
 \tag{296}
\] Indeed convergence on a fixed ball follows by multiplying strongly
convergent vectors by uniformly convergent bounded multipliers there.
First taking \(g\downarrow0\), then \(R\to\infty\), controls the full
error using (294) and the polynomial times Gaussian tail of the limit.
In particular, with \(k=3/2\) and
\(d\rho(q)=q^{k-1}e^{-q}dq/\Gamma(k)\), \[
 \langle R_g\rangle_{\psi_g}\to\lambda k,\qquad
 \langle R_g^n\rangle_{\psi_g}\to\lambda^n(k)_n
 \quad(n\ge1).
 \tag{297}
\] One obtains each expectation by (296) and the strong unit vacuum
limit.

\subsubsection{22.5. Weighted transport on actual carrier time-orbit
vectors}\label{weighted-transport-on-actual-carrier-time-orbit-vectors}

For each fixed box and real \(\beta\), set \[
 m_g(\beta)=\langle\psi_g,e^{i\beta Q_g}\psi_g\rangle,\qquad
 v_g(\beta)=(e^{i\beta Q_g}-m_g(\beta))\psi_g,\qquad
 h_g(\beta,t)=e^{-tA}v_g(\beta),\quad t\ge0.
\] The fixed-box spectral theorem and bounded phase-vector convergence
give strong transport to the exact oscillator time vector \[
 h_0(\beta,t)=e^{-tA_0}
       (e^{i\beta q}-(1-i\beta)^{-k})\Phi_0.
 \tag{298}
\] This convergence is uniform for \(\beta,t\) in compact parameter
sets. A direct proof uses compactness of the comparison phase-vector
family in \(L^2\), a finite uniform net and finite oscillator spectral
cutoffs, convergence of the corresponding actual projections, and
uniform convergence of the finitely many exponential energy factors. The
initial phase vectors converge uniformly even over all phase parameters:
the actual multiplier equals \(e^{i\beta q}\) on a ball whose radius
tends to infinity, while both multipliers have modulus one.

The original positive heat semigroup satisfies \[
 |e^{-tA}(u\psi_g)|\le\|u\|_\infty\psi_g.
 \tag{299}
\] For completeness, the product-group electric heat kernels are
positive; the full real bounded Wilson multiplication has positive
exponential. The Trotter product and strong convergence preserve
positivity and \(|e^{-tA}u|\le e^{-tA}|u|\). The identity
\(e^{-tA}\psi_g=\psi_g\) then proves (299). Gauge restriction preserves
it. Thus \[
 |h_g(\beta,t)|\le2\psi_g,
 \tag{300}
\] uniformly for all real \(\beta\) and \(t\ge0\).

Multiply the squared tail bound (294) by four and use (300). The same
local convergence argument as in (296) proves, for every fixed integer
\(m\ge1\), \[
 \boxed{\mathcal B_g(R_g^m h_g(\beta,t))
       \longrightarrow(\lambda q)^m h_0(\beta,t)
       \quad\hbox{strongly, uniformly on compact }(\beta,t)\hbox{ sets}.}
 \tag{301}
\] The local part is uniform by (298) and the fixed-ball multiplier
bound; the tail is uniform by (294),(300). Finite linear combinations
inherit the same result. This proves the required polynomial powers on
these time-orbit vectors. It makes no assertion about arbitrary
interspersed unbounded products.

\subsubsection{22.6. A positive-coupling diagonal retaining the weighted
tests}\label{a-positive-coupling-diagonal-retaining-the-weighted-tests}

For \(j\ge2\), retain \[
 L_j=j^2,\quad a_j=\frac1{100j},\quad
 \sigma_j=\sqrt8\sin\frac{\pi}{4j^2+2},\quad
 \lambda_j=\frac{2\sigma_j}{a_j},\quad
 c_j=j\lambda_j\longrightarrow c_\Gamma=100\sqrt2\pi.
\] The preexisting diagonal selected only bounded phase/time kernels and
configuration flatness. Those tests alone do not establish (301) after
multiplication by a growing power of \(j\). Refine the positive dyadic
choice at stage \(j\) as follows. In addition to every previous
stage-\(j\) test and \(g<j^{-5}\), require \[
 \sup_{|\beta|\le j^2,\ 0\le t\le j}
 \left\|\mathcal B_g((jR_g)^m h_g(\beta,t))
                   -c_j^m q^m h_0(\beta,t)\right\|<\frac1j,
 \qquad 0\le m\le j,
 \tag{302}
\] The zero power explicitly includes unweighted vector transport; the
compact-family convergence of (298) proves this case. Also impose the
corresponding vacuum-vector tests \[
 \left\|\mathcal B_g((jR_g)^m\psi_g)
                   -c_j^m q^m\Phi_0\right\|<\frac1j
 \quad(1\le m\le j),\qquad
 |\langle R_g\rangle-\lambda_j k|<j^{-5}.
 \tag{303}
\] At each fixed \(j\), there are finitely many powers and each compact
supremum tends to zero by (296),(301). Thus all tests hold at every
sufficiently small positive \(g\); a dyadic exists. Take the least
positive dyadic exponent satisfying the entire list. This is a definite
refined sequence, not an assertion that an earlier least dyadic already
satisfies the new conditions. Its original coefficients remain
\(\kappa_j=200jg_j^2\), \(b_j=50j/g_j^2\), \(\xi_j=1/(4g_j^4)\), with
every face and scalar Wilson term present. Spectral tests for the
centered \(R_g\) vacuum vector used below can be included in this same
finite list, by the fixed-box argument in Section 22.8.

\subsubsection{22.7. The exact nonconstant carrier
operator}\label{the-exact-nonconstant-carrier-operator}

Fix \(b_0\ne0\), a real offset \(\alpha\), and put \(s_j=b_0\sqrt j\),
\(\beta_j=s_j+\alpha\). The exact oscillator radial map sends the vacuum
to \(1\) in \(L^2(\rho)\), and \[
 A_{0,j}=\lambda_jN,\qquad
 N=-q\partial_q^2-(k-q)\partial_q,\qquad
 U_s f=e^{isq}f,\qquad R_{0,j}=\lambda_jQ,\quad Qf=qf.
 \tag{304}
\] The identity on exponential functions \[
 e^{-uN}e^{i\beta q}
 =[1-i\beta(1-e^{-u})]^{-k}
  \exp\left(\frac{i\beta e^{-u}q}
                  {1-i\beta(1-e^{-u})}\right)
 \tag{305}
\] can be verified by differentiating in \(u\) and using (304). At
\(u=0\) it has the specified initial value. The differentiated functions
lie in the operator domain at finite \(u\): their derivatives are
polynomials times an exponential of nonpositive real part, hence
square-integrable for the Gamma density. More explicitly, multiplying by
the radial oscillator ground Gaussian and returning to the three
original selected-mode coordinates gives a smooth Schwartz function of
those coordinates, in the original oscillator operator domain; this
fixes the radial endpoint domain rather than inferring it from
integrability alone. Uniqueness of the self-adjoint semigroup evolution
proves (305). Its prefactor has modulus at most one, and the real part
of its exponential coefficient is \(-\beta^2e^{-u}(1-e^{-u})/
[1+\beta^2(1-e^{-u})^2]\le0\).

Take \(u=t\lambda_j\) in (305) and multiply by \(e^{-is_jq}\). Its
prefactor tends to one, while its exponential coefficient tends to
\(i\alpha-c_\Gamma b_0^2t\), because \[
 1-e^{-t\lambda_j}\sim t\lambda_j,\quad
 \beta_j(1-e^{-t\lambda_j})\to0,\quad
 s_j\beta_j(1-e^{-t\lambda_j})\to c_\Gamma b_0^2t.
 \tag{306}
\] For \(t=0\) the same limits are exact directly. The subtracted
centered term is \((1-i\beta_j)^{-k}e^{-is_jq}\); its \(q^m\)-weighted
norm is \((1+\beta_j^2)^{-k/2}\sqrt{(k)_{2m}}\to0\). The uncentered
expression and its limit have modulus at most one. Dominated convergence
against \(q^{2m}\rho\), whose integral is \((k)_{2m}\), proves \[
 q^m U_{s_j}^{-1}e^{-t\lambda_jN}
       \{e^{i\beta_jq}-(1-i\beta_j)^{-k}\}
 \longrightarrow q^m e^{i\alpha q-c_\Gamma b_0^2tq}
 \quad\hbox{in }L^2(\rho),\qquad m\ge0.
 \tag{307}
\] The estimates are uniform for bounded \(\alpha,t\). Equivalently one
can combine ordinary strong convergence and the domination bound with
the split \(R^{2m}\|\text{unweighted error}\|^2
+4\int_{q>R}q^{2m}d\rho\).

Combining (302),(307), the exact radial coordinate map and demodulation
therefore identifies the actual operator \(jR_g\) on finite carrier
time-orbit vectors with \[
 R_\infty=c_\Gamma Q,\qquad
 H_{b_0}=c_\Gamma b_0^2Q=b_0^2R_\infty .
 \tag{308}
\] For a precise common-coordinate statement, first transport by
\(\mathcal B_{g_j}\), then use the orthogonal mode/Gaussian radial
isometry on its comparison vector and demodulate by \(U_{s_j}^{-1}\).
The norm error between the actual transported vector and this comparison
vector tends to zero by (302). The other modes remain in their
normalized vacuum factors. Thus the varying number of other modes
introduces no unidentified Hilbert-space map in the stated Gram and
operator limits.

For nonnegative times \(u,v\), offsets \(\alpha,\delta\), and an integer
\(m\ge0\), the complete limiting insertion is \[
 \boxed{
 \left\langle e^{-uH_{b_0}}e^{i\alpha q},
            R_\infty^m e^{-vH_{b_0}}e^{i\delta q}\right\rangle
 =\frac{c_\Gamma^m(k)_m}
 {[1+c_\Gamma b_0^2(u+v)+i(\alpha-\delta)]^{k+m}} .
 }
 \tag{309}
\] This follows by integrating \(c_\Gamma^mq^m\) against the original
Gamma density. Its real denominator part is positive, fixing the complex
power. The same limit for actual insertions follows from (302),(307) and
Cauchy--Schwarz. In particular it retains their imaginary parts.

The multiplication operator \(c_\Gamma Q\) is nonconstant, positive and
self-adjoint on \(\{f:\int q^2|f|^2d\rho<\infty\}\): its nonreal
resolvents are the bounded multipliers \((c_\Gamma q-z)^{-1}\), and
testing the adjoint on compactly supported functions proves maximality
of this domain. The carrier time-orbit vectors, together with their
modulations, are dense. For example the functions \(e^{i\alpha q}\) are
total: orthogonality makes the Fourier transform of the finite density
\(\overline f\,d\rho\) zero. Extend that integrable density by zero to
the negative real axis, and convolve it with a Gaussian. The Gaussian
Fourier integral and Fubini express the convolution as the integral of
the zero Fourier transform times a Gaussian multiplier, so every such
convolution vanishes. Gaussian convolutions approach the original
density in \(L^1\): uniform continuity and Gaussian tails prove this
first for continuous compactly supported functions, and density with the
\(L^1\) contraction bound proves the general case. Thus \(f=0\). The
same argument for the finite measure \((1+q^{2m})d\rho\) proves that the
modulation span is a core for the maximal multiplication domain of
\(Q^m\). Formula (309) specifies the operator correspondence on the
tested polynomial domains, rather than only its expectation in one
vector. It does not assert convergence of arbitrary interspersed
unbounded products or identify a continuum local field algebra.

For the actual first carrier energy, (289),(303) give \[
 q_{A_j}[v_{g_j}(\beta_j)]
 =\beta_j^2\langle R_{g_j}\rangle
 \longrightarrow k c_\Gamma b_0^2.
 \tag{310}
\] Indeed \(\beta_j^2\lambda_j\to c_\Gamma b_0^2\), while
\(\beta_j^2|\langle R_{g_j}\rangle-\lambda_j k|
\le\beta_j^2j^{-5}\to0\). This proves the previously unavailable
unbounded first moment at physical time zero.

Finally, for \(s_j=b_0\sqrt j\), (288) gives the unchanged-regulator
realization \[
 jR_{g_j}
 =\frac{B_{s_j}^{-1}A_jB_{s_j}
        +B_{-s_j}^{-1}A_jB_{-s_j}-2A_j}{2b_0^2}.
 \tag{311}
\] It explains exactly in which sense this nonconstant multiplier is a
rescaled electric-energy observable.

\subsubsection{22.8. The vacuum image and its entire raw limiting
measure}\label{the-vacuum-image-and-its-entire-raw-limiting-measure}

The vacuum cannot be discarded when extending the observable algebra. At
each fixed \(j\), (296),(297) prove \[
 \mathcal B_g
   \{j(R_g-\langle R_g\rangle)\psi_g\}
 \longrightarrow c_j(q-k)\Phi_0.
 \tag{312}
\] The exact Laguerre identity \(N(q-k)=q-k\) and \(\int(q-k)^2d\rho=k\)
show that the complete comparison raw excitation measure of this vector
is \[
 k c_j^2\,\delta_{\lambda_j}.
 \tag{313}
\] This is a spectral-vector calculation: all other oscillator factors
are in their vacuum, and \(q-k\) is exactly its first radial Laguerre
eigenvector up to sign and norm.

The strong vector convergence (312) and finite spectral projection
convergence imply weak convergence of the actual raw measure to (313) at
fixed \(j\), including its total mass. To check tightness, choose a
finite comparison spectral cutoff containing this eigenvector. Its
limiting omitted mass is zero, while total masses converge by (312). The
actual omitted mass then tends to zero. Finite cluster projections and
their converging energies determine every bounded continuous test. This
repeats the exact argument used for the bounded-phase vectors, now with
the separately proved weighted strong vector convergence.

Add to stage \(j\)'s finite tests a bounded-Lipschitz error below
\(1/j\) between this actual measure and (313), and the corresponding
norm error below \(1/j\). These thresholds exist at every fixed box by
the preceding proof and can be included in the same positive dyadic
choice. Since \(\lambda_j\to0\) and \(c_j\to c_\Gamma\), the actual
centered vacuum images
\(w_j=j(R_{g_j}-\langle R_{g_j}\rangle)\psi_{g_j}\) then satisfy \[
 \boxed{\nu_{w_j}\Longrightarrow k c_\Gamma^2\delta_0,\qquad
 \|w_j\|^2\to k c_\Gamma^2>0,\qquad
 \langle w_j,e^{-tA_j}w_j\rangle\to k c_\Gamma^2\quad(t\ge0).}
 \tag{314}
\] Every finite \(w_j\) is exactly vacuum-orthogonal. It is nonzero:
\(R_g\) vanishes on an open region outside the cutoff support and is
strictly positive on an open set near a point where the selected
quadratic has nonzero derivative. Since \(\psi_g>0\), \(R_g\) cannot
have zero variance. This argument does not require a limiting norm to
establish finite-regulator nonzero states.

A representation retaining the vacuum \(\Omega\), this operator's
centered vacuum image \(w\), its inner products and its correlations
must have \(\langle\Omega,w\rangle=0\) and \(\|w\|^2=k c_\Gamma^2\). If
its Hamiltonian is nonnegative, (314) forces \(w\) into its kernel: for
any \(t>0\), \(\int(1-e^{-tE})d\nu_w(E)=0\), and positivity makes the
measure supported at \(E=0\). Thus it has a second independent
zero-energy vector. Constant autocorrelation without retained
orthogonality would not imply this conclusion; orthogonality is part of
the operator/vacuum map just specified.

There is an exact relation to the previously constructed origin
modulation sector. In \(L^2(\rho)\), \[
 \frac{e^{i\alpha q}-(1-i\alpha)^{-k}}{i\alpha}
       \longrightarrow q-k\qquad(\alpha\to0).
 \tag{315}
\] The inequality \(|(e^{i\alpha q}-1)/\alpha|\le q\) and the finite
second Gamma moment prove strong convergence of its first term;
differentiating the Gamma integral gives the scalar derivative \(k\).
Hence the centered vacuum image of \(R_\infty=c_\Gamma Q\) is the
precise derivative of the origin modulation state. Its zero-energy limit
is consistent with that sector's already computed vanishing
physical-time generator.

The positive carrier result (308)--(310) is therefore an actual
nonconstant observable extension beyond the scalar unscaled
configuration cylinders. Its vacuum extension is equally explicit and
has the zero-energy consequence (314). Neither statement alone
constructs the desired local interacting four-dimensional theory; they
specify the maps, domains, moments and the exact obstruction for this
particular retained observable/state family.

\subsubsection{22.9. Exact relation to the unchanged local
cylinders}\label{exact-relation-to-the-unchanged-local-cylinders}

The retained configuration bridge uses the same original face-filling
map from Section 13. If a fixed coarse index \(k_0\) divides \(j\),
replace a coarse edge by the ordered product of \(j/k_0\) fine links.
Its physical endpoints agree because \(a_j(j/k_0)n=a_{k_0}n\). No
noncommuting factors are reordered. On the cofinal subsequence \(j=n!\),
this gives every fixed coarse index eventually. The vacuum estimate in
(107), with \(m=j/k_0\) and \(\sqrt{N_jM_j}\le18\sqrt5j^6\), gives \[
\int\|I-U_p\|_{\rm HS}^2\,d(p_{k_0,j})_*\mu_j
\le\frac{216\sqrt5}{k_0^2}g_j^2j^8
\le\frac{216\sqrt5}{k_0^2j^2}.
\tag{316}
\] All elementary coarse faces therefore become flat in every
subsequential compact configuration limit. Equation (110) identifies all
such flat fields as \(U_e=q_{s(e)}^{-1}q_{t(e)}\), with one fixed root
value. They form one compact vertex-gauge orbit. Haar averaging on that
orbit is its unique invariant probability: the gauge average of a
continuous function is constant on the transitive orbit, and integration
against any invariant probability gives that same constant. Thus all
subsequential vacuum measures have the same limit. Every fixed
continuous gauge-invariant cylinder \(F_j=F\circ p_{k_0,j}\)
consequently satisfies \[
\|(F_j-F(I))\psi_j\|\longrightarrow0 .
\tag{317}
\] The full nonlinear semigroup domination (299) gives, for a finite
carrier time-orbit combination
\(w_j=\sum_\ell a_\ell e^{-t_\ell A_j}v_j(\beta_{j,\ell})\), \[
\|(F_j-F(I))w_j\|
\le2\sum_\ell|a_\ell|\,\|(F_j-F(I))\psi_j\|
\longrightarrow0 .
\tag{318}
\] The same inequality includes any finite collection of bounded phase
carriers. Their limiting cylinder representation is exactly
\(F\mapsto F(I)I\); this is a unital star representation because
evaluation preserves sums, products and conjugation. Its norm bound by
\(\|F\|_\infty\) extends it to the uniform closure.

Equations (308)--(315) now give a concrete observable extension beyond
(318): \(jR_g\) acts by the nonconstant \(c_\Gamma Q\) on a carrier
sector, has the complete polynomial insertion kernel (309), and its
centered vacuum image has the positive zero-energy raw measure (314).
The original phase coordinate uses a selected global transverse mode, so
this result does not identify \(R_g\) with a spatial local electric or
curvature field. No interchange of an unbounded multiplier with (318)
was used; its tails and weighted limits were proved separately. The
original native magnetic covariance and either prescribed
running-coupling trajectory remain different, unresolved continuum
calculations.

\subsection{23. Exact correspondence between local energy and radial
vacuum
images}\label{exact-correspondence-between-local-energy-and-radial-vacuum-images}

The complete
\href{sources/local_band_inputs/LOCAL_ENERGY_BAND_TRANSFER.md}{local-energy
first-band proof} is retained together with its
\href{sources/local_band_inputs/extensive_quantum_blocking.md}{preceding
definitions and proofs}. Its equation numbers are explicitly called
local-band source numbers below; the numbered equations of this
manuscript continue from (318). That source proves the actual fixed-box
local-energy graph limit, the seven-dimensional spectral frame, and the
local spatial Riemann limits. Section 22 supplies the phase-square
weighted limits. We now prove their exact common operator and state
maps, including a common positive-coupling sequence and every displayed
raw amplitude.

\subsubsection{23.1. One original finite-regulator
space}\label{one-original-finite-regulator-space}

Fix \(L\ge2\), \(a>0\), and \(g>0\). Retain the entire open box, every
edge and face, and the physical invariant Haar Hilbert space. Put \[
 H=\kappa\sum_e E_e+b\sum_p(2-\operatorname{tr}U_p),\qquad
 E_e=-\sum_{\alpha=1}^3X_{e,\alpha}^2,\qquad
 \kappa=\frac{2g^2}{a},\quad b=\frac1{2g^2a}.
\] Let \(\psi\) be the actual positive unit vacuum, with energy
\(\mathcal E\), and \(A=H-\mathcal E\). Write \(d\mu=\psi^2dU\). For
real edge weights \(f_e\) set \[
 f_p=\frac14\sum_{e\in\partial p}f_e,\qquad
 D_f=\kappa\sum_e f_eE_e+b\sum_p f_p(2-\operatorname{tr}U_p),
 \quad C_f=D_f-\langle\psi,D_f\psi\rangle,\quad \Xi_f=C_f\psi.
 \tag{319}
\] The signed electric sum has its Peter--Weyl self-adjoint domain: in
the joint Casimir expansion, the sum of the squared absolute weighted
eigenvalues times the squared coefficients must be finite. The bounded
self-adjoint magnetic perturbation makes \(D_f\) self-adjoint on that
same domain; no invariance of the domain under the multiplier is
asserted or needed. All identities below initially use smooth physical
functions; \(\psi\) and every finite spectral eigenvector are smooth.
Thus all displayed products on those vectors are defined, including when
unbounded operators occur.

Let \(Q=g^{-2}F\) be the globally smooth real invariant function of
Section 22, with its original cutoff. Put \[
 B_s=e^{isQ},\qquad
 R_g=\kappa\sum_{e,\alpha}(X_{e,\alpha}Q)^2,\qquad
 z_g=(R_g-\langle R_g\rangle_\psi)\psi.
 \tag{320}
\] At every fixed regulator \(R_g\) is a bounded smooth nonnegative
multiplier. It is not the local weighted differential operator \(D_f\).

There is nevertheless an exact common-operator relation. Define \[
 R_{g,f}^{Q}=\kappa\sum_{e,\alpha}f_e(X_{e,\alpha}Q)^2,\qquad
 J_f^Q=i[D_f,Q].
\] Applying each \(X\) twice to \(B_s v\) gives \[
 B_s^*D_fB_s=D_f+sJ_f^Q+s^2R_{g,f}^{Q},\qquad
 R_{g,f}^{Q}=\tfrac12[Q,[D_f,Q]],
 \tag{321}
\] where \[
 J_f^Q=-i\kappa\sum_{e,\alpha}f_e
 \{2(X_{e,\alpha}Q)X_{e,\alpha}+X_{e,\alpha}^2Q\}.
\] The face multiplication and its entire scalar commute with \(Q\) and
\(B_s\); this is why they contribute no commutator, not a deletion from
\(D_f\). For \(s\ne0\), (321) also proves \[
 R_{g,f}^{Q}
 =\frac{B_s^*D_fB_s+B_{-s}^*D_fB_{-s}-2D_f}{2s^2}.
 \tag{322}
\] For \(0\le f_e\le M\), \[
 0\le R_{g,f}^{Q}\le M R_g
\] pointwise. For arbitrary real weights its absolute value is at most
\(\|f\|_\infty R_g\). These follow term by term from the nonnegative
squares. When \(f_e=1\), \(D_f=H\) and \(R_{g,f}^{Q}=R_g\) exactly.

The additional exact commutator is \[
 [D_f,R_g]=-\kappa\sum_{e,\alpha}f_e
 \{(X_{e,\alpha}^2R_g)+2(X_{e,\alpha}R_g)X_{e,\alpha}\}.
 \tag{323}
\] It follows by the same product rule and again retains the full
\(D_f\). Equations (321)--(323) are concrete finite-regulator operator
maps between the two constructions, before any oscillator comparison.

For a regular compact fluid profile \(u(s,x)\), the retained local-band
source uses exactly
\(h_s(x)=\lambda_{\rm fl}^2|\operatorname{curl}u(s,x)|^2\), with
calibration \(\lambda_{\rm fl}\). The edge weights are
\(f_e=h_s(a\,m(e))\), with the original face averages. They are
nonnegative. Equations (319)--(323) apply to those very weights. The
fluid parameter \(s\) has not been identified with the phase parameter,
the quantum heat time, or the coupling regulator.

\subsubsection{23.2. Exact finite-frame overlap, with the Gram tensor
retained}\label{exact-finite-frame-overlap-with-the-gram-tensor-retained}

Use the retained local-band source's actual spectral projection \[
 P_g=\mathbf1_{(\Delta/2,\,11\Delta/10)}(H-\mathcal E),\qquad
 \Pi_g=|\psi\rangle\langle\psi|+P_g,\qquad
 \Delta=2\sigma/a .
\]

For sufficiently small \(g>0\) at a fixed box, \(P_g\) has rank six and
\(\Pi_g\) has rank seven. Retain \[
 {\cal K}=\mathbb C\oplus\operatorname{Sym}_3(\mathbb C),\qquad
 \langle(c,B),(d,E)\rangle_{\cal K}=\bar c d+6\operatorname{tr}(B^*E),
\] the exact frame \(M_g\), \(G_g=M_g^*M_g\), and inverse
\(G_g^{-1}M_g^*\) on the range, as in local-band source (133). Set
\(e_{\rm vac}=(1,0)\), and define \[
 x_g=M_g^{-1}P_g z_g,\qquad
 y_{g,f}=M_g^{-1}P_g\Xi_f.
\] Since \(M_ge_{\rm vac}=\psi\), \(C_f\) has zero vacuum expectation,
and \(\Pi_g\) projects onto the frame, \[
 y_{g,f}={\cal D}_{g,f}e_{\rm vac},\qquad
 \langle z_g,P_g\Xi_f\rangle
       =\langle x_g,G_g y_{g,f}\rangle_{\cal K}.
 \tag{324}
\] This is exact, because the part of \(z_g\) orthogonal to the band has
zero inner product with \(P_g\Xi_f\), and \(M_g^*M_g=G_g\). Likewise \[
 \|P_g z_g\|^2=\langle x_g,G_gx_g\rangle_{\cal K},\qquad
 \|z_g\|^2=\|P_gz_g\|^2+\|(I-\Pi_g)z_g\|^2.
 \tag{325}
\] The last equality uses \(\langle\psi,z_g\rangle=0\), which follows
from its actual expectation subtraction. No vector is divided by its
norm.

One can also write the complete unprojected raw overlap in the original
vacuum measure. Since \(\psi\) is strictly positive and smooth on a
compact space, \(d_f=(D_f\psi)/\psi\) is a smooth real physical
function. Then \[
 \langle z_g,\Xi_f\rangle
 =\int (R_g-\mu(R_g))(d_f-\mu(d_f))\,d\mu.
 \tag{326}
\] This gives its exact covariance; it does not replace \(d_f\) by the
classical fluid configuration or by a sampled classical energy.

\subsubsection{23.3. Exact identification of the selected oscillator
direction}\label{exact-identification-of-the-selected-oscillator-direction}

The two sources retain the same tree linear map \(T\), \(G=TT^*\),
\(C=d_1\jmath\), and \(B=G^{1/2}C^*CG^{1/2}\). To avoid a conflict with
the phase-square \(R_g\), call the edge-space isometry \[
 {\cal R}=T^*G^{-1/2}.
\] Direct multiplication gives \({\cal R}^*{\cal R}=I\). Its range is
\((\ker T)^\perp=(\operatorname{im}d_0)^\perp\), the original transverse
edge space. The exact identity \(d_1T^*G^{-1}=C\) gives
\(d_1{\cal R}=CG^{1/2}\), hence \[
 {\cal R}^*d_1^*d_1{\cal R}=B.
 \tag{327}
\] For the selected lowest eigenmode \(\nu\) of Section 22, its original
edge vector is \(v_\nu=\mathcal R Oe_\nu\). Thus it belongs to the same
lowest three-dimensional eigenspace spanned by the retained local-band
source's explicit orthonormal edge cochains \(V_1,V_2,V_3\). Define \[
 n_i=\langle V_i,v_\nu\rangle,\qquad
 n\in\mathbb R^3,\quad n^{\mathsf T}n=1,\qquad P_n=nn^{\mathsf T}.
 \tag{328}
\] Equations (327)--(328) prove the mode alignment without replacing an
unspecified orthogonal basis by an invented one. On each of the three
color copies,
\(a_{\nu,\alpha}^\dagger=\sum_i n_i a_{i,\alpha}^\dagger\).
Orthogonality of this mode change preserves the oscillator measure and
creation commutators.

Keep \(\lambda_{\rm osc}=2\sigma/a\) and \(k=3/2\);
\(\lambda_{\rm osc}\) is an energy, not \(\lambda_{\rm fl}\). With
\(q=\sigma|z_\nu|^2/4\) and
\(z_{\nu,\alpha}=\sqrt{2/\sigma}(a_{\nu,\alpha}+a_{\nu,\alpha}^\dagger)\),
direct expansion gives \[
 q=\tfrac12\sum_\alpha
 \{a_{\nu,\alpha}^2+(a_{\nu,\alpha}^\dagger)^2
                    +2a_{\nu,\alpha}^\dagger a_{\nu,\alpha}+1\}.
\] Therefore, with the retained local-band source's creation-vector map
\(\Phi\), \[
 (q-k)\Phi_0=\tfrac12\Phi(P_n),\qquad
 \|\tfrac12\Phi(P_n)\|^2=\tfrac64\operatorname{tr}(P_n^2)=\tfrac32.
 \tag{329}
\] The source's factor six is retained, including all three colors. This
exact identity is the relation between the radial Gamma vacuum-image
direction and the local first band. The phase direction is one rank-one
matrix in that six-dimensional symmetric-matrix band.

At fixed \(L,a\), the two source vector limits and projection
convergence give \[
 x_g\longrightarrow (0,\tfrac{\lambda_{\rm osc}}2P_n),\qquad
 y_{g,f}\longrightarrow(0,\tfrac1{4a}{\mathsf R}_f^{\min}),\qquad
 \|(I-\Pi_g)z_g\|\longrightarrow0.
 \tag{330}
\] The last assertion follows by applying the converging band projection
to the strong vector limit \(\lambda_{\rm osc}(q-k)\Phi_0\), which is
entirely in that band by (329).

Using the original inner product in (324) now proves \[
 \langle z_g,P_g\Xi_f\rangle
 \longrightarrow\frac{3\lambda_{\rm osc}}{4a}
                  n^{\mathsf T}{\mathsf R}_f^{\min}n,\qquad
 \|z_g\|^2\longrightarrow\tfrac32\lambda_{\rm osc}^2.
 \tag{331}
\] The unprojected overlap has the same limit: \(\Xi_f\) has a bounded
fixed-box norm by its proved strong graph limit, and the omitted part of
\(z_g\) tends strongly to zero. This controls the missing cross term by
Cauchy--Schwarz; it does not declare all local-energy leakage zero.

\subsubsection{23.4. The phase-square operator on the whole finite
spectral
frame}\label{the-phase-square-operator-on-the-whole-finite-spectral-frame}

A further fixed-box result follows from the retained local-band source's
excited-vector moment identity (131). Let \(u_g\) be an eigenvector in
the bounded finite spectral range, with bounded norm and
\(H\)-eigenvalue \(e_g\). For \(U_g=W/g^2\), positivity of the electric
part gives the initial moment bound below. Division of the moment
recursion by \(g^{2n+2}\) gives the subsequent bounds for integers
\(n\ge1\): \[
 \begin{aligned}
 \int U_g|u_g|^2&\le2ae_g\|u_g\|^2,\\
 \int U_g^{n+1}|u_g|^2
 &\le2ae_g\int U_g^n|u_g|^2+
       16n^2\int U_g^{n-1}|u_g|^2,\qquad n\ge1.
 \end{aligned}
 \tag{332}
\] Consequently every fixed moment is bounded at the fixed box,
uniformly over the bounded eigenvalue range. The global inequality
\(0\le R_g\le D_SU_g\) from Section 22 is pointwise and does not depend
on positivity of \(u_g\). For any fixed integer \(m\ge1\), it therefore
gives exactly the same squared tail estimate as (294), with its vacuum
moment \(B_{2m+1}\) replaced by the eigenvector moment bound in (332):
\[
 \int_{|x|>R}|{\cal B}_g(R_g^m u_g)|^2
 \le D_S^{2m}B^{\rm eig}_{2m+1}
       \left(\frac1{c_*R^2}+\frac{g^2}{c_\rho}\right).
 \tag{333}
\] Indeed, in the inner chart \(U_g\ge c_*R^2\); outside it,
\(U_g\ge c_\rho/g^2\). On each region, \(R_g^{2m}\) is bounded by
\(D_S^{2m}U_g^{2m}\). One extra power of \(U_g\) proves (333).

On every fixed ball the transported multiplier tends uniformly to
\(\lambda_{\rm osc}q\). Strong eigenvector convergence, multiplication
by these locally bounded coefficients, (333), and the polynomial
Gaussian tail of the comparison eigenvector give \[
 {\cal B}_g(R_g^m u_g)\longrightarrow
             (\lambda_{\rm osc}q)^m u_0 .
 \tag{334}
\] For the entire finite frame, choose an eigenbasis along a
subsequence, use the bounded finite rank and projection convergence to
extract convergent basis vectors, and apply (334) to each. The computed
limiting matrix is independent of this extraction; otherwise a sequence
failing matrix convergence would have a subsequence with the computed
limit, a contradiction. This proves finite-frame matrix convergence,
with all constants still depending on the fixed box. It is not a
uniform-volume theorem.

Write \(Z_g=R_g-\langle R_g\rangle_\psi\) and compress by the exact
frame. The resulting comparison operator on \(\mathcal K\) is \[
 {\cal Z}_0(c,B)=\lambda_{\rm osc}
 \left(3\operatorname{tr}(P_nB),
       \tfrac c2P_n+P_nB+BP_n\right).
 \tag{335}
\] Pair annihilation gives \(6\operatorname{tr}(P_nB)\), multiplied by
\(\lambda_{\rm osc}/2\). Pair creation of the vacuum gives
\(c\lambda_{\rm osc}P_n/2\). The number term has commutator
\(\Phi(P_nB+BP_n)\). These computations prove (335), including its
self-adjointness for the factor-six metric.

The compression has nonzero leakage. For any band vector \(\Phi(B)\), \[
 (I-\Pi_0)\lambda_{\rm osc}(q-k)\Phi(B)
 =\tfrac{\lambda_{\rm osc}}2\,{\cal S}_{P_n}{\cal S}_B\Phi_0,
 \tag{336}
\] where \(\mathcal S_B\) is the pair-creation polynomial defined in the
retained local-band source. The pair-annihilation and number terms stay
in \(\Pi_0\) and there are no higher spatial modes from \(P_n\). For
\(B=P_n\), this is nonzero: with three color oscillators, \[
 \left\|\left(\sum_\alpha a_\alpha^{\dagger2}\right)^2\Phi_0\right\|^2
 =4^2\,2!\,(3/2)_2=120.
\] For verification, its three terms \((a_\alpha^\dagger)^4\Phi_0\) have
total squared norm \(3\cdot4!=72\), and its three terms
\(2(a_\alpha^\dagger)^2(a_\beta^\dagger)^2\Phi_0\), \(\alpha<\beta\),
have total squared norm \(3\cdot4\cdot2!\cdot2!=48\). Different
occupation vectors are orthogonal. The leakage squared norm in (336) is
thus \(30\lambda_{\rm osc}^2\) for \(B=P_n\).

At finite regulator, inserting \(I=\Pi_g+(I-\Pi_g)\) gives \[
 \Pi_g Z_g C_f\Pi_g
 =\Pi_g Z_g\Pi_g C_f\Pi_g+
 ((I-\Pi_g)Z_g\Pi_g)^*((I-\Pi_g)C_f\Pi_g).
 \tag{337}
\] All range vectors are smooth, \(Z_g\) is a bounded smooth multiplier,
and \(C_f\) preserves smoothness, so every displayed product is defined.
The graph limits (334) and local-band source (132) identify each
fixed-box leakage matrix element by inner products of their full
limiting vectors. Thus neither compressed multiplication nor a
comparison of its matrices discards the term in (337).

\subsubsection{23.5. Ordered local-profile limit: one-third overlap and
its
constants}\label{ordered-local-profile-limit-one-third-overlap-and-its-constants}

Now use exactly \[
 L_j=j^2,\qquad a_j=\frac1{100j},\qquad
 \ell_j=a_j(2j^2+1),\qquad \lambda_j=\frac{2\sigma_j}{a_j},
 \qquad c_j=j\lambda_j\longrightarrow c_\Gamma=100\sqrt2\pi.
\] For a fixed real \(h\in C_c(\mathbb R^3)\), set
\(H_h=\int_{\mathbb R^3}h(x)\,dx\) and \(a_h=\sqrt2\pi H_h\). The
retained local-band source's proved matrix limit is \[
 \frac{\ell_j^4}{a_j}{\mathsf R}_{h,j}^{\min}
       \longrightarrow 4\sqrt2\pi H_h I_3.
 \tag{338}
\] Thus the two comparison frame vectors, with their displayed raw scale
factors, are \[
 jz_{0,j}=(0,\tfrac{c_j}{2}P_{n_j}),\qquad
 \ell_j^4v_{h,0,j}=(0,\tfrac{\ell_j^4}{4a_j}
                      {\mathsf R}_{h,j}^{\min})
                 =(0,a_h I_3+o(1)).
 \tag{339}
\] The unit vector \(n_j\) need not converge:
\(\operatorname{tr}(P_{n_j})=\operatorname{tr}(P_{n_j}^2)=1\) for every
\(j\), so all the following scalar conclusions are independent of that
basis freedom. Direct factor-six contractions give \[
 \begin{split}
 \|jz_{0,j}\|^2&\longrightarrow\tfrac32c_\Gamma^2
                         =30000\pi^2,\\
 \|\ell_j^4v_{h,0,j}\|^2&\longrightarrow36\pi^2H_h^2,\\
 \langle jz_{0,j},\ell_j^4v_{h,0,j}\rangle
                  &\longrightarrow3c_\Gamma a_h
                         =600\pi^2H_h.
 \end{split}
 \tag{340}
\] For \(H_h\ne0\), the squared cosine of their angle therefore tends to
\(1/3\). This is a computed fraction of raw Gram entries, not a
replacement of either original state by a unit vector.

The exact comparison projection of \(v_{h,0,j}\) onto \(z_{0,j}\) has
coefficient \[
 \frac{n_j^{\mathsf T}{\mathsf R}_{h,j}^{\min}n_j}
               {2a_j\lambda_j};
\] for the scaled vectors in (339) that coefficient tends to \(H_h/50\).
Consequently \[
 \|\ell_j^4v_{h,0,j}-(H_h/50)jz_{0,j}\|^2
                   \longrightarrow24\pi^2H_h^2.
 \tag{341}
\] Equations (340)--(341) also hold as the ordered actual limits
\(\lim_{j\to\infty}\lim_{g\downarrow0}\), by (330)--(331) at every fixed
\(j\). They have not been asserted along either independently selected
coupling sequence in the two sources.

For the actual regular NS profile, substitute
\(H_h=\lambda_{\rm fl}^2\|\operatorname{curl}u(s)\|_2^2\). This is
positive for a nonzero compact smooth incompressible profile: the
compact divergence/curl identity makes zero curl imply zero gradient,
hence \(u=0\). The leading local curvature state thus has a nonzero
overlap with the phase direction, and has the positive residual in
(341). This is an actual proved relation, rather than an inference of
nonrelation from the different presentations.

The residual has a concrete completion within the same lowest modes.
Define three phase quadratics using each of the original \(V_i\), with
the same admissible cutoff, and write \(z_g^{(i)}\) for their centered
phase-square vacuum vectors. Their fixed-box comparison matrices are
\((\lambda_{\rm osc}/2)e_i e_i^{\mathsf T}\). Their three pair vectors
are orthogonal. Summing them gives \((\lambda_{\rm osc}/2)I_3\). At
comparison level, \[
 \ell_j^4v_{h,0,j}-(H_h/50)\sum_{i=1}^3jz_{0,j}^{(i)}
                         \longrightarrow0.
 \tag{342}
\] No new spatial-locality identification is involved: these are three
global selected-mode phases. The same complete compressed operator
relation holds in the ordered comparison:
\(\sum_i j\mathcal Z_{0,j}^{(i)}\) acts as
\(c_j(3\operatorname{tr}B,cI_3/2+2B)\). Multiplication by \(H_h/50\)
makes its limit exactly the retained local-band source's \[
 \mathcal T_h(c,B)=\sqrt2\pi H_h
                 (6\operatorname{tr}B,cI_3+4B).
\] The proof is substitution into (335), with \(\sum_iP_i=I_3\). This
relates the entire seven-dimensional compressed operators, not merely
their vacuum columns. The leakage (337) remains present.

\subsubsection{23.6. Exact nonidentity and the joint state
errors}\label{exact-nonidentity-and-the-joint-state-errors}

For nonnegative local \(h\) with at least one positive sampled edge
weight, \(C_f\) is not any scalar multiple of the bounded multiplier
\(Z_g\) as an operator on the smooth physical core. Choose a plaquette
\(p\) incident to an edge \(e\) with \(f_e>0\), and set
\(\phi=\operatorname{tr}U_p\). There are configurations with
\(X_{e,\alpha}\phi\ne0\), as follows by varying that edge with all other
plaquette links fixed to the identity. The smooth function \(\phi\) is
gauge invariant. Applying \(D_f\) to \(e^{it\phi}\) and dividing by that
nonzero multiplier gives a polynomial in \(t\) whose quadratic
coefficient is \(\kappa\sum_{e,\alpha}f_e(X_{e,\alpha}\phi)^2\),
strictly positive somewhere. A fixed multiplication operator has no
\(t\)-dependence. Equality on every \(e^{it\phi}\) is therefore
impossible. The same argument applies after subtracting the vacuum
scalar.

This establishes the specific full-operator nonidentity while
(321)--(342) prove the strongest explicit relations derived here. The
three-phase completion does not remove the full-operator nonidentity or
the leakage, and does not identify a local field.

The two retained constructions choose different finite lists for their
least dyadic exponent. Fixed-box convergence alone does not prove that
either existing least exponent satisfies the other's new scaled tests.
The concrete joint state errors, which Section 23.8 controls on its
defined refined sequence, are in the common comparison coordinates \[
 \|{\cal B}_{g_j}(jz_{g_j,j})
                 -(c_j/2)\Phi(P_{n_j})\|\to0,\qquad
 \|\ell_j^4{\cal B}_{g_j}P_{g_j}\Xi_{f_{h,j}}
                 -\Phi(a_hI_3)\|\to0,
 \tag{343}
\] together with the already retained frame Gram/projection control.
These are the precise errors whose product estimates by Cauchy--Schwarz
would promote (340) to a single actual diagonal. Each test is attainable
at every fixed box by the preceding proofs, so their finite union can be
imposed in a newly defined refined dyadic selection. Section 23.8 below
defines that additional sequence and proves the common state map on it.
Neither preexisting sequence is declared unchanged by this refinement,
and no uniform-volume bound is inferred.

For products without intermediate projection one additionally needs the
joint scaled limit of the exact leakage in (337), including the retained
local-band source's higher spatial modes and four-creation terms. Their
fixed-box maps are explicit here and in local-band source (147)--(148);
their simultaneous volume limit is not supplied by the two source
diagonals. None of the calculations above assumes it.

\subsubsection{23.7. Unchanged bounded cylinders on the weighted phase
core}\label{unchanged-bounded-cylinders-on-the-weighted-phase-core}

This additional estimate uses the already selected sequence of Section
22; it requires no local-energy-band error estimate. Let \(F_j\) be the
pullback of one fixed continuous gauge-invariant coarse cylinder, let
\(\Delta F_j=F_j-F(I)\), and retain (317): \(\|\Delta F_j\psi_j\|\to0\).
Its uniform bound is \(|\Delta F_j|\le2\|F\|_\infty\). Therefore \[
 \int|\Delta F_j|^4\psi_j^2
 \le4\|F\|_\infty^2\|\Delta F_j\psi_j\|^2\longrightarrow0.
 \tag{344}
\] For each fixed nonnegative integer \(m\), (303) with power \(2m\)
gives a uniform bound on \(\|(jR_g)^{2m}\psi_j\|\); for \(m=0\), use the
unit vacuum norm instead. Write \(h_j\) for an actual carrier time-orbit
vector. Equation (300) gives \(|h_j|\le2\psi_j\). All the multipliers
commute, so Cauchy--Schwarz in the original measure gives \[
 \begin{split}
 \|\Delta F_j(jR_g)^m h_j\|^2
 &\le4\int|\Delta F_j|^2(jR_g)^{2m}\psi_j^2\\
 &\le4\left(\int|\Delta F_j|^4\psi_j^2\right)^{1/2}
               \|(jR_g)^{2m}\psi_j\|\\
 &\le8\|F\|_\infty\|\Delta F_j\psi_j\|\,
               \|(jR_g)^{2m}\psi_j\|\longrightarrow0.
 \end{split}
 \tag{345}
\] For a finite carrier combination, the bound
\(|h_j|\le2(\sum_\ell|a_\ell|)\psi_j\) replaces the factor four by
\(4(\sum_\ell|a_\ell|)^2\). For a weighted vacuum vector, the same proof
starts with factor one. Since \(j\langle R_g\rangle_\psi\) is bounded by
(303), it also proves the scalar action on
\(j(R_g-\langle R_g\rangle_\psi)\psi_j\) after subtracting its bounded
multiple of \(\psi_j\).

Thus the unchanged bounded cylinder \(F\) acts by exactly \(F(I)\) in
all the retained polynomial phase-square vector limits. This is a proved
extension of the scalar cylinder representation to those polynomial
cores, with the actual moment bounds used in (345). It is compatible
with (340)--(342): the local operator \(D_f\) is a differential operator
with coefficients \(\kappa f_e\) and the magnetic term \(bf_p\), and its
scaled first-band action is not a fixed bounded configuration cylinder.
The unbounded coefficient and projection operations cannot be replaced
by the scalar action in (345). Equations (321), (335), and (337) give
the exact relationships to that different operator class.

\subsubsection{23.8. A defined common dyadic sequence for every fixed
local
profile}\label{a-defined-common-dyadic-sequence-for-every-fixed-local-profile}

Define one common sequence using the finite union of the two source
requirements and the additional state tests below. It does not identify
the numerical least dyadic of either preexisting sequence with this new
one.

At fixed \(j\), retain the original finite edge set \(E_j\) and let
\(e^{(r)}\) be the edge weight that is one at \(r\) and zero elsewhere,
always with the prescribed quarter-face averages. Define the actual
comparison error map \[
 {\cal T}_{g,j}f
  ={\cal B}_gP_g C_{g,f}\psi_g
          -\Phi({\mathsf R}_{f,j}^{\min}/(4a_j)).
 \tag{346}
\] Here both terms are in the same full comparison Hilbert space,
including the original constant-Jacobian kinetic coordinate isometry
used to express \(\Phi\). The map is linear in the finite real weight
vector \(f\): \(D_f\) is linear, its expectation is linear, and
\(P_g,\psi_g\) do not depend on \(f\). The oscillator matrix is also
linear. The source graph and projection convergence imply
\(\|\mathcal T_{g,j}e^{(r)}\|\to0\) for every \(r\). Thus \[
 \sup_{\|f\|_\infty\le j}\|{\cal T}_{g,j}f\|
 \le j\sum_{r\in E_j}\|{\cal T}_{g,j}e^{(r)}\|
 \longrightarrow0\quad(g\downarrow0,\ j\hbox{ fixed}).
 \tag{347}
\] Every sum in (347) is finite. This proves the required uniformity
over the entire finite weight cube without assuming uniformity as its
dimension grows.

At each fixed \(j\), let \(\nu_1,\nu_2,\nu_3\) be the three original
lowest column indices of \(O\), in their existing order. They include
the original selected \(\nu_j\) of Section 22. Put \[
 n_{i,j}=V^*{\cal R}Oe_{\nu_i},\qquad
 P_{i,j}=n_{i,j}n_{i,j}^{\mathsf T},\qquad
 n_{i,j}^{\mathsf T}n_{r,j}=\delta_{ir},\qquad
 \sum_{i=1}^3P_{i,j}=I_3.
\] The identities follow from (327), the orthonormality of both
lowest-mode bases, and their equal three-dimensional ranges. For each
\(i=1,2,3\), use the same original cutoff and the quadratic
\(\sigma_j|z_{\nu_i}|^2/4\) for that original \(O\)-column to define
\(Q_g^{(i)}\), \(R_g^{(i)}\), and
\(z_g^{(i)}=(R_g^{(i)}-\langle R_g^{(i)}\rangle)\psi_g\). These
constructions preserve the original tree and Gram metric: \(O\) is not
changed. Its exact comparison in the \(V\)-basis is (327)--(328), acting
identically on all three colors. The proof of their fixed-box weighted
vacuum limits uses the same established steps, with this selected
quadratic: it and its first differential vanish at the unique minimum,
its derivative square \(S_i\) is bounded by a fixed-box constant times
\(W\), and its local quadratic coefficient is \(\lambda_jq_i\). The
source moment recursion and tail proof then give \[
 {\cal B}_g(jz_g^{(i)})
             \longrightarrow(c_j/2)\Phi(P_{i,j}).
 \tag{348}
\] This is a proved application to three explicitly specified
quadratics; it is not an assumption of new weighted convergence. One is
exactly the original selected-mode quadratic, so its phase-square
direction has not been replaced. The earlier sum in the \(V\)-basis in
(342) has this same comparison because both sums of the three rank-one
projectors are \(I_3\).

At stage \(j\) choose the least positive integer \(n_j^\dagger\) for
which \(g_j^\dagger=2^{-n_j^\dagger}<j^{-5}\) satisfies all original
stage-\(j\) tests in both sources and, in addition, \[
 \begin{split}
 \ell_j^4 j\sum_{r\in E_j}
       \|{\cal T}_{g_j^\dagger,j}e^{(r)}\|&<1/j,\\
 \|{\cal B}_{g_j^\dagger}(jz_{g_j^\dagger}^{(i)})
            -(c_j/2)\Phi(P_{i,j})\|&<1/j
                 \quad(i=1,2,3),\\
 j\|{\cal Z}_{g_j^\dagger}^{(i)}
                -{\cal Z}_{0,j}^{(i)}\|_{\mathcal K}&<1/j
                 \quad(i=1,2,3),\\
 \|(I-P_{g_j^\dagger})jz_{g_j^\dagger}^{(i)}\|&<1/j
                 \quad(i=1,2,3).
 \end{split}
 \tag{349}
\] The original source compact-family tests are functions of \(g\) that
tend to zero at fixed \(j\), as proved there; the number of powers, edge
basis vectors, and additional quadratics is finite at that \(j\). Here
\[
 {\cal Z}_g^{(i)}
 =G_g^{-1}M_g^*\Pi_g
       (R_g^{(i)}-\langle R_g^{(i)}\rangle)\Pi_gM_g,
\] and \({\cal Z}_{0,j}^{(i)}\) is (335) with \(P_n=P_{i,j}\). Equations
(334)--(335) prove the finite-frame test tends to zero. Equations
(347)--(348) prove the vector tests tend to zero. For the last test,
(348)'s limit lies wholly in the comparison first band, and the
transported actual band projection converges in operator norm at fixed
\(j\). Applying the two projections to the strongly converging vector
proves that its omitted norm tends to zero. Thus the last test is
attainable as well. Their finite union therefore holds for every
sufficiently small positive \(g\). The dyadics tend to zero, so the set
of admissible positive exponents is nonempty and has a least element.
This proves existence of the sequence just defined. No uniform-volume
estimate has been assumed. Its coefficients are precisely
\(\kappa_j=200j(g_j^\dagger)^2\) and \(b_j=50j/(g_j^\dagger)^2\),
including every face and the original Wilson scalar.

Use daggers to distinguish all actual states on this sequence: \[
 v_{h,j}^\dagger=P_{g_j^\dagger}\Xi_{f_{h,j}},\qquad
 w_{i,j}^\dagger=jz_{g_j^\dagger}^{(i)},\qquad
 w_{n,j}^\dagger=jz_{g_j^\dagger}.
\] For any fixed real \(h\in C_c(\mathbb R^3)\),
\(\|f_{h,j}\|_\infty\le\|h\|_\infty\) is at most \(j\) eventually. The
first test in (349) and (347) then bound its scaled actual-to-comparison
error by \(1/j\). No finite dense list of profiles is required: the
entire finite weight cube was controlled on the same sequence.

The source's Riemann matrix limit (338), (348)--(349), and the isometry
of the common coordinate maps prove \[
 \boxed{\left\|\ell_j^4v_{h,j}^\dagger
          -\frac{H_h}{50}\sum_{i=1}^3 w_{i,j}^\dagger
         \right\|\longrightarrow0
          \quad\hbox{for every fixed }h\in C_c(\mathbb R^3).}
 \tag{350}
\] Explicitly its transported norm is at most \((1+3|H_h|/50)/j\) plus
the norm of
\(\Phi(\ell_j^4\mathsf R_{h,j}^{\min}/(4a_j)-H_hc_jI_3/100)\). The last
matrix tends to zero because \(H_hc_\Gamma/100=\sqrt2\pi H_h\). The
factor-six formula makes matrix convergence equivalent to convergence of
this final creation-vector norm. This proves (350) completely.

The same actual-to-comparison estimates prove on this defined single
sequence \[
 \begin{split}
 \|w_{n,j}^\dagger\|^2&\longrightarrow30000\pi^2,\\
 \|\ell_j^4v_{h,j}^\dagger\|^2&\longrightarrow36\pi^2H_h^2,\\
 \langle w_{n,j}^\dagger,\ell_j^4v_{h,j}^\dagger\rangle
                                  &\longrightarrow600\pi^2H_h,\\
 \|\ell_j^4v_{h,j}^\dagger-(H_h/50)w_{n,j}^\dagger\|^2
                                  &\longrightarrow24\pi^2H_h^2.
 \end{split}
 \tag{351}
\] Indeed each comparison vector in (339) has bounded norm, so the error
in every inner product is bounded by the two vector errors times those
bounds, plus their product. The computed constants (340)--(341) then
give (351). For \(H_h\ne0\), the squared angular fraction is \(1/3\).
For \(H_h=0\), (350) gives the vanishing scaled local vector; no
division by that norm is made.

All these statements retain the raw original local state: it has norm of
order \(\ell_j^{-4}\) when \(H_h\ne0\). The explicit scale \(\ell_j^4\)
records its comparison amplitude; it is not a unit normalization. Each
\(w_{i,j}^\dagger\) likewise retains its definition
\(j(R_g^{(i)}-\langle R_g^{(i)}\rangle)\psi_g\) and its nonzero raw
limiting norm. No equality of the fluid time and quantum time was
introduced. The original phase sequence's bounded-cylinder tests are in
the finite union, so the polynomial-core scalar action (345) persists on
this defined refined sequence. Mixed products without \(\Pi_g\) still
retain the leakage (337); (350) by itself does not determine its
simultaneous volume limit.

The added finite-frame test also promotes the ordered compressed
operator relation following (342) to this same sequence: \[
 \left\|\ell_j^4{\cal D}_{g_j^\dagger,f_{h,j}}
       -\frac{H_h}{50}\sum_{i=1}^3
                       j{\cal Z}_{g_j^\dagger}^{(i)}
 \right\|_{\mathcal K}\longrightarrow0.
 \tag{352}
\] The first term tends to \(\mathcal T_h\) by the retained local-band
source's edge-basis and Gram tests, which are included in (349). For the
second term its error from the comparison sum is at most
\(3|H_h|/(50j)\), by (349). Equations (335) and \(\sum_iP_{i,j}=I\) make
that comparison sum \((H_hc_j/50)(3\operatorname{tr}B,cI_3/2+2B)\),
which tends exactly to \(\mathcal T_h\) because
\(c_\Gamma=100\sqrt2\pi\). This proves (352), retaining the full actual
Gram frame. It concerns compressed operators; the exact full-product
term (337) is still present.

\subsection{24. The complete local directional electric
spectrum}\label{the-complete-local-directional-electric-spectrum}

The following proof retains the actual Hamiltonian and vacuum while
calculating the original direction-one electric observable. It gives its
anisotropic spectral density and every positive-time moment on one
strengthened positive-coupling sequence.

\subsubsection{24.1 Exact original objects and fixed-box graph
map}\label{exact-original-objects-and-fixed-box-graph-map}

For \(L\ge2\), the original vertices are \(\{-L,\ldots,L\}^3\), with
every contained positive edge and face. For \(a,g>0\), write \[
 H_{L,a,g}=\kappa\sum_eE_e+bW,\qquad
 E_e=-\sum_{\alpha=1}^3X_{e,\alpha}^2,\qquad
 W=\sum_p(2-\operatorname{tr}U_p),\qquad
 \kappa=\frac{2g^2}{a},\quad b=\frac1{2g^2a}.
 \tag{353}
\] Here \(T_\alpha=-i\sigma_\alpha/2\) and the \(X_{e,\alpha}\) are the
original link derivatives. The full Wilson scalar is included in \(bW\).
The actual positive unit physical vacuum is \(\psi_{L,a,g}\), with
actual energy \(\mathcal E_{L,a,g}\); put
\(A_g=H_{L,a,g}-\mathcal E_{L,a,g}\ge0\). The original vertex gauge law
and its inverse-parameter presentation have the same invariant Hilbert
subspace.

Fix real \(h\in C_c^\infty(\mathbb R^3)\) and define \[
 f_{h,L,a}(n,i)=\mathbf1_{\{i=1\}}h(a(n+\mathbf e_1/2)),\qquad
 \Gamma_h=\sum_{e\parallel1}h(am(e))E_e,
 \tag{354}
\] \[
 D^E_h=\kappa\Gamma_h,\qquad
 \Xi^E_{h,g}=(D^E_h-\langle\psi_g,D^E_h\psi_g\rangle)\psi_g.
 \tag{355}
\] All sums use the original finite edge set, including the boundary.
There is no face-average term in \(D^E_h\). For signed weights, the
joint Peter--Weyl decomposition realizes \(\Gamma_h\) as the real
multiplier \(\sum_e f_h(e)j_e(j_e+1)\), on the domain defined by the
squares of those eigenvalues. This gives its self-adjoint realization;
smooth functions lie in its domain. The actual vacuum is smooth, so
(355) is smooth, physical, exactly vacuum-orthogonal, and in every power
domain of the full \(H\).

Retain the exact additive tree-to-chord map \(T\), with \(G=TT^*>0\),
the chord insertion \(\iota\), and the original oriented curl
\(C=d_1\iota\). If \[
 O^*G^{1/2}C^*CG^{1/2}O=\Sigma^2,\qquad
 \Sigma=\operatorname{diag}(\sigma_\nu),\qquad
 V_\nu=T^*G^{-1/2}Oe_\nu,
 \tag{356}
\] then the \(V_\nu\) are a real orthonormal transverse edge basis. Set
\[
 D^h_{\nu\eta}=\sum_{e\parallel1}h(am(e))V_\nu(e)V_\eta(e),
 \qquad \omega_\nu=\frac{\sigma_\nu}{a}.
 \tag{357}
\] No kinetic metric or counting inner product has changed.

The exact Haar/log map on based chord functions is \[
 (\mathcal B_gF)(x)
 =g^{3r/2}\mathcal J(gx)^{1/2}F(\exp(gx))
 \quad(x\in\Omega_g),\qquad
 \Omega_g=\{x:|x_c|<2\pi/g\},
 \tag{358}
\] and is zero outside \(\Omega_g\), with \(\mathcal J(y)=(16\pi^2)^{-r}
\prod_c[\sin(|y_c|/2)/(|y_c|/2)]^2\). The kinetic-coordinate unitary is
\[
 f(x)\longmapsto(\det G)^{3/4}f(G^{1/2}Oz).
 \tag{359}
\] We include both (358) and (359) whenever comparison vectors are
written in \(z\) coordinates, and denote their composition by
\(\mathcal C_g\). These are isometries with their exact image domains;
neither identifies the finite-\(g\) vacuum with a trial Gaussian.

Here is the Section 17 graph argument needed for (355). Its actual
vacuum moment calculation gives \[
 \begin{gathered}
 \int W^n\psi_g^2\,dU\le B_ng^{2n},\qquad
 B_0=1,\quad B_1=2aK,\\
 B_{n+1}=2aKB_n+16n^2B_{n-1}\quad(n\ge1),\qquad
 K=3\sqrt{N_EM_F}/a .
 \end{gathered}
 \tag{360}
\] where \(N_E,M_F\) are the original edge and face counts. It follows
from the full ground-state equation, the bound
\(\sum_{e,\alpha}|X_{e,\alpha}W|^2\le16W\), and the ground-state form
identity. At fixed box the potential has its unique zero at the identity
chords, with \(W(y)\ge c|y|^2\) locally and \(W\ge c_\rho>0\) outside
that neighborhood. The third scaled moment in (360) controls the squared
tail of \((W/g^2)\psi_g\) by \(B_3/(cR^2)+g^2B_3/c_\rho\). The full
Taylor coefficient and strong actual-vacuum convergence therefore give
\[
 \mathcal B_g((W/g^2)\psi_g)
 \longrightarrow \tfrac14\sum_\alpha\|Cx^\alpha\|^2\Phi_0.
 \tag{361}
\] The unchanged eigen-equation is \[
 g^2\sum_eE_e\psi_g
 =\frac{a\mathcal E_g}{2}\psi_g-\frac{W}{4g^2}\psi_g.
 \tag{362}
\] Together with (361), it proves convergence in the total electric
graph. Since the separate positive link Casimirs strongly commute, \[
 \|\Gamma_hF\|\le\|f_h\|_\infty\|\textstyle\sum_eE_eF\|.
 \tag{363}
\] Apply (363) to the difference between the actual vacuum and the
inverse image of a compact invariant cutoff times its limiting Gaussian.
On that compact chart the exact differentiated coefficients converge.
Let the cutoff radius tend to infinity using the Gaussian graph norm.
This proves the arbitrary signed weight graph limit, with all constants
fixed before \(g\to0\): \[
 \mathcal C_g(g^2\Gamma_h\psi_g)
 \longrightarrow
 -\sum_{\nu,\eta,\alpha}D^h_{\nu\eta}
          \partial_{z_\nu^\alpha}\partial_{z_\eta^\alpha}\Phi_0.
 \tag{364}
\] It uses the full potential moments, rather than inferring unbounded
operator convergence from vacuum norm convergence.

Use the original creation operators and variance, \[
 \partial_{z_\nu^\alpha}
 =\sqrt{\sigma_\nu/8}\,(a_{\nu\alpha}-a_{\nu\alpha}^\dagger),
 \quad
 z_\nu^\alpha=\sqrt{2/\sigma_\nu}
                    (a_{\nu\alpha}+a_{\nu\alpha}^\dagger),
 \quad
 A_0=\sum_{\nu,\alpha}\omega_\nu
                      a_{\nu\alpha}^\dagger a_{\nu\alpha}.
 \tag{365}
\] Subtracting the expectation in (364) removes its scalar term. Since
\(\kappa/g^2=2/a\), the exact fixed-box vector limit is \[
 \boxed{\mathcal C_g\Xi^E_{h,g}\longrightarrow
 X^E_{h,0}
 =-\frac14\sum_{\nu,\eta,\alpha}
     \sqrt{\omega_\nu\omega_\eta}\,D^h_{\nu\eta}
          a_{\nu\alpha}^\dagger a_{\eta\alpha}^\dagger\Phi_0.}
 \tag{366}
\] The sign is negative. Indeed
\((a_\nu-a_\nu^\dagger)(a_\eta-a_\eta^\dagger)\Phi_0
=-\delta_{\nu\eta}\Phi_0+
a_\nu^\dagger a_\eta^\dagger\Phi_0\), and the electric operator has the
preceding minus sign.

Define the real symmetric pair matrix \[
 B^{E,L,a}_{h;\nu\eta}
 =-\sqrt{\omega_\nu\omega_\eta}
       \sum_{e\parallel1}h(am(e))V_\nu(e)V_\eta(e).
 \tag{367}
\] There are three equal-colour contractions. The norm of a symmetric
two-creation kernel is twice its squared Hilbert--Schmidt norm:
commuting its two annihilators through the two creators gives the two
ordered contractions. Consequently the complete fixed-box raw measure is
\[
 \boxed{\nu^{E,L,a}_{h,0}
 =\frac38\sum_{\nu,\eta}|B^{E,L,a}_{h;\nu\eta}|^2
          \delta_{\omega_\nu+\omega_\eta}.}
 \tag{368}
\] The sum is ordered, including both occurrences of distinct modes. The
coefficient is exactly \(3\cdot2/4^2=3/8\). The strong vector theorem
(366), convergence of all finite physical spectral projections and
energies, and convergence of the total vector norms prove weak
convergence of actual finite measures to (368) at fixed box. A finite
comparison spectral cutoff contains all but a chosen tail of the
comparison vector; projection and norm convergence give the same bound
for the actual vector. Within it, only finitely many clusters occur and
their projections and energies converge. This proves all bounded
continuous tests.

\subsubsection{24.2 Every open-box mode and the directional parity
calculation}\label{every-open-box-mode-and-the-directional-parity-calculation}

Put \(N=2L+1\) and \(\ell=aN\). The physical endpoints are still
\(\pm aL\). On vertices and positive edges respectively use \[
 v_j(n)=\sqrt{\frac{2-\delta_{j0}}N}
       \cos\frac{\pi j(n+L+1/2)}N,\quad
 s_j=2\sin\frac{\pi j}{2N},\quad 0\le j<N,
 \tag{369}
\] \[
 w_j(n)=-\sqrt{\frac2N}
       \sin\frac{\pi j(n+L+1)}N,\qquad j>0.
 \tag{370}
\] There is no \(w_0\). Finite geometric sums give orthonormality on the
unchanged vertex and edge ranges. Direct differencing gives
\(d_0v_j=s_jw_j\), and transposition gives \(d_0^*w_j=s_jv_j\).

For \(\boldsymbol j=(j_1,j_2,j_3)\), let \(I=\{i:j_i>0\}\) and
\(s=(s_{j_1},s_{j_2},s_{j_3})\). Each real unit
\(\epsilon\in\mathbb R^I\) with \(s\cdot\epsilon=0\) gives \[
 V_{\boldsymbol j,\epsilon}(n,i)
 =\epsilon_iw_{j_i}(n_i)\prod_{d\ne i}v_{j_d}(n_d),
 \qquad \sigma=|s|.
 \tag{371}
\] There are \(|I|-1\) polarizations for \(|I|\ge2\) and none otherwise.
The tensor one-form space at this index is \(\mathbb R^I\); its gradient
line is \(\mathbb Rs\). The exact curl squared form is
\(|s|^2|\epsilon|^2-|s\cdot\epsilon|^2\). Thus its transverse complement
is exactly (371), proving completeness and the original frequencies.
Orthogonal changes inside eigenspaces give the same ordered spectral sum
(368).

Each mode component is bounded by \(2\sqrt2N^{-3/2}\). Choose \(R_h\) so
that a fixed unit neighborhood of the support of \(h\) lies in
\([-R_h,R_h]^3\). For \(a\le1\), the number of direction-one edge
samples meeting this support is at most \((2R_h/a+3)^3\). Thus (367)
satisfies the global bound \[
 |B^{E,L,a}_{h;\nu\eta}|
 \le K_h\ell^{-3}\sqrt{\omega_\nu\omega_\eta},\qquad
 K_h=8(2R_h+3)^3\|h\|_\infty .
 \tag{372}
\] This constant does not depend on \(L,a\) once the box contains the
fixed neighborhood. Every edge sample and zero weight remains in the
original sum.

The sine bounds give \[
 \frac{2|\boldsymbol j|}{\ell}
      \le\omega_{\boldsymbol j}
      \le\frac{\pi|\boldsymbol j|}{\ell}.
 \tag{373}
\] For \(\ell\ge1\) and each \(u>0\), grouping triples in shells
\(m\le|\boldsymbol j|<m+1\) proves \[
 \ell^{-3}\sum_\nu\omega_\nu e^{-u\omega_\nu}\le C_u.
 \tag{374}
\] The shell has at most \(C(m+1)^2\) triples, and each has at most two
polarizations. One justification for this shell count is to surround
each nonnegative integer triple by its disjoint unit cube: their union
lies between spheres whose radii differ by a fixed constant. Their
volume difference is bounded by \(C(m+1)^2\). The resulting series is
bounded by a constant times
\(\ell^{-4}\sum_{m\ge1}(m+1)^3e^{-2um/\ell}\); integral comparison
bounds it uniformly for \(\ell\ge1\).

Combining (372),(374), the contribution of \(\omega_\nu+\omega_\eta>R\)
to (368) after multiplication by \(e^{-t(\omega_\nu+\omega_\eta)}\) is
at most \(C_{h,t}e^{-tR/2}\), uniformly in the regulators. Every fixed
additional power of the pair energy is handled by using a smaller
positive time. The same shell proof gives \[
 \limsup_{\ell\to\infty}\ell^{-3}
       \sum_{\omega_\nu<\varepsilon}\omega_\nu
       \le C\varepsilon^4.
 \tag{375}
\] In a bounded frequency region the triples with a zero coordinate
number \(O_R(\ell^2)\), so their paired contribution is
\(O_{h,R}(\ell^{-1})\). Strips of width \(\delta\) adjoining a
coordinate plane have limiting contribution \(O_{h,R}(\delta)\). These
bounds justify retaining a bounded momentum region away from its
coordinate planes and zero before taking its Riemann sum, then removing
these restrictions.

For positive index coordinates set \(k_i=\pi j_i/\ell\),
\(\beta_i=j_i\bmod2\) and
\(\chi_{\boldsymbol j}=(-1)^{\sum_i\lfloor j_i/2\rfloor}\). At the
physical location of each mode component the exact standing-wave
identity is \[
 \ell^{3/2}a^{-3/2}V_{\boldsymbol j,\epsilon}(x)
 =\chi_{\boldsymbol j}\sum_{\rho\in\{-1,1\}^3}
       U_{\beta\rho}\epsilon^\rho e^{i(\rho k)\cdot x},
 \quad
 U_{\beta\rho}=\frac{i}{\sqrt8}e^{i\pi\rho\cdot\beta/2},
 \quad \epsilon_i^\rho=\rho_i\epsilon_i.
 \tag{376}
\] The common \(x\) denotes the component expressions restricted to
their original edge positions in the sum. No sample is moved. Expand
each cosine and negative sine in (369)--(371) to verify (376); the
differentiated component contributes \(i\rho_i\) relative to the two
undifferentiated components. Also \[
 \sum_{\beta\in\{0,1\}^3}
 e^{i\pi(\rho-\rho')\cdot\beta/2}
 =8\mathbf1_{\{\rho=\rho'\}},
 \tag{377}
\] so the eight-by-eight matrix \(U\) is unitary.

On bounded frequencies, \(s/a\to k\) uniformly. The directional edge sum
in (367) is a Riemann sum for the product of the two direction-one
components and \(h\); smoothness and compact support bound its sample
error by \(O_{h,R}(a)\) after its volume factor. There are
\(O_R(\ell^6)\) mode pairs, and the coefficient error is
\(o(1)\ell^{-3}\) per pair. Equation (372) therefore bounds the error in
the squared sum by \(o(1)\). One may sum polarizations using
\(I-kk^*/|k|^2\), avoiding any singular choice of frame; local
continuous polarization frames on the retained compact sets give the
same statement.

For each fixed parity the momentum grid has step \(2\pi/\ell\). Align
its eight parity grids with one grid; their shifts are at most
\(\pi/\ell\) in each coordinate and compact-set uniform continuity makes
the change tend to zero. Equation (377), applied to both mode indices,
then turns the two standing-wave parity sums into the Hilbert--Schmidt
sum over the two momentum signs. The signed factors
\(\chi_{\boldsymbol j}\) cancel in the modulus squared, not in the
original individual modes. The two grid densities and the squared
coefficient give precisely \(\ell^{-6}(\ell/(2\pi))^6=(2\pi)^{-6}\).
Their sign grids exhaust momentum space after removing only the
controlled coordinate strips.

For any measurable real transverse polarization frames \(e_r(p)\), the
resulting plane-wave coefficient is \[
 \boxed{
 \mathcal R^{E,rs}_h(p,q)
 =-\frac{\widehat h(-p-q)}{(2\pi)^3}
       \sqrt{|p||q|}\,(e_r(p))_1(e_s(q))_1,\qquad
 \widehat h(k)=\int_{\mathbb R^3}e^{-ik\cdot x}h(x)\,dx .
 }
 \tag{378}
\] The minus sign is the electric pair sign in (366). Polarization
completeness yields exactly \[
 \sum_r(e_r(p))_1^2=1-\frac{p_1^2}{|p|^2},\qquad
 \sum_{r,s}|(e_r(p))_1(e_s(q))_1|^2
 =\left(1-\frac{p_1^2}{|p|^2}\right)
  \left(1-\frac{q_1^2}{|q|^2}\right).
 \tag{379}
\] Values at zero momentum are irrelevant to the integral. The preferred
direction one has been retained. In particular the coefficient at
\(q=-p\) is generally nonzero. There is no magnetic pair term in this
observable to cancel it.

It follows, with all cutoffs removed by (372)--(375), that for any
\(a\downarrow0\), \(\ell=a(2L+1)\to\infty\) and \(t>0\), \[
 \boxed{
 C^{E,L,a}_{0;h,h}(t)\longrightarrow C^E_h(t)
 =\frac{3}{8(2\pi)^6}
 \int_{\mathbb R^3}\!\int_{\mathbb R^3}
 |p||q|\left(1-\frac{p_1^2}{|p|^2}\right)
       \left(1-\frac{q_1^2}{|q|^2}\right)
 |\widehat h(p+q)|^2e^{-t(|p|+|q|)}\,dp\,dq .
 }
 \tag{380}
\] All energy moments with the same positive-time damping converge by
the proved uniform tails. Real mixed covariances follow by polarization.
This is the directional electric counterpart of the incoming all-mode
result, with a different retained pair kernel.

\subsubsection{24.3 Exact anisotropic
density}\label{exact-anisotropic-density}

Write \(k=p+q\), \(K=|k|\), and \(k_\perp^2=K^2-k_1^2\). The raw locally
finite measure in (380) is absolutely continuous: \[
 \boxed{
 \rho^E_h(\omega)=\frac1{2560\pi^5}
 \int_{|k|<\omega}|\widehat h(k)|^2
 \left[(\omega^2-K^2)^2+
       (\omega^2-K^2)k_\perp^2+k_\perp^4\right]\,dk,\qquad
 d\nu^E_h(\omega)=\rho^E_h(\omega)\,d\omega .
 }
 \tag{381}
\] There is no atom at zero or at the threshold. Here is a direct
evaluation of the full inner angular integral, including the preferred
axis.

Fix \(K>0\), put \(r=|p|\), \(s=|k-p|\), and \(\omega=r+s\). Align a
temporary polar coordinate axis with \(k\), without moving the fixed
vector \(\mathbf e_1\). The polar-angle change to \(s\) gives measure
\((rs/K)\,dr\,d\phi\) after the energy delta function is integrated. For
\(\omega>K\), set \[
 x=\frac{r-s}{K}\in[-1,1],\quad
 r=\frac{\omega+Kx}{2},\quad s=\frac{\omega-Kx}{2},
 \quad D=\omega^2-K^2,
 \tag{382}
\] \[
 A=\frac{K+\omega x}{2},\quad B=\frac{K-\omega x}{2},
 \quad R^2=\frac{D(1-x^2)}4,\quad
 \mu=\frac{k_1}{K},\quad \eta=\mu^2,\quad e=1-\eta.
 \tag{383}
\] The variables \(A,B\) are the axial components of \(p,q\) in this
calculation, not the Hamiltonian or its coefficients. The transverse
parts are opposite, with length \(R\). Thus \[
 p_1=\mu A+\sqrt e\,R\cos\phi,\qquad
 q_1=\mu B-\sqrt e\,R\cos\phi,\qquad
 r^2=A^2+R^2,\quad s^2=B^2+R^2.
 \tag{384}
\] The original angular integrand times its Jacobian reduces to
\(\tfrac12(r^2-p_1^2)(s^2-q_1^2)\,dx\,d\phi\). Using the exact averages
of \(\cos\phi,\cos^2\phi,\cos^4\phi\), the azimuthal mean of that
product before its factor \(1/2\) is \[
 P=e^2A^2B^2+
 \left[\frac{e(1+\eta)}2(A^2+B^2)-2\eta e AB\right]R^2
 +\left(\eta+\frac{3e^2}{8}\right)R^4.
 \tag{385}
\] All four needed polynomial integrals are \[
 \begin{split}
 \int_{-1}^1A^2B^2\,dx
  &=\frac{K^4-\frac23K^2\omega^2+\frac15\omega^4}{8},\\
 \int_{-1}^1(A^2+B^2)R^2\,dx
  &=\frac{D(5K^2+\omega^2)}{30},\\
 \int_{-1}^1ABR^2\,dx
  &=\frac{D(5K^2-\omega^2)}{60},\\
 \int_{-1}^1R^4\,dx&=\frac{D^2}{15}.
 \end{split}
 \tag{386}
\] Each follows by expanding (383) and using
\(\int_{-1}^1x^{2m}dx=2/(2m+1)\). Substituting \(\eta=1-e\) and
\(\omega^2=D+K^2\) into (385)--(386) gives \[
 \int_{-1}^1P\,dx=\frac{D^2+DeK^2+e^2K^4}{15}.
 \tag{387}
\] The coefficients of \(D^2,DeK^2,e^2K^4\) are all \(1/15\); the
remaining expanded terms cancel. Consequently the original inner
momentum integral is \[
 \begin{split}
 &\int_{\mathbb R^3}|p||k-p|
 \left(1-\frac{p_1^2}{|p|^2}\right)
 \left(1-\frac{(k_1-p_1)^2}{|k-p|^2}\right)
 \delta(\omega-|p|-|k-p|)\,dp\\
 &\qquad=\frac{\pi}{15}
       [(\omega^2-K^2)^2+(\omega^2-K^2)k_\perp^2+k_\perp^4]
       \mathbf1_{\{\omega>K\}}.
 \end{split}
 \tag{388}
\] The triangle inequality gives zero below threshold. The degenerate
boundary has zero measure in the joint integral. At \(K=0\) the
expression has limit \(\pi\omega^4/15\); that outer null set requires no
special assignment. Tonelli applies to the nonnegative integrand, and
the retained constant is \([3/(8(2\pi)^6)](\pi/15)=1/(2560\pi^5)\),
proving (381).

The density is nonnegative and locally integrable. Its bracket is
strictly positive for \(K<\omega\), since \(D>0\). If nonzero compact
smooth \(h\) had \(\widehat h=0\) on a ball, its convergent entire
Taylor expansion, obtained by integrating the exponential series over
its compact support, would vanish everywhere. Fourier uniqueness then
forces \(h=0\). For a direct uniqueness proof, multiply the zero
transform by a Gaussian and use its elementary inverse Fourier integral
and Fubini; every Gaussian convolution of \(h\) is zero, and these
convolutions converge to \(h\). Thus \(\rho^E_h(\omega)>0\) for every
\(\omega>0\) when \(h\ne0\). In particular the measure gives positive
mass to every \((0,\varepsilon)\) and has no zero atom.

Integrating the polynomial in (381) over \(\omega\ge K\) also evaluates
the exact physical-time correlation: \[
 \boxed{
 C^E_h(t)=\frac1{2560\pi^5}
 \int_{\mathbb R^3}e^{-tK}|\widehat h(k)|^2
 \left[
 \frac{k_\perp^4}{t}+
 \frac{2Kk_\perp^2}{t^2}+
 \frac{10K^2-2k_1^2}{t^3}+
 \frac{24K}{t^4}+\frac{24}{t^5}
 \right]\,dk,\qquad t>0 .
 }
 \tag{389}
\] To check every coefficient, expand the bracket of (381) as \[
 \omega^4-(K^2+k_1^2)\omega^2+
                    (K^4-K^2k_1^2+k_1^4)
 \tag{390}
\] and integrate each monomial by repeated integration by parts from
\(K\) to infinity. All coefficients in (389) are nonnegative; in
particular \(10K^2-2k_1^2\ge8K^2\).

\subsubsection{24.4 Infrared coefficients and the raw large-time
amplitudes}\label{infrared-coefficients-and-the-raw-large-time-amplitudes}

Let \(m\) be the least degree of a nonzero homogeneous Taylor polynomial
\(P_m\) of \(\widehat h\) at zero. Its explicit value is \[
 P_m(k)=\frac{(-i)^m}{m!}\int_{\mathbb R^3}(k\cdot x)^m h(x)\,dx.
 \tag{391}
\] For \(h\ne0\) such a degree exists by the preceding analyticity and
uniqueness argument. Uniformly for \(|z|\le1\),
\(\widehat h(\omega z)=\omega^mP_m(z)+O_h(\omega^{m+1})\). Scaling
\(k=\omega z\) in (381) proves \[
 \rho^E_h(\omega)=A^E_{h,m}\omega^{7+2m}
                   +O_h(\omega^{8+2m}),\qquad\omega\downarrow0,
 \tag{392}
\] where the full positive anisotropic coefficient is \[
 \begin{split}
 A^E_{h,m}
 &=\frac1{2560\pi^5}
 \int_{|z|<1}|P_m(z)|^2
 \bigl[(1-|z|^2)^2+
       (1-|z|^2)(|z|^2-z_1^2)+(|z|^2-z_1^2)^2\bigr]\,dz\\
 &=\frac1{2560\pi^5}\int_{S^2}|P_m(n)|^2
 \left[
 \frac1{2m+3}-\frac{1+n_1^2}{2m+5}
             +\frac{1-n_1^2+n_1^4}{2m+7}
 \right]\,d\Omega(n)>0 .
 \end{split}
 \tag{393}
\] The second line is the elementary radial integration of the first,
using homogeneity. Positivity follows from the first line: its bracket
is positive inside the ball and a nonzero polynomial cannot vanish on an
open ball.

For \(H_h=\int h\ne0\), \(m=0\) and \(P_0=H_h\). Using
\(\int_{S^2}1=4\pi\), \(\int n_1^2=4\pi/3\), \(\int n_1^4=4\pi/5\) in
(393) gives \[
 \boxed{\rho^E_h(\omega)\sim
       \frac{H_h^2}{3360\pi^4}\omega^7.}
 \tag{394}
\] The sphere integrals follow by a polar axis along \(\mathbf e_1\) and
\(2\pi\int_{-1}^1u^{2r}du\).

Equations (381),(392), the rapid decay of \(\widehat h\), and the
substitution \(v=t\omega\) give \[
 C^E_h(t)\sim
 A^E_{h,m}\Gamma(8+2m)t^{-(8+2m)},\qquad
 -{C^E_h}'(t)\sim
 A^E_{h,m}\Gamma(9+2m)t^{-(9+2m)}
 \quad(t\to\infty).
 \tag{395}
\] For rigor, split the frequency integral at a fixed small positive
number. The complement is exponentially small in large \(t\); inside,
the remainder in (392) is bounded by its next power and its gamma
integral. Repeated integration by parts evaluates each integer gamma
integral. In particular, if \(H_h\ne0\), \[
 \boxed{
 C^E_h(t)\sim\frac{3H_h^2}{2\pi^4t^8},\qquad
 -{C^E_h}'(t)\sim\frac{12H_h^2}{\pi^4t^9},\qquad
 -{C^E_h}'(t)/C^E_h(t)\sim\frac8t.
 }
 \tag{396}
\] For general \(m\), the quotient is \((8+2m)/t+o(t^{-1})\). The norms
in (395)--(396) tend to zero with their full displayed amplitudes. These
are not unit-vector limits.

An exact continuity bound follows directly from (380): \[
 \boxed{C^E_h(t)\le\frac{3}{2\pi^4t^8}\|h\|_1^2.}
 \tag{397}
\] Indeed \(|\widehat h|\le\|h\|_1\) and \[
 \int_{\mathbb R^3}|p|
 \left(1-\frac{p_1^2}{|p|^2}\right)e^{-t|p|}\,dp
 =\frac{8\pi}{3}\int_0^\infty r^3e^{-tr}\,dr
 =\frac{16\pi}{t^4}.
 \tag{398}
\] Squaring (398) and multiplying by the exact prefactor in (380) proves
(397).

\subsubsection{24.5 Complete ultraviolet polynomial and equal-time
divergence}\label{complete-ultraviolet-polynomial-and-equal-time-divergence}

Since \(h\) is compact smooth, integration by parts makes \(\widehat h\)
decrease faster than any inverse power. Extending the integral in
(381),(390) from \(K<\omega\) to all \(k\) therefore makes an error
\(O_{h,N}(\omega^{-N})\) for every fixed integer \(N\). On the omitted
region \(K\ge\omega\), all terms of (390) are bounded by a constant
times \(K^4\); the stated remainder follows by choosing a sufficiently
high Fourier decay power and integrating its radial tail.

Put \(\nabla_\perp=(\partial_2,\partial_3)\) and
\(\Delta_\perp=\partial_2^2+\partial_3^2\), and define \[
 A_h=\|h\|_2^2,\qquad
 B_h=\|\nabla h\|_2^2+\|\partial_1h\|_2^2,\qquad
 D_h=\|\Delta_\perp h\|_2^2+
        \|\partial_1\nabla_\perp h\|_2^2+
        \|\partial_1^2h\|_2^2 .
 \tag{399}
\] Here \(D_h\) denotes this scalar only in (399)--(402); the
directional electric operator remains \(D^E_h\) in (355). The full
ultraviolet expansion is \[
 \boxed{
 \rho^E_h(\omega)=
 \frac{A_h\omega^4-B_h\omega^2+D_h}{320\pi^2}
          +O_{h,N}(\omega^{-N})\quad(\omega\to\infty).
 }
 \tag{400}
\] Every coefficient follows from (390) and Fourier norm identities. In
particular \[
 K^4-K^2k_1^2+k_1^4
 =k_\perp^4+k_1^2k_\perp^2+k_1^4,
 \tag{401}
\] which gives the three nonnegative terms in \(D_h\). With the Fourier
convention (378), Parseval contributes \((2\pi)^3\), giving
\(8\pi^3/(2560\pi^5)=1/(320\pi^2)\). This identity can be justified
directly here: insert a Gaussian factor in \(\int|\widehat h(k)|^2dk\),
apply Fubini and the elementary Gaussian Fourier integral, and let its
width tend to zero. The resulting approximate identity gives
\((2\pi)^3\int|h(x)|^2dx\). Integration by parts yields the same
statement for the derivatives in (399).

For \(h\ne0\), \(A_h>0\), so the locally finite raw measure has infinite
total mass. Positive-time damping gives finite mass and all energy
moments, since the density is bounded by a constant times
\(1+\omega^4\). Its exact small-time singular expansion is \[
 C^E_h(t)=\frac1{320\pi^2}
    \left(\frac{24A_h}{t^5}-\frac{2B_h}{t^3}
                          +\frac{D_h}{t}\right)+O_h(1)
 \quad(t\downarrow0).
 \tag{402}
\] To see that the remainder is bounded, subtract the polynomial in
(400) on the whole half-line. The resulting function is integrable: it
is bounded near zero and rapidly decreasing at infinity. Its Laplace
integral is therefore bounded. In particular \[
 C^E_h(t)\sim\frac{3\|h\|_2^2}{40\pi^2t^5},\quad
 -{C^E_h}'(t)\sim\frac{3\|h\|_2^2}{8\pi^2t^6},\quad
 -{C^E_h}'(t)/C^E_h(t)\sim\frac5t.
 \tag{403}
\] The first-moment statement follows either by differentiating the
integral with its polynomial remainder, or repeating (402) with one
additional power of \(\omega\). Spatial smearing alone therefore does
not supply an equal-time Hilbert vector in this limiting two-creation
representation. This does not assert a divergence at any finite
regulator, where (355) is smooth.

\subsubsection{24.6 The explicit common two-creation state
map}\label{the-explicit-common-two-creation-state-map}

Let the one-particle space be the complex \(L^2\) transverse vector
fields on momentum space with three colour coordinates, using Lebesgue
measure. Take its symmetric Fock space and
\(A_{\mathrm{fr}}=d\Gamma(|p|)\). On the \(n\)-particle space this
multiplies by \(\sum_{i=1}^n|p_i|\), with maximal squared energy domain;
the direct sum of these maximal real multiplication operators is
nonnegative and self-adjoint. Its nonreal resolvents are the
corresponding bounded multipliers, which also verify these domain
statements.

For \(t>0\) define the equal-colour symmetric two-particle vector \[
 J^E_t h=\frac14\sum_{\alpha=1}^3\sum_{r,s=1}^2
 \int e^{-t(|p|+|q|)/2}\mathcal R^{E,rs}_h(p,q)
       a^\dagger_{r\alpha}(p)a^\dagger_{s\alpha}(q)
                       \Omega\,dp\,dq.
 \tag{404}
\] Precisely, its two-particle wavefunction is \(\sqrt2/4\) times the
displayed symmetric kernel with equal colours, so (404) requires no
distributional operator product. Equation (380) proves square
integrability. It is invariant under the original simultaneous adjoint
rotations of the three colour coordinates. The trace metric
\(-2\operatorname{tr}(T_\alpha T_\beta)=\delta_{\alpha\beta}\) makes
those rotations orthogonal; their contraction in (404) is exactly
preserved.

The full raw identities are \[
 \|J^E_t h\|^2=C^E_h(t),\qquad
 \langle J^E_t h,A_{\mathrm{fr}}^nJ^E_t h\rangle
   =\int\omega^ne^{-t\omega}\rho^E_h(\omega)\,d\omega,
 \qquad
 e^{-sA_{\mathrm{fr}}/2}J^E_t h=J^E_{t+s}h.
 \tag{405}
\] They follow directly from multiplication by the pair energy and the
factor \(3/8\) already proved. For nonzero \(h\) the norm is positive by
(381), proving injectivity on the real test profiles. Complex-linear
extension has the same Hermitian norm formula. Changing a transverse
polarization frame applies pointwise orthogonal matrices and their
tensor squares to (378),(404), preserving the Hilbert vector and all
norms. Spatial translation \(h(x)\mapsto h(x-b)\) multiplies the kernel
by \(e^{i(p+q)\cdot b}\), the exact two-particle translation
representation. The direction-one axis remains fixed; full spatial
rotation covariance would also rotate that prescribed axis.

This gives a typed common Hilbert and spectral map for the positive-time
vectors. Its Hamiltonian is the explicitly defined free transverse
multiplication operator. The correspondence does not identify it with
the full interacting continuum theory or establish products of several
local observables.

\subsubsection{24.7 A common actual positive dyadic
sequence}\label{a-common-actual-positive-dyadic-sequence}

Retain \(L_j=j^2\), \(a_j=1/(100j)\), and \(\ell_j=a_j(2L_j+1)\), for
integers \(j\ge2\). For actual finite-regulator vectors write \[
 C^E_{g;h,k}(t)=
 \langle\Xi^E_{h,g},e^{-tA_g}\Xi^E_{k,g}\rangle,\qquad
 y^E_{h,g}(t)=e^{-tA_g/2}\Xi^E_{h,g}.
 \tag{406}
\] Enumerate the positive rational times \(t_1,t_2,\ldots\). Let
\(\mathsf E_1\) be the original direction-one edges and
\(n_1=|\mathsf E_1|\). For each \(e\in\mathsf E_1\), use the original
single-edge electric state \[
 \Xi^E_{e,g}=\kappa(E_e-\langle E_e\rangle_{\psi_g})\psi_g.
 \tag{407}
\] These states have no added quarter-face observable term. The full
Hamiltonian used to evolve them is still (353). Set
\(\varepsilon_j=j^{-40}(1+\ell_j)^{-10}\). At stage \(j\), impose all
previously proved stage-\(j\) conditions that are to be retained in the
programme, \(0<g<j^{-5}\), and the additional finite tests \[
 \left|C^E_{g;e,e'}(t_n)
     -C^{E,L_j,a_j}_{0;e,e'}(t_n)\right|
       <\frac{\varepsilon_j}{n_1^2}
 \quad(e,e'\in\mathsf E_1,\ 1\le n\le j).
 \tag{408}
\] At fixed \(j\), (366)--(368) prove convergence for each such test.
Mixed terms follow by real polarization, since the Hamiltonian, weights
and vacuum are real. The list is finite at each stage, so it holds for
every sufficiently small positive coupling. The previously retained
finite tests likewise hold on tails by their proved fixed-box
convergence. Take the least positive integer exponent \(k_j\) for which
the entire list holds at \(g_j=2^{-k_j}\). Existence follows from these
tails and well ordering. This refines the coupling selection; it does
not assert that an earlier least dyadic already met these extra electric
tests.

Linearity in the original direction-one weights gives \[
 |C^E_{g_j;h,k}(t_n)-C^{E,L_j,a_j}_{0;h,k}(t_n)|
 \le\varepsilon_j\|h\|_\infty\|k\|_\infty
 \quad(n\le j).
 \tag{409}
\] Thus one sequence gives every fixed compact smooth profile at each
positive rational time. The all-mode limit (380) identifies these
limits. For every real \(t>0\), the actual diagonal covariance
\(C^E_{g_j;h,h}(t)\) is decreasing in \(t\) because its spectral measure
is nonnegative. Squeeze it between two rational times and use the
continuity of (389); then polarize to obtain \[
 C^E_{g_j;h,k}(t)\longrightarrow C^E_{h,k}(t)
 \quad(t>0)
 \tag{410}
\] for all real compact smooth \(h,k\).

All original couplings remain \[
 \kappa_j=200jg_j^2,\qquad
 b_j=50j/g_j^2,\qquad
 \xi_j=1/(4g_j^4),
 \tag{411}
\] with every nonlinear Wilson term and its scalar, and the actual
vacuum at each positive \(g_j\). No uniform dependence of the fixed-box
approximation on \(L,a\) was assumed. This sequence is not identified
with the prescribed logarithmic or fixed-electric-coefficient coupling
paths.

For precision, let \(\nu^E_{h,g_j}\) be the actual unfiltered raw
spectral measure of (355). For fixed \(t>0\), define the finite measure
\(\mu_{j,t}=e^{-t\omega}\nu^E_{h,g_j}\). Its mass tends to \(C^E_h(t)\)
and its Laplace transform at \(s>0\) tends to \(C^E_h(t+s)\). The exact
bound \[
 (1-e^{-sR})\mu_{j,t}([R,\infty))
 \le C^E_{g_j;h,h}(t)-C^E_{g_j;h,h}(t+s)
 \tag{412}
\] proves tightness: first choose \(s>0\) small using continuity at
\(t\), then \(R\) large so its denominator is bounded away from zero.
Finitely many initial measures are controlled separately. On a compact
energy interval, polynomials in \(e^{-\omega}\) approximate all
continuous functions by the change \(z=e^{-\omega}\) and polynomial
approximation on its compact range. The Laplace limits and tightness
therefore identify the full weak limit: \[
 \boxed{\mu_{j,t}\Longrightarrow
           e^{-t\omega}\rho^E_h(\omega)\,d\omega.}
 \tag{413}
\] Applying (413) at \(t/2\) to the bounded continuous test
\(\omega^n e^{-t\omega/2}\) proves every positive-time energy moment,
including \[
 \|y^E_{h,g_j}(t)\|^2\to C^E_h(t),\qquad
 \langle y^E_{h,g_j}(t),A_{g_j}y^E_{h,g_j}(t)\rangle
                   \to-{C^E_h}'(t).
 \tag{414}
\] The damped vectors lie in every excitation-power domain: the spectral
multiplier \(\omega^m e^{-t\omega/2}\) is bounded at each fixed \(t>0\).

For finitely many profiles and positive times, their actual Gram and
semigroup matrix elements converge to those of (404). The mixed time is
the sum of the two half-times and any intervening nonnegative evolution
time, so (410) applies. This proves the common state correspondence at
the level claimed, without asserting an isometry between the entire
varying Hilbert spaces.

\subsubsection{24.8 Locally finite raw convergence and unfiltered
escape}\label{locally-finite-raw-convergence-and-unfiltered-escape}

Equation (413) also proves vague convergence of the actual unfiltered
measures to (381): multiply any compactly supported continuous test by
\(e^{t\omega}\). Because the limiting density has no atoms, the same
argument by continuous upper and lower approximation gives \[
 \nu^E_{h,g_j}([0,\Omega])
 \longrightarrow\int_0^\Omega\rho^E_h(\omega)\,d\omega
 \quad(0<\Omega<\infty).
 \tag{415}
\] The limit is finite and positive for \(h\ne0\). For the damped
vectors their raw low-energy mass is instead
\(\int_0^\Omega e^{-t\omega}\rho^E_h(\omega)d\omega>0\). These limits
retain the original raw spectral amplitudes.

On the same actual sequence the unfiltered norms diverge: \[
 \boxed{\|\Xi^E_{h,g_j}\|^2\longrightarrow\infty
        \quad(h\ne0).}
 \tag{416}
\] Indeed their total spectral masses dominate \(C^E_{g_j;h,h}(t)\). For
any prescribed \(M\), choose a fixed small \(t>0\) with \(C^E_h(t)>2M\),
possible by (402). Equation (410) makes that lower bound exceed \(M\) at
every sufficiently late regulator. Since \(M\) was arbitrary, this
proves (416), with no unsupported rate in the regulator.

After these raw facts have been retained, their consequence for the
corresponding probabilities is \[
 \frac{\nu^E_{h,g_j}([0,\Omega])}
      {\|\Xi^E_{h,g_j}\|^2}\longrightarrow0
 \quad(\Omega<\infty).
 \tag{417}
\] This does not replace any original state by a unit vector; it records
what division by its explicitly diverging mass would do. Positive raw
low-energy mass in (415) and probability escape in (417) both follow
from the full calculation.

Equations (395),(414) give the physical large-time filtered raw
amplitudes at each fixed positive time before any time sequence is
taken. If an actual time-growing sequence is desired, set \(t=n\), and
choose the least increasing regulator index \(j_n\) for which the two
errors in (414) are each smaller than \(1/n\) times their respective
strictly positive limits. These indices exist by (414). Then the actual
vectors \(y^E_{h,g_{j_n}}(n)\) retain the norm and energy amplitudes of
(395), with ratio \((8+2m)/n+o(n^{-1})\). This is an explicit further
selection, not an exchange of the positive-time and regulator limits.

The original all-box covariance \(\Gamma=\sum_{e\parallel1}n_2^2E_e\)
has a different noncompact spatial coefficient and its raw scaling is
not determined by replacing it with a fixed \(h\). The native magnetic
\(K_\theta\) state requires its already specified separate covariance
comparison. This note supplies the exact direction-one local electric
coefficient, full positive-time spectral map, anisotropic density,
ultraviolet and infrared raw amplitudes, and a common positive-coupling
sequence. It neither identifies that sequence with a prescribed
running-coupling path nor constructs an interacting four-dimensional
mass-gap counterexample.

\subsection{25. The original weighted covariance at every physical
energy
scale}\label{the-original-weighted-covariance-at-every-physical-energy-scale}

This calculation retains the original noncompact coordinate weight and
both open boundaries. The full mode trace supplies an exact leading raw
mass, finite-energy fraction and native magnetic transfer on an
explicitly strengthened cusp.

\subsubsection{25.1 Exact original operator and the component
trace}\label{exact-original-operator-and-the-component-trace}

Let \(N=2L+1\), \(\ell=aN\), and let \(\mathsf E_L\) be every oriented
positive edge of \(\{-L,\ldots,L\}^3\). Retain \[
 H_g=\frac{2g^2}{a}\sum_e E_e+
 \frac1{2g^2a}\sum_p(2-W_p),\quad
 A_g=H_g-E_g,\quad
 \Gamma=\sum_{e=(n,1)}n_2^2E_e,\quad
 v_{\Gamma,g}=(\Gamma-\langle\Gamma\rangle_g)\psi_g .
 \tag{418}
\] Here \(\psi_g\) is the actual positive unit physical vacuum. The
original period, covering, coordinate, and Haar maps are those of the
cumulative manuscript. If
\(D^E_h=(2g^2/a)\sum_{e\parallel1}h(am(e))E_e\), then on this entire
finite box the physical weight \(h(x)=x_2^2\) gives the exact identity
\[
 D^E_{x_2^2}=2g^2a\Gamma,\qquad
 g^2v_{\Gamma,g}=\frac1{2a}
       (D^E_{x_2^2}-\langle D^E_{x_2^2}\rangle_g)\psi_g .
 \tag{419}
\] This growing-box weight is not a fixed compact test function. We
prove its limit directly below.

Let \(\mathscr R\), \(\Lambda=\operatorname{diag}\sigma_\nu\), and
\(D=\mathscr R^{\mathsf T}\mathsf W\mathscr R\) have exactly their
Section 18 meanings, with \(\mathsf W_{(n,1)}=n_2^2\) and its other
entries zero. The ordered form of its full comparison measure is \[
 m^\Gamma_{L,a}=\frac3{32}\sum_{\nu,\eta}
   \sigma_\nu\sigma_\eta D_{\nu\eta}^2
              \delta_{(\sigma_\nu+\sigma_\eta)/a}.
 \tag{420}
\] The factor \(3/32\) includes three colours, both pair contractions
and the original coefficient \(-1/8\) in the pair vector. For \(t\ge0\),
set \[
 T_t=\mathscr R\Lambda e^{-t\Lambda/a}\mathscr R^{\mathsf T},
 \qquad C^\Gamma_{L,a}(t)=\int e^{-t\omega}\,dm^\Gamma_{L,a}(\omega)
           =\frac3{32}\operatorname{Tr}(\mathsf W T_t\mathsf W T_t).
 \tag{421}
\] Moving the finite rectangular matrices cyclically proves this
equality without a missing projection term.

Write the original exact one-dimensional vertex and edge modes as \[
 \begin{split}
 v_r(n)&=\sqrt{\frac{2-\delta_{r0}}N}
   \cos\frac{\pi r(n+L+1/2)}N,\quad 0\le r<N,\\
 s_r&=2\sin\frac{\pi r}{2N},\\
 w_r(n)&=-\sqrt{\frac2N}
   \sin\frac{\pi r(n+L+1)}N,\quad 1\le r<N .
 \end{split}\tag{422}
\] Finite geometric sums and \(d_0v_r=s_rw_r\) give orthonormality and
the full tensor decomposition. In a frequency block \(\boldsymbol r\),
its one-form space has components only where \(r_i>0\); the gradient is
its line spanned by \(s=(s_{r_1},s_{r_2},s_{r_3})\). The transverse
projection is \(I-ss^{\mathsf T}/|s|^2\) on this space. Consequently,
the direction-one to direction-one block of \(T_t\), in the exact basis
\[
 w_{r_1}(n_1)v_{r_2}(n_2)v_{r_3}(n_3),
 \quad 1\le r_1<N,\quad 0\le r_2,r_3<N,
\] is diagonal, with scalar \[
 f_t(\boldsymbol r)=
 \frac{s_{r_2}^2+s_{r_3}^2}{\sigma_{\boldsymbol r}}\,
            e^{-t\sigma_{\boldsymbol r}/a},\qquad
 \sigma_{\boldsymbol r}=(s_{r_1}^2+s_{r_2}^2+s_{r_3}^2)^{1/2}.
 \tag{423}
\] The scalar is zero if \(r_2=r_3=0\). This includes the absent
transverse block rather than counting a spurious polarization. Since
\(\mathsf W\) has only direction-one entries, (421) uses this component
block twice. All original transverse modes and their exact polarization
sums have therefore been retained.

\subsubsection{25.2 Discrete weight estimates including both
boundaries}\label{discrete-weight-estimates-including-both-boundaries}

Let \(Z\) be multiplication by \((n_2/N)^2\) on the vertex coordinate
\(n_2=-L,\ldots,L\). In its cosine basis denote
\(Z_{rs}=\langle v_r,Zv_s\rangle\). There is a constant independent of
odd \(N\ge3\) such that \[
 |Z_{rs}|\le\frac{C}{(1+|r-s|)^2},
 \quad \sum_{|r-s|>M}|Z_{rs}|^2\le CM^{-3},
 \quad \sum_s|Z_{rs}|^2\le\frac1{16}.
 \tag{424}
\] We give the boundary proof. For \(d=2,4\), define \[
 q^{(d)}_h=\frac1N\sum_{n=-L}^L
       (n/N)^d\cos\frac{\pi h(n+L+1/2)}N .
\] Reflect the \(N\) samples evenly at the two endpoint half-cells to a
\(2N\)-periodic sequence. In its interior the absolute second
differences of the polynomial samples are at most \(C_d/N^2\). At each
of the two reflection junctions the second difference is bounded by
\(C_d/N\), by the bound on the first derivative of \(x^d\) on
\([-1/2,1/2]\). Thus the sum of the absolute second differences around
the periodic sequence is at most \(C_d/N\). The Fourier multiplier of
minus the periodic second difference is \(4\sin^2(\pi h/(2N))\). Even
reflection identifies its Fourier coefficient, up to a unit phase, with
\(q^{(d)}_h\). Division by the \(2N\) in that coefficient gives \[
 |q^{(d)}_h|\le
 \frac{C_d}{N^2\sin^2(\pi h/(2N))}
 \le\frac{C'_d}{\operatorname{dist}(h,2N\mathbb Z)^2}
 \quad(h\notin2N\mathbb Z).
 \tag{425}
\] The second inequality uses \(\sin u\ge 2u/\pi\) for
\(0\le u\le\pi/2\), after reflection. The zero coefficient is bounded
directly by \(2^{-d}\). The product-of-cosines identity gives \[
 Z_{rs}=\frac{\sqrt{(2-\delta_{r0})(2-\delta_{s0})}}2
                  (q^{(2)}_{r-s}+q^{(2)}_{r+s}).
 \tag{426}
\] For \(0\le r,s<N\), both \(r+s\) and \(2N-r-s\) are at least
\(|r-s|\). Thus (425) proves the first inequality in (424), including
the high-frequency boundary alias \(r+s\) near \(2N\). Summing the
fourth-power tail proves the second; the last follows from
\(Z^2\le I/16\).

The diagonal of \(Z^2\) is \[
 A_r:=\langle v_r,Z^2v_r\rangle
 =\begin{cases}\mu_N,&r=0,\\
 \mu_N+q^{(4)}_{2r},&1\le r<N,\end{cases}
 \qquad
 \mu_N=\frac{L(L+1)(3L^2+3L-1)}{15N^4}
                 \longrightarrow\frac1{80}.
 \tag{427}
\] Indeed \(\sum_{n=-L}^Ln^4=L(L+1)(2L+1)(3L^2+3L-1)/15\), obtained by
summing the fourth finite difference of a fifth-degree polynomial. The
formula for \(v_r^2\) proves the diagonal identity. In particular \[
 |A_r-\mu_N|\le
 \frac{C}{(1+\min(r,N-r))^2}\quad(1\le r<N).
 \tag{428}
\]

\subsubsection{25.3 Full positive-time
trace}\label{full-positive-time-trace}

For every simultaneous \(a\downarrow0\), \(\ell=aN\to\infty\), and fixed
\(t>0\), we prove \[
 C^\Gamma_{L,a}(t)
   \sim\frac{3\ell^7}{12800\pi^2a^2t^5}.
 \tag{429}
\] Let \(S_t\) be the direction-one diagonal matrix whose entries are
\(f_t/a\), and extend \(Z\) on that component edge space as the identity
in coordinates one and three. For finite self-adjoint matrices, \[
 \operatorname{Tr}(Z^2S_t^2)-\operatorname{Tr}(ZS_tZS_t)
              =\frac12\|[Z,S_t]\|_{\rm HS}^2.
 \tag{430}
\] Expansion of the squared Hilbert--Schmidt norm and cyclicity of the
finite trace prove the equality; its right side is nonnegative.

Put \(p_i=s_{r_i}/a\), \(\omega_{\boldsymbol r}=|p|\). The scalar of
\(S_t\) is \[
 d_t(p)=\frac{p_2^2+p_3^2}{|p|}e^{-t|p|},\qquad d_t(0)=0.
 \tag{431}
\] This is a continuous Lipschitz function. Away from zero,
differentiating its homogeneous degree-one prefactor bounds its gradient
by \(C(1+t|p|)e^{-t|p|}\le C_t e^{-t|p|/2}\). The same bound along
segments follows by integration, since a point at the origin has
one-dimensional measure zero and the function is Lipschitz there. Also
\[
 \frac{2|\boldsymbol r|}{\ell}\le\omega_{\boldsymbol r}
 \le\frac{\pi|\boldsymbol r|}{\ell},\qquad
 \left|\frac{\partial (s_r/a)}{\partial r}\right|\le\frac\pi\ell.
 \tag{432}
\] The number of nonnegative integer triples in a shell
\(m\le|\boldsymbol r|<m+1\) is at most \(C(m+1)^2\): associate their
disjoint unit cubes to an annulus whose radii differ by at most
\(1+2\sqrt3\), and bound its volume. Equations (432), followed by
comparison of the resulting polynomial-exponential series with its
integral on each unit interval, prove \[
 \sum_{\boldsymbol r}\omega_{\boldsymbol r}^k
                 e^{-u\omega_{\boldsymbol r}}\le C_{k,u}\ell^3,
 \quad \ell\ge1,\quad k=0,1,2,\ldots,\quad u>0 .
 \tag{433}
\] Excluding \(\boldsymbol r=0\) when necessary only decreases these
sums. The same argument after factoring out \(e^{-uR/2}\) proves a tail
bound \(C_{k,u}\ell^3e^{-uR/2}\) on \(\omega_{\boldsymbol r}>R\), with
an adjustment of the constant and a smaller exponent if needed.

In the commutator, only \(r_2,s_2\) differ. Split at
\(1\le |r_2-s_2|\le M\), where \(M/\ell\le1\). Integration of the
derivative bound and (432) gives \[
 |d_t(p_{\boldsymbol r})-d_t(p_{r_1,s_2,r_3})|
       \le C_t\frac M\ell e^{-t\omega_{\boldsymbol r}/4}.
\] Here the frequency changes by at most \(\pi M/\ell\le\pi\), so
replacing the minimum frequency along the segment by the initial one
costs only a constant depending on \(t\). Using the last bound of (424)
and then (433), the squared commutator contribution of these entries is
at most \(C_t\ell M^2\). On \(|r_2-s_2|>M\), use
\(|d-d'|^2\le2d^2+2(d')^2\), the symmetry of \(Z\), and its tail bound
in (424). The result is at most \(C_t\ell^3M^{-3}\). Choose
\(M=\lfloor\sqrt\ell\rfloor\) for sufficiently large \(\ell\). Then \[
 \|[Z,S_t]\|_{\rm HS}^2=o(\ell^3).
 \tag{434}
\] Both boundaries and all high-frequency modes have been controlled in
this estimate.

Next,
\(\operatorname{Tr}(Z^2S_t^2)=\sum_{\boldsymbol r}A_{r_2}d_t(p_{\boldsymbol r})^2\).
On \(\omega_{\boldsymbol r}\le R\), each coordinate index is at most
\(R\ell/2\), and for sufficiently small \(a\) this is less than \(N/2\).
Equation (428), summation in \(r_2\), and the \(O_R(\ell^2)\) possible
pairs \((r_1,r_3)\) bound the error in replacing \(A_{r_2}\) by
\(\mu_N\) by \(C_R\ell^2\). The error on the complement divided by
\(\ell^3\) tends uniformly to zero as \(R\to\infty\), since \(A_{r_2}\)
and \(\mu_N\) are bounded and (433) gives the exponential tail. Thus \[
 \operatorname{Tr}(Z^2S_t^2)
       =\mu_N\sum_{\boldsymbol r}d_t(p_{\boldsymbol r})^2+o(\ell^3).
 \tag{435}
\] On every fixed bounded physical frequency region,
\(s_{r_i}/a-(\pi r_i/\ell)\to0\) uniformly, using the sine Taylor
remainder bounded by its cubic term. The continuous function (431), the
grid step \(\pi/\ell\), and the tail estimate give \[
 \begin{split}
 \ell^{-3}\sum_{\boldsymbol r}d_t(p_{\boldsymbol r})^2
 &\longrightarrow
 \frac1{\pi^3}\int_{\mathbb R_+^3}
     |p|^2(1-p_1^2/|p|^2)^2 e^{-2t|p|}\,dp\\
 &=\frac1{5\pi^2t^5}.
 \end{split}\tag{436}
\] Coordinate planes contribute \(O_R(\ell^2)\) in each bounded region
and have zero limiting measure. The angular integral on the full sphere
is \(2\pi\int_{-1}^1(1-u^2)^2du=32\pi/15\), and the positive octant
gives one eighth of it. Four integrations by parts give
\(\int_0^\infty r^4e^{-2tr}dr=24/(2t)^5=3/(4t^5)\). These constants
prove the second line. Finally, \[
 C^\Gamma_{L,a}(t)
       =\frac3{32}N^4a^2\operatorname{Tr}(ZS_tZS_t).
\] Insert (427), (430), (434)--(436) and \(N=\ell/a\) to prove (429).

\subsubsection{25.4 Entire raw mass and the exact finite-energy
fraction}\label{entire-raw-mass-and-the-exact-finite-energy-fraction}

Define the positive finite lattice integral \[
 I_{\rm lat}=\frac1{\pi^3}\int_{[0,\pi]^3}
 \frac{(s_2^2+s_3^2)^2}{s_1^2+s_2^2+s_3^2}\,dk_1\,dk_2\,dk_3,
 \qquad s_i=2\sin(k_i/2),
 \tag{437}
\] with the integrand set to zero at the single origin. Then \[
 C_L^\Gamma=C^\Gamma_{L,a}(0)
       \sim \frac{3N^7}{2560}I_{\rm lat},\qquad
                       \frac83<I_{\rm lat}<4.
 \tag{438}
\] For the trace proof use the direction-one matrix with scalar
\(d_0(s)=(s_2^2+s_3^2)/|s|\), extended continuously at zero. It is
globally Lipschitz on the bounded cube. Repeat (430) with index
derivative at most \(C/N\). The near-diagonal bound is \(CNM^2\), and
the far bound is \(CN^3M^{-3}\), since the dimension is less than
\(N^3\) and the scalar is bounded. Choosing \(M=\lfloor\sqrt N\rfloor\)
makes the commutator squared \(o(N^3)\). Equation (428) gives total
diagonal-weight error \(O(N^2)\), including \(r_2\) near \(N\). The
Riemann sum of the continuous symbol squared divided by \(N^3\) tends to
(437). Multiplication by \(3N^4/32\) and by \(\mu_N\to1/80\) proves the
asymptotic.

To prove the bounds, take expectation for uniform \(k\) on the cube, put
\(X=s_1^2\), \(S=s_1^2+s_2^2+s_3^2\), and use \(\mathbb EX=2\),
\(\mathbb ES=6\). The integral is
\(\mathbb E(S-2X+X^2/S)=2+\mathbb E(X^2/S)\). Cauchy--Schwarz applied to
\(X/\sqrt S\) and \(\sqrt S\) gives \(\mathbb E(X^2/S)\ge 4/6\), with
strict inequality because \(X/S\) is not constant almost everywhere.
Also \(X^2/S<X\) almost everywhere since the two other coordinates have
positive squares almost everywhere. This proves both strict bounds
without replacing the integral by an estimated coefficient.

Let \(a\to0,\ell\to\infty\). The locally finite raw spectral limit, with
its full dimensional factor recorded, is \[
 \frac{a^2}{\ell^7}m^\Gamma_{L,a}
      \longrightarrow \frac{\omega^4}{102400\pi^2}\,d\omega
 \quad\hbox{on bounded energy intervals}.
 \tag{439}
\] For a complete measure proof, multiply these measures by
\(e^{-t\omega}\) for a fixed \(t>0\). Their masses and Laplace
transforms converge by (429). If their masses are \(M_j(t)\), then \[
 (1-e^{-sR})\mu_j([R,\infty))
       \le M_j(t)-M_j(t+s).
\] Continuity of \(3/(12800\pi^2t^5)\), first choosing small \(s>0\) and
then large \(R\), proves tightness; finitely many initial measures cause
no difficulty. On a bounded interval, \(z=e^{-\omega}\) and Bernstein
polynomial approximation show that constants and finite linear
combinations of \(e^{-n\omega}\) approximate every continuous function.
Tightness and the convergent transforms therefore identify the weak
limit. Its density is the one in (439) times \(e^{-t\omega}\), because
\(\int_0^\infty\omega^4e^{-t\omega}d\omega=24/t^5\). Multiplication by
\(e^{t\omega}\) on a bounded interval proves (439), including its
interval masses since this density has no atoms. Applying weak
convergence at time \(t/2\) to the bounded continuous function
\(\omega^k e^{-t\omega/2}\) also proves every positive-time moment: \[
 \frac{a^2}{\ell^7}\int\omega^k e^{-t\omega}dm^\Gamma_{L,a}
       \longrightarrow
         \frac{(k+4)!}{102400\pi^2t^{k+5}},
                  \quad k=0,1,2,\ldots .
 \tag{440}
\] For every fixed \(\Omega>0\), (438)--(439) give the full leading raw
mass and probability, \[
 \begin{split}
 m^\Gamma_{L,a}([0,\Omega])
       &\sim\frac{\ell^7}{a^2}\frac{\Omega^5}{512000\pi^2},\\
 \rho^\Gamma_{L,a}([0,\Omega])
       &\sim a^5\frac{\Omega^5}{600\pi^2 I_{\rm lat}} .
 \end{split}\tag{441}
\] Here the probability uses precisely the positive raw mass (438); the
first line remains part of the assertion. On \(L_j=j^2,a_j=1/(100j)\),
the first mass grows as a fixed constant times \(j^9\), while the whole
mass grows as a fixed constant times \(j^{14}\). Their vanishing
fraction has exact order \(j^{-5}\). This refines, rather than
contradicts, the previous uniform bound.

\subsubsection{25.5 Transport to one actual positive-coupling
sequence}\label{transport-to-one-actual-positive-coupling-sequence}

Retain every earlier finite-stage test from Sections 18 and 22--24,
including the single-edge directional tests, and the supplied full
local-spectrum sequence, and retain \(g<j^{-5}\). Enumerate the positive
rational times \(t_1,t_2,\ldots\). At fixed \(j\), set \(N_j=2j^2+1\),
\(a_j=1/(100j)\), \(\ell_j=a_jN_j\), and in addition require \[
 \begin{split}
 |g^4\|v_{\Gamma,g}\|^2-C^\Gamma_{L_j}|&<N_j^7/j,\\
 \left|g^4\langle v_{\Gamma,g},e^{-t_rA_g}v_{\Gamma,g}\rangle
                  -C^\Gamma_{L_j,a_j}(t_r)\right|
                         &<\ell_j^7/(a_j^2j),\quad r\le j .
 \end{split}\tag{442}
\] The fixed-box graph and finite spectral projection theorem (231)
proves convergence of the total mass and every bounded heat multiplier
as \(g\downarrow0\). Each prior finite-stage test likewise holds
throughout a sufficiently small positive interval by its proved
fixed-box limit. Intersect these finitely many intervals. Define \(k_j\)
as the least positive integer such that every dyadic \(g=2^{-q}\),
\(q\ge k_j\), satisfies this finite list, and set \(g_j=2^{-k_j}\). The
intersection proves existence. This is a refinement of the selected
actual sequence, with all earlier results retained, and not a claim
about either prescribed coupling path.

Every coefficient remains \[
 \kappa_j=200j\,2^{-2k_j},\quad
 b_j=50j\,2^{2k_j},\quad
 \xi_j=2^{4k_j-2},\quad
 2b_jM_j=100j\,2^{2k_j}M_j .
 \tag{443}
\] Let \(\nu_{\Gamma,j}\) be the unchanged raw spectral measure of
\(v_{\Gamma,g_j}\), so \(m^\Gamma_{L_j,a_j,g_j}=g_j^4\nu_{\Gamma,j}\).
Equations (442) and the trace theorem give the heat limits at rational
times. Monotonicity of positive spectral heat integrals squeezes every
real positive time between rational times, and continuity of the limit
gives the same conclusion there. The preceding damped-measure and moment
proof applies unchanged. Consequently \[
 \begin{split}
 g_j^4\|v_{\Gamma,g_j}\|^2&\sim(3I_{\rm lat}/2560)N_j^7,\\
 \nu_{\Gamma,j}([0,\Omega])
   &\sim g_j^{-4}\frac{\ell_j^7}{a_j^2}
                         \frac{\Omega^5}{512000\pi^2},\\
 \zeta_{\Gamma,j}([0,\Omega])
   &\sim a_j^5\frac{\Omega^5}{600\pi^2I_{\rm lat}},\\
 \|e^{-tA_{g_j}/2}v_{\Gamma,g_j}\|^2
   &\sim g_j^{-4}\frac{3\ell_j^7}{12800\pi^2a_j^2t^5}.
 \end{split}\tag{444}
\] All positive-time energy moments carry the coefficient in (440),
multiplied by \(g_j^{-4}\ell_j^7/a_j^2\). In particular the quotient of
first energy and squared norm of the displayed actual filtered vector
tends to \(5/t\).

\subsubsection{25.6 Exact native magnetic transfer on the strengthened
original cusp
family}\label{exact-native-magnetic-transfer-on-the-strengthened-original-cusp-family}

To transfer the leading finite-energy fraction, the old relative error
\(O(j^{-4})\) is insufficient compared with \(a_j^5\). We explicitly
strengthen the external cusp depth while retaining the full finite
Hamiltonian, vacuum, coordinate order and cover: \[
 \begin{split}
 F(\xi)&=\xi^{3/4}(6+64\xi),\\
 T_j^{\ddagger}
    &=2C+j^5 e^{8\pi\xi_j}F(\xi_j)^{1/4},\\
 D_j^{\ddagger}
    &=L_{T_j^{\ddagger}}q_{T_j^{\ddagger}}
                             +6m_{T_j^{\ddagger}}^2,\\
 \theta_j^{\ddagger}&=\frac{2\pi}{10^4j^2D_j^{\ddagger}},\quad
 v_j^{\ddagger}=e^{-2\pi T_j^{\ddagger}/j},\quad
 t_{c,j}^{\ddagger}=(v_j^{\ddagger})^j .
 \end{split}\tag{445}
\] The cover remains \(j\), its matrix is the original
\(P_T\operatorname{diag}(1,1,j,j)\), its pullback metric is unchanged in
form, its signed determinant is \(-j^2D_j^\ddagger\), and its positive
volume is \(j^2D_j^\ddagger\). The original nonwrapping connection is
\(h_1(n)=\exp(-\theta_j^\ddagger n_2H_c)\), \(h_2=h_3=I\). The
previously proved embedding and systole estimates apply with this larger
original cusp parameter.

Relative to (242), the fourth power of \(T-C\) has an extra factor
\(j^4\) in the denominator of the angle error. Thus the all-angle
estimate, with its original constant \(A_{\rm rel}=138240\pi^2/10^8\),
proves \[
 \eta_j^\ddagger\le A_{\rm rel}j^{-8},
                  \qquad \eta_j^\ddagger/a_j^5\longrightarrow0 .
 \tag{446}
\] Let
\(\chi_j^\ddagger=(K_{\theta_j^\ddagger}-\langle K_{\theta_j^\ddagger}\rangle)\psi_{g_j}\),
and let \(n_j^\ddagger\) be its unchanged raw excitation measure. The
exact norm remainder and raw measure comparison from Section 15 are \[
 \begin{split}
 \|\chi_j^\ddagger+(2/3)(\theta_j^\ddagger)^2v_{\Gamma,g_j}\|
  &\le\eta_j^\ddagger(2/3)(\theta_j^\ddagger)^2
                    \|v_{\Gamma,g_j}\|,\\
 \|n_j^\ddagger-(4/9)(\theta_j^\ddagger)^4\nu_{\Gamma,j}\|_{\rm TV,1}
  &\le\eta_j^\ddagger(2+\eta_j^\ddagger)
          (4/9)(\theta_j^\ddagger)^4\|v_{\Gamma,g_j}\|^2 .
 \end{split}\tag{447}
\] Multiplying the second line by
\(9g_j^4a_j^2/[4(\theta_j^\ddagger)^4\ell_j^7]\) makes its bound
asymptotic to a fixed constant times \(\eta_j^\ddagger/a_j^5\), which
tends to zero by (444) and (446). It follows, with every raw factor
retained, that for \(\Omega>0\), \(t>0\), \[
 \begin{split}
 \|\chi_j^\ddagger\|^2
    &\sim\frac{I_{\rm lat}}{1920}
       (\theta_j^\ddagger)^4g_j^{-4}N_j^7,\\
 n_j^\ddagger([0,\Omega])
    &\sim\frac{(\theta_j^\ddagger)^4g_j^{-4}\ell_j^7}{a_j^2}
                                      \frac{\Omega^5}{1152000\pi^2},\\
 \frac{n_j^\ddagger([0,\Omega])}{\|\chi_j^\ddagger\|^2}
    &\sim a_j^5\frac{\Omega^5}{600\pi^2I_{\rm lat}},\\
 \|e^{-tA_{g_j}/2}\chi_j^\ddagger\|^2
    &\sim
      \frac{(\theta_j^\ddagger)^4g_j^{-4}\ell_j^7}
                       {9600\pi^2a_j^2t^5},\\
 \frac{\langle e^{-tA_{g_j}/2}\chi_j^\ddagger,
          A_{g_j}e^{-tA_{g_j}/2}\chi_j^\ddagger\rangle}
             {\|e^{-tA_{g_j}/2}\chi_j^\ddagger\|^2}
    &\longrightarrow 5/t .
 \end{split}\tag{448}
\] For the energy assertion, every \(\omega^ke^{-t\omega}\) is a bounded
multiplier; use (447) with its supremum norm and (440). No
unbounded-moment conclusion is taken from unweighted weak convergence.

The raw native mass still vanishes. For example the original cusp bounds
give \(D_j^\ddagger\ge c(T_j^\ddagger-C)^2\) for a fixed \(c>0\) on the
proved tail. Thus \((\theta_j^\ddagger)^4g_j^{-4}N_j^7\) is at most a
fixed power of \(j,\xi_j\) times \(e^{-64\pi\xi_j}\), since
\(g_j^{-4}=4\xi_j\). This tends to zero; expansion of the exponential
into its positive power series bounds it against every fixed inverse
power. The finite-energy and time-filtered masses in (448) also vanish,
either from that bound or because they are bounded above by the total
raw mass.

For completeness there is a rigorously selected low-quotient family of
actual filtered native vectors. For each integer \(n\ge1\), apply the
proved fixed-time norm and first-moment limits at \(t=n\), and choose
the least \(j_n>j_{n-1}\) for which both relative errors in (448) are
less than \(1/n\). Existence follows from the two limits and their
positive leading constants. The actual vectors
\(e^{-nA_{g_{j_n}}/2}\chi_{j_n}^\ddagger\) then have quotient
\(5/n+o(n^{-1})\), with their raw masses given by the fourth line of
(448) on this subsequence. The filter is an explicit operation, not an
identification with the original unfiltered state.

This theorem establishes the exact leading finite-energy fraction of the
original weighted covariance and its proved native magnetic transfer on
an explicitly strengthened cusp and coupling selection. It does not
identify that selection with \(g_j^2=1/\log j\) or
\(g_j^2=\kappa_*/(200j)\); it supplies no uniform-volume rate on those
prescribed paths. It also does not identify the depth \(T_j^\ddagger\)
with the initial depth \(T_j=j^2\), or prove survival of a nonabelian
interaction in a spatial continuum. Those original targets remain part
of the active assignment.

\subsection{26. Original finite-angle native states at weak
coupling}\label{original-finite-angle-native-states-at-weak-coupling}

This proof retains the original unmodified cusp angle at each spatial
regulator. It first proves a fixed-box theorem for every fixed
nontrivial angle. The subsequent simultaneous sequence retains the
original depth \(T_j=j^2\), periods, metric, cover and native state. Its
strictly positive couplings are selected from proved limits; they are
not identified with either prescribed coupling trajectory.

\subsubsection{26.1 Exact operator, gauge space and translated
faces}\label{exact-operator-gauge-space-and-translated-faces}

Fix \(L\ge2\), \(a>0\), and all positive edges and elementary faces of
the original open box \(\{-L,\ldots,L\}^3\). Write \[
 N_E=3(2L)(2L+1)^2,\quad M=12L^2(2L+1),\quad
 W(U)=\sum_p(2-\operatorname{tr}U_p),
\] \[
 H_g=\frac{2g^2}{a}H_0+\frac{W}{2g^2a},\quad
 H_0=\sum_eE_e,\quad E_e=-\sum_{\alpha=1}^3X_{e,\alpha}^2,
 \quad T_\alpha=-i\sigma_\alpha/2 .
 \tag{449}
\] The full potential includes every original face and its scalar
\(2bM\). All spectral projections below are on the physical subspace
fixed by every original vertex gauge transformation. The operator and
form domains are the intersections of that subspace with the product
Sobolev spaces \(H^2,H^1\). Its actual positive unit vacuum is
\(\psi_g\), with energy \(E_g\), and \(A_g=H_g-E_g\). The previously
proved variational estimate gives \[
 0\le E_g\le K:=\frac{3\sqrt{N_EM}}a,\qquad
 \|H_0^{1/2}\psi_g\|^2\le\frac{aK}{2g^2}.
 \tag{450}
\] The second inequality follows directly from (449), since \(W\ge0\).

The native links are \(h_1(n)=\exp(-\theta n_2H_c)\), \(h_2=h_3=I\),
\(H_c=-2T_3=i\sigma_3\). Since \(n_2\) is integral, their period in
\(\theta\) is exactly \(2\pi\). Choose a representative with
\(0<|\theta|\le\pi\). For every edge \(e=(n,1)\) with \(n_2\ne0\), let
\(q_e\in SU(2)\) be independent with Haar probability. Put
\(h_e^q=q_eh_eq_e^{-1}\); identity edges remain identity. There are \[
 m=4L^2(2L+1)
 \tag{451}
\] such edges. The exact physical average established earlier is the
positive-kernel self-adjoint contraction \[
 K_\theta F(U)=\int F(((h_e^q)^{-1}U_e)_e)\,dq,\qquad
 \chi_{\theta,g}=K_\theta\psi_g-c_{\theta,g}\psi_g,\quad
 c_{\theta,g}=\langle\psi_g,K_\theta\psi_g\rangle>0.
 \tag{452}
\] Each direction-one translated edge has a different source vertex,
which is why the source gauge averages give precisely these independent
conjugacy averages. Their central convolution factors commute with every
edge Casimir and with the full gauge action. Each vector in (452) is
therefore smooth and physical, and \(\chi_{\theta,g}\perp\psi_g\).

For a fixed realization \(q\), write \(V_e=(h_e^q)^{-1}U_e\). Consider a
\(12\) face with \(n_2=0\). Its bottom direction-one link is unchanged,
and only its top direction-one link is translated. In the original
positive face word, \[
 V_p=U_p\bigl(U_2(n)h_{(n+e_2,1)}^qU_2(n)^{-1}\bigr).
 \tag{453}
\] For a \(12\) face with \(n_2=-1\), only its bottom direction-one link
is translated, and instead \[
 V_p=(h_{(n,1)}^q)^{-1}U_p .
 \tag{454}
\] The extra factor in either formula has eigenvalues
\(e^{i\theta},e^{-i\theta}\), possibly interchanged. If \(B\) is such a
factor and \(S\in SU(2)\), write \(B=B_{1/2}^2\) with eigenvalues of
\(B_{1/2}\) equal to \(e^{\pm i\theta/2}\). Then \[
 I+B=2\cos(\theta/2)B_{1/2},\qquad
 \operatorname{tr}S+\operatorname{tr}(SB)
       =2\cos(\theta/2)\operatorname{tr}(SB_{1/2})
       \le4\cos(\theta/2).
 \tag{455}
\] Here \(\cos(\theta/2)\ge0\), and every \(SU(2)\) trace is at most
two. Cyclicity handles the left factor in (454). There are exactly
\(2L(2L+1)\) faces in each selected plane. Summing their nonnegative
defects, and retaining all other faces in \(W\), proves the global
nonlinear separation \[
 \boxed{W(U)+W(((h_e^q)^{-1}U_e)_e)
    \ge\Delta_{L,\theta}:=
       16L(2L+1)(1-\cos(\theta/2))>0.}
 \tag{456}
\] It holds for every configuration and every twirl realization, not
just near a flat configuration.

Let \(\delta=\Delta_{L,\theta}/3\) and \(Q=\mathbf1_{\{W\le\delta\}}\),
as a bounded multiplication projection. It preserves the physical
subspace because \(W\) is gauge invariant. Equation (456) implies the
exact support identity \[
 QK_\theta Q=0.
 \tag{457}
\] Indeed two arguments with potential at most \(\delta\) would have sum
at most \(2\Delta/3\), contradicting (456). In particular \(K_\theta Q\)
has support in \(W\ge2\delta\).

\subsubsection{26.2 All potential moments on the actual finite-energy
space}\label{all-potential-moments-on-the-actual-finite-energy-space}

For a unit eigenfunction \(u_g\) of the full physical \(H_g\), with
eigenvalue \(\lambda_g\le E_*\), set \(I_n=\int W^n|u_g|^2dU\). The
original derivative estimate is
\(\sum_{e,\alpha}|X_{e,\alpha}W|^2\le16W\). Haar integration by parts
gives \[
 \operatorname{Re}\langle W^nu_g,H_0u_g\rangle
 =\sum_{e,\alpha}\|X_{e,\alpha}(W^{n/2}u_g)\|^2
       -\frac{n^2}{4}\int W^{n-2}
                  \sum_{e,\alpha}|X_{e,\alpha}W|^2|u_g|^2
 \ge-4n^2I_{n-1}.
 \tag{458}
\] The identity is valid for complex eigenfunctions. At zeros of \(W\),
apply it first with \(W+\varepsilon\). Its derivative obeys the same
bound by \(16(W+\varepsilon)\); dominated convergence and the form norm
give (458). This also supplies the domain justification for odd \(n\).

Using (449), \(I_0=1\), and nonnegativity of the kinetic energy for
\(I_1\), we obtain \[
 I_n\le B_n(E_*)g^{2n},\qquad
 B_0=1,\quad B_1=2aE_*,\quad
 B_{n+1}=2aE_*B_n+16n^2B_{n-1}\quad(n\ge1).
 \tag{459}
\] Every constant is fixed before \(g\to0\). For the vacuum, take
\(E_*=K\).

Fix \(\Omega>0\), and let \(P_g=\mathbf1_{(0,\Omega]}(A_g)\) on the
physical subspace. Its rank is bounded by an integer \(R\) for all
sufficiently small positive \(g\). Here is precisely the fixed-box
input: Section 17 proves convergence of each ordered physical eigenvalue
of \(H_g\) to the corresponding eigenvalue of its compact-resolvent
oscillator on the adjoint-invariant chord space. Choose a comparison
eigenvalue strictly greater than \(K+\Omega+1\). Its actual eigenvalue
is greater than \(K+\Omega\) for all sufficiently small \(g\). Equation
(450) then bounds the number of physical eigenvalues below
\(E_g+\Omega\) by this fixed index. This argument is not a claim about
the rank on the full nongauge-invariant configuration Hilbert space.

Choose a physical orthonormal eigenbasis of \(\operatorname{ran}P_g\)
and apply (459) with \(E_*=K+\Omega\). The squared Hilbert--Schmidt
bound obtained by summing its at most \(R\) columns yields \[
 \|(I-Q)P_g\|
 \le\sqrt{R B_n(K+\Omega)}\,\delta^{-n/2}g^n,\qquad
 \|(I-Q)\psi_g\|\le\sqrt{B_n(K)}\,\delta^{-n/2}g^n.
 \tag{460}
\] These hold for every integer \(n\ge1\), after possibly reducing the
fixed-box positive coupling threshold for the rank assertion.

The support identity (457), contraction of \(K_\theta\), and adjoints in
(460) give \[
 \|P_gK_\theta\psi_g\|
 \le\left[\sqrt{R B_n(K+\Omega)}+\sqrt{B_n(K)}\right]
                  \delta^{-n/2}g^n .
 \tag{461}
\] To see the decomposition explicitly, write the input as
\(Q\psi+(I-Q)\psi\); the first image equals \((I-Q)K_\theta Q\psi\), and
the second has norm at most \(\|(I-Q)\psi\|\). Since \(P_g\psi=0\), the
left side equals \(\|P_g\chi_{\theta,g}\|\). Also, by splitting the
first argument and then the second in the scalar overlap and using
(457), \[
 0<c_{\theta,g}\le2\|(I-Q)\psi_g\|
       \le2\sqrt{B_n(K)}\,\delta^{-n/2}g^n .
 \tag{462}
\] Thus the raw low-energy vector and the uncentered vacuum overlap
decrease faster than every power of \(g\), at each fixed box, angle and
energy cutoff.

\subsubsection{26.3 A polynomial lower bound for the unchanged raw
norm}\label{a-polynomial-lower-bound-for-the-unchanged-raw-norm}

A relative spectral conclusion requires a lower bound for the same
native state. Positivity of the actual vacuum provides it without
replacing that vacuum by a trial function. Let
\(T_hF(U)=F((h_e^{-1}U_e)_e)\), so \(\|T_h\psi_g\|=1\). Put \[
 \eta=\frac1{2(1+\sqrt{2maK})},\qquad
 b_*=\frac{\eta^3}{3\pi^3},\qquad c_*=\frac12 b_*^m .
 \tag{463}
\] Restrict \(0<g\le1\). For each twirl variable restrict
\(q_e=\exp(Y_e)\), where \(|Y_e|\le\eta g\) in the metric making
\(T_\alpha\) orthonormal. The curve \(q_e(t)=\exp(tY_e)\) conjugates
\(h_e\). Its left or right logarithmic velocity has norm at most
\(2|Y_e|\), because it is a difference of two orthogonal adjoint images.
Differentiating the product translation along this curve and applying
Cauchy--Schwarz in its edge and Lie indices gives \[
 \|T_{h^q}\psi_g-T_h\psi_g\|
 \le2\eta g\sqrt m\,\|H_0^{1/2}\psi_g\|
 \le\eta\sqrt{2maK}\le\frac12 .
 \tag{464}
\] Every product left translation preserves the full electric form,
since each edge Laplacian is bi-invariant. This proves the uniform
gradient bound along the entire curve used in (464).

Both translated vacua are positive real unit vectors. Their inner
product equals one minus half the squared distance, so it is at least
\(7/8\) in this neighborhood. At all other twirl parameters their inner
product remains nonnegative. The exact Haar density in exponential
coordinates is \[
 \frac1{16\pi^2}
       \left(\frac{\sin(|Y|/2)}{|Y|/2}\right)^2d^3Y.
 \tag{465}
\] For \(r\le1\), the inequality \(\sin u/u\ge2/\pi\) on
\(0\le u\le\pi/2\) and the Euclidean ball volume give Haar measure at
least \(r^3/(3\pi^3)\). Thus the product neighborhood in (464) has
probability at least \(b_*^m g^{3m}\). Integrating only its positive
contributions and using Cauchy--Schwarz with the unit test vector
\(T_h\psi_g\) proves \[
 \boxed{\|K_\theta\psi_g\|
 \ge \langle T_h\psi_g,K_\theta\psi_g\rangle
 \ge c_*g^{3m}.}
 \tag{466}
\] The factor \(1/2\) in \(c_*\) is smaller than the proved \(7/8\); it
provides a uniform displayed lower bound.

Choose an integer \(n>3m\) in (462). For all sufficiently small positive
\(g\), that overlap is at most \(c_*g^{3m}/2\). Orthogonal centering
gives the exact equality \[
 d_{\theta,g}:=\|\chi_{\theta,g}\|^2
       =\|K_\theta\psi_g\|^2-c_{\theta,g}^2
       \ge\frac34c_*^2g^{6m}>0.
 \tag{467}
\] In particular, every raw low-energy mass from (461) is being compared
with the actual nonzero norm, not an assumed unit excitation.

Combining (461) and (467), for each fixed real \(p>0\) choose an integer
\(n\) with \(2n-6m\ge p\). The full physical native probability obeys \[
 \boxed{
 \frac{\langle\chi_{\theta,g},
      \mathbf1_{(0,\Omega]}(A_g)\chi_{\theta,g}\rangle}
      {d_{\theta,g}}
       =O_{L,a,\theta,\Omega,p}(g^p),\qquad g\downarrow0 .}
 \tag{468}
\] The raw numerator bound, with its explicit constants, remains (461)
squared, and the raw denominator lower bound remains (467). No statement
uniform in the growing box is hidden in (468).

\subsubsection{26.4 The unfiltered raw mass and the central
endpoint}\label{the-unfiltered-raw-mass-and-the-central-endpoint}

For \(0<|\theta|<\pi\), the raw norm also tends to zero. We prove a rate
that suffices for later simultaneous selection. Choose the original edge
\(e=((0,1,0),1)\). Its central factor in \(K_\theta\) is
\(C_{e,\theta}\), with exact spin-\(s\) multiplier \[
 c_s(\theta)=\frac{\sin((2s+1)\theta)}
                    {(2s+1)\sin\theta},\qquad
 |c_s(\theta)|\le\frac1{(2s+1)|\sin\theta|}.
 \tag{469}
\] This is used only when \(\sin\theta\ne0\). All other central factors
commute with it and are contractions.

Let \(P_d^e\) be the projection onto edge spins with \(2s+1\le d\),
\(d\) a positive integer. Its one-edge projection kernel has constant
diagonal \(D_d=\sum_{r=1}^dr^2=d(d+1)(2d+1)/6\le d^3\). For fixed other
links, the original face \(p=((0,0,0),12)\) has word
\(U_p=A B U_e^{-1}D_0^{-1}\), where \(A=U_1(0,0,0)\), \(B=U_2(1,0,0)\),
and \(D_0=U_2(0,0,0)\) are fixed by the exterior. Cyclicity gives
\(\operatorname{tr}U_p=\operatorname{tr}(D_0^{-1}ABU_e^{-1})\). Haar
invariance and inversion therefore make the volume of its defect cap
independent of the exterior. If \(W\le u\), its face defect is at most
\(u\). The Haar class angle \(\phi\in[0,\pi]\) has density
\(2\sin^2\phi/\pi\), so the exact defect \(w=2-2\cos\phi\) has density
\[
 \frac1\pi\sqrt{w-w^2/4}\,dw,\qquad 0\le w\le4 .
\] Consequently the allowed subset of \(U_e\) has Haar measure at most
\[
 C_{\rm cap}u^{3/2},\qquad C_{\rm cap}=\frac2{3\pi},
                         \quad u>0.
 \tag{470}
\] Indeed integrate the displayed density over \(0\le w\le\min(u,4)\),
bound \(\sqrt{w-w^2/4}\le\sqrt w\), and enlarge the upper integration
limit to \(u\). This proves the estimate also when \(u>4\).

For any such fiber subset \(B\), the Hilbert--Schmidt norm squared of
\(P_d^e\mathbf1_B\) is \(\int_BD_d\,dU_e\). Apply this bound fiberwise
to \(\mathbf1_{\{W\le u\}}\psi_g\), and apply (459) to its complementary
part. The triangle inequality squared gives the full edge-spin
cumulative probability \[
 F_g(d):=\|P_d^e\psi_g\|^2
       \le2C_{\rm cap}d^3u^{3/2}
                  +2B_n(K)g^{2n}u^{-n}.
 \tag{471}
\] The fiberwise argument uses the full original Haar space; its input
is the actual physical vacuum. It requires no factorization of the
vacuum measure.

Use \(n=4\), and choose \(u_d=g^2(dg)^{-6/11}>0\) separately for each
spin cutoff. This cutoff is used only in (471), which is valid for every
positive \(u\); it does not change the fixed separation cutoff in (457).
Substitution gives \[
 F_g(d)\le A_4(dg)^{24/11},\qquad
 A_4=\frac4{3\pi}+2B_4(K).
\] The recurrence (459), with \(x=2aK\), gives exactly \(B_2=x^2+16\),
\(B_3=x^3+80x\), and \(B_4=x^4+224x^2+2304\). If \(\mathsf d=2s_e+1\)
denotes the positive integer edge-spin dimension in the actual squared
coefficient distribution, Tonelli's theorem gives the exact identity \[
 \mathbb E_{\psi_g}\mathsf d^{-2}
 =\int_1^\infty2t^{-3}F_g(\lfloor t\rfloor)\,dt .
 \tag{472}
\] For each integer \(d\), the contribution of its probability mass is
\(\int_d^\infty2t^{-3}\,dt=d^{-2}\), proving the identity before
summing. For \(0<g\le1\), split at \(g^{-1}\). On the first part use the
displayed cumulative bound and \(\lfloor t\rfloor\le t\), and on the
second use \(F_g\le1\). This gives \[
 \mathbb E_{\psi_g}\mathsf d^{-2}
 \le2A_4g^{24/11}\int_1^{g^{-1}}t^{24/11-3}\,dt+g^2
 =11A_4(g^2-g^{24/11})+g^2.
\] Parseval and (469), with the commuting contraction factors retained,
therefore prove \[
 d_{\theta,g}\le\|K_\theta\psi_g\|^2
 \le\frac{1+44/(3\pi)+22B_4(K)}{\sin^2\theta}\,g^2,
 \qquad 0<g\le1.
 \tag{473}
\] For every sufficiently small positive \(g\), (467) supplies the
simultaneous strictly positive lower bound for this same raw norm. The
angle-dependent constant and every original norm remain explicit.

At \(\theta=\pi\), every direction-one link is central,
\(h_1(n)=(-1)^{n_2}I\). Then \(K_\pi=T_h\) is a unitary involution.
Equations (462) and (467) are still valid, but now \[
 d_{\pi,g}=1-c_{\pi,g}^2\longrightarrow1 .
 \tag{474}
\] The complete low-energy probability still satisfies (468). At
\(\theta=0\) modulo \(2\pi\), \(K_\theta=I\) and the centered state is
exactly zero. These branches preserve the discrete center endpoint and
the original continuous path; they do not identify the initial cusp
exhaustion with its endpoint.

\subsubsection{26.5 Original geometry on one common actual
sequence}\label{original-geometry-on-one-common-actual-sequence}

Let \(j\) range over an eventual dyadic tail \(j=2^s\), on which the
original cusp domain is defined and \(0<\theta_j<\pi\). Retain exactly
\[
 T_j=j^2,\quad L_j=j^2,\quad a_j=\frac1{100j},\quad
 D_j=L_{j^2}q_{j^2}+6m_{j^2}^2,\quad
 \theta_j=\frac{2\pi}{10^4j^2D_j},
 \quad v_j=e^{-2\pi j},\quad t_{c,j}=v_j^j .
 \tag{475}
\] The original period entries, \(B_T\), coordinate permutation, cover
\(j\), pullback metric and signed determinant \(-j^2D_j\) remain those
of Section 2. Its proved estimate \(D_j=j^4+O(j^2)\) is used only after
the exact angle in (475); in particular
\(\theta_j\sim(2\pi/10^4)j^{-6}\), and \(0<\theta_j<\pi\) on an eventual
tail.

For a fixed \(j\), (468),(473) are actual \(g\downarrow0\) results at
that box and its original angle. Set \[
 R_j=\frac{4\sqrt3}{a_j}+1.
 \tag{476}
\] The comparison covariance measure (228) has pair energies at most
\(4\sqrt3/a_j\), because every \(\sigma_\nu\le\sqrt{12}\). The fixed-box
convergence (231) therefore implies that the actual covariance
probability of \((R_j,\infty)\) tends to zero.

Keep all finite-stage tests defining Section 25's actual coupling
selection, including Sections 18 and 22--24 and the retained full
local-energy tests. Add \[
 \begin{split}
 &\zeta_{\theta_j,g}([0,R_j])<a_j^5/j,\qquad
    \tfrac34c_{*,j}^2g^{6m_j}\le d_{\theta_j,g}<j^{-2},\\
 &\zeta_{\Gamma,L_j,a_j,g}((R_j,\infty))<j^{-2},
                       \qquad 0<g<j^{-5}.
 \end{split}\tag{477}
\] Here \(\zeta_{\theta,g}\) denotes precisely the probability in (468);
it has no vacuum atom. Here \(m_j=4L_j^2(2L_j+1)\) and \(c_{*,j}\) is
exactly (463) at \(L_j,a_j\); it does not depend on \(g\). The first two
new inequalities hold throughout a sufficiently small positive interval
by (467),(468),(473). The third holds by the fixed-box covariance
convergence and its stated comparison support. Every previous stage test
also holds on a sufficiently small interval by its proved fixed-box
limit. Their finite intersection is nonempty. Define \(k_j\) as the
least positive integer such that all dyadics \(2^{-q}\), \(q\ge k_j\),
meet this list, and set \(g_j=2^{-k_j}\). This gives a single exact
actual sequence, retaining all original physical coefficients \[
 \kappa_j=200j\,2^{-2k_j},\quad
 b_j=50j\,2^{2k_j},\quad
 \xi_j=2^{4k_j-2},\quad
 2b_jM_j=100j\,2^{2k_j}M_j.
 \tag{478}
\] There is no claimed numerical uniform-volume rate for the least
exponent.

Write \(\chi_j=\chi_{\theta_j,g_j}\), retaining the original, unchanged
depth (475). The same physical Hamiltonians have closing gaps by the
retained Section 18 tests. Meanwhile, for each fixed \(\Omega>0\), \[
 \begin{split}
 &0<\|\chi_j\|^2<j^{-2},\qquad
   \zeta_{\theta_j,g_j}([0,\Omega])\le a_j^5/j
                                     =o(a_j^5),\\
 &\zeta_{\Gamma,L_j,a_j,g_j}([0,\Omega])
       \sim a_j^5\frac{\Omega^5}{600\pi^2I_{\rm lat}},\\
 &\frac{\zeta_{\theta_j,g_j}([0,\Omega])}
           {\zeta_{\Gamma,L_j,a_j,g_j}([0,\Omega])}
                                  \longrightarrow0 .
 \end{split}\tag{479}
\] The first estimate uses \(\Omega<R_j\) eventually; the second is the
retained all-mode theorem of Section 25 on this refined sequence. The
explicit raw lower bound for the first line is (467) at its actual
\(L_j,a_j,\theta_j,g_j\), with its full \(m_j,c_{*,j}\); the raw
spectral numerator is bounded by \(\|\chi_j\|^2a_j^5/j\). Thus both raw
factors accompany the probability assertion.

There is a stronger exact spectral relation at the growing cutoff (476).
By (477), \[
 \sup_{B\ {\rm Borel}}|\zeta_{\theta_j,g_j}(B)
                  -\zeta_{\Gamma,L_j,a_j,g_j}(B)|
       \ge1-j^{-2}-a_j^5/j\longrightarrow1 .
 \tag{480}
\] The supremum is at most one, so its limit is one; the total variation
mass norm tends to two. Splitting the scalar product at the same
physical projection and using Cauchy--Schwarz gives, with all raw norms
displayed, \[
 \frac{|\langle\chi_j,v_{\Gamma,g_j}\rangle|}
           {\|\chi_j\|\,\|v_{\Gamma,g_j}\|}
       \le\sqrt{a_j^5/j}+j^{-1}\longrightarrow0 .
 \tag{481}
\] Thus the native state and its electric angular coefficient become
orthogonal in their common actual Hilbert space on this selected
original geometry, despite both probabilities escaping fixed physical
energies.

\subsubsection{26.6 The angular comparison and its precise
nonuniformity}\label{the-angular-comparison-and-its-precise-nonuniformity}

Define the actual relative Taylor remainder at the original angle by \[
 \eta_j^{\rm orig}
  =\frac{\|\chi_j+(2/3)\theta_j^2v_{\Gamma,g_j}\|}
          {(2/3)\theta_j^2\|v_{\Gamma,g_j}\|}.
 \tag{482}
\] The retained all-mode mass theorem and \(N_j=2j^2+1\) give \[
 \frac23\theta_j^2\|v_{\Gamma,g_j}\|
 \sim
 \frac{8\pi^2}{3\cdot10^8}
       \sqrt{\frac{3I_{\rm lat}}{20}}\,
               g_j^{-2}j^{-5}\longrightarrow\infty,
 \tag{483}
\] since \(g_j<j^{-5}\). Every factor follows from the exact original
\(D_j\), with its asymptotic applied at this final step. Because
\(\|\chi_j\|\le1\), the reverse triangle inequality proves \[
 |\eta_j^{\rm orig}-1|
       \le\frac{\|\chi_j\|}
                 {(2/3)\theta_j^2\|v_{\Gamma,g_j}\|}
           \longrightarrow0 .
 \tag{484}
\] This establishes actual failure of the small relative Taylor
approximation on this selected sequence, rather than merely divergence
of one certified upper bound. It makes no assertion of that failure on
the logarithmic or fixed-electric-coefficient paths.

The strengthened cusp of Section 25 is still defined on these same
Hamiltonians and couplings, using its exact \(\theta_j^\ddagger\) and
actual centered state \(\chi_j^\ddagger\). Its proved relative error is
\(o(a_j^5)\); the exact probability comparison (167) therefore bounds
its spectral probability distance from the covariance probability by
that error. Equations (480)--(481) therefore imply \[
 \sup_B|\zeta_{\theta_j,g_j}(B)-\zeta_j^\ddagger(B)|
       \longrightarrow1,\qquad
 \frac{|\langle\chi_j,\chi_j^\ddagger\rangle|}
                 {\|\chi_j\|\,\|\chi_j^\ddagger\|}
       \longrightarrow0 .
 \tag{485}
\] For the second assertion, the normalized strengthened state converges
in norm to the negative normalized covariance vector by the exact
relative norm remainder in (447); applying (481) and the triangle
inequality proves the limit. This comparison records ratios only after
retaining each original raw vector and the raw estimates
(448),(467),(479).

There is also an exact unscaled map between the two angular
presentations. Let \(\mathcal R_\theta=K_\theta-I+(2/3)\theta^2\Gamma\),
on \(D(\Gamma^2)\), and let \(P_\psi^\perp=I-|\psi\rangle\langle\psi|\).
For any two nonzero real angles, \[
 \begin{split}
 \chi_\theta-\frac{\theta^2}{\phi^2}\chi_\phi
 &=P_\psi^\perp
       \left(\mathcal R_\theta-\frac{\theta^2}{\phi^2}
                       \mathcal R_\phi\right)\psi,\\
 \left\|\chi_\theta-\frac{\theta^2}{\phi^2}\chi_\phi\right\|
 &\le\frac29\theta^2(\theta^2+\phi^2)\|\Gamma^2\psi\|.
 \end{split}\tag{486}
\] Subtract the two exact centered expansions and use the all-angle
graph remainder from Section 15 to prove these identities. They retain
both amplitude factors and identify the map whose relative behavior
(484)--(485) now determines on the selected sequence.

The original cusp state thus has a proved full-Hamiltonian spectral
escape sequence at weak positive coupling without changing its external
cusp depth. This result extends the original-state calculation beyond
the earlier strengthened-cusp theorem. The prescribed coupling paths and
an interacting spatial continuum with nonzero original native low-energy
weight remain unresolved; (477) is an explicitly selected sequence, not
a substitute definition of those remaining targets.

\subsection{26.7 A sharper actual-vacuum fourth electric
moment}\label{a-sharper-actual-vacuum-fourth-electric-moment}

This continuation records a bound needed when testing the original angle
on prescribed coupling paths. It is a fixed finite box statement, with
every constant and the original Hamiltonian retained. It does not assert
a volume-uniform continuum estimate.

Fix \(L\ge2\), \(a>0\), \(0<g\le1\), and use \[
 H=\kappa H_0+bW,\quad \kappa=2g^2/a,\quad b=(2g^2a)^{-1},
\quad W=\sum_p(2-\operatorname{tr}U_p),\quad 0\le W\le4M .
\] Let \(\psi\) be the actual positive unit vacuum, \(H\psi=E\psi\),
\(E\le K=3\sqrt{N_EM}/a\). The moment recurrence already proved for this
same state is \[
 \langle W^n\rangle_\psi\le B_n(K)g^{2n},\qquad
 B_0=1,\ B_1=2aK,\ B_{n+1}=2aK B_n+16n^2B_{n-1}.
 \tag{487}
\] In particular \(B_2=x^2+16\) and \(B_4=x^4+224x^2+2304\), \(x=2aK\).

The eigen-equation is an equality of smooth functions: \[
 H_0\psi=f\psi,\qquad f=(E-bW)/\kappa. \tag{488}
\] Applying \(H_0\) once more and using
\(E_e=-\sum_\alpha X_{e,\alpha}^2\) gives the exact product identity \[
 H_0^2\psi=f^2\psi-\frac b\kappa[H_0,W]\psi,\qquad
 [H_0,W]\psi=(H_0W)\psi-2\sum_{e,\alpha}(X_{e,\alpha}W)(X_{e,\alpha}\psi).
 \tag{489}
\] The sign follows by expanding
\(H_0(f\psi)=fH_0\psi+(H_0f)\psi-2\sum(Xf)(X\psi)\), with
\(H_0f=-(b/\kappa)H_0W\) and \(Xf=-(b/\kappa)XW\).

Every fundamental face trace is an \(H_0\)-eigenfunction with eigenvalue
\(3\). Hence \[
 H_0W=3W-6M. \tag{490}
\] The exact full-face gradient estimate is
\(\sum_{e,\alpha}|X_{e,\alpha}W|^2\le16W\). Since
\(\langle H_0\rangle_\psi=(E-b\langle W\rangle_\psi)/\kappa\le E/\kappa\),
Cauchy--Schwarz and (487) imply \[
\begin{aligned}
\|[H_0,W]\psi\|
&\le3\|W\psi\|+6M
 +2\Big(\int|\nabla W|^2\psi^2\Big)^{1/2}
       \Big(\int|\nabla\psi|^2\Big)^{1/2}\\
&\le3\sqrt{B_2(K)}\,g^2+6M+8aK=:C_{LW}(L,a).
\end{aligned}\tag{491}
\] Indeed the product under the square root is at most
\((16B_1(K)g^2)(Ka/(2g^2))=16aK\,B_1(K)/2\le16a^2K^2\), and the
prefactor \(2\) gives \(8aK\) after using \(B_1=2aK\). (Keeping the
unsimplified product gives the same displayed bound.)

For real \(r,s\), \((r+s)^4\le8(r^4+s^4)\). Therefore (487)--(488) give
\[
\|f^2\psi\|
\le\frac{\sqrt{8(K^4+b^4B_4(K)g^8)}}{\kappa^2}
=\frac{a^2}{4g^4}
\sqrt{8\left(K^4+\frac{B_4(K)}{16a^4}\right)}.
\tag{492}
\] Combining (489), (491), and \(b/\kappa=1/(4g^4)\) proves the explicit
actual-vacuum bound \[
\boxed{\ \|H_0^2\psi\|\le C_2(L,a)\,g^{-4},\quad
C_2(L,a)=\frac{a^2}{4}\sqrt{8\left(K^4+\frac{B_4(K)}{16a^4}\right)}
+\frac14C_{LW}(L,a).\ }\tag{493}
\] No replacement state or unlisted boundary face is used.

Finally \(\Gamma=\sum_{(n,1)}n_2^2E_e\) is a positive joint spectral
multiplier of the commuting edge Casimirs and \(0\le n_2^2\le4L^2\). The
joint spectral theorem gives \[
\|\Gamma^2\psi\|\le16L^4\|H_0^2\psi\|
\le16L^4C_2(L,a)\,g^{-4}. \tag{494}
\] This is sharper in \(g\) than the earlier \(g^{-7}\) estimate at each
fixed box. Its explicit \(L,a\)-dependence still grows with the
regulator, so PM8 alone does not prove the prescribed logarithmic or
fixed-electric-coefficient continuum path. It is an exact input for the
centered spin/remainder estimate now under construction.

\section{Centered spin remainder
advance}\label{centered-spin-remainder-advance}

\subsection{26.8 A polynomial covariance lower
bound}\label{a-polynomial-covariance-lower-bound}

This subsection records an exact centered scalar expansion and the
strongest finite-box covariance estimate obtained without replacing the
actual positive vacuum. It is a fixed-box result, with the remaining
volume-uniform centered-remainder estimate stated explicitly below.

\subsection{Exact joint-spin fourth
coefficient}\label{exact-joint-spin-fourth-coefficient}

On a joint spin block, with \(\gamma=\sum_e n_2^2 \lambda_e\),
\(\lambda_e=s_e(s_e+1)\), and independent uniform magnetic weights as in
(138)--(140), write

\(r_\theta(\gamma,lambda)=k(\theta)-1+(2/3)\theta^2 \gamma\).

The exact fourth term from the cosine expansion is

\(r_\theta = (2/9) \theta^4 \gamma^2 -(4/45)\theta^4 \sum_e n_2^4 \lambda_e^2 -(2/45)\theta^4 \sum_e n_2^4 \lambda_e + R_6\).

Here \(R_6\) is the finite block remainder obtained by replacing
\(cos x\) by its sixth-order Taylor polynomial; the displayed
coefficient follows directly from (140), without suppressing any
diagonal spin term. After centering, the deterministic part of
\((2/9)\theta^4 \gamma^2\) is removed by \(P_psi^perp\), leaving the
exact covariance contribution
\((2/9)\theta^4 P_psi^perp((\Gamma^2-<\Gamma^2>)psi)\) plus the two
diagonal terms and the sixth remainder. This identifies precisely what a
variance-sensitive estimate must control; the old bound
\((2/9)\theta^4\|\Gamma^2 psi\|\) is valid but does not exploit this
cancellation.

No bound that replaces \(\|\Gamma^2 psi\|\) by a variance has been
asserted: the joint spin blocks carry arbitrary actual-vacuum
coefficients, and a scalar variance inequality alone is false without an
additional moment/concentration input.

\subsection{Polynomial covariance lower bound from the exact
vacuum}\label{polynomial-covariance-lower-bound-from-the-exact-vacuum}

Let \(W=\sum_p(2-tr U_p)\), and let \(u=log psi\). The exact operator
identity is

\(\Gamma W = (3/4) \sum_p S_p (w_p-2)\),

where \(S_p=sum_{e \in p, direction 1} n_2(e)^2\), and
\(A_L=\sum_p S_p=sum_{e direction 1} n_2(e)^2 r_e\). In particular
\(A_L\) is of order \(L^5\), with all boundary incidence factors
retained. The actual energy bound gives

\(A_L = (4 L^2(2L+1)/3)(4L^2+2L+1)\).

\(<W> \le  E/b \le  6 g^2 \sqrt(N M) = 6 g^2 \sqrt(r) M\),
\(r=N/M\le 5/4\).

Since \(S_p \le  2 L^2\),

\(\sum_p S_p <w_p> \le  2 L^2 <W>\) and therefore

\(<\Gamma W> \le  -(3/2)A_L + (3/2)L^2 <W>\).

Using the displayed energy estimate, the explicit condition
\(6 g^2 \sqrt(r) L^2 M \le  A_L/2\) implies \(<\Gamma W> \le  -3A_L/4\);
hence \(|<\Gamma W>| \ge  3A_L/4\) on that domain. This is an ordinary
inequality, with no lower bound on individual plaquette expectations
assumed.

Integration by parts gives
\(<\Gamma W\ge 2 Re int (\Gamma-gradient W) dot (\Gamma-gradient u) psi^2\).
Therefore

\(|<\Gamma W>|^2 \le  4 S_\Gamma <\Gamma>\),
\(S_\Gamma=<sum_{e dir1} n_2^2 |grad_e W|^2>\), while
\(<\Gamma> \le  L^2 <H_0> \le  3 L^2 \sqrt(NM) \sqrt(xi) = 3 L^2 M \sqrt(r xi)\).

Consequently, whenever the preceding small-coupling inequality holds,

\(S_\Gamma \ge  (A_L^2)/(12 L^2 M \sqrt(r xi))\).

This is a polynomial, volume-explicit replacement for the previous
local-density bound \(S_\Gamma \ge  (3/4)e^{-32 pi xi} mathscr S_L\).
For an explicit denominator, use the retained exact identity
\([H_0,W]psi=(3W-6M)psi-2 grad W dot grad psi\). The actual-vacuum graph
estimate gives \(\|[H_0,W]psi\| \le  6M+3 \sqrt(M_2)+8 \sqrt(I)\), where
\(I=int W |grad log psi|^2 psi^2 \le  a^2 K^2+3M = 9rM^2+3M\) and
\(M_2\le B_2(K)g^4\). Since \(M\ge 240\), \(r\le 5/4\), and \(g\le 1\),
one has \(B_2(K)g^4 \le  36rM^2+16\), so this denominator is \(<60M\)
(retaining this coarse integer only for a transparent bound). Therefore
the preceding inequalities imply the fully explicit covariance estimate

\(\| (\Gamma-<\Gamma>)psi \| \ge  \sqrt(xi) A_L^2 /(720 L^2 M^2 \sqrt(r))\),

whenever \(6g^2 \sqrt(r)L^2M \le  A_L/2\). This is polynomial in \(L\)
and \(xi\) and applies to the actual interacting vacuum; it replaces the
former exponential-in-\(xi\) lower bound. It still does not control the
centered fourth remainder, whose diagonal and sixth-order terms require
a separate weighted fourth-moment estimate, so no prescribed-path
spectral conclusion is claimed here.

\subsection{Consequence and non-claim}\label{consequence-and-non-claim}

The centered expansion is an exact morphism on each finite joint-spin
block and the all-block graph domain. The polynomial covariance lower
bound holds for the actual interacting ground state on its displayed
small-coupling condition. It does not establish nonzero native
low-energy weight in the prescribed continuum paths, nor does it prove
an interacting four-dimensional continuum; those remain open.

\section{Fixed-box native angle scaled as theta\_g = g
tau}\label{fixed-box-native-angle-scaled-as-theta_g-g-tau}

This note proves the finite spatial regulator limit of the actual vacuum
native state when the background angle is \(\theta_g=g\tau+o(g)\), with
\(L\ge2\), \(a>0\), and fixed nonzero real \(\tau\). It is a fixed-box
theorem; no interchange with \(L\to\infty\) is asserted.

\subsection{1. Exact finite operator and
chart}\label{exact-finite-operator-and-chart}

Keep the full physical Hamiltonian \[
 H_g=\frac{2g^2}{a}H_0+\frac{1}{2g^2a}W,\qquad
 H_0=\sum_eE_e,\quad E_e=-\sum_{\alpha=1}^3X_{e,\alpha}^2,
\] with every edge and face, and let \(\psi_g\) be its strictly positive
unit physical vacuum. Let \(K_\theta\) be the exact independent
conjugacy twirl on the \(m=4L^2(2L+1)\) direction-one edges with
\(n_2\ne0\), as in (452), and put
\(\chi_{\theta,g}=K_\theta\psi_g-c_{\theta,g}\psi_g\),
\(c_{\theta,g}=\langle\psi_g,K_\theta\psi_g\rangle\).

Use the exact maximal-tree, chord, logarithm and Haar-density isometry
\(\mathcal B_g\) of (194)--(197), followed by \(x=G^{1/2}Oz\) in every
colour. Write \(R=T^*G^{-1/2}O\), an isometry from chord modes onto the
original transverse edge space. The actual-vacuum theorem gives \[
 f_g:=\mathcal B_g\psi_g\to\Phi_0,\qquad
 \Phi_0(z)=\prod_\nu(\sigma_\nu/4\pi)^{3/4}e^{-\sigma_\nu|z_\nu|^2/8}
\] strongly in \(L^2\), together with convergence of every fixed finite
physical spectral projection and of the required weighted graph vectors.

For each twirled edge let \(\eta_e(q)=n_2\,\operatorname{Ad}_{q_e}T_3\)
for direction one and \(\eta_e=0\) otherwise. Since \(H_c=-2T_3\) and
\(h_e=\exp(-\theta n_2H_c)=\exp(2\theta n_2T_3)\), the inverse
translated link is \[
 (h_e^q)^{-1}U_e=\exp(-2\theta n_2\operatorname{Ad}_{q_e}T_3)U_e.
\] Thus the tangent displacement in original edge coordinates is
\(-2\tau\eta(q)\) when \(\theta_g/g\to\tau\). The transverse chord
displacement is exactly \[
 \delta(q):=-2\tau R^*\eta(q). \tag{S1}
\] The sign follows from the inverse left action.

\subsection{2. Strong convergence of the twirl, including
tails}\label{strong-convergence-of-the-twirl-including-tails}

Let \(u\in C_c^\infty(\mathbb R^{3r})\) be invariant under the residual
simultaneous adjoint color rotation; such functions form a dense
physical core. For each fixed \(q\), apply the exact inverse chart to
the configuration with tree links equal to the identity and chord links
\(Z=\exp(gG^{1/2}Oz)\). Every translated chord word is a product of
finitely many links. Bi-invariance of the SU(2) metric and
\(d(\exp(Y),I)\le |Y|\) give, uniformly in \(q\) and on the support of
\(u\), \[
 z'_g(z,q)=z+\delta(q)+o(1),\qquad
 \partial_z^\alpha z'_g(z,q)\to\partial_z^\alpha(z+\delta(q)). \tag{S2}
\] The exact words retain all source/target tree factors; the only
first-order term is the edge displacement \(-2\tau\eta\), transported by
\(R^*\). Haar factors satisfy
\(\mathcal J(gG^{1/2}Oz)^{1/2}/\mathcal J(gG^{1/2}Oz'_g)^{1/2}\to1\)
uniformly on this compact set.

The same word estimate supplies tails rather than merely local
convergence. A path through the rooted box has at most \(6L\) tree
links, so each chord word in (194) has at most \(\ell_*:=12L+1\)
factors. Every translated factor has \(d(h_e^q,I)\le2|g\tau+o(g)|L\).
Hence, with \(C_L^{\rm edge}:=2L\ell_*\), bi-invariance gives a uniform
word displacement at most \(C_L^{\rm edge}|g\tau+o(g)|\). If
\(|z|\le R\), the unshifted chord logarithms are bounded by
\(g\|G^{1/2}\|R\), and the translated ones by
\(g(\|G^{1/2}\|R+C_L^{\rm edge}(|\tau|+1))\) for small \(g\). The
inverse translation has the same bound. Since \(R^*\) is an isometry and
\(\|G^{-1/2}\|\) is fixed, this gives a common rescaled support ball of
radius \(C(R+1)\), uniformly in \(q\). Thus the exact twirl of
\(\mathcal B_g^*u\) has compact support independent of sufficiently
small \(g\). Dominated convergence in the finite Haar average proves \[
 \mathcal B_gK_{\theta_g}\mathcal B_g^*u\to A_\tau u,\qquad
 (A_\tau u)(z)=\int u(z+\delta(q))\,dq \tag{S3}
\] in \(L^2\). This uses the original-chart pullback and word tails; it
does not assert that a fixed-\(q\) quotient pullback is a contraction.

The operators \(K_{\theta_g}\) are self-adjoint positive contractions on
the exact physical Hilbert space. Choose smooth compact cutoffs \(u_R\)
with \(u_R\to\Phi_0\) in Gaussian graph norm. The contraction bounds,
(203), and (S3), first with fixed \(R\) and then with \(R\to\infty\),
imply \[
 \mathcal B_gK_{\theta_g}\psi_g\to A_\tau\Phi_0\quad\hbox{strongly},\qquad
 c_{\theta_g,g}\to c_\tau:=\langle\Phi_0,A_\tau\Phi_0\rangle. \tag{S4}
\] Consequently \[
 \mathcal B_g\chi_{\theta_g,g}\to w_\tau:=A_\tau\Phi_0-c_\tau\Phi_0,\quad
 d_{\theta_g,g}:=\|\chi_{\theta_g,g}\|^2\to d_\tau=\|w_\tau\|^2. \tag{S5}
\]

\subsection{3. Exact Gaussian translation and centered spectral
kernel}\label{exact-gaussian-translation-and-centered-spectral-kernel}

Define \(\beta(q)=-\sqrt{\Lambda/8}\,\delta(q)\), where
\(\Lambda=\operatorname{diag}(\sigma_\nu)\) acts on each colour.
Translation of the product Gaussian is the coherent vector \[
 U_{\delta(q)}\Phi_0(z):=\Phi_0(z+\delta(q)),\qquad
 U_{\delta(q)}\Phi_0=e^{-\|\beta(q)\|^2/2}\sum_{N\in\mathbb N^{3r}}
 \frac{\beta(q)^N}{\sqrt{N!}}\Phi_N . \tag{S6}
\] The oscillator excitation energy of \(\Phi_N\) is
\(E_N=\sum_{\nu,\alpha}\sigma_\nu N_{\nu\alpha}/a\). Therefore the
uncentered limiting spectral measure of \(A_\tau\Phi_0\) is \[
 \widetilde\mu_\tau=\sum_N |b_N(\tau)|^2\delta_{E_N},\qquad
 b_N(\tau)=\int e^{-\|\beta(q)\|^2/2}\frac{\beta(q)^N}{\sqrt{N!}}\,dq. \tag{S7}
\] The centered measure is \[
 \mu_\tau=\widetilde\mu_\tau-c_\tau^2\delta_0,\qquad
 c_\tau=b_0(\tau)=\int e^{-\|\beta(q)\|^2/2}dq. \tag{S8}
\] The subtraction is exact because all nonzero oscillator levels are
orthogonal to \(\Phi_0\). Its total mass is \[
 d_\tau=\iint e^{-\frac1{16}(\delta(q)-\delta(q'))^T\Lambda(\delta(q)-\delta(q'))}dq\,dq'-c_\tau^2. \tag{S9}
\] The exact two-twirl Euclidean-time kernel, for \(t\ge0\), is \[
 \langle A_\tau\Phi_0,e^{-tA_{\rm osc}}A_\tau\Phi_0\rangle
 =\iint\exp\!\left[-\frac{\|\beta(q)\|^2+\|\beta(q')\|^2}{2}
 +e^{-t\Lambda/a}\beta(q)\!\cdot\!\beta(q')\right]dq\,dq'. \tag{S10}
\] For the centered state subtract \(c_\tau^2\). These formulas retain
all modes and colours.

If \(\tau\ne0\), then \(d_\tau>0\). The random vector \(R^*\eta(q)\) is
not almost surely constant: vary one independent \(q_e\) on an edge with
\(n_2\ne0\), and pair with a plaquette curl containing that edge; the
transverse projection cannot vanish for every such variation. Hence
\(\beta\) is nonconstant and strict Cauchy--Schwarz in (S9) gives
\(d_\tau>0\). Equivalently, reflection invariance kills linear
coefficients while \[
 \sum_{\nu,\alpha}|b_{2e_{\nu\alpha}}|^2>0,
\] because \(\int\|\Lambda^{1/2}R^*\eta(q)\|^2dq>0\). For \(\tau=0\),
\(A_0=I\), \(w_0=0\).

\subsection{4. Transfer of the full physical spectral
measure}\label{transfer-of-the-full-physical-spectral-measure}

Let \(A_g=H_g-E_g\). Fixed-box spectral-projection convergence from
Section 17 applies to bounded intervals whose boundaries avoid the
oscillator eigenvalues. Combining it with (S5), for every bounded
continuous \(F\) on \([0,\infty)\), \[
 \langle\chi_{\theta_g,g},F(A_g)\chi_{\theta_g,g}\rangle
 \longrightarrow\int F(\lambda)\,d\mu_\tau(\lambda). \tag{S11}
\] First use a finite oscillator spectral cutoff, where projection
convergence is in norm and vectors converge strongly; bound the
complementary vector norm by the coherent Gaussian tail uniformly in
\(q\), and send the cutoff to infinity. For \(F\equiv1\) this is (S5),
so normalized native probabilities converge weakly to
\(\mu_\tau/d_\tau\) whenever \(\tau\ne0\). Every finite oscillator
interval has the exact mass obtained by summing (S7)--(S8), and some
finite positive level has nonzero mass whenever \(\tau\ne0\). No
simultaneous infinite-volume conclusion is made.

\subsection{26.10 Prescribed paths: the centered fourth and sixth
remainders remain volume
uncontrolled}\label{prescribed-paths-the-centered-fourth-and-sixth-remainders-remain-volume-uncontrolled}

This note isolates what the exact centered spin expansion and the proved
actual-vacuum moments do, and do not, imply on the two prescribed
coupling paths. The state is always the positive ground vector of the
full finite-box Hamiltonian; no Gaussian replacement is made.

Write \[
 v_\Gamma=(\Gamma-\langle\Gamma\rangle_\psi)\psi,
 \qquad
 \chi_\theta=P_\psi^\perp(K_\theta-I)\psi .
\] The exact joint-spin expansion gives \[
 \chi_\theta=-\frac23\theta^2v_\Gamma+P_\psi^\perp R_{4,\theta}\psi,
 \qquad
 \|R_{4,\theta}\psi\|\le\frac29\theta^4\|\Gamma^2\psi\| .
 \tag{P1}
\] The centered fourth coefficient is more precise blockwise: \[
 r_\theta=\frac29\theta^4\gamma^2
 -\frac4{45}\theta^4\sum_e n_2^4\lambda_e^2
 -\frac2{45}\theta^4\sum_e n_2^4\lambda_e+R_{6,\theta},
 \tag{P2}
\] where the deterministic part of the first term disappears under
\(P_\psi^\perp\), but the variance of \(\Gamma^2\), the two diagonal
spin terms, and the sixth remainder remain. Thus (1) is a valid graph
estimate, while replacing its \(\Gamma^2\)-norm by the variance of
\(\Gamma\) is not a consequence of scalar variance alone.

For the original regulator take \(L=j^2\), \(a=(100j)^{-1}\), and
\(\theta_j=2\pi/(10^4j^2D_j)\), with \(D_j=j^4+O(j^2)\). Hence
\(\theta_j\asymp j^{-6}\). Put \[
 M=12L^2(2L+1),\qquad
 A_L=\frac{4L^2(2L+1)}3(4L^2+2L+1).
\] The exact covariance estimate from the actual vacuum is \[
 \|v_\Gamma\|\ge
 \frac{\sqrt\xi\,A_L^2}{720L^2M^2\sqrt r},
 \qquad r=1+\frac1{2L},\qquad \sqrt\xi=\frac1{2g^2},
 \tag{P3}
\] provided \(6g^2\sqrt r\,L^2M\le A_L/2\). This condition holds
eventually on both prescribed paths, since its left/right ratio is
\(O(g^2)\).

The proved fourth electric moment estimate (the preceding (C\_2(L,a))
bound) gives \[
 \|\Gamma^2\psi\|\le16L^4C_2(L,a)g^{-4},
 \tag{P4}
\] with that explicit \(C_2\). Along the displayed regulator,
\(K=3\sqrt{NM}/a=O(L^{7/2})\), \(a=O(L^{-1/2})\), and the exact
recurrence for \(B_4\) gives \(C_2(L,a)=O(L^6)\). Combining (1), (3),
and (4), the certified relative fourth-order error obeys \[
 \frac{\|P_\psi^\perp R_{4,\theta_j}\psi\|}
      {(2/3)\theta_j^2\|v_\Gamma\|}
 \le
 7680\,\frac{\theta_j^2L^6M^2C_2(L,a)\sqrt r}{A_L^2}\,g^{-2}
 =O(j^4g^{-2}).
 \tag{P5}
\] The power follows from \(L^6M^2C_2/A_L^2=O(L^8)=O(j^{16})\) and
\(\theta_j^2=O(j^{-12})\). Consequently, on \[
 g_j^2=\frac1{\log j},
 \qquad\text{or}\qquad
 g_j^2=\frac{\kappa_*}{200j},
 \tag{P6}
\] the right side of (5) is respectively \(O(j^4\log j)\) and
\(O(j^5)\). Both diverge. Therefore the available exact bounds do not
prove that the centered fourth term is small relative to the covariance
term; they become quantitatively vacuous precisely in the simultaneous
volume/coupling limits under consideration.

The sixth term is strictly less controlled. Taylor's scalar remainder is
of order \(\theta^6\gamma^3\) on a spin block, so a global estimate
requires a uniform bound for \(\|\Gamma^3\psi\|\) (and the centered
diagonal terms in (2) require corresponding mixed moments). The
recurrence for moments of \(W\) at a fixed box supplies finite constants
for every fixed order, but its constants grow with \(L\); it supplies no
volume-uniform concentration of the joint-spin coefficients. In
particular, a bound on \(\operatorname{Var}(\Gamma)\) alone cannot
control \(\operatorname{Var}(\Gamma^2)\) or the sixth moment:
probability measures with the same second moment can have arbitrarily
large fourth and sixth moments. Hence the exact cancellation in (2) does
not close without a new, volume-uniform weighted fourth/sixth-moment
theorem for the actual interacting vacuum.

This is a rigorous obstruction to the prescribed-path Taylor argument,
not a claim that the native low-energy weight vanishes there. The
finite-box scaled-angle theorem and the selected original-cusp sequence
remain valid, while nonzero native low-energy weight on either path and
an interacting four-dimensional continuum remain unproved.

\subsection{26.11 Volume scaling of the fixed-box scaled-angle
kernel}\label{volume-scaling-of-the-fixed-box-scaled-angle-kernel}

This note extracts a uniform consequence of the fixed-box limit
\(\theta_g/g\to\tau\) from Section 26.9. It does not assert a joint
\(L\to\infty\) theorem. All estimates below concern the exact Gaussian
kernel obtained after the fixed-box limit.

Let \(P=\mathscr R\mathscr R^{\mathsf T}\) and
\(K=\mathscr R\Lambda\mathscr R^{\mathsf T}\), acting separately on each
of the three colour components. Thus \(0\le K\le\sqrt{12}\,P\), and
\(\operatorname{tr}(KW)=3\operatorname{tr}(\Lambda D)\), where
\(W=\operatorname{diag}(n_2^2)\) on direction-one edges. Write \[
 \eta(q)=\sum_e n_2\,\operatorname{Ad}_{q_e}T_3\,\mathbf e_e,
 \qquad Y_\tau(q)=\delta(q)^{\mathsf T}\Lambda\delta(q)
        =4\tau^2\eta(q)^{\mathsf T}K\eta(q).
\] All expectations below use normalized product Haar probability. With
the normalization \(\|T_3\|=1\), Haar invariance gives \[
 \mathbb E\eta_e=0,
 \qquad \mathbb E(\eta_{e,\alpha}\eta_{e,\beta})
     =\frac{n_2^2}{3}\,\delta_{\alpha\beta},
 \qquad \|\eta_e\|^2=n_2^2.
\] Consequently \[
 \mathbb E Y_\tau=4\tau^2\operatorname{tr}(\Lambda D). \tag{V1}
\]

The exact trace count from (232) and \(\sigma_\nu<\sqrt{12}\) imply \[
 \frac{\mathcal A_L}{\sqrt{12}}\le \operatorname{tr}(\Lambda D)
 \le\sqrt{12}\operatorname{tr}D,
 \qquad
 \mathcal A_L=\frac{4}{3}L^2(2L+1)(4L^2+2L+1). \tag{V2}
\] For \(L\ge2\), \(\mathcal A_L\ge(32/3)L^5\) and
\(\sum_e n_2^4\le18L^7\). In particular \[
 \mathbb E Y_\tau\ge \frac{128}{3\sqrt{12}}\tau^2L^5,
 \qquad \operatorname{tr}(\Lambda D)\le72L^5. \tag{V3}
\]

We use the following elementary quadratic-form estimate. If \(\xi_e\)
are independent centred vectors with \(\|\xi_e\|^2\le w_e\),
\(\mathbb E\xi_e\xi_e^{\mathsf T}=(w_e/3)I_3\), and
\(0\le K\le\sqrt{12}I\), then expansion into diagonal and off-diagonal
terms, followed by independence and Cauchy--Schwarz, gives \[
 \operatorname{Var}(\sum_{e,f}\xi_e^{\mathsf T}K_{ef}\xi_f)
 \le 768\,\operatorname{tr}(K W K W)
 \le 27648\sum_e w_e^2. \tag{V4}
\] The constant is deliberately enlarged: the first inequality uses only
\(\mathbb E\|\xi_e\|^4\le w_e^2\), and the second uses
\(\|K\|\le\sqrt{12}\) and \(\operatorname{tr}(W^2)=3\sum_e w_e^2\) for
the three-colour block. Applying (V4) to \(\eta\) and multiplying by
\(16\tau^4\) yields \[
 \operatorname{Var}Y_\tau\le 442368\tau^4\sum_e n_2^4
 \le 7{,}962{,}624\,\tau^4L^7. \tag{V5}
\] Chebyshev and (V3) therefore imply, for every fixed \(\tau\ne0\), \[
 \mathbb P\left(Y_\tau\le \frac{1}{2}\mathbb EY_\tau\right)
 \le 10^6L^{-3}. \tag{V6}
\]

Let \(q,q'\) be independent and put
\(Y_\tau^\Delta=(\delta(q)-\delta(q'))^{\mathsf T}\Lambda(\delta(q)-\delta(q'))\).
The same argument applies to \(\eta(q)-\eta(q')\): its covariance is
doubled and its fourth moments are at most \(16n_2^4\). Enlarging the
constant in (V6) by a factor 32 gives \[
 \mathbb P\left(Y_\tau^\Delta\le \frac{1}{2}\mathbb E Y_\tau^\Delta\right)
 \le 3.2\cdot10^7L^{-3},
 \qquad
 \mathbb E Y_\tau^\Delta=8\tau^2\operatorname{tr}(\Lambda D). \tag{V7}
\]

The two-twirl kernel in (S9) is
\(J_{L,\tau}=\mathbb E\exp(-Y_\tau^\Delta/16)\), and
\(d_{L,\tau}=J_{L,\tau}-c_{L,\tau}^2\le J_{L,\tau}\). Since
\(e^{-x}\le1\) and \(e^{-x}\le e^{-\mathbb EY_\tau^\Delta/32}\) on the
complementary event, \[
 J_{L,\tau}
 \le 3.2\cdot10^7L^{-3}
   +\exp\!\left[-\frac{\tau^2\operatorname{tr}(\Lambda D)}{4}\right]
 \le 3.2\cdot10^7L^{-3}+e^{-\tau^2L^5/(\sqrt{12})}. \tag{V8}
\] Thus the fixed-\(\tau\) Gaussian native mass satisfies \[
 d_{L,\tau}\longrightarrow0\qquad(L\to\infty,\ \tau\ne0\ \text{fixed}). \tag{V9}
\] This is a genuine volume obstruction: the fixed-box limit has
positive mass for every \(L\), but that mass is not uniformly positive
in volume.

There is also a small-angle bound useful for the original cusp scales.
From \(1-e^{-x}\le x\), \(c_{L,\tau}\ge1-\mathbb EY_\tau/16\), hence \[
 d_{L,\tau}\le1-c_{L,\tau}^2
 \le\frac{\mathbb EY_\tau}{8}
 \le36\tau^2L^5. \tag{V10}
\] On the retained sequence \(L_j=j^2\),
\(\theta_j\sim(2\pi/10^4)j^{-6}\). Therefore the ideal oscillator
parameter \(\tau_j=\theta_j/g_j\) obeys \[
 \tau_j^2L_j^5=O(j^{-2}\log j)\quad(g_j^2=1/\log j),\qquad
 \tau_j^2L_j^5=O(j^{-1})\quad(g_j^2=\kappa_*/(200j)). \tag{V11}
\] The right side of (V10) consequently tends to zero on both prescribed
trajectories. Equation (V11) is an oscillator-kernel statement only; the
fixed-box convergence (S11) supplies no uniform error in \(L\), so it
cannot by itself be promoted to a theorem for the exact simultaneous
lattice sequence. Any such promotion would require a new uniform,
centred, finite-volume graph estimate.

\subsection{26.12 Exact finite-box small-angle expansion of the centered
Gaussian native
mass}\label{exact-finite-box-small-angle-expansion-of-the-centered-gaussian-native-mass}

This is a fixed-box consequence of the exact Gaussian kernel in Section
26.9. Fix \(L\ge2\) and \(a>0\), and write \(\delta(q)=\tau d(q)\),
where \(d(q)=-2R^*\eta(q)\). Put \[
u(q)=d(q)^{\mathsf T}\Lambda d(q),\quad z(q,q')=d(q)^{\mathsf T}\Lambda d(q'),\quad \mu=\mathbb E u,\quad \nu=\mathbb E u^2,\quad \zeta=\mathbb E z^2.
\] Since \(R^*\) has operator norm \(1\), \(\|\Lambda\|<\sqrt{12}\), and
\(\sum_{e=(n,1)}n_2^2=:S_{2,L}=\frac23L^2(L+1)(2L+1)^2\), we may take
the fully explicit bound \[
 U_L:=4\sqrt{12}\,S_{2,L},
\] so \(0\le u\le U_L\), \(|z|\le U_L\), and
\(0\le u(q)+u(q')-2z(q,q')\le4U_L\). For \(x\ge0\),
\(|e^{-x}-(1-x+x^2/2)|\le x^3/6\). Applying this bound to the two-twirl
kernel and to \(c(\tau)=\mathbb E e^{-\tau^2u/16}\), uniformly for
\(|\tau|\le1\), gives \[
J(\tau)=1-\frac{\mu}{8}\tau^2+\frac{\nu+\mu^2+2\zeta}{256}\tau^4+R_J,
\qquad |R_J|\le\frac{U_L^3}{384}|\tau|^6,
\] \[
c(\tau)=1-\frac{\mu}{16}\tau^2+\frac{\nu}{512}\tau^4+r_c,
\qquad |r_c|\le\frac{U_L^3}{24576}|\tau|^6.
\] Writing \(a_0=\mu/16\), \(b_0=\nu/512\), and \(R_c=U_L^3/24576\),
direct multiplication yields \[
\left|c(\tau)^2-\left(1-2a_0\tau^2+(a_0^2+2b_0)\tau^4\right)\right|
\le C_{c,L}|\tau|^6,
\] where \(C_{c,L}=2a_0b_0+b_0^2+2(1+a_0+b_0)R_c+R_c^2\). Therefore \[
d_{L,\tau}=J(\tau)-c(\tau)^2=\frac{\zeta}{128}\tau^4+R_{d,L}(\tau),
\qquad |R_{d,L}(\tau)|\le\left(\frac{U_L^3}{384}+C_{c,L}\right)|\tau|^6.
\]

All expectations use normalized product Haar probability. In the
retained \(T_\alpha\)-coordinate norm \(-2\operatorname{tr}(XY)\),
\(\|T_3\|=1\); the physical metric \(c=-\operatorname{tr}(XY)/2\) is
unchanged. Haar invariance gives covariance
\(\mathbb E[(\operatorname{Ad}_qT_3)_\alpha(\operatorname{Ad}_qT_3)_\beta]=\delta_{\alpha\beta}/3\).
Thus each colour block of \(d\) has covariance \((4/3)D\); summing the
three colour blocks, whose cross-colour covariances vanish, gives \[
\zeta=\frac{16}{3}\operatorname{tr}(\Lambda D\Lambda D).
\] Consequently \[
\boxed{d_{L,\tau}=\frac{\tau^4}{24}\operatorname{tr}(\Lambda D\Lambda D)+R_{d,L}(\tau)},
\qquad |R_{d,L}(\tau)|\le\left(\frac{U_L^3}{384}+C_{c,L}\right)|\tau|^6.
\] Since \(C_L^\Gamma=(3/32)\operatorname{tr}(\Lambda D\Lambda D)\), the
leading term is \((4/9)\tau^4C_L^\Gamma\). On \(L_j=j^2\) and
\(\theta_j\asymp j^{-6}\), its leading scale \(\tau_j^4L_j^7\) is
\(O(j^{-10}(\log j)^2)\) for \(g_j^2=1/\log j\), and \(O(j^{-8})\) for
\(g_j^2=\kappa_*/(200j)\). The preceding Taylor remainder is
regulator-dependent. A sharper exact finite-box identity gives a useful
uniform relative statement. Choose a fixed \(r\in\mathrm{SU}(2)\) with
\(\operatorname{Ad}_rT_3=-T_3\) and right-multiply every \(q'_e\) by
\(r\). Product Haar is invariant, \(u'\) is unchanged and
\(z\mapsto-z\); hence, with \(t=\tau^2\), \[
 d_{L,\tau}=\mathbb E\!\left[e^{-t(u(q)+u(q'))/16}
 \left(\cosh\!\frac{t\,z(q,q')}{8}-1\right)\right].
\] Writing \(b=t(u+u')/16\) and \(v=tz/8\), Cauchy--Schwarz gives
\(|v|\le b\). Therefore \[
 e^{-tU_L/8}\frac{t^2\zeta}{128}
 \le d_{L,\tau}\le\frac{t^2\zeta}{128},
 \qquad
 0\ge R_{d,L}(\tau)\ge
 -\frac{\tau^6U_L}{192}\operatorname{tr}(\Lambda D\Lambda D).
\] Thus on \(L_j=j^2\), whenever \(\tau_j^2L_j^5\to0\), the oscillator
kernel satisfies \[
 d_{L_j,\tau_j}\sim\frac{\tau_j^4}{24}
 \operatorname{tr}(\Lambda D\Lambda D).
\] For the two retained paths this condition follows from (V11), but the
fixed-box convergence (S11) still has no uniform-in-\(L\) error.
Consequently this remains an oscillator-kernel result and gives no
simultaneous exact-lattice or interacting-continuum theorem.

\subsection{26.13 Exact prescribed-path rates for the fixed-box
oscillator
kernel}\label{exact-prescribed-path-rates-for-the-fixed-box-oscillator-kernel}

This source derives only the exact finite-box Gaussian scaled-angle
kernel from Sections 26.9, 26.11 and 26.12. It does not identify that
kernel with the simultaneous exact interacting lattice state.

\subsection{Data retained exactly}\label{data-retained-exactly}

Set \(L_j=j^2\), \(a_j=(100j)^{-1}\), \[
 c_*:=\frac{2\pi}{10^4},\qquad
 D_j=L_{j^2}q_{j^2}+6m_{j^2}^2,\qquad
 \theta_j=\frac{c_*}{j^2D_j},
\] so the proved cusp asymptotic is \(D_j/j^4\to1\). Let \[
 \mathsf Q_L:=\operatorname{tr}(\Lambda D\Lambda D),
 \qquad
 C_L^\Gamma=\frac3{32}\mathsf Q_L.
\] The exact bounds already proved in Sections 24--26 imply \[
 \frac{128}{297}L^7\le \mathsf Q_L\le264L^7. \tag{PK1}
\] The exact small-angle bound in Section 26.12 uses \[
 U_L=4\sqrt{12}\,S_{2,L}
 =\frac{16\sqrt3}{3}L^2(L+1)(2L+1)^2,
\] therefore \(U_L/L^5\to64\sqrt3/3\).

For any \(\tau\in\mathbb R\), put \(t=\tau^2\). The sign-controlled cosh
identity gives the exact two-sided inequality \[
 e^{-tU_L/8}\frac{t^2\mathsf Q_L}{128}
 \le d_{L,\tau}
 \le\frac{t^2\mathsf Q_L}{128}. \tag{PK2}
\] Thus, for \(\tau\ne0\), \(d_{L,\tau}>0\), and with \(x=tU_L/8\), \[
0\le1-\frac{d_{L,\tau}}{t^2\mathsf Q_L/128}
 \le1-e^{-x}\le x. \tag{PK3}
\]

\subsection{Substitution along the two prescribed
paths}\label{substitution-along-the-two-prescribed-paths}

Use the ideal oscillator parameter \(\tau_j=\theta_j/g_j\), with
\(g_j>0\). Since \(L_j^7=j^{14}\), (PK1)--(PK2) give the exact
finite-\(j\) bounds \[
 e^{-x_j}\,\frac{c_*^4g_j^{-4}j^6}{297D_j^4}
 \le d_{L_j,\tau_j}
 \le \frac{33}{16}\,\frac{c_*^4g_j^{-4}j^6}{D_j^4},
 \qquad
 x_j:=\frac{\theta_j^2U_{L_j}}{8g_j^2}. \tag{PK4}
\] No replacement \(D_j=j^4\) is made in (PK4). Also \[
 x_j=\frac{2\sqrt3}{3}\,
 \frac{c_*^2(L_j+1)(2L_j+1)^2}{D_j^2g_j^2}. \tag{PK5}
\]

\subsubsection{Logarithmic path}\label{logarithmic-path}

For \(g_j^2=1/\log j\), (PK5) and \(D_j/j^4\to1\) give \[
 x_j=\frac{8\sqrt3}{3}c_*^2\frac{\log j}{j^2}(1+o(1))
 \longrightarrow0. \tag{PK6}
\] Consequently (PK3)--(PK4) prove the relative oscillator-kernel
estimate \[
 d_{L_j,\tau_j}
 =\frac{c_*^4(\log j)^2}{128j^8D_j^4}Q_{L_j}
 \,[1+O(\log j/j^2)], \tag{PK7}
\] where the \(O(\cdot)\) is an explicit one-sided bound from (PK3).
Using only (PK1), the strongest asymptotic constants currently justified
are \[
 \frac{c_*^4}{297}
 \le\liminf_{j\to\infty}\frac{j^{10}d_{L_j,\tau_j}}{(\log j)^2},
 \qquad
 \limsup_{j\to\infty}\frac{j^{10}d_{L_j,\tau_j}}{(\log j)^2}
 \le\frac{33}{16}c_*^4. \tag{PK8}
\] The bounds are strict-positive at every finite \(j\) and show
\(d_{L_j,\tau_j}=\Theta\big(j^{-10}(\log j)^2\big)\) in the
two-sided-bound sense.

\subsubsection{Fixed-electric-coefficient
path}\label{fixed-electric-coefficient-path}

For \(g_j^2=\kappa_*/(200j)\), with fixed \(\kappa_*>0\), (PK5) gives \[
 x_j=\frac{1600\sqrt3}{3\kappa_*}c_*^2\frac1j(1+o(1))
 \longrightarrow0. \tag{PK9}
\] The exact relative estimate and (PK1) imply \[
 d_{L_j,\tau_j}
 =\frac{40000c_*^4}{128\kappa_*^2j^6D_j^4}Q_{L_j}
 \,[1+O(1/j)], \tag{PK10}
\] and hence \[
 \frac{40000c_*^4}{297\kappa_*^2}
 \le\liminf_{j\to\infty}j^8d_{L_j,\tau_j},
 \qquad
 \limsup_{j\to\infty}j^8d_{L_j,\tau_j}
 \le\frac{82500c_*^4}{\kappa_*^2}. \tag{PK11}
\] Here \(82500=40000\cdot(264/128)=40000\cdot(33/16)\). Thus
\(d_{L_j,\tau_j}=\Theta(j^{-8})\) for this oscillator kernel in the
two-sided-bound sense.

\subsection{What this does and does not prove for the exact
lattice}\label{what-this-does-and-does-not-prove-for-the-exact-lattice}

Equations (PK4)--(PK11) are exact finite-box/oscillator conclusions.
They prove positivity and the leading small-\(\tau_j\) scaling after
taking the fixed-box \(g\to0\) limit. They do \textbf{not} imply the
same estimates for the exact positive-vacuum vector at the simultaneous
pair \((L_j,g_j)\): (S11) is pointwise in fixed \(L,a\), with no rate
uniform in \(L\), while \(L_j\to\infty\) and \(g_j\to0\) together. In
particular, no bound currently controls the difference between the exact
normalized native spectral measure and the oscillator measure after
division by the vanishing mass \(d_{L_j,\tau_j}\). The independent
centered Taylor estimate (P5) is volume-vacuous, \(O(j^4\log j)\) and
\(O(j^5)\), and the sixth/mixed moments have no volume-uniform bound.
Thus the new cosh inequality closes the oscillator-kernel relative error
but leaves the simultaneous exact-lattice and interacting-continuum
bridge obstructed by the missing uniform fixed-box convergence rate.

This is a precise obstruction, not a claim that either prescribed path
has zero native low-energy weight.

\subsection{27. Exact collapse of the finite-box weak-coupling disk at
fixed
coupling}\label{exact-collapse-of-the-finite-box-weak-coupling-disk-at-fixed-coupling}

This bounded section isolates the precise limitation of the
small-coupling argument on a growing open box. It retains the original
interacting \(SU(2)\) Wilson Hamiltonian, every face, the physical
spacing \(a\), and the positive coupling \(g=g_{\mathrm{YM}}\). It
proves neither a continuum mass gap nor its absence.

For \(L\ge2\), put \(m=2L\), and let \(P_L\) be the elementary faces of
the open box. Then \[
 M_L:=|P_L|=3m^2(m+1)=12L^2(2L+1). \tag{FC1}
\] On \(L^2(SU(2)^{E_L},d\lambda)\), with \[
 H_{L,g,a}=\frac{2g^2}{a}H_{0,L}
 +\frac{1}{2g^2a}\sum_{p\in P_L}(2-W_p),\qquad
 H_{0,L}=\sum_{e\in E_L}E_e, \tag{FC2}
\] define \(W_L=\sum_{p\in P_L}W_p\) and \(\xi=1/(4g^4)\). The exact
decomposition is \[
 H_{L,g,a}=\frac{2g^2}{a}(H_{0,L}-\xi W_L)
 +\frac{M_L}{g^2a}I. \tag{FC3}
\] The final term is the complete scalar \(2bM_L\), \(b=(2g^2a)^{-1}\);
it is retained and only shifts all eigenvalues equally.

Since \(|W_p|\le2\), \(\|W_L\|_\infty\le2M_L\). At the configuration
\(U_e=I\) for every edge all \(W_p=2\), so equality is attained: \[
 \|W_L\|_{\mathcal B}=2M_L. \tag{FC4}
\] With \(T_\alpha=-i\sigma_\alpha/2\), the first nonzero one-link
Casimir is \(3/4\), hence \[
 \operatorname{dist}\bigl(0,\operatorname{spec}(H_{0,L})\setminus\{0\}\bigr)
 =\frac34. \tag{FC5}
\] Consequently \[
 \|\xi W_L\|=\frac{M_L}{2g^4},\qquad
 \frac{\|\xi W_L\|}{3/4}
 =\frac{2M_L}{3g^4}
 =\frac{8L^2(2L+1)}{g^4}. \tag{FC6}
\] For every fixed finite \(g>0\), this ratio diverges with \(L\).

Take the circle \(\Gamma=\{|z|=3/8\}\). Equation (FC5) gives
\(\|(H_{0,L}-z)^{-1}\|\le8/3\) on \(\Gamma\), while \[
 H_{0,L}-\xi W_L-z
 =\bigl(I-\xi W_L(H_{0,L}-z)^{-1}\bigr)(H_{0,L}-z). \tag{FC7}
\] The global operator-norm Neumann argument therefore certifies
invertibility on \(\Gamma\), and the corresponding rank-one Riesz
projection, only when \[
 2|\xi|M_L\frac83<1,\qquad\text{i.e.}\qquad
 \boxed{|\xi|<\frac{3}{16M_L}
 =\frac{1}{64L^2(2L+1)}.} \tag{FC8}
\] This is a sufficient disk from this proof, not a claim that the true
analyticity radius equals its boundary.

Along \(L=j^2\), the right side of (FC8) is \[
 \frac{1}{64j^4(2j^2+1)}\sim\frac{1}{128j^6}. \tag{FC9}
\] At fixed \(g\), \(\xi=1/(4g^4)\) is constant, so for all sufficiently
large \(j\) it lies outside this sufficient disk. Equivalently,
remaining inside would require \(16j^4(2j^2+1)<g^4\), impossible for an
unbounded sequence at fixed finite \(g\).

Thus (FC1)--(FC9) prove a regulator-level obstruction to this
volume-uniform perturbative certification route: the norm of the
retained Wilson interaction grows exactly with the number of faces,
while the unperturbed gap remains \(3/4\). They do not prove that the
interacting eigenvalues lack a continuum limit, and they do not prove or
refute a Yang--Mills mass gap. A fixed-coupling continuum proof must
provide a different volume-uniform nonperturbative estimate while
retaining the full Hamiltonian and gauge-invariant subspace. The
arithmetic replay is the accompanying file
check\_fixed\_coupling\_analytic\_disk\_obstruction.py; it verifies
(FC1), (FC6), (FC8), and eventual fixed-coupling failure for
representative finite couplings.

\subsection{28. Current mathematical
conclusion}\label{current-mathematical-conclusion}

The fixed-coupling diagonal (12) retains the exact nonlinear geometric
map, actual state and measure refinement, configuration compactness and
the centered norm bound (22). The new volume-uniform proof (62)--(76)
determines its physical spectral consequence for every fixed
\(g_0^4\ge12288\): the actual magnetic excitation probability has
exactly zero weight on every bounded physical energy interval
eventually. This extends the former total-interaction restriction (47)
to a coupling interval independent of volume.

For every fixed \(g_0>0\), the unchanged fixed physical Wilson loop
states have the positive high-energy raw mass (82), threshold (83), and
failure of strong continuity (85) in any direct
configuration/correlation limit with those observables. This result
proves loss of a positive spectral fraction; outside (62) it does not
assert zero low-energy loop weight or determine the native magnetic
state's full spectral measure. Tightness for these loop states requires
\(g_j\to0\), as proved in (84).

The controlled varying-coupling diagonal (38) still has the full
concentration and amplitude estimates (36)--(43), projective-Haar limit
(59), and complete spectral escape (42). The explicit alternative (88)
retains a vanishing actual magnetic state bound and configuration
compactness while leaving both proved whole-spectrum escape domains. The
fixed-electric-coefficient alternative (91) retains a growing magnetic
interaction. The full weak-coupling variational calculation in Section
12 now improves its earlier constant state bound (92) to the vanishing
estimate (104), so its actual uncentered translated and vacuum
configuration measures have the same subsequential limits. Section 16
gives an explicit corrected cusp for both alternatives and proves
comparison of their full native spectral probabilities with the
corresponding actual electric covariance probabilities. Neither
probability has yet been proved tight or convergent to an interacting
spatial continuum.

The unresolved target is the actual local interacting continuum
construction and nonzero original native low-energy spectral weight
outside the proved exclusions, followed by an exact determination of
whether that construction gives the proposed four-dimensional mass-gap
counterexample. The repository supplies full proofs of the obtained
estimates, their constants, and all stated exact maps. No missing
covariance, tightness, or dynamical convergence statement has been used
as an assumption.

The exact trial integrals and gauge averaging in (93)--(99) bound the
actual mean plaquette defect by \(3\sqrt5\,g^2\). The retained
eigen-equation improves its full electric second moment in (100)--(101),
yielding the new native magnetic-state estimates (102)--(104) at all
positive couplings. Ordered non-Abelian face filling then proves (107).
On the explicit trajectory \(g_j^2=j^{-10}\), all actual vacuum
cylinders converge to the consistently constructed pure-gauge measure
(110)--(111), with zero centered variance for gauge-invariant bounded
cylinders. That exact limit describes the loss of those unchanged
observables; it supplies neither nonzero magnetic spectral probability
nor an interacting continuum field theory.

Sections 14--16 prove a local exponential lower bound for the exact
electric covariance, a quantitative all-angle lower bound for the
original native magnetic mass, an all-angle fourth-order graph
remainder, and a full physical spectral probability comparison on an
explicit corrected cusp. The correction preserves each finite
Hamiltonian, its actual vacuum, the original coordinate and metric
dictionary, and all raw amplitudes. The proved total variation error
tends to zero on both decreasing-coupling trajectories. Consequently
finite-energy weights and tightness of those corrected native
probabilities are exactly linked to those of the unchanged covariance
vectors; neither tightness nor nonzero low-energy covariance weight is
assumed or established by this comparison.

Sections 17--18 derive the full actual weighted-electric graph and
spectral map at fixed box from potential moments and the retained
eigen-equation. They calculate all physical pair-state weights and prove
a uniform native comparison-mass bound of order L\^{}7, with
finite-energy probability at most the exact rank bound (237). On the
original box and mesh diagonal that probability is O(j\^{}-5). The
precisely selected positive-coupling diagonal and strengthened original
cusp then give an actual native magnetic probability escaping every
bounded physical energy interval while the same actual finite-box gaps
close. Its full state and period maps, raw masses and physical-time
limits are proved. This excludes that particular selected state sequence
as the sought interacting continuum; it does not determine the
prescribed logarithmic or fixed-electric-coefficient paths, or exclude
all other native-state continuations.

Section 20 computes a second physical observable's complete raw
Gamma/Laguerre measure. Its separately selected macroscopic limit is
\(d\delta_0\), \(d=1-2^{-3/2}>0\); preservation of the original vacuum
orthogonality then precludes a unique-vacuum realization. Section 21
proves the exact fluid-to-curvature-and-current map, including finite
magnetic spacetime integral, invariant concentration, and the necessary
divergence (274). The displayed source core also gives the
shrinking-loop trace cluster set \([-2,2]\) by (275)--(277). The new
source's fluid existence proof has not been independently verified here,
and none of these classical maps constructs the remaining interacting
quantum continuum or a nonzero native low-energy spectral measure on
either prescribed coupling path.

Section 22 adds a different radial observable with actual raw Gamma
spectral mass on every \((0,\epsilon)\), retaining positive finite
couplings and the original physical energy scale. Its positive
double-commutator multiplier has an independently proved weighted limit
\(jR_g\mapsto c_\Gamma Q\), complete carrier insertion kernels, and
convergence of the actual first energy moment. The centered vacuum image
of this same observable has nonzero raw limit
\((3/2)c_\Gamma^2\delta_0\). Thus retaining that observable action on
the vacuum produces a precise second zero-energy vector, while the
original bounded local cylinders act by scalars on the carrier sector.
These results do not identify the radial family with the native magnetic
translation, either prescribed coupling path, or the desired local
interacting four-dimensional theory.

Section 23 proves an exact correspondence on one refined positive dyadic
sequence, retaining both earlier families of convergence tests. For
every fixed real compact profile \(h\), with \(H_h=\int h\), the
original local first-band vector and the three original radial
vacuum-image directions obey \[
 \left\|\ell_j^4v_{h,j}^{\dagger}
       -\frac{H_h}{50}\sum_{i=1}^3w_{i,j}^{\dagger}\right\|\to0.
\] The full raw cross inner product with one radial direction tends to
\(600\pi^2H_h\); its two squared norms tend to \(30000\pi^2\) and
\(36\pi^2H_h^2\), respectively. The complete seven-dimensional
compressed operator correspondence is (352). The orthogonal coordinate
morphism, polynomial-core cylinder map, and nonzero terms leaving the
compressed band are all proved. This connects the regular compact
curvature-profile map to the radial phase construction without
identifying their full local operator algebras or assuming a joint limit
for the retained higher-mode terms. The original interacting continuum
and native prescribed-coupling calculations remain unfinished.

Sections 24--25 prove the complete directional-electric spectral density
and the exact original weighted trace. The local density in (381) has
positive raw mass on every finite positive energy interval, while its
total unfiltered mass diverges; all positive-time vectors and energy
moments have their explicit common free transverse representation. The
noncompact original weight is treated by its exact component trace,
including the reflected cosine boundary terms. Its whole raw comparison
mass is asymptotic to \(3N^7I_{\rm lat}/2560\), its raw finite-energy
mass is asymptotic to \(\ell^7\Omega^5/(512000\pi^2a^2)\), and its
finite-energy probability is asymptotic to
\(a^5\Omega^5/(600\pi^2I_{\rm lat})\). These results are transferred to
one actual positive-coupling sequence with all previous stage tests
retained. The further original cusp depth in (445) makes the native
relative error \(o(a^5)\), proving the same precise native probability
coefficient and every raw factor in (448). Section 26 now treats the
original depth \(T=j^2\) on a common selected weak-coupling sequence.
The prescribed coupling trajectories and survival of a nonabelian
continuum interaction remain unresolved.

Section 26 proves the original unmodified native state's finite-energy
escape directly. At each fixed box and nonzero angle, the exact support
separation (456), all eigenfunction moments (459), and raw lower bound
(467) imply that its low-energy probability is smaller than every power
of the positive coupling. The noncentral raw mass also has the explicit
upper bound (473) proportional to \(g^2/\sin^2\theta\); the central
endpoint retains raw mass tending to one. The refined common sequence
(477)--(478) preserves the original cusp depth and all preceding stage
tests, and gives native low-energy probability \(o(a_j^5)\), compared
with the covariance coefficient \(a_j^5\Omega^5/(600\pi^2I_{\rm lat})\).
Equations (480)--(485) prove maximal spectral separation, asymptotic
vector orthogonality, and actual relative Taylor failure, together with
the exact two-angle map (486). No conclusion about either prescribed
coupling path or an interacting continuum is inferred from this selected
sequence.

\end{document}
